\input{preamble}

% OK, start here.
%
\begin{document}

\title{Variétés}


\maketitle

\phantomsection
\label{section-phantom}

\tableofcontents

\section{Introduction}
\label{section-introduction}

\noindent
Dans ce chapitre, nous commençons à étudier les variétés et, plus
généralement, les schémas sur un corps. Une référence fondamentale est \cite{EGA}.








\section{Notation}
\label{section-notation}

\noindent
Tout au long de ce chapitre, nous utilisons la lettre $k$ pour désigner le corps de base.










\section{Variétés}
\label{section-varieties}

\noindent
Dans le projet Stacks, nous utiliserons ce
qui suit comme notre définition d'une variété.

\begin{definition}
\label{definition-variety}
Supposons que $k$ soit un corps. Une {\it variété } est un
schéma $X$ sur $k$ tel que $X$ soit intègre et que le
morphisme structural $X \to \Spec(k)$ soit séparé et de type fini.
\end{definition}

\noindent
Cette définition présente le inconvénient suivant.
Supposons que $k'/k$ soit une extension de corps. Supposons
que $X$ soit une variété sur $k$. Alors le changement de base
$X_{k'} = X \times_{\Spec(k)} \Spec(k')$ n'est pas
nécessairement une variété sur $k'$. Ce phénomène (dans une
plus grande généralité) sera discuté en détail dans les sections
suivantes. Le produit de deux variétés n'est pas nécessairement
une variété (c'est vraiment le même phénomène). Voici un exemple.

\begin{example}
\label{example-product-not-a-variety}
Soit $k = \mathbf{Q}$. Soit $X =
\Spec(\mathbf{Q}(i))$ et $Y = \Spec(\mathbf{Q}(i))$.
Alors le produit $X \times_{\Spec(k)} Y$ des variétés $X$
et $Y$ n'est pas une variété, car il est réductible. (Il
est isomorphe à l'union disjointe de deux copies de $X$.)
\end{example}

\noindent
Si le corps de base est algébriquement clos cependant, alors le produit de
variétés est une variété. Cela découle des résultats du chapitre sur
l'algèbre, mais là-bas nous traitons des situations beaucoup plus générales.
Il existe également une preuve directe simple de cela que nous présentons ici.

\begin{lemma}
\label{lemma-product-varieties}
\begin{slogan}
Les produits de variétés sont des variétés sur des corps algébriquement clos.
\end{slogan}
Supposons que $k$ soit un corps algébriquement clos.
Supposons que $X$, $Y$ soient des variétés sur $k$.
Alors $X \times_{\Spec(k)} Y$ est une variété sur $k$.
\end{lemma}

\begin{proof}
Le morphisme $X \times_{\Spec(k)} Y \to \Spec(k)$
est de type fini et séparé parce qu'il est la
composition des morphismes $X \times_{\Spec(k)} Y
\to Y \to \Spec(k)$ qui sont séparés et de type
fini, voir morphismes, Lemme
\ref{morphisms-lemma-base-change-finite-type} et
\ref{morphisms-lemma-composition-finite-type} et
schémas, Lemme
\ref{schemes-lemma-separated-permanence}. Pour terminer
la démonstration, il suffit de montrer que $X
\times_{\Spec(k)} Y$ est intégral. Supposons que $X =
\bigcup_{i = 1, \ldots, n} U_i$, $Y = \bigcup_{j = 1,
\ldots, m} V_j$ soient des recouvrements affines ouverts
finis. Si nous pouvons montrer que chaque $U_i
\times_{\Spec(k)} V_j$ est intégral, alors nous avons terminé
par Propriétés, Lemme
\ref{properties-lemma-characterize-reduced},
\ref{properties-lemma-characterize-irreducible} et \ref{properties-lemma-characterize-integral}. Cela nous ramène au cas affine.

\medskip\noindent
Le cas affine se traduit par l'énoncé algébrique suivant : Supposons que $A$,
$B$ soient des anneaux intègres et des $k$ -algèbres de type fini. Alors $A
\otimes_k B$ est un anneau intègre. Pour arriver à une contradiction, supposons que
$$
(\sum\nolimits_{i = 1, \ldots, n} a_i \otimes b_i)
(\sum\nolimits_{j = 1, \ldots, m} c_j \otimes d_j) = 0
$$
dans $A \otimes_k B$ avec les deux facteurs non nuls dans $A
\otimes_k B$. Nous pouvons supposer que $b_1, \ldots, b_n$ sont
$k$-linéairement indépendants dans $B$, et que $d_1, \ldots, d_m$ sont
$k$-linéairement indépendants dans $B$. Bien sûr, nous pouvons aussi
supposer que $a_1$ et $c_1$ sont non nuls dans $A$. Par conséquent,
$D(a_1c_1) \subset \Spec(A)$ n'est pas vide. Par le Nullstellensatz de
Hilbert (Algèbre, Théorème \ref{algebra-theorem-nullstellensatz}) nous
pouvons trouver un idéal maximal $\mathfrak m \subset A$ contenu dans
$D(a_1c_1)$ et $A/\mathfrak m = k$ puisque $k$ est algébriquement
clos. Désignons par $\overline{a}_i, \overline{c}_j$ les classes
résiduelles de $a_i, c_j$ dans $A/\mathfrak m = k$. L'équation ci-dessus devient
$$
(\sum\nolimits_{i = 1, \ldots, n} \overline{a}_i b_i)
(\sum\nolimits_{j = 1, \ldots, m} \overline{c}_j d_j) = 0
$$
ce qui est une contradiction avec $\mathfrak m \in
D(a_1c_1)$, l'indépendance linéaire de $b_1, \ldots, b_n$
et $d_1, \ldots, d_m$, et le fait que $B$ est un domaine.
\end{proof}






\section{Variétés et applications rationnelles}
\label{section-varieties-rational-maps}

\noindent
Supposons que $k$ soit un corps. Supposons que $X$ et $Y$ soient des variétés
sur $k$. Nous utiliserons l'expression {\it application rationnelle de
variétés de $X$ vers $Y$ } pour désigner une $\Spec(k)$ -application
rationnelle du schéma $X$ vers le schéma $Y$ comme défini dans morphismes,
Définition \ref{morphisms-definition-rational-map}. Comme il est d'usage,
l'expression `` application rationnelle de variétés '' ne se réfère pas au corps de
base (commun) des variétés, même si pour les schémas généraux nous faisons la
distinction entre applications rationnelles et applications rationnelles sur une base donnée.

\medskip\noindent
Le titre de cette section fait référence au théorème fondamental suivant.

\begin{theorem}
\label{theorem-varieties-rational-maps}
Supposons que $k$ soit un corps. La catégorie des variétés
et des applications rationnelles dominantes est équivalente
à la catégorie des extensions de corps de type fini $K/k$.
\end{theorem}

\begin{proof}
Supposons que $X$ et $Y$ soient des variétés avec des points
génériques $x \in X$ et $y \in Y$. Rappelons que les applications
rationnelles dominantes de $X$ vers $Y$ sont exactement celles des
applications rationnelles qui envoient $x$ vers $y$ (morphismes,
Définition \ref{morphisms-definition-dominant-rational} et discussion
qui suit). Ainsi, étant donnée une application rationnelle dominante $X
\supset U \to Y$, nous obtenons une application des corps des fonctions
$$
k(Y) = \kappa(y) = \mathcal{O}_{Y, y}
\longrightarrow
\mathcal{O}_{X, x} = \kappa(x) = k(X)
$$
Inversement, une telle application de $k$ -algèbre (qui est
automatiquement locale puisque la source et le but sont des corps)
détermine (unique) une application rationnelle dominante par morphismes,
Lemme \ref{morphisms-lemma-rational-map-finite-presentation}. De cette façon,
nous obtenons un foncteur pleinement fidèle. Pour terminer la démonstration,
il suffit de montrer que toute extension de corps de type fini $K/k$ se trouve
dans l'image essentielle. Comme $K/k$ est de type fini, il existe une $k$
-algèbre de type fini $A \subset K$ telle que $K$ est le corps des fractions de
$A$. Alors $X = \Spec(A)$ est une variété dont le corps des fonctions est $K$.
\end{proof}

\noindent
Supposons que $k$ soit un corps. Supposons que $X$ et $Y$ soient des
variétés sur $k$. Nous utiliserons l'expression {\it $X$ et $Y$ sont
des variétés birationnelles } pour signifier que $X$ et $Y$ sont
$\Spec(k)$ -birationnelles comme défini dans morphismes, Définition
\ref{morphisms-definition-birational}. Comme d'habitude, l'expression ``
variétés birationnelles '' ne se réfère pas au corps de base (commun) des
variétés, même si pour les schémas irréductibles généraux nous faisons la
distinction entre être birationnel et être birationnel sur une base donnée.

\begin{lemma}
\label{lemma-birational-varieties}
Supposons que $X$ et $Y$ soient des variétés
sur un corps $k$. Ce qui suit est équivalent
\begin{enumerate}
\item $X$ et $Y$ sont des variétés
birationnelles, \item le corps des fonctions $k(X)$
et $k(Y)$ sont isomorphes, \item il existe des
ouverts non vides de $X$ et $Y$ qui sont isomorphes en
tant que variétés, \item il existe un ouvert $U \subset
X$ et un morphisme birationnel $U \to Y$ de variétés.
\end{enumerate}
\end{lemma}

\begin{proof}
Ceci est un cas spécial de morphismes, Lemme
\ref{morphisms-lemma-criterion-birational-finite-presentation}.
\end{proof}







\section{Changement de corps et anneaux locaux}
\label{section-local-rings}

\noindent
Quelques résultats préliminaires sur ce qui arrive aux
anneaux locaux lors d'une extension des corps de base.

\begin{lemma}
\label{lemma-change-fields-flat}
Supposons que $K/k$ soit une extension de corps. Supposons que $X$ soit un
schéma sur $k$ et définissons $Y = X_K$. Si $y \in Y$ avec image $x \in X$, alors
\begin{enumerate}
\item $\mathcal{O}_{X, x} \to \mathcal{O}_{Y, y}$
est un homomorphisme local d'anneaux fidèlement
plat, \item avec $\mathfrak p_0 = \Ker(\kappa(x)
\otimes_k K \to \kappa(y))$ nous avons $\kappa(y) =
\kappa(\mathfrak p_0)$, \item $\mathcal{O}_{Y, y} =
(\mathcal{O}_{X, x} \otimes_k K)_\mathfrak p$ où
$\mathfrak p \subset \mathcal{O}_{X, x} \otimes_k K$
est l'image réciproque de $\mathfrak p_0$. \item
nous avons $\mathcal{O}_{Y, y}/\mathfrak
m_x\mathcal{O}_{Y, y} = (\kappa(x) \otimes_k K)_{\mathfrak p_0}$
\end{enumerate}
\end{lemma}

\begin{proof}
Nous pouvons supposer que $X = \Spec(A)$ est affine. Alors $Y =
\Spec(A \otimes_k K)$. Comme $K$ est plat sur $k$, nous voyons
que $A \to A \otimes_k K$ est plat. Par conséquent, $Y \to X$
est plat et nous obtenons la première affirmation si nous
utilisons également Algèbre, Lemme
\ref{algebra-lemma-local-flat-ff}. La deuxième affirmation découle de
schémas, Lemme \ref{schemes-lemma-points-fibre-product}. Maintenant,
$y$ correspond à un idéal premier $\mathfrak q \subset A \otimes_k
K$ et $x$ à $\mathfrak r = A \cap \mathfrak q$. Alors $\mathfrak p_0$
est le noyau de l'application induite $\kappa(\mathfrak r) \otimes_k
K \to \kappa(\mathfrak q)$. L'application sur les anneaux locaux est
$$
A_\mathfrak r \longrightarrow (A \otimes_k K)_\mathfrak q
$$
Nous pouvons factoriser cette application
à travers $A_\mathfrak r \otimes_k K = (A
\otimes_k K)_{\mathfrak r}$ pour obtenir
$$
A_\mathfrak r \longrightarrow A_\mathfrak r \otimes_k K
\longrightarrow (A \otimes_k K)_\mathfrak q
$$
et puis la deuxième flèche est une localisation en un certain
nombre premier. Cet idéal premier est l'image réciproque de
$\mathfrak p_0$ (détails omis) et cela prouve (3). Pour voir (4) utilisez
(3) et que la localisation et $- \otimes_k K$ sont des foncteurs exacts.
\end{proof}

\begin{lemma}
\label{lemma-change-fields-algebraic-dim}
Notation comme dans le Lemme \ref{lemma-change-fields-flat}.
Supposons que $X$ soit localement de type fini sur $k$. Ensuite
$$
\dim(\mathcal{O}_{Y, y}/\mathfrak m_x\mathcal{O}_{Y, y}) =
\text{trdeg}_k(\kappa(x)) - \text{trdeg}_K(\kappa(y)) =
\dim(\mathcal{O}_{Y, y}) - \dim(\mathcal{O}_{X, x})
$$
\end{lemma}

\begin{proof}
Ceci est une reformulation de l'algèbre, Lemme
\ref{algebra-lemma-inequalities-under-field-extension}.
\end{proof}

\begin{lemma}
\label{lemma-change-fields-algebraic-unramified}
Notation comme dans le Lemme
\ref{lemma-change-fields-flat}. Supposons que $X$ soit
localement de type fini sur $k$, que $\dim(\mathcal{O}_{X, x}) =
\dim(\mathcal{O}_{Y, y})$ et que $\kappa(x) \otimes_k K$ soit
réduit (par exemple si $\kappa(x)/k$ est séparable ou que $K/k$ est
séparable). Alors $\mathfrak m_x \mathcal{O}_{Y, y} = \mathfrak m_y$.
\end{lemma}

\begin{proof}
(La phrase entre parenthèses découle de l'algèbre, Lemme
\ref{algebra-lemma-separable-extension-preserves-reducedness}.)
En combinant les lemmes \ref{lemma-change-fields-flat} et
\ref{lemma-change-fields-algebraic-dim}, nous voyons que
$\mathcal{O}_{Y, y}/\mathfrak m_x \mathcal{O}_{Y, y}$ a une
dimension $0$ et est réduit. Par conséquent, c'est un corps.
\end{proof}





\section{Schémas géométriquement réduits}
\label{section-geometrically-reduced}

\noindent
Si $X$ est un schéma réduit sur un corps, il peut arriver que
$X$ devienne non réduit après l'extension du corps de base.
Cela ne se produit pas pour les schémas géométriquement réduits.

\begin{definition}
\label{definition-geometrically-reduced}
Supposons que $k$ est un corps.
Supposons que $X$ est un schéma sur $k$.
\begin{enumerate}
\item Supposons que $x \in X$ soit un point. Nous
disons que $X$ est {\it géométriquement réduit en $x$
} si pour toute extension de corps $k'/k$ et tout
point $x' \in X_{k'}$ au-dessus de $x$ l'anneau local
$\mathcal{O}_{X_{k'}, x'}$ est réduit. \item Nous
disons que $X$ est {\it géométriquement réduit } sur $k$
si $X$ est géométriquement réduit en tout point de $X$.
\end{enumerate}
\end{definition}

\noindent
Cela peut sembler un peu mystérieux au début, mais c'est en
réalité la même chose que la notion discutée dans le chapitre sur
l'algèbre. Voici quelques résultats de base expliquant la
connexion. Voir également le Lemme
\ref{lemma-geometrically-reduced-dense-smooth-open}, qui implique que ``
géométriquement réduit '' est équivalent à `` génériquement lisse '' pour des variétés sur un corps.

\begin{lemma}
\label{lemma-geometrically-reduced-at-point}
\begin{slogan}
La réduction géométrique peut être vérifiée sur les anneaux locaux.
\end{slogan}
Supposons que $k$ soit un corps.
Supposons que $X$ soit un schéma
sur $k$. Soit $x \in X$. Les
énoncés suivants sont équivalents
\begin{enumerate}
\item $X$ est géométriquement réduit à $x$, et
\item l’anneau $\mathcal{O}_{X, x}$ est
géométriquement réduit sur $k$ (voir Algèbre,
Définition \ref{algebra-definition-geometrically-reduced}).
\end{enumerate}
\end{lemma}

\begin{proof}
Supposons (1). Cela implique en particulier que
$\mathcal{O}_{X, x}$ est réduit. Supposons que $k'/k$ soit une
extension de corps finie purement inséparable. Considérons
l'anneau $\mathcal{O}_{X, x} \otimes_k k'$. D'après l'algèbre,
le lemme \ref{algebra-lemma-p-ring-map}, son spectre est le
même que le spectre de $\mathcal{O}_{X, x}$. Par conséquent,
c'est également un anneau local (Algèbre, lemme
\ref{algebra-lemma-characterize-local-ring}). Donc, il existe un
point unique $x' \in X_{k'}$ au-dessus de $x$ et de
$\mathcal{O}_{X_{k'}, x'} \cong \mathcal{O}_{X, x} \otimes_k k'$, voir
le lemme \ref{lemma-change-fields-flat}. Par hypothèse, il s'agit d'un
anneau réduit. Ainsi, nous déduisons (2) d'après l'algèbre, lemme
\ref{algebra-lemma-geometrically-reduced-finite-purely-inseparable-extension}.

\medskip\noindent
Supposons (2). Supposons que $k'/k$ soit une extension
de corps. Puisque $\Spec(k') \to \Spec(k)$ est surjectif,
$X_{k'} \to X$ est aussi surjectif (morphismes, Lemme
\ref{morphisms-lemma-base-change-surjective}). Supposons
que $x' \in X_{k'}$ soit un point quelconque situé
au-dessus de $x$. L'anneau local $\mathcal{O}_{X_{k'}, x'}$
est une localisation de l'anneau $\mathcal{O}_{X, x}
\otimes_k k'$ selon le Lemme \ref{lemma-change-fields-flat}.
Par conséquent, il est réduit par hypothèse et (1) est prouvé.
\end{proof}

\noindent
La notion n'est pas intéressante en caractéristique zéro.

\begin{lemma}
\label{lemma-perfect-reduced}
Supposons que $X$ soit un schéma sur un corps parfait $k$ (par
exemple \ $k$ a une caractéristique nulle). Soit $x \in X$. Si
$\mathcal{O}_{X, x}$ est réduit, alors $X$ est géométriquement réduit
en $x$. Si $X$ est réduit, alors $X$ est géométriquement réduit sur $k$.
\end{lemma}

\begin{proof}
La première affirmation découle du Lemme
\ref{lemma-geometrically-reduced-at-point} et de
l'Algèbre, du Lemme
\ref{algebra-lemma-separable-extension-preserves-reducedness} et de
la définition d'un corps parfait (Algèbre, Définition
\ref{algebra-definition-perfect}). La seconde affirmation découle de la première.
\end{proof}

\begin{lemma}
\label{lemma-geometrically-reduced}
Supposons que $k$ soit un corps de caractéristique $p > 0$. Supposons
que $X$ soit un schéma sur $k$. Les énoncés suivants sont équivalents
\begin{enumerate}
\item $X$ est géométriquement réduit, \item
$X_{k'}$ est réduit pour toute extension de
corps $k'/k$, \item $X_{k'}$ est réduit pour
toute extension de corps finie purement
inséparable $k'/k$, \item $X_{k^{1/p}}$ est
réduit, \item $X_{k^{perf}}$ est réduit, \item
$X_{\bar k}$ est réduit, \item pour tout ouvert
affine $U \subset X$ l'anneau $\mathcal{O}_X(U)$
est géométriquement réduit (voir Algèbre,
Définition \ref{algebra-definition-geometrically-reduced}).
\end{enumerate}
\end{lemma}

\begin{proof}
Supposons (1). Alors pour chaque extension de corps $k'/k$ et
chaque point $x' \in X_{k'}$, l'anneau local de $X_{k'}$ en $x'$
est réduit. En d'autres termes, $X_{k'}$ est réduit. D'où (2).

\medskip\noindent
Supposons (2). Supposons que $U \subset X$ soit un ouvert affine.
Alors, pour toute extension de corps $k'/k$, le schéma $X_{k'}$
est réduit, donc $U_{k'} = \Spec(\mathcal{O}(U)\otimes_k k')$ est
réduit, donc $\mathcal{O}(U)\otimes_k k'$ est réduit (voir
Propriétés, Section \ref{properties-section-integral}). En d'autres
termes, $\mathcal{O}(U)$ est géométriquement réduit, donc (7) est vérifié.

\medskip\noindent
Supposons (7). Pour toute extension de corps $k'/k$, le
changement de base $X_{k'}$ est obtenu en collant les spectres
des anneaux $\mathcal{O}_X(U) \otimes_k k'$ où $U$ est un
ouvert affine dans $X$ (voir schémas, Section
\ref{schemes-section-fibre-products}). Donc $X_{k'}$ est réduit. Ainsi (1) est vrai.

\medskip\noindent
Cela prouve que (1), (2) et (7) sont équivalents. Ceux-ci sont
équivalents à (3), (4), (5) et (6) car nous pouvons appliquer
l'algèbre, le Lemme
\ref{algebra-lemma-geometrically-reduced-finite-purely-inseparable-extension} à $\mathcal{O}_X(U)$ pour $U \subset X$ ouvert affine.
\end{proof}

\begin{lemma}
\label{lemma-check-only-finite-inseparable-extensions}
Supposons que $k$ soit un corps de caractéristique $p > 0$. Supposons que $X$
soit un schéma sur $k$. Soit $x \in X$. Les énoncés suivants sont équivalents
\begin{enumerate}
\item $X$ est géométriquement réduit en $x$, \item
$\mathcal{O}_{X_{k'}, x'}$ est réduit pour chaque
extension de corps finie purement inséparable $k'$ de
$k$ et $x' \in X_{k'}$ le point unique au-dessus de $x$,
\item $\mathcal{O}_{X_{k^{1/p}}, x'}$ est réduit pour
$x' \in X_{k^{1/p}}$ le point unique au-dessus de $x$,
et \item $\mathcal{O}_{X_{k^{perf}}, x'}$ est réduit pour
$x' \in X_{k^{perf}}$ le point unique au-dessus de $x$.
\end{enumerate}
\end{lemma}

\begin{proof}
Notez que si $k'/k$ est purement inséparable, alors $X_{k'}
\to X$ induit un homéomorphisme sur les espaces
topologiques sous-jacents, voir Algèbre, Lemme
\ref{algebra-lemma-p-ring-map}. D'où l'unicité de $x'$
au-dessus de $x$ mentionnée dans l'énoncé. De plus, dans ce cas
$\mathcal{O}_{X_{k'}, x'} = \mathcal{O}_{X, x} \otimes_k k'$.
Par conséquent, le lemme découle du Lemme
\ref{lemma-geometrically-reduced-at-point} ci-dessus et du Lemme
\ref{algebra-lemma-geometrically-reduced-finite-purely-inseparable-extension} d'Algèbre.
\end{proof}

\begin{lemma}
\label{lemma-geometrically-reduced-upstairs}
Supposons que $k$ soit un corps. Supposons que $X$
soit un schéma sur $k$. Supposons que $k'/k$ soit
une extension de corps. Supposons que $x \in X$ soit
un point, et soit $x' \in X_{k'}$ un point situé
au-dessus de $x$. Les énoncés suivants sont équivalents
\begin{enumerate}
\item $X$ est géométriquement réduit à $x$,
\item $X_{k'}$ est géométriquement réduit à $x'$.
\end{enumerate}
En particulier, $X$ est géométriquement réduit sur $k$ si et
seulement si $X_{k'}$ est géométriquement réduit sur $k'$.
\end{lemma}

\begin{proof}
Il est clair que (1) implique (2). Supposons (2). Supposons que
$k''/k$ soit une extension de corps finie purement inséparable et
soit $x'' \in X_{k''}$ un point situé au-dessus de $x$ (en fait, il
est unique). Choisissez une extension de corps commune $k'''/k$ de
$k''/k$ et $k'/k$, voir corps, Lemme
\ref{fields-lemma-common-extension-field}. Choisissez un point $x''' \in X_{k'''}$
situé au-dessus à la fois de $x'$ et $x''$. Considérons l'application des anneaux locaux
$$
\mathcal{O}_{X_{k''}, x''} \longrightarrow \mathcal{O}_{X_{k'''}, x'''}.
$$
Ceci est un homomorphisme local plat d'anneaux et donc fidèlement
plat. Par (2), nous voyons que l'anneau local à droite est réduit.
Ainsi, d'après L'Algèbre, Lemme
\ref{algebra-lemma-descent-reduced}, nous concluons que
$\mathcal{O}_{X_{k''}, x''}$ est réduit. Ainsi, d'après le Lemme
\ref{lemma-check-only-finite-inseparable-extensions}, nous concluons que $X$ est géométriquement réduit en $x$.
\end{proof}

\begin{lemma}
\label{lemma-geometrically-reduced-any-base-change}
Supposons que $k$ soit un corps. Supposons que $X$, $Y$ soient des schémas sur $k$.
\begin{enumerate}
\item Si $X$ est géométriquement réduit à $x$, et
$Y$ est réduit, alors $X \times_k Y$ est réduit
en chaque point au-dessus de $x$. \item Si $X$ est
géométriquement réduit sur $k$ et $Y$ est réduit,
alors $X \times_k Y$ est réduit. \item Si $X$ et
$Y$ sont géométriquement réduits sur $k$, alors $X
\times_k Y$ est géométriquement réduit. \item Si
$k$ est parfait et $X$ et $Y$ sont réduits, alors
$X \times_k Y$ est réduit. \item Ajouter plus ici.
\end{enumerate}
\end{lemma}

\begin{proof}
Pour prouver (1), combinez le Lemme
\ref{lemma-geometrically-reduced-at-point} avec le Lemme
d'Algèbre
\ref{algebra-lemma-geometrically-reduced-any-reduced-base-change}.
Pour prouver (2), combinez le Lemme
\ref{lemma-geometrically-reduced} avec le Lemme d'Algèbre
\ref{algebra-lemma-geometrically-reduced-any-reduced-base-change}.
Pour prouver (3), notez que $(X \times_k Y)_{\overline{k}} =
X_{\overline{k}} \times_{\overline{k}} Y_{\overline{k}}$ et utilisez
(2) ainsi que le Lemme \ref{lemma-geometrically-reduced}. Pour prouver
(4), utilisez (3) combiné avec le Lemme \ref{lemma-perfect-reduced}.
\end{proof}

\begin{lemma}
\label{lemma-generic-points-geometrically-reduced}
Supposons que $k$ est un corps.
Supposons que $X$ est un schéma sur $k$.
\begin{enumerate}
\item Si $x' \leadsto x$ est une spécialisation et $X$
est géométriquement réduit en $x$, alors $X$ est
géométriquement réduit en $x'$. \item Si $x \in X$ tel
que (a) $\mathcal{O}_{X, x}$ est réduit, et (b) pour
chaque spécialisation $x' \leadsto x$ où $x'$ est un
point générique d'une composante irréductible de $X$, le
schéma $X$ est géométriquement réduit en $x'$, alors $X$
est géométriquement réduit en $x$. \item Si $X$ est
réduit et géométriquement réduit en tous les points
génériques des composantes irréductibles de $X$, alors $X$
est géométriquement réduit. \item Si $X$ est une variété sur
$k$, alors $X$ est géométriquement réduit si et seulement si
son corps des fonctions est une extension séparable de $k$.
\end{enumerate}
\end{lemma}

\begin{proof}
La partie (1) découle du Lemme
\ref{lemma-geometrically-reduced-at-point} et du fait que
si $A$ est une $k$-algèbre géométriquement réduite, alors
$S^{-1}A$ est une $k$-algèbre géométriquement réduite pour
tout sous-ensemble multiplicatif $S$ de $A$, voir Algèbre,
Lemme \ref{algebra-lemma-geometrically-reduced-permanence}.

\medskip\noindent
Soit $A = \mathcal{O}_{X, x}$. Les hypothèses (a) et (b) de (2)
impliquent que $A$ est réduit, et que $A_{\mathfrak q}$ est
géométriquement réduit sur $k$ pour chaque idéal premier minimal
$\mathfrak q$ de $A$. Par conséquent, $A$ est géométriquement
réduit sur $k$, voir Algèbre, Lemme
\ref{algebra-lemma-generic-points-geometrically-reduced}. Ainsi, $X$ est
géométriquement réduit en $x$, voir Lemme \ref{lemma-geometrically-reduced-at-point}.

\medskip\noindent
La partie (3) découle trivialement de la partie (2).
Enfin, (4) découle de (3) par Algèbre, Lemme
\ref{algebra-lemma-characterize-separable-field-extensions}.
\end{proof}

\begin{lemma}
\label{lemma-Noetherian-geometrically-reduced-at-point}
Supposons que $k$ soit un corps. Supposons
que $X$ soit un schéma sur $k$. Soit $x \in
X$. Supposons que $X$ soit localement
noethérien et géométriquement réduit en $x$.
Alors il existe un voisinage ouvert $U \subset
X$ de $x$ qui est géométriquement réduit sur $k$.
\end{lemma}

\begin{proof}
Supposons que $X$ soit localement noethérien et
géométriquement réduit en $x$. D'après le Lemme des
propriétés
\ref{properties-lemma-ring-affine-open-injective-local-ring}, nous
pouvons trouver un voisinage affine ouvert $U \subset X$ de $x$
tel que $R = \mathcal{O}_X(U) \to \mathcal{O}_{X, x}$ soit
injectif. Selon le Lemme \ref{lemma-geometrically-reduced-at-point},
l'hypothèse signifie que $\mathcal{O}_{X, x}$ est géométriquement
réduit sur $k$. D'après le Lemme d'algèbre
\ref{algebra-lemma-subalgebra-separable}, cela implique que $R$ est
géométriquement réduit sur $k$, ce qui implique à son tour que $U$ est géométriquement réduit.
\end{proof}

\begin{example}
\label{example-not-geometrically-reduced}
Supposons que $k = \mathbf{F}_p(s, t)$, c’est-à-dire une
extension purement transcendantale du corps premier.
Considérons la variété $X = \Spec(k[x, y]/(1 + sx^p +
ty^p))$. Soit $k'/k$ une extension quelconque telle que $s$
et $t$ possèdent toutes deux une racine $p$ dans $k'$.
Alors le changement de base $X_{k'}$ n’est pas réduit. À
savoir, l’anneau $k'[x, y]/(1 + s x^p + ty^p)$ contient
l’élément $1 + s^{1/p}x + t^{1/p}y$ dont la puissance $p$ est
zéro mais qui n’est pas nul (puisque l’idéal $(1 + sx^p + ty^p)$
ne contient certainement aucun élément non nul de degré $< p$).
\end{example}

\begin{lemma}
\label{lemma-finite-extension-geometrically-reduced}
Supposons que $k$ est un corps. Supposons que $X \to
\Spec(k)$ est localement de type fini. Supposons que $X$ a
un nombre fini de composantes irréductibles. Alors il
existe une extension finie purement inséparable $k'/k$
telle que $(X_{k'})_{red}$ est géométriquement réduit sur $k'$.
\end{lemma}

\begin{proof}
Pour prouver ce lemme, nous pouvons remplacer $X$ par sa
réduction $X_{red}$. Ainsi, nous pouvons supposer que $X$
est réduit et localement de type fini sur $k$. Supposons que
$x_1, \ldots, x_n \in X$ soient les points génériques des
composantes irréductibles de $X$. Notez que pour toute
extension algébrique purement inséparable $k'/k$, le
morphisme $(X_{k'})_{red} \to X$ est un homéomorphisme, voir
Algèbre, Lemme \ref{algebra-lemma-p-ring-map}. Ainsi, les
points $x'_1, \ldots, x'_n$ situés au-dessus de $x_1, \ldots,
x_n$ sont les points génériques des composantes irréductibles
de $(X_{k'})_{red}$. Comme $X$ est réduit, les anneaux locaux
$K_i = \mathcal{O}_{X, x_i}$ sont des corps, voir Algèbre,
Lemme \ref{algebra-lemma-minimal-prime-reduced-ring}. Comme
$X$ est localement de type fini sur $k$, les extensions de
corps $K_i/k$ sont des extensions de corps de type fini.
Enfin, les anneaux locaux $\mathcal{O}_{(X_{k'})_{red}, x'_i}$
sont les corps $(K_i \otimes_k k')_{red}$. Par l'Algèbre,
Lemme \ref{algebra-lemma-make-separable}, nous pouvons trouver
une extension finie purement inséparable $k'/k$ telle que $(K_i
\otimes_k k')_{red}$ soient des extensions séparables de corps
de $k'$. En particulier, chaque $(K_i \otimes_k k')_{red}$ est
géométriquement réduit sur $k'$ par l'Algèbre, Lemme
\ref{algebra-lemma-characterize-separable-field-extensions}. À ce
stade, le Lemme \ref{lemma-generic-points-geometrically-reduced}
partie (3) implique que $(X_{k'})_{red}$ est géométriquement réduit.
\end{proof}







\section{Schémas géométriquement connexes}
\label{section-geometrically-connected}

\noindent
Si $X$ est un schéma connexe sur un corps, il peut arriver
que $X$ devienne déconnecté après l'extension du corps de
base. Cela n'arrive pas pour les schémas géométriquement connexes.

\begin{definition}
\label{definition-geometrically-connected}
Supposons que $X$ soit un schéma sur le corps $k$. On dit que
$X$ est {\it géométriquement connexe } sur $k$ si le schéma
$X_{k'}$ est connecté pour toute extension de corps $k'$ de $k$.
\end{definition}

\noindent
Par convention, un espace topologique connexe est non vide ; par
conséquent a fortiori, les schémas géométriquement connexes sont non
vides. Voici un exemple d'une variété qui n'est pas géométriquement connexe.

\begin{example}
\label{example-not-geometrically-irreducible}
Soit $k = \mathbf{Q}$. Le schéma $X =
\Spec(\mathbf{Q}(i))$ est une variété sur
$\Spec(\mathbf{Q})$. Mais le changement de base
$X_{\mathbf{C}}$ est le spectre de $\mathbf{C}
\otimes_{\mathbf{Q}} \mathbf{Q}(i) \cong \mathbf{C} \times \mathbf{C}$
qui est l'union disjointe de deux copies de $\Spec(\mathbf{C})$. Donc
en fait, c'est un exemple d'une variété non géométriquement connexe.
\end{example}

\begin{lemma}
\label{lemma-geometrically-connected-check-after-extension}
Supposons que $X$ soit un schéma sur le corps $k$.
Supposons que $k'/k$ soit une extension de corps.
Alors $X$ est géométriquement connexe sur $k$ si et
seulement si $X_{k'}$ est géométriquement connexe sur $k'$.
\end{lemma}

\begin{proof}
Si $X$ est géométriquement connexe sur $k$, alors il est clair
que $X_{k'}$ est géométriquement connexe sur $k'$. Pour la
réciproque, pour toute extension de corps $k''/k$, il existe
une extension de corps commune $k'''/k'$ et $k'''/k''$, voir
corps, Lemme \ref{fields-lemma-common-extension-field}. Comme
le morphisme $X_{k'''} \to X_{k''}$ est surjectif (en tant que
changement de base d'un morphisme surjectif entre spectres de
corps), nous voyons que la connexité de $X_{k'''}$ implique la
connexité de $X_{k''}$. Ainsi, si $X_{k'}$ est géométriquement
connexe sur $k'$, alors $X$ est géométriquement connexe sur $k$.
\end{proof}

\begin{lemma}
\label{lemma-bijection-connected-components}
Supposons que $k$ soit un corps. Supposons
que $X$, $Y$ soient des schémas sur $k$.
Supposons que $X$ soit géométriquement
connexe sur $k$. Alors le morphisme de projection
$$
p : X \times_k Y \longrightarrow Y
$$
induit une bijection entre les composantes connexes.
\end{lemma}

\begin{proof}
Les fibres schématiques de $p$ sont connectées,
puisque ce sont des changements de base du
schéma géométriquement connexe $X$ par
extensions de corps. De plus, les fibres
schématiques sont homéomorphes aux fibres
ensemblistes, voir schémas, Lemme
\ref{schemes-lemma-fibre-topological}. Par morphismes,
Lemme
\ref{morphisms-lemma-scheme-over-field-universally-open},
l'application $p$ est ouverte. Ainsi, nous pouvons appliquer Topologie,
Lemme \ref{topology-lemma-connected-fibres-connected-components} pour conclure.
\end{proof}

\begin{lemma}
\label{lemma-affine-geometrically-connected}
Supposons que $k$ soit un corps. Supposons que $A$
soit une $k$ -algèbre. Alors $X = \Spec(A)$ est
géométriquement connexe sur $k$ si et seulement si $A$
est géométriquement connexe sur $k$ (voir Algèbre,
Définition \ref{algebra-definition-geometrically-connected} ).
\end{lemma}

\begin{proof}
Immédiat d'après les définitions.
\end{proof}

\begin{lemma}
\label{lemma-separably-closed-field-connected-components}
Supposons que $k'/k$ soit une extension de corps.
Supposons que $X$ soit un schéma sur $k$. Supposons
que $k$ soit séparablement algébriquement clos. Alors
le morphisme $X_{k'} \to X$ induit une bijection des
composantes connexes. En particulier, $X$ est
géométriquement connexe sur $k$ si et seulement si $X$ est connexe.
\end{lemma}

\begin{proof}
Puisque $k$ est séparablement algébriquement clos, nous
voyons que $k'$ est géométriquement connexe sur $k$, voir
Algèbre, Lemme
\ref{algebra-lemma-separably-closed-connected-implies-geometric}.
Par conséquent, $Z = \Spec(k')$ est géométriquement connexe sur
$k$ d'après le Lemme \ref{lemma-affine-geometrically-connected}
ci-dessus. Puisque $X_{k'} = Z \times_k X$, le résultat est un cas
particulier du Lemme \ref{lemma-bijection-connected-components}.
\end{proof}

\begin{lemma}
\label{lemma-characterize-geometrically-connected}
Supposons que $k$ soit un corps. Supposons que $X$
soit un schéma sur $k$. Supposons que $\overline{k}$
soit une clôture algébrique séparablement de $k$.
Alors $X$ est géométriquement connexe si et seulement
si le changement de base $X_{\overline{k}}$ est connexe.
\end{lemma}

\begin{proof}
Supposons que $X_{\overline{k}}$ soit connecté. Supposons
que $k'/k$ soit une extension de corps. Il existe une
extension de corps $\overline{k}'/\overline{k}$ telle que
$k'$ s'embarque dans $\overline{k}'$ comme une extension de
$k$. D'après le Lemme
\ref{lemma-separably-closed-field-connected-components}, nous voyons que
$X_{\overline{k}'}$ est connecté. Puisque $X_{\overline{k}'} \to X_{k'}$
est surjectif, nous concluons que $X_{k'}$ est connecté comme souhaité.
\end{proof}

\begin{lemma}
\label{lemma-descend-open}
Supposons que $k$ soit un corps. Supposons
que $X$ soit un schéma sur $k$. Supposons
que $A$ soit une $k$-algèbre. Supposons que
$V \subset X_A$ soit un ouvert
quasi-compact. Alors il existe une
$k$-subalgèbre de type fini $A' \subset A$ et un ouvert
quasi-compact $V' \subset X_{A'}$ tels que $V = V'_A$.
\end{lemma}

\begin{proof}
Nous remarquons que si $X$ est également quasi-séparé, cela
découle de Limites, Lemme \ref{limits-lemma-descend-opens}.
Supposons que $U_1, \ldots, U_n$ soit constitué de finitement
nombreux ouverts affines de $X$ tels que $V \subset \bigcup U_{i,
A}$. Disons $U_i = \Spec(R_i)$. Puisque $V$ est quasi-compact,
nous pouvons trouver finitement nombreux $f_{ij} \in R_i
\otimes_k A$, $j = 1, \ldots, n_i$ tels que $V = \bigcup_i
\bigcup_{j = 1, \ldots, n_i} D(f_{ij})$ où $D(f_{ij}) \subset U_{i,
A}$ est l'ouvert standard correspondant. (Nous ne prétendons pas
que $V \cap U_{i, A}$ est l'union des $D(f_{ij})$, $j = 1, \ldots,
n_i$.) Il est clair que nous pouvons trouver une $k$-sous-algèbre
de type fini $A' \subset A$ telle que $f_{ij}$ soit l'image de
certains $f_{ij}' \in R_i \otimes_k A'$. Posons $V' = \bigcup
D(f_{ij}')$ qui est un ouvert quasi-compact de $X_{A'}$. Notons $\pi
: X_A \to X_{A'}$ le morphisme canonique. Nous avons $\pi(V) \subset
V'$ comme $\pi(D(f_{ij})) \subset D(f_{ij}')$. Si $x \in X_A$ avec
$\pi(x) \in V'$, alors $\pi(x) \in D(f_{ij}')$ pour certains $i, j$
et nous voyons que $x \in D(f_{ij})$ en tant que $f_{ij}'$ s'envoie
vers $f_{ij}$. Ainsi, nous voyons que $V = \pi^{-1}(V')$ comme souhaité.
\end{proof}

\noindent
Supposons que $k$ soit un corps.
Supposons que $\overline{k}/k$ soit
une extension galoisienne
(éventuellement infinie). Par
exemple, $\overline{k}$ pourrait
être la clôture algébrique séparable
de $k$. Pour tout $\sigma \in
\text{Gal}(\overline{k}/k)$, nous
obtenons un automorphisme
correspondant $ \Spec(\sigma) :
\Spec(\overline{k}) \longrightarrow
\Spec(\overline{k}) $. Notez que
$\Spec(\sigma) \circ \Spec(\tau) = \Spec(\tau \circ
\sigma)$. Par conséquent, nous obtenons une action
$$
\text{Gal}(\overline{k}/k)^{opp} \times \Spec(\overline{k})
\longrightarrow
\Spec(\overline{k})
$$
du groupe opposé sur le schéma
$\Spec(\overline{k})$. Supposons que $X$ est un
schéma sur $k$. Puisque $X_{\overline{k}} =
\Spec(\overline{k}) \times_{\Spec(k)} X$ par définition,
nous voyons que l'action ci-dessus induit une action canonique
\begin{equation}
\label{equation-galois-action-base-change-kbar}
\text{Gal}(\overline{k}/k)^{opp} \times X_{\overline{k}}
\longrightarrow
X_{\overline{k}}.
\end{equation}

\begin{lemma}
\label{lemma-Galois-action-quasi-compact-open}
Supposons que $k$ soit un corps. Supposons que $X$ soit un
schéma sur $k$. Supposons que $\overline{k}$ soit une
extension de Galois (éventuellement infinie) de $k$. Supposons que
$V \subset X_{\overline{k}}$ soit un ouvert quasi-compact. Alors
\begin{enumerate}
\item il existe une sous-extension finie
$\overline{k}/k'/k$ et un ouvert quasi-compact $V'
\subset X_{k'}$ tels que $V = (V')_{\overline{k}}$, \item il
existe un sous-groupe ouvert $H \subset
\text{Gal}(\overline{k}/k)$ tels que $\sigma(V) = V$ pour tout $\sigma \in H$.
\end{enumerate}
\end{lemma}

\begin{proof}
D'après le Lemme \ref{lemma-descend-open}, il existe une sous-extension
finie $k'/k \subset \overline{k}$ et un ouvert $V' \subset X_{k'}$ qui se
rétracte sur $V$. Cela prouve (1). Puisque $\text{Gal}(\overline{k}/k')$ est
ouvert dans $\text{Gal}(\overline{k}/k)$, la partie (2) est également claire.
\end{proof}

\begin{lemma}
\label{lemma-closed-fixed-by-Galois}
Supposons que $k$ soit un corps. Supposons que $\overline{k}/k$ soit une
extension de Galois (éventuellement infinie). Supposons que $X$ soit un schéma
sur $k$. Que $\overline{T} \subset X_{\overline{k}}$ ait les propriétés suivantes
\begin{enumerate}
\item $\overline{T}$ est un sous-ensemble fermé de
$X_{\overline{k}}$, \item pour chaque $\sigma \in
\text{Gal}(\overline{k}/k)$ nous avons $\sigma(\overline{T}) = \overline{T}$.
\end{enumerate}
Alors il existe un sous-ensemble fermé $T \subset X$ dont
l'image réciproque dans $X_{\overline{k}}$ est $\overline{T}$.
\end{lemma}

\begin{proof}
Ce lemme se réduit immédiatement au cas où $X = \Spec(A)$
est affine. Dans ce cas, soit $\overline{I} \subset A
\otimes_k \overline{k}$ l'idéal radical correspondant à
$\overline{T}$. L'hypothèse (2) implique que
$\sigma(\overline{I}) = \overline{I}$ pour tout $\sigma \in
\text{Gal}(\overline{k}/k)$. Choisissez $x \in \overline{I}$. Il
existe une extension galoisienne finie $k'/k$ contenue dans
$\overline{k}$ telle que $x \in A \otimes_k k'$. Posez $G = \text{Gal}(k'/k)$. Posez
$$
P(T) = \prod\nolimits_{\sigma \in G} (T - \sigma(x)) \in (A \otimes_k k')[T]
$$
Il est clair que $P(T)$ est unitaire et est en fait un élément
de $(A \otimes_k k')^G[T] = A[T]$ (selon la théorie de Galois de
base). De plus, si nous écrivons $P(T) = T^d + a_1T^{d - 1} +
\ldots + a_d$ alors nous voyons que $a_i \in I := A \cap
\overline{I}$. En combinant $P(x) = 0$ et $a_i \in I$ nous
trouvons $x^d = - a_1 x^{d - 1} - \ldots - a_d \in I(A \otimes_k
\overline{k})$. Ainsi, $x$ est contenu dans le radical de $I(A
\otimes_k \overline{k})$. Par conséquent, $\overline{I}$ est le radical
de $I(A \otimes_k \overline{k})$ et en posant $T = V(I)$ est une solution.
\end{proof}

\begin{lemma}
\label{lemma-characterize-geometrically-disconnected}
Supposons que $k$ soit un corps.
Supposons que $X$ soit un schéma sur
$k$. Les éléments suivants sont équivalents
\begin{enumerate}
\item $X$ est géométriquement connexe,
\item pour toute extension de corps finie
séparable $k'/k$ le schéma $X_{k'}$ est connecté.
\end{enumerate}
\end{lemma}

\begin{proof}
Il découle immédiatement de la définition que (1)
implique (2). Supposons que $X$ ne soit pas
géométriquement connexe. Supposons que $k \subset
\overline{k}$ soit une clôture algébrique séparable de $k$.
D'après le lemme
\ref{lemma-characterize-geometrically-connected}, il s'ensuit que
$X_{\overline{k}}$ est déconnecté. Disons $X_{\overline{k}} = \overline{U} \amalg
\overline{V}$ avec $\overline{U}$ et $\overline{V}$ ouverts, fermés et non vides.

\medskip\noindent
Supposons que $W \subset X$ soit un ouvert quasi-compact
quelconque. Alors $W_{\overline{k}} \cap \overline{U}$ et
$W_{\overline{k}} \cap \overline{V}$ sont ouverts et
fermés dans $W_{\overline{k}}$. En particulier,
$W_{\overline{k}} \cap \overline{U}$ et $W_{\overline{k}}
\cap \overline{V}$ sont quasi-compacts, et d'après le Lemme
\ref{lemma-Galois-action-quasi-compact-open},
$W_{\overline{k}} \cap \overline{U}$ et $W_{\overline{k}} \cap
\overline{V}$ sont définis sur une sous-extension finie et
invariants sous un sous-groupe ouvert de
$\text{Gal}(\overline{k}/k)$. Nous utiliserons cela sans autre mention dans ce qui suit.

\medskip\noindent
Choisissez $W_0 \subset X$ quasi-compact ouvert de
telle sorte que $W_{0, \overline{k}} \cap
\overline{U}$ et $W_{0, \overline{k}} \cap
\overline{V}$ soient tous deux non vides. Choisissez
une sous-extension finie $\overline{k}/k'/k$ et une
décomposition $W_{0, k'} = U_0' \amalg V_0'$ en
parties ouvertes et fermées de sorte que $W_{0,
\overline{k}} \cap \overline{U} =
(U'_0)_{\overline{k}}$ et $W_{0, \overline{k}} \cap
\overline{V} = (V'_0)_{\overline{k}}$. Soit $H =
\text{Gal}(\overline{k}/k') \subset \text{Gal}(\overline{k}/k)$.
En particulier $\sigma(W_{0, \overline{k}} \cap \overline{U}) =
W_{0, \overline{k}} \cap \overline{U}$ et de même pour $\overline{V}$.

\medskip\noindent
Ayant choisi $W_0$, $k'$ comme ci-dessus, pour
chaque ouvert quasi-compact $W \subset X$ nous posons
$$
U_W =
\bigcap\nolimits_{\sigma \in H} \sigma(W_{\overline{k}} \cap \overline{U}),
\quad
V_W =
\bigcup\nolimits_{\sigma \in H} \sigma(W_{\overline{k}} \cap \overline{V}).
$$
Maintenant, puisque $W_{\overline{k}} \cap \overline{U}$ et
$W_{\overline{k}} \cap \overline{V}$ sont fixés par un sous-groupe ouvert
de $\text{Gal}(\overline{k}/k)$, nous voyons que l'union et
l'intersection ci-dessus sont finies. Par conséquent, $U_W$ et $V_W$ sont à la
fois ouverts et fermés. De plus, par construction $W_{\bar k} = U_W \amalg V_W$.

\medskip\noindent
Nous affirmons que si $W \subset W' \subset X$ sont
quasi-compacts ouverts, alors $W_{\overline{k}} \cap U_{W'}
= U_W$ et $W_{\overline{k}} \cap V_{W'} = V_W$. Vérification
omise. Ainsi, nous voyons qu'en définissant $U = \bigcup_{W
\subset X} U_W$ et $V = \bigcup_{W \subset X} V_W$ nous
obtenons $X_{\overline{k}} = U \amalg V$ est une union
disjointe de ouverts et fermés. Il est évident que $V$ est
non vide car il est construit en prenant des unions
(localement). D'autre part, $U$ est non vide puisqu'il contient
$W_0 \cap \overline{U}$ par construction. Enfin, $U, V \subset
X_{\bar k}$ sont fermés et $H$-invariants par construction.
Ainsi, par le Lemme \ref{lemma-closed-fixed-by-Galois} nous
avons $U = (U')_{\bar k}$, et $V = (V')_{\bar k}$ pour un certain
$U', V' \subset X_{k'}$ fermé. Il est clair que $X_{k'} = U'
\amalg V'$ et nous voyons que $X_{k'}$ est déconnecté comme souhaité.
\end{proof}

\begin{lemma}
\label{lemma-tricky}
Supposons que $k$ soit un corps. Supposons que $\overline{k}/k$ soit une
extension de Galois (éventuellement infinie). Supposons que $f : T \to X$
soit un morphisme de schémas sur $k$. Supposons que $T_{\overline{k}}$
soit connexe et $X_{\overline{k}}$ soit déconnecté. Alors $X$ est déconnecté.
\end{lemma}

\begin{proof}
Écrire $X_{\overline{k}} = \overline{U} \amalg \overline{V}$
avec $\overline{U}$ et $\overline{V}$ ouverts et fermés. Noter
$\overline{f} : T_{\overline{k}} \to X_{\overline{k}}$ le
changement de base de $f$. Puisque $T_{\overline{k}}$ est connexe,
nous voyons que $T_{\overline{k}}$ est contenu soit dans
$\overline{f}^{-1}(\overline{U})$, soit dans
$\overline{f}^{-1}(\overline{V})$. Dire $T_{\overline{k}} \subset \overline{f}^{-1}(\overline{U})$.

\medskip\noindent
Fixez un ouvert quasi-compact $W \subset X$. Il
existe une sous-extension galoisienne finie
$\overline{k}/k'/k$ telle que $\overline{U} \cap
W_{\overline{k}}$ et $\overline{V} \cap
W_{\overline{k}}$ proviennent d'ouverts quasi-compacts $U', V'
\subset W_{k'}$. Alors aussi $W_{k'} = U' \amalg V'$. Considérons
$$
U'' = \bigcap\nolimits_{\sigma \in \text{Gal}(k'/k)} \sigma(U'),
\quad
V'' = \bigcup\nolimits_{\sigma \in \text{Gal}(k'/k)} \sigma(V').
$$
Ce sont des invariants de Galois, ouverts et
fermés, et $W_{k'} = U'' \amalg V''$.
D'après le Lemme
\ref{lemma-closed-fixed-by-Galois}, nous obtenons des
ouverts et fermés $U_W, V_W \subset W$ tels que $U'' =
(U_W)_{k'}$, $V'' = (V_W)_{k'}$ et $W = U_W \amalg V_W$.

\medskip\noindent
Nous affirmons que si $W \subset W' \subset X$ sont
quasi-compacts ouverts, alors $W \cap U_{W'} = U_W$ et $W
\cap V_{W'} = V_W$. Vérification omise. Par conséquent,
nous voyons qu'en définissant $U = \bigcup_{W \subset X}
U_W$ et $V = \bigcup_{W \subset X} V_W$, nous obtenons $X =
U \amalg V$. Il est clair que $V$ est non vide car il est
construit en prenant des unions (localement). D'autre part,
$U$ est non vide puisqu'il contient $f(T)$ par construction.
\end{proof}

\begin{lemma}
\label{lemma-geometrically-connected-criterion}
\begin{reference}
\cite[IV Corollary 4.5.13.1(i)]{EGA}
\end{reference}
Supposons que $k$ est un corps. Supposons que $T \to X$ est un
morphisme de schémas sur $k$. Supposons que $T$ est
géométriquement connexe et $X$ connecté. Alors $X$ est géométriquement connexe.
\end{lemma}

\begin{proof}
Ceci est une reformulation
du Lemme \ref{lemma-tricky}.
\end{proof}

\begin{lemma}
\label{lemma-geometrically-connected-if-connected-and-point}
Supposons que $k$ soit un corps. Supposons que $X$ soit
un schéma sur $k$. Supposons que $X$ soit connexe et
possède un point $x$ tel que $k$ soit algébriquement clos
dans $\kappa(x)$. Alors $X$ est géométriquement connexe.
En particulier, si $X$ possède un point $k$-rationnel et
que $X$ est connexe, alors $X$ est géométriquement connexe.
\end{lemma}

\begin{proof}
Posons $T = \Spec(\kappa(x))$. Supposons que
$\overline{k}$ soit une clôture algébrique séparablement
de $k$. L'hypothèse sur $\kappa(x)/k$ implique que
$T_{\overline{k}}$ est irréductible, voir Algèbre, Lemme
\ref{algebra-lemma-field-extension-geometrically-irreducible}.
Donc, d'après le Lemme
\ref{lemma-geometrically-connected-criterion}, nous voyons que
$X_{\overline{k}}$ est connexe. D'après le Lemme
\ref{lemma-characterize-geometrically-connected}, nous concluons que $X$ est géométriquement connexe.
\end{proof}

\begin{lemma}
\label{lemma-inverse-image-connected-component}
Supposons que $K/k$ soit une extension de corps.
Supposons que $X$ soit un schéma sur $k$. Pour chaque
composante connexe $T$ de $X$, l'image inverse $T_K
\subset X_K$ est une union de composantes connexes de $X_K$.
\end{lemma}

\begin{proof}
Ceci est une déclaration purement topologique. On désigne par $p
: X_K \to X$ le morphisme de projection. Supposons que $T \subset
X$ est une composante connexe de $X$. Supposons que $t \in T_K =
p^{-1}(T)$. Soit $C \subset X_K$ une composante connexe
contenant $t$. Alors $p(C)$ est un sous-ensemble connexe de $X$ qui
rencontre $T$, donc $p(C) \subset T$. Par conséquent $C \subset T_K$.
\end{proof}

\noindent
Le lemme suivant sera supplanté par le lemme plus
fort \ref{lemma-image-connected-component} ci-dessous.

\begin{lemma}
\label{lemma-image-connected-component-finite-extension}
Supposons que $K/k$ soit une extension finie de corps et
soit $X$ un schéma sur $k$. On note par $p : X_K \to X$ le
morphisme de projection. Pour chaque composante connexe $T$
de $X_K$, l'image $p(T)$ est une composante connexe de $X$.
\end{lemma}

\begin{proof}
L'image $p(T)$ est contenue dans une certaine composante
connexe $X'$ de $X$. Considérons $X'$ comme un
sous-schéma fermé de $X$ de quelque manière que ce soit.
Alors $T$ est également une composante connexe de $X'_K =
p^{-1}(X')$ et nous pouvons donc supposer que $X$ est
connexe. Le morphisme $p$ est ouvert (morphismes, Lemme
\ref{morphisms-lemma-scheme-over-field-universally-open}),
fermé (morphismes, Lemme
\ref{morphisms-lemma-integral-universally-closed}) et les fibres de
$p$ sont des ensembles finis (morphismes, Lemme
\ref{morphisms-lemma-finite-quasi-finite}). Ainsi, nous pouvons appliquer
Topologie, Lemme \ref{topology-lemma-finite-fibre-connected-components} pour conclure.
\end{proof}

\begin{lemma}[Gabber]
\label{lemma-image-connected-component}
\begin{reference}
Email de Ofer Gabber daté du 4 juin 2016
\end{reference}
Supposons que $K/k$ soit une extension de corps. Supposons que $X$
soit un schéma sur $k$. Notons $p : X_K \to X$ le morphisme de
projection. Supposons que $\overline{T} \subset X_K$ soit une
composante connexe. Alors $p(\overline{T})$ est une composante connexe de $X$.
\end{lemma}

\begin{proof}
Lorsque $K/k$ est fini, c'est le Lemme
\ref{lemma-image-connected-component-finite-extension}.
En général, la preuve est plus difficile.

\medskip\noindent
Supposons que $T \subset X$ soit la composante connexe de $X$
contenant l'image de $\overline{T}$. Nous pouvons remplacer $X$ par
$T$ (avec la structure de sous-schéma réduit induite). Ainsi, nous
pouvons supposer que $X$ est connexe. Supposons que $A = H^0(X,
\mathcal{O}_X)$. Soit $L \subset A$ l'algèbre $k$ maximale faiblement
\' étale, voir More on Algèbre, Lemme
\ref{more-algebra-lemma-max-weakly-etale-subalgebra}. Comme $A$ n'a pas
d'idempotents non triviaux, nous voyons que $L$ est un corps et une
extension algébrique séparable de $k$ d'après More on Algèbre, Lemme
\ref{more-algebra-lemma-class-weakly-etale-over-field}. Observons que $L$
est également l'algèbre $L$ maximale faiblement \' étale de $A$ (car toute
algèbre faiblement \' étale $L$ est faiblement \' étale sur $k$ d'après
More on Algèbre, Lemme \ref{more-algebra-lemma-composition-weakly-etale}).
D'après schémas, Lemme \ref{schemes-lemma-morphism-into-affine}, nous
obtenons une factorisation $X \to \Spec(L) \to \Spec(k)$ du morphisme de structure.

\medskip\noindent
Supposons que $L'/L$ soit une extension
finie séparable. D'après la cohomologie des
schémas, Lemme
\ref{coherent-lemma-finite-locally-free-base-change-cohomology} nous avons
$$
A \otimes_L L' =
H^0(X \times_{\Spec(L)} \Spec(L'), \mathcal{O}_{X \times_{\Spec(L)} \Spec(L')})
$$
La $L'$-sous-algèbre faiblement \' etale maximale de $A \otimes_L
L'$ est $L \otimes_L L' = L'$, d'après Plus sur l'algèbre, lemme
\ref{more-algebra-lemma-change-fields-max-weakly-etale-subalgebra}.
En particulier, $A \otimes_L L'$ n'a pas d'idempotents non
triviaux (un tel idempotent engendrerait une sous-algèbre faiblement
\' etale) et nous concluons que $X \times_{\Spec(L)} \Spec(L')$ est
connexe. Par le lemme
\ref{lemma-characterize-geometrically-disconnected}, nous concluons que $X$ est géométriquement connexe sur $L$.

\medskip\noindent
Donnons à $\overline{T}$ la structure de
schéma induit réduite et considérons la composition
$$
\overline{T} \xrightarrow{i} X_K = X \times_{\Spec(k)} \Spec(K)
\xrightarrow{\pi}
\Spec(L \otimes_k K)
$$
L'image est contenue dans une composante connexe de $\Spec(L
\otimes_k K)$. Puisque $K \to L \otimes_k K$ est intégral,
nous voyons que les composantes connexes de $\Spec(L \otimes_k
K)$ sont des points et tous les points sont fermés, voir
Algèbre, Lemme \ref{algebra-lemma-integral-over-field}. Ainsi,
nous obtenons un corps quotient $L \otimes_k K \to E$ tel que
$\overline{T}$ se mappe dans $\Spec(E) \subset \Spec(L \otimes_k
K)$. Par conséquent $i(\overline{T}) \subset \pi^{-1}(\Spec(E))$. Mais
$$
\pi^{-1}(\Spec(E)) =
(X \times_{\Spec(k)} \Spec(K)) \times_{\Spec(L \otimes_k K)} \Spec(E) =
X \times_{\Spec(L)} \Spec(E)
$$
qui est connexe parce que $X$ est géométriquement connexe sur
$L$. Ensuite, nous obtenons l'égalité $\overline{T} = X
\times_{\Spec(L)} \Spec(E)$ (au sens des ensembles) et nous
concluons que $\overline{T} \to X$ est surjective comme souhaité.
\end{proof}

\noindent
Supposons que $X$ soit un schéma. Nous désignons
$\pi_0(X)$ l'ensemble des composantes connexes de $X$.

\begin{lemma}
\label{lemma-galois-action-connected-components}
Supposons que $k$ soit un corps, avec clôture
algébrique séparable $\overline{k}$. Supposons
que $X$ soit un schéma sur $k$. Il y a une action
$$
\text{Gal}(\overline{k}/k)^{opp} \times \pi_0(X_{\overline{k}})
\longrightarrow
\pi_0(X_{\overline{k}})
$$
avec les propriétés suivantes :
\begin{enumerate}
\item Un élément $\overline{T} \in
\pi_0(X_{\overline{k}})$ est fixé par l'action si et
seulement s'il existe une composante connexe $T \subset X$, qui
est géométriquement connexe sur $k$, telle que
$T_{\overline{k}} = \overline{T}$. \item Pour toute extension de corps
$k'/k$ avec clôture algébrique séparable $\overline{k}'$ le diagramme
$$
\xymatrix{
\text{Gal}(\overline{k}'/k') \times \pi_0(X_{\overline{k}'})
\ar[r] \ar[d] &
\pi_0(X_{\overline{k}'}) \ar[d] \\
\text{Gal}(\overline{k}/k) \times \pi_0(X_{\overline{k}})
\ar[r] &
\pi_0(X_{\overline{k}})
}
$$
est commutatif (où la flèche verticale droite est une bijection
selon le Lemme \ref{lemma-separably-closed-field-connected-components}).
\end{enumerate}
\end{lemma}

\begin{proof}
L'action (\ref{equation-galois-action-base-change-kbar})
de $\text{Gal}(\overline{k}/k)$ sur $X_{\overline{k}}$
induit une action sur ses composantes connexes. Les
composantes connexes sont toujours fermées (Topologie,
Lemme \ref{topology-lemma-connected-components}). Donc si
$\overline{T}$ est comme dans (1), alors par le Lemme
\ref{lemma-closed-fixed-by-Galois} il existe un sous-ensemble
fermé $T \subset X$ tel que $\overline{T} =
T_{\overline{k}}$. Notez que $T$ est géométriquement connexe
sur $k$, voir le Lemme
\ref{lemma-characterize-geometrically-connected}. Pour voir que $T$ est
une composante connexe de $X$, supposons que $T \subset T'$, $T \not =
T'$ où $T'$ est une composante connexe de $X$. Dans ce cas, $T'_{k'}$
contient strictement $\overline{T}$ et est donc déconnecté. Par le Lemme
\ref{lemma-tricky} cela signifie que $T'$ est déconnecté ! Contradiction.

\medskip\noindent
Nous omettons la preuve de la fonctorialité en (2).
\end{proof}

\begin{lemma}
\label{lemma-galois-action-connected-components-continuous}
Supposons que $k$ soit un corps, avec clôture
algébrique séparable $\overline{k}$.
Supposons que $X$ soit un schéma sur $k$. Supposons
\begin{enumerate}
\item $X$ est quasi-compact, et \item les
composantes connexes de $X_{\overline{k}}$ sont ouvertes.
\end{enumerate}
Alors
\begin{enumerate}
\item [(a)] $\pi_0(X_{\overline{k}})$ est fini,
et \item [(b)] l'action de
$\text{Gal}(\overline{k}/k)$ sur $\pi_0(X_{\overline{k}})$ est continue.
\end{enumerate}
De plus, les hypothèses (1) et (2) sont
satisfaites lorsque $X$ est de type fini sur $k$.
\end{lemma}

\begin{proof}
Puisque les composantes connexes sont ouvertes, couvrent
$X_{\overline{k}}$ (Topologie, Lemme
\ref{topology-lemma-connected-components}) et
$X_{\overline{k}}$ est quasi-compact, nous concluons qu'il
n'y en a qu'un nombre fini. Ainsi, (a) est vérifié. D'après
le Lemme \ref{lemma-descend-open}, ces composantes connexes
sont chacune définies sur une sous-extension finie de
$\overline{k}/k$ et nous obtenons (b). Si $X$ est de type
fini sur $k$, alors $X_{\overline{k}}$ est de type fini sur
$\overline{k}$ (morphismes, Lemme
\ref{morphisms-lemma-base-change-finite-type}). Par conséquent,
$X_{\overline{k}}$ est un schéma noethérien (morphismes, Lemme
\ref{morphisms-lemma-finite-type-noetherian}). Ainsi,
$X_{\overline{k}}$ a un nombre fini de composantes irréductibles
(Propriétés, Lemme \ref{properties-lemma-Noetherian-irreducible-components})
et a fortiori un nombre fini de composantes connexes (qui sont donc ouvertes).
\end{proof}









\section{Schémas géométriquement irréductibles}
\label{section-geometrically-irreducible}

\noindent
Si $X$ est un schéma irréductible sur un corps, il peut arriver
que $X$ devienne réductible après l'extension du corps de base.
Cela n'arrive pas pour les schémas géométriquement irréductibles.

\begin{definition}
\label{definition-geometrically-irreducible}
Supposons que $X$ soit un schéma sur le corps $k$. On dit que
$X$ est {\it géométriquement irréductible } sur $k$ si le
schéma $X_{k'}$ est irréductible \footnote { Un espace
irréductible est non vide. } pour toute extension de corps $k'$ de $k$.
\end{definition}

\begin{lemma}
\label{lemma-geometrically-irreducible-check-after-extension}
Supposons que $X$ soit un schéma sur le corps $k$.
Supposons que $k'/k$ soit une extension de corps. Alors
$X$ est géométriquement irréductible sur $k$ si et
seulement si $X_{k'}$ est géométriquement irréductible sur $k'$.
\end{lemma}

\begin{proof}
Si $X$ est géométriquement irréductible sur $k$, alors il est
clair que $X_{k'}$ est géométriquement irréductible sur $k'$.
Pour la réciproque, pour toute extension de corps $k''/k$, il
existe une extension de corps commune $k'''/k'$ et $k'''/k''$,
voir corps, Lemme \ref{fields-lemma-common-extension-field}. Comme
le morphisme $X_{k'''} \to X_{k''}$ est surjectif (en tant que
changement de base d'un morphisme surjectif entre spectres de
corps), nous voyons que l'irréductibilité de $X_{k'''}$ implique
l'irréductibilité de $X_{k''}$. Ainsi, si $X_{k'}$ est géométriquement
irréductible sur $k'$, alors $X$ est géométriquement irréductible sur $k$.
\end{proof}

\begin{lemma}
\label{lemma-separably-closed-irreducible}
Supposons que $X$ soit un schéma sur un corps
séparablement clos $k$. Si $X$ est irréductible, alors
$X_K$ est irréductible pour toute extension de corps
$K/k$. C.-à-d., $X$ est géométriquement irréductible sur $k$.
\end{lemma}

\begin{proof}
Utilisez les propriétés, Lemme
\ref{properties-lemma-characterize-irreducible} et Algèbre, Lemme \ref{algebra-lemma-separably-closed-irreducible}.
\end{proof}

\begin{lemma}
\label{lemma-bijection-irreducible-components}
Supposons que $k$ soit un corps. Supposons
que $X$, $Y$ soient des schémas sur $k$.
Supposons que $X$ soit géométriquement
irréductible sur $k$. Ensuite, le morphisme de projection
$$
p : X \times_k Y \longrightarrow Y
$$
induit une bijection entre les composantes irréductibles.
\end{lemma}

\begin{proof}
Tout d'abord, notez que les fibres schématiques de $p$
sont irréductibles, puisqu'elles sont des changements de
base du schéma géométriquement irréductible $X$ par des
extensions de corps. De plus, les fibres schématiques
sont homéomorphes aux fibres ensemblistes, voir schémas,
Lemme \ref{schemes-lemma-fibre-topological}. Selon
morphismes, Lemme
\ref{morphisms-lemma-scheme-over-field-universally-open}, l'application
$p$ est ouverte. Ainsi, nous pouvons appliquer Topologie, Lemme
\ref{topology-lemma-irreducible-fibres-irreducible-components} pour conclure.
\end{proof}

\begin{lemma}
\label{lemma-geometrically-irreducible-local}
\begin{slogan}
L'irréductibilité géométrique est Zariski locale modulo la connexité.
\end{slogan}
Supposons que $k$ soit un corps. Supposons que $X$ soit
un schéma sur $k$. Les éléments suivants sont équivalents
\begin{enumerate}
\item $X$ est géométriquement irréductible sur $k$, \item
pour chaque ouvert affine non vide $U$ l'algèbre $k$
$\mathcal{O}_X(U)$ est géométriquement irréductible sur $k$
(voir Algèbre, Définition
\ref{algebra-definition-geometrically-irreducible}), \item $X$
est irréductible et il existe un recouvrement ouvert affine $X =
\bigcup U_i$ tel que chaque algèbre $k$ $\mathcal{O}_X(U_i)$ est
géométriquement irréductible, et \item il existe un recouvrement
ouvert $X = \bigcup_{i \in I} X_i$ avec $I \not = \emptyset$
tel que $X_i$ est géométriquement irréductible pour chaque $i$ et
tel que $X_i \cap X_j \not = \emptyset$ pour tout $i, j \in I$.
\end{enumerate}
De plus, si $X$ est géométriquement irréductible, il en
est de même pour tout sous-schéma ouvert non vide de $X$.
\end{lemma}

\begin{proof}
Un schéma affine $\Spec(A)$ sur $k$ est géométriquement
irréductible si et seulement si $A$ est géométriquement
irréductible sur $k$ ; cela découle immédiatement des
définitions. Rappelons que si un schéma est irréductible, alors
tout sous-schéma ouvert non vide de $X$ l'est aussi, et que deux
ouverts non vides ont une intersection non vide. De plus, si tout
ouvert affine est irréductible, alors le schéma est irréductible,
voir Propriétés, Lemme
\ref{properties-lemma-characterize-irreducible}. Ainsi, l'énoncé final du
lemme est clair, ainsi que les implications (1) $\Rightarrow$ (2), (2)
$\Rightarrow$ (3), et (3) $\Rightarrow$ (4). Si (4) est vrai, alors pour
toute extension de corps $k'/k$, le schéma $X_{k'}$ possède un recouvrement
par des ouverts irréductibles dont les intersections deux à deux sont non
vides. Par conséquent, $X_{k'}$ est irréductible. Ainsi, (4) implique (1).
\end{proof}

\begin{lemma}
\label{lemma-geometrically-irreducible-function-field}
Supposons que $X$ soit un schéma irréductible sur le corps $k$. Supposons que
$\xi \in X$ soit son point générique. Les énoncés suivants sont équivalents :
\begin{enumerate}
\item $X$ est géométriquement irréductible sur $k$ ; \item
$\kappa(\xi)$ est géométriquement irréductible sur $k$ ; \item la
clôture algébrique séparable de $k$ dans $\kappa(\xi)$ est égale à $k$.
\end{enumerate}
\end{lemma}

\begin{proof}
Par l'algèbre, le lemme
\ref{algebra-lemma-geometrically-irreducible-separable-elements},
(2) et (3) sont équivalents.

\medskip\noindent
Supposons (1). Rappelons que $\mathcal{O}_{X, \xi}$ est
la limite colimitée filtrée de $\mathcal{O}_X(U)$ où $U$
parcourt les sous-schémas affines ouverts non vides de
$X$. En combinant le Lemme
\ref{lemma-geometrically-irreducible-local} et le Lemme
d'Algèbre
\ref{algebra-lemma-subalgebra-geometrically-irreducible}, nous voyons
que $\mathcal{O}_{X, \xi}$ est géométriquement irréductible sur $k$.
Puisque $\mathcal{O}_{X, \xi} \to \kappa(\xi)$ est une surjection avec
noyau localement nilpotent (voir le Lemme d'Algèbre
\ref{algebra-lemma-minimal-prime-reduced-ring}), il s'ensuit que $\kappa(\xi)$ est
géométriquement irréductible, voir le Lemme d'Algèbre \ref{algebra-lemma-p-ring-map}.

\medskip\noindent
Supposons (2). Nous pouvons supposer que $X$ est réduit. Supposons que
$U \subset X$ est un ouvert affine non vide. Alors $U = \Spec(A)$ où $A$
est un domaine avec corps des fractions $\kappa(\xi)$. Ainsi $A$ est
une $k$-sous-algèbre d'une $k$-algèbre géométriquement irréductible. Par
conséquent, d'après le Lemme d'Algèbre
\ref{algebra-lemma-subalgebra-geometrically-irreducible}, nous voyons que $A$ est
géométriquement irréductible sur $k$. D'après le Lemme
\ref{lemma-geometrically-irreducible-local}, nous en concluons que $X$ est géométriquement irréductible sur $k$.
\end{proof}

\begin{lemma}
\label{lemma-geometrically-irreducible-self-product}
Un schéma $X$ sur un corps $k$ est géométriquement
irréductible sur $k$ si et seulement si $X\times_k X$ est irréductible.
\end{lemma}

\begin{proof}
Si $X$ est géométriquement irréductible sur $k$,
alors $X\times_k X$ est irréductible selon le
Lemme \ref{lemma-bijection-irreducible-components}.

\medskip\noindent
Inversement, supposons que $X \times_k X$ soit
irréductible. En remplaçant $X$ par le schéma réduit
sous-jacent, nous pouvons supposer que $X$ est réduit sans
changer le problème. Puisque la première projection $X
\times_k X \to X$ est surjective, $X$ est également
irréductible, donc intégral. Supposons qu'il existe un élément
$\alpha$ dans le corps des fonctions $k(X)$ qui soit
algébrique séparable sur $k$ mais pas dans $k$. Soit $E = k(\alpha)$.

\medskip\noindent
Puis il existe un sous-schéma affine ouvert non vide $U$ de $X$
tel que $k \subset E \subset O(U)$, et donc $E \otimes_k E
\subset O(U)\otimes_k O(U) = O(U\times_k U)$. Ainsi nous avons
un morphisme dominant $U \times_k U \to \Spec(E\otimes_k E)$.
Puisque $X\times_k X$ est irréductible, $U \times_k U$ l'est
aussi, et donc $\Spec(E \otimes_k E)$ est connexe. Si $E$ est
strictement plus grand que $k$, alors $E \otimes_k E$ est un
produit de corps incluant $E$ comme un facteur. Il a une
dimension $(\dim_k(E))^2>\dim_k(E)$ comme espace vectoriel sur
$k$, et donc $E\otimes_k E$ est un produit d'au moins deux corps,
ce qui contredit que $\Spec(E \otimes_k E)$ est connexe. Donc en
fait $E$ est égal à $k$. Selon le Lemme
\ref{lemma-geometrically-irreducible-function-field}, $X$ est géométriquement irréductible sur $k$.
\end{proof}

\begin{lemma}
\label{lemma-separably-closed-field-irreducible-components}
Supposons que $k'/k$ est une extension de corps. Supposons que
$X$ est un schéma sur $k$. Posons $X' = X_{k'}$. Supposons que
$k$ est séparablement algébriquement clos. Alors le morphisme
$X' \to X$ induit une bijection des composantes irréductibles.
\end{lemma}

\begin{proof}
Puisque $k$ est séparablement algébriquement clos, nous voyons que
$k'$ est géométriquement irréductible sur $k$, voir Algèbre, Lemme
\ref{algebra-lemma-separably-closed-irreducible-implies-geometric}.
Par conséquent, $Z = \Spec(k')$ est géométriquement irréductible
sur $k$, selon le Lemme
\ref{lemma-geometrically-irreducible-local} ci-dessus. Puisque $X' = Z \times_k X$, le
résultat est un cas particulier du Lemme \ref{lemma-bijection-irreducible-components}.
\end{proof}

\begin{lemma}
\label{lemma-characterize-geometrically-irreducible}
\begin{slogan}
L'irréductibilité géométrique peut être testée sur
une clôture algébrique séparable du corps de base.
\end{slogan}
Supposons que $k$ soit un corps. Supposons que $X$
soit un schéma sur $k$. Ce qui suit est équivalent :
\begin{enumerate}
\item $X$ est géométriquement irréductible sur $k$,
\item pour toute extension de corps finie séparables
$k'/k$ le schéma $X_{k'}$ est irréductible, et \item
$X_{\overline{k}}$ est irréductible, où $k \subset
\overline{k}$ est une clôture algébrique séparables de $k$.
\end{enumerate}
\end{lemma}

\begin{proof}
Supposons que $X_{\overline{k}}$ soit irréductible, c'est-à-dire,
supposons (3). Supposons que $k'/k$ soit une extension de corps.
Il existe une extension de corps $\overline{k}'/\overline{k}$
telle que $k'$ s'implante dans $\overline{k}'$ comme une extension
de $k$. D'après le Lemme
\ref{lemma-separably-closed-field-irreducible-components}, nous voyons que
$X_{\overline{k}'}$ est irréductible. Puisque $X_{\overline{k}'} \to X_{k'}$ est
surjective, nous concluons que $X_{k'}$ est irréductible. Par conséquent, (1) est vrai.

\medskip\noindent
Supposons que $k \subset \overline{k}$ soit une clôture
algébrique séparable de $k$. Supposons que non (3),
c'est-à-dire supposons que $X_{\overline{k}}$ soit
réductible. Notre but est de montrer que $X_{k'}$ est
également réductible pour une certaine sous-extension
finie $\overline{k}/k'/k$. Supposons que $X = \bigcup_{i
\in I} U_i$ soit un recouvrement affine ouvert avec $U_i$
non vide. Si pour un certain $i$ le schéma $U_i$ est
réductible, ou si pour une certaine paire $i \not = j$
l'intersection $U_i \cap U_j$ est vide, alors $X$ est
réductible (Propriétés, Lemme
\ref{properties-lemma-characterize-irreducible}) et nous
avons terminé. En particulier, nous pouvons supposer que
$U_{i, \overline{k}} \cap U_{j, \overline{k}}$ pour tout $i, j
\in I$ est non vide et nous en concluons que $U_{i,
\overline{k}}$ doit être réductible pour un certain $i$. Selon
Algèbre, Lemme \ref{algebra-lemma-geometrically-irreducible},
cela signifie que $U_{i, k'}$ est réductible pour une certaine
extension de corps finie séparable $k'/k$. Par conséquent, $X_{k'}$
est également réductible. Ainsi, nous voyons que (2) implique (3).

\medskip\noindent
L'implication (1) $\Rightarrow$ (2)
est immédiate. Cela prouve le lemme.
\end{proof}

\begin{lemma}
\label{lemma-inverse-image-irreducible}
Supposons que $K/k$ soit une extension de corps.
Supposons que $X$ soit un schéma sur $k$. Pour chaque
composante irréductible $T$ de $X$, l'image inverse $T_K
\subset X_K$ est une union de composantes irréductibles de $X_K$.
\end{lemma}

\begin{proof}
Supposons que $T \subset X$ soit une composante irréductible de
$X$. Le morphisme $T_K \to T$ est plat, donc les généralisations
se relèvent le long de $T_K \to T$. Ainsi, chaque $\xi \in T_K$
qui est un point générique d'une composante irréductible de
$T_K$ se mappe au point générique $\eta$ de $T$. Si $\xi'
\leadsto \xi$ est une spécialisation dans $X_K$, alors $\xi'$ se
mappe à $\eta$ puisqu'il n'y a pas de points se spécialisant en
$\eta$ dans $X$. Ainsi $\xi' \in T_K$ et nous concluons que $\xi =
\xi'$. En d'autres termes, $\xi$ est le point générique d'une
composante irréductible de $X_K$. Cela signifie que les composantes
irréductibles de $T_K$ sont toutes des composantes irréductibles de $X_K$.
\end{proof}

\noindent
Pour un schéma $X$, nous notons
$\text{IrredComp}(X)$ l'ensemble des composantes irréductibles de $X$.

\begin{lemma}
\label{lemma-image-irreducible}
Supposons que $K/k$ soit une extension de corps.
Supposons que $X$ soit un schéma sur $k$. Pour chaque
composante irréductible $\overline{T} \subset X_K$,
l'image de $\overline{T}$ dans $X$ est une composante
irréductible dans $X$. Cela définit une application canonique
$$
\text{IrredComp}(X_K)
\longrightarrow
\text{IrredComp}(X)
$$
qui est surjective.
\end{lemma}

\begin{proof}
Considérez le diagramme
$$
\xymatrix{
X_K \ar[d] & X_{\overline{K}} \ar[d] \ar[l] \\
X & X_{\overline{k}} \ar[l]
}
$$
où $\overline{K}$ est la clôture algébrique séparable de $K$,
et où $\overline{k}$ est la clôture algébrique séparable de
$k$. D'après le lemme
\ref{lemma-separably-closed-field-irreducible-components}, le
morphisme $X_{\overline{K}} \to X_{\overline{k}}$ induit une bijection
entre les composantes irréductibles. Il suffit donc de démontrer le
lemme pour les morphismes $X_{\overline{k}} \to X$ et $X_{\overline{K}}
\to X_K$. En d'autres termes, nous pouvons supposer que $K = \overline{k}$.

\medskip\noindent
Le morphisme $p : X_{\overline{k}} \to X$ est intégral, plat et
surjectif. La platitude implique que les généralisations se relèvent
le long de $p$, voir morphismes, Lemme
\ref{morphisms-lemma-generalizations-lift-flat}. Par conséquent, les
points génériques des composantes irréductibles de $X_{\overline{k}}$ se
projettent sur des points génériques des composantes irréductibles de
$X$. L'intégralité implique que $p$ est universellement fermé, voir
morphismes, Lemme \ref{morphisms-lemma-integral-universally-closed}.
Ainsi, nous concluons que l'image $p(\overline{T})$ d'une composante
irréductible est un sous-ensemble fermé irréductible qui contient un point
générique d'une composante irréductible de $X$, donc $p(\overline{T})$ est
une composante irréductible de $X$. Cela prouve la première assertion. Si
$T \subset X$ est une composante irréductible, alors $p^{-1}(T) =T_K$ est
une union non vide de composantes irréductibles, voir Lemme
\ref{lemma-inverse-image-irreducible}. Chacune d'elles se projette nécessairement
sur $T$ selon la première partie. Par conséquent, l'application est surjective.
\end{proof}

\begin{lemma}
\label{lemma-irreducible-dense-rational-points-geometrically-irreducible}
Supposons que $k$ soit un corps. Supposons que $X$ soit un schéma
sur $k$. Si $X$ est irréductible et possède un ensemble dense de
points $k$-rationnels, alors $X$ est géométriquement irréductible.
\end{lemma}

\begin{proof}
Supposons que $k'/k$ soit une extension finie de corps et soit
$Z, Z' \subset X_{k'}$ des composantes irréductibles. Il
suffit de montrer $Z = Z'$, voir le Lemme
\ref{lemma-characterize-geometrically-irreducible}. Par le Lemme
\ref{lemma-image-irreducible}, nous avons $p(Z) = p(Z') = X$ où
$p : X_{k'} \to X$ est la projection. Si $Z \not = Z'$ alors $Z
\cap Z'$ est nulle part dense dans $X_{k'}$ et donc $p(Z \cap
Z')$ n'est pas dense par morphismes, Lemme
\ref{morphisms-lemma-image-nowhere-dense-finite} ; ici nous
utilisons aussi que $p$ est un morphisme fini comme le changement
de base du morphisme fini $\Spec(k') \to \Spec(k)$, voir
morphismes, Lemme \ref{morphisms-lemma-base-change-finite}. Ainsi
nous pouvons choisir un point rationnel $k$ $x \in X$ avec $x \not \in
p(Z \cap Z')$. Puisque le corps résiduel de $x$ est $k$, nous voyons
que $p^{-1}(\{x\}) = \{x'\}$ où $x' \in X_{k'}$ est un point dont le
corps résiduel est $k'$. Puisque $x \in p(Z) = p(Z')$, nous concluons
que $x' \in Z \cap Z'$, ce qui est la contradiction que nous cherchions.
\end{proof}

\begin{lemma}
\label{lemma-galois-action-irreducible-components}
Supposons que $k$ soit un corps, avec clôture
algébrique séparable $\overline{k}$. Supposons
que $X$ soit un schéma sur $k$. Il y a une action
$$
\text{Gal}(\overline{k}/k)^{opp} \times \text{IrredComp}(X_{\overline{k}})
\longrightarrow
\text{IrredComp}(X_{\overline{k}})
$$
avec les propriétés suivantes :
\begin{enumerate}
\item Un élément $\overline{T} \in
\text{IrredComp}(X_{\overline{k}})$ est fixé par l'action si et
seulement s'il existe une composante irréductible $T \subset X$,
qui est géométriquement irréductible sur $k$, telle que
$T_{\overline{k}} = \overline{T}$. \item Pour toute extension de corps
$k'/k$ avec clôture algébrique séparable $\overline{k}'$, le diagramme
$$
\xymatrix{
\text{Gal}(\overline{k}'/k') \times \text{IrredComp}(X_{\overline{k}'})
\ar[r] \ar[d] &
\text{IrredComp}(X_{\overline{k}'}) \ar[d] \\
\text{Gal}(\overline{k}/k) \times \text{IrredComp}(X_{\overline{k}})
\ar[r] &
\text{IrredComp}(X_{\overline{k}})
}
$$
est commutatif (où la flèche verticale droite est une bijection selon
le Lemme \ref{lemma-separably-closed-field-irreducible-components}).
\end{enumerate}
\end{lemma}

\begin{proof}
L'action ( \ref{equation-galois-action-base-change-kbar} ) de
$\text{Gal}(\overline{k}/k)$ sur $X_{\overline{k}}$ induit
une action sur ses composantes irréductibles. Les composantes
irréductibles sont toujours fermées (Topologie, Lemme
\ref{topology-lemma-connected-components}). Par conséquent, si
$\overline{T}$ est comme dans (1), alors par le Lemme
\ref{lemma-closed-fixed-by-Galois} il existe un sous-ensemble
fermé $T \subset X$ tel que $\overline{T} = T_{\overline{k}}$.
Notez que $T$ est géométriquement irréductible sur $k$, voir le
Lemme \ref{lemma-characterize-geometrically-irreducible}. Pour
voir que $T$ est une composante irréductible de $X$, supposons
que $T \subset T'$, $T \not = T'$ où $T'$ est une composante
irréductible de $X$. Supposons que $\overline{\eta}$ est le
point générique de $\overline{T}$. Il se mappe au point
générique $\eta$ de $T$. Ensuite, le point générique $\xi \in
T'$ se spécialise en $\eta$. Comme $X_{\overline{k}} \to X$ est
plat, il existe un point $\overline{\xi} \in X_{\overline{k}}$
qui se mappe vers $\xi$ et se spécialise en $\overline{\eta}$. Il
s'ensuit que l'adhérence du singleton $\{\overline{\xi}\}$ est un
sous-ensemble fermé irréductible de $X_{\overline{\xi}}$ qui
contient strictement $\overline{T}$. C'est la contradiction souhaitée.

\medskip\noindent
Nous omettons la preuve de la fonctorialité en (2).
\end{proof}

\begin{lemma}
\label{lemma-orbit-irreducible-components}
Supposons que $k$ soit un corps, avec clôture
algébrique séparable $\overline{k}$. Supposons que
$X$ soit un schéma sur $k$. Les fibres de l'application
$$
\text{IrredComp}(X_{\overline{k}})
\longrightarrow
\text{IrredComp}(X)
$$
des lemmes \ref{lemma-image-irreducible}
sont exactement les orbites de
$\text{Gal}(\overline{k}/k)$ sous l'action du lemme
\ref{lemma-galois-action-irreducible-components}.
\end{lemma}

\begin{proof}
Supposons que $T \subset X$ soit une composante
irréductible de $X$. Supposons que $\eta \in T$ soit
son point générique. D'après les lemmes
\ref{lemma-inverse-image-irreducible} et
\ref{lemma-image-irreducible}, les points génériques des
composantes irréductibles de $\overline{T}$ qui se
projettent sur $T$ se projettent sur $\eta$. Selon le lemme
d'algèbre \ref{algebra-lemma-Galois-orbit}, le groupe de
Galois agit transitivement sur tous les points de
$X_{\overline{k}}$ se projetant sur $\eta$. D'où le lemme en découle.
\end{proof}

\begin{lemma}
\label{lemma-galois-action-irreducible-components-locally-finite-type}
Supposons que $k$ soit un corps.
Supposons que $X \to \Spec(k)$ soit
localement de type fini. Dans ce cas
\begin{enumerate}
\item l'action
$$
\text{Gal}(\overline{k}/k)^{opp} \times \text{IrredComp}(X_{\overline{k}})
\longrightarrow
\text{IrredComp}(X_{\overline{k}})
$$
est continue si nous donnons à
$\text{IrredComp}(X_{\overline{k}})$ la topologie
discrète, \item chaque composante irréductible de
$X_{\overline{k}}$ peut être définie sur une extension finie
de $k$, et \item étant donnée une composante irréductible
$T \subset X$, le schéma $T_{\overline{k}}$ est une union
finie de composantes irréductibles de $X_{\overline{k}}$ qui
sont toutes dans le même $\text{Gal}(\overline{k}/k)$-orbite.
\end{enumerate}
\end{lemma}

\begin{proof}
Supposons que $\overline{T}$ soit une composante irréductible
de $X_{\overline{k}}$. Nous pouvons choisir un ouvert affine $U
\subset X$ tel que $\overline{T} \cap U_{\overline{k}}$ ne soit
pas vide. Écrivons $U = \Spec(A)$, ainsi $A$ est une
$k$-algèbre de type fini, voir morphismes, Lemme
\ref{morphisms-lemma-locally-finite-type-characterize}. Par
conséquent, $A_{\overline{k}}$ est une $\overline{k}$-algèbre de
type fini, et en particulier noethérienne. Supposons que $\mathfrak
p = (f_1, \ldots, f_n)$ soit l'idéal premier correspondant à
$\overline{T} \cap U_{\overline{k}}$. Comme $A_{\overline{k}} = A
\otimes_k \overline{k}$, nous voyons qu'il existe une sous-extension
finie $\overline{k}/k'/k$ telle que chaque $f_i \in A_{k'}$. Il est clair
que $\text{Gal}(\overline{k}/k')$ fixe $\overline{T}$, ce qui prouve (1).

\medskip\noindent
La partie (2) suit en appliquant le Lemme
\ref{lemma-galois-action-irreducible-components} (1) à
la situation sur $k'$, ce qui implique que la composante
irréductible $\overline{T}$ est de la forme
$T'_{\overline{k}}$ pour un certain irréductible $T' \subset X_{k'}$.

\medskip\noindent
Pour prouver (3), soit $T \subset X$ une composante
irréductible. Choisissez une composante irréductible
$\overline{T} \subset X_{\overline{k}}$ qui se mappe sur
$T$, voir le Lemme \ref{lemma-image-irreducible}. D'après ce
qui précède, l'orbite de $\overline{T}$ est finie, disons
qu'elle est $\overline{T}_1, \ldots, \overline{T}_n$. Alors
$\overline{T}_1 \cup \ldots \cup \overline{T}_n$ est un
sous-ensemble fermé invariant par $\text{Gal}(\overline{k}/k)$
de $X_{\overline{k}}$ et donc de la forme $W_{\overline{k}}$
pour un certain $W \subset X$ fermé selon le Lemme
\ref{lemma-closed-fixed-by-Galois}. Clairement $W = T$ et nous avons gagné.
\end{proof}

\begin{lemma}
\label{lemma-finite-extension-geometrically-irreducible-components}
Supposons que $k$ soit un corps. Supposons que $X
\to \Spec(k)$ soit localement de type fini.
Supposons que $X$ ait un nombre fini de composantes
irréductibles. Alors il existe une extension finie
séparable $k'/k$ telle que chaque composante irréductible
de $X_{k'}$ soit géométriquement irréductible sur $k'$.
\end{lemma}

\begin{proof}
Supposons que $\overline{k}$ est une clôture algébrique
séparable de $k$. L'hypothèse que $X$ a un nombre fini de
composantes irréductibles combinée avec le Lemme
\ref{lemma-galois-action-irreducible-components-locally-finite-type}
(3) montre que $X_{\overline{k}}$ a un nombre fini de
composantes irréductibles $\overline{T}_1, \ldots,
\overline{T}_n$. Par le Lemme
\ref{lemma-galois-action-irreducible-components-locally-finite-type} (2), il existe
une extension finie $\overline{k}/k'/k$ et des composantes irréductibles $T_i
\subset X_{k'}$ telles que $\overline{T}_i = T_{i, \overline{k}}$ et nous avons terminé.
\end{proof}

\begin{lemma}
\label{lemma-irreducible-components-geometrically-irreducible}
Supposons que $X$ soit un schéma sur le corps
$k$. Supposons que $X$ ait un nombre fini de
composantes irréductibles, toutes géométriquement
irréductibles. Alors $X$ a un nombre fini de
composantes connexes, chacune étant géométriquement connexe.
\end{lemma}

\begin{proof}
Ceci est clair parce qu'une composante connexe est
une union de composantes irréductibles. Détails omis.
\end{proof}







\section{Schémas géométriquement intègres}
\label{section-geometrically-integral}

\noindent
Si $X$ est un schéma intègre sur un corps, il peut arriver que $X$
devienne soit non réduit soit réductible après extension du corps de
base. Cela n'arrive pas pour les schémas géométriquement intègres.

\begin{definition}
\label{definition-geometrically-integral}
Supposons que $X$ soit un schéma sur le corps $k$.
\begin{enumerate}
\item Soit $x \in X$. Nous disons que $X$ est {\it
géométriquement ponctuellement intègre en $x$ } si pour toute
extension de corps $k'/k$ et tout $x' \in X_{k'}$ au-dessus de $x$,
l'anneau local $\mathcal{O}_{X_{k'}, x'}$ est intègre. \item Nous
disons que $X$ est {\it géométriquement ponctuellement intègre } si
$X$ est géométriquement ponctuellement intègre en chaque point. \item
Nous disons que $X$ est {\it géométriquement intègre } sur $k$ si le
schéma $X_{k'}$ est intègre pour toute extension de corps $k'$ de $k$.
\end{enumerate}
\end{definition}

\noindent
La distinction entre les notions (2) et (3) est nécessaire.
Par exemple, si $k = \mathbf{R}$ et $X =
\Spec(\mathbf{C}[x])$, alors $X$ est géométriquement point par point
intégral sur $\mathbf{R}$ mais bien sûr pas géométriquement intègre.

\begin{lemma}
\label{lemma-geometrically-integral}
Supposons que $k$ est un corps. Supposons
que $X$ est un schéma sur $k$. Alors $X$ est
géométriquement intègre sur $k$ si et
seulement si $X$ est à la fois géométriquement
réduit et géométriquement irréductible sur $k$.
\end{lemma}

\begin{proof}
Voir les propriétés, Lemme \ref{properties-lemma-characterize-integral}.
\end{proof}

\begin{lemma}
\label{lemma-proper-geometrically-reduced-global-sections}
Supposons que $k$ soit un corps. Supposons que $X$ soit un schéma propre sur $k$.
\begin{enumerate}
\item $A = H^0(X, \mathcal{O}_X)$ est une $k$ -algèbre
de dimension finie, \item $A = \prod_{i = 1, \ldots,
n} A_i$ est un produit d'$k$ -algèbres locales
artiniennes, un facteur pour chaque composante connexe
de $X$, \item si $X$ est réduit, alors $A = \prod_{i =
1, \ldots, n} k_i$ est un produit de corps, chacun
étant une extension finie de $k$, \item si $X$ est
géométriquement réduit, alors $k_i$ est fini séparable
sur $k$, \item si $X$ est géométriquement connexe, alors
$A$ est géométriquement irréductible sur $k$, \item si
$X$ est géométriquement irréductible, alors $A$ est
géométriquement irréductible sur $k$, \item si $X$ est
géométriquement réduit et géométriquement connexe, alors $A =
k$, et \item si $X$ est géométriquement intègre, alors $A = k$.
\end{enumerate}
\end{lemma}

\begin{proof}
Par la cohomologie des schémas, le lemme
\ref{coherent-lemma-proper-over-affine-cohomology-finite},
nous voyons que $A = H^0(X, \mathcal{O}_X)$ est
une $k$-algèbre de dimension finie. Cela prouve (1).

\medskip\noindent
Alors $A$ est un produit des $k$-algèbres artiniennes locales par
l'algèbre, le Lemme
\ref{algebra-lemma-finite-dimensional-algebra} et la Proposition
\ref{algebra-proposition-dimension-zero-ring}. Si $X = Y \amalg Z$
avec $Y$ et $Z$ ouverts dans $X$, alors nous obtenons un idempotent $e
\in A$ en prenant la section de $\mathcal{O}_X$ qui est $1$ sur $Y$ et
$0$ sur $Z$. Réciproquement, si $e \in A$ est un idempotent, alors nous
obtenons une décomposition correspondante de $X$. Enfin, comme $X$ a un
espace topologique sous-jacent noethérien, ses composantes connexes
sont ouvertes. Par conséquent, les composantes connexes de $X$
correspondent de $1$ à $1$ avec des idempotents primitifs de $A$. Cela prouve (2).

\medskip\noindent
Si $X$ est réduit, alors $A$ est réduit. Par
conséquent, les anneaux locaux $A_i = k_i$ sont réduits
et donc corps (par exemple d'après Algèbre, Lemme
\ref{algebra-lemma-minimal-prime-reduced-ring}). Cela prouve (3).

\medskip\noindent
Si $X$ est géométriquement réduit, alors $A
\otimes_k \overline{k} = H^0(X_{\overline{k}},
\mathcal{O}_{X_{\overline{k}}})$ (égalité par la
cohomologie des schémas, Lemme
\ref{coherent-lemma-flat-base-change-cohomology})
est réduit. Cela implique que $k_i \otimes_k
\overline{k}$ est un produit de corps et donc $k_i/k$
est séparable, par exemple d'après Algèbre, Lemmata
\ref{algebra-lemma-characterize-separable-field-extensions}
et
\ref{algebra-lemma-geometrically-reduced-finite-purely-inseparable-extension}. Cela prouve (4).

\medskip\noindent
Si $X$ est géométriquement connexe, alors $A \otimes_k
\overline{k} = H^0(X_{\overline{k}},
\mathcal{O}_{X_{\overline{k}}})$ est un anneau local de
dimension zéro d'après la partie (2) et donc son spectre a un
point, en particulier il est irréductible. Ainsi $A$ est
géométriquement irréductible. Cela prouve (5). Bien sûr, (5) implique (6).

\medskip\noindent
Si $X$ est géométriquement réduit et géométriquement connexe, alors
$A = k_1$ est un corps et l'extension $k_1/k$ est finie séparable et
géométriquement irréductible. Cependant, alors $k_1 \otimes_k
\overline{k}$ est un produit de $[k_1 : k]$ copies de $\overline{k}$ et
nous concluons que $k_1 = k$. Cela prouve (7). Bien sûr, (7) implique (8).
\end{proof}

\noindent
Voici une version simplifiée de la factorisation de Stein ; la
factorisation de Stein proprement dite sera abordée dans Plus sur les
morphismes, Section \ref{more-morphisms-section-stein-factorization}.

\begin{lemma}
\label{lemma-baby-stein}
Supposons que $X$ soit un schéma propre sur un corps $k$.
Posons $A = H^0(X, \mathcal{O}_X)$. Les fibres du morphisme
canonique $X \to \Spec(A)$ sont géométriquement connexes.
\end{lemma}

\begin{proof}
Posons $S = \Spec(A)$. Le morphisme canonique $X \to S$
est le morphisme correspondant à $\Gamma(S, \mathcal{O}_S)
= A = \Gamma(X, \mathcal{O}_X)$ via les schémas, Lemme
\ref{schemes-lemma-morphism-into-affine}. L'algèbre
$k$-$A$ est un produit fini $A = \prod A_i$ d'algèbres
locales artiniennes $k$ finies sur $k$, voir Lemme
\ref{lemma-proper-geometrically-reduced-global-sections}.
Désignons par $s_i \in S$ le point correspondant à l'idéal
maximal de $A_i$. Choisissons une clôture algébrique
$\overline{k}$ de $k$ et posons $\overline{A} = A \otimes_k
\overline{k}$. Choisissons un plongement $\kappa(s_i) \to
\overline{k}$ sur $k$ ; cela détermine un morphisme d'algèbres $\overline{k}$.
$$
\sigma_i : \overline{A} = A \otimes_k \overline{k} \to
\kappa(s_i) \otimes_k \overline{k} \to \overline{k}
$$
Considérez le changement de base
$$
\xymatrix{
\overline{X} \ar[r] \ar[d] & X \ar[d] \\
\overline{S} \ar[r] & S
}
$$
de $X$ à $\overline{S} = \Spec(\overline{A})$. Par la
cohomologie des schémas, Lemme
\ref{coherent-lemma-flat-base-change-cohomology} nous avons
$\Gamma(\overline{X}, \mathcal{O}_{\overline{X}}) =
\overline{A}$. Si $\overline{s}_i \in \Spec(\overline{A})$
désigne le point $\overline{k}$-rationnel correspondant à
$\sigma_i$, alors nous voyons que $\overline{s}_i$ se mappe à $s_i
\in S$ et $\overline{X}_{\overline{s}_i}$ est le changement de
base de $X_{s_i}$ par $\Spec(\sigma_i)$. Ainsi, nous voyons qu'il
suffit de prouver le lemme dans le cas où $k$ est algébriquement clos.

\medskip\noindent
Supposons que $k$ est algébriquement clos. Dans ce cas,
$\kappa(s_i)$ est algébriquement clos et nous devons montrer
que $X_{s_i}$ est connexe. La décomposition en produit $A =
\prod A_i$ correspond à une décomposition en union disjointe
$\Spec(A) = \coprod \Spec(A_i)$, voir Algèbre, Lemme
\ref{algebra-lemma-spec-product}. Notons $X_i$ l'image
inverse de $\Spec(A_i)$. Il découle du Lemme
\ref{lemma-proper-geometrically-reduced-global-sections},
partie (2), que $A_i = \Gamma(X_i, \mathcal{O}_{X_i})$. Il faut
observer que $X_{s_i} \to X_i$ est une immersion fermée
induisant un isomorphisme sur les espaces topologiques
sous-jacents (parce que $\Spec(A_i)$ est un singleton). Ainsi, si
$X_{s_i}$ n'est pas connexe, alors $X_i$ ne l'est pas non plus.
Donc soit $X_i$ est vide et $A_i = 0$ ou $X_i$ peut être écrit
comme $U \amalg V$ avec $U$ et $V$ ouverts et non vides, ce qui
impliquerait que $A_i$ possède un idempotent non trivial. Puisque
$A_i$ est local, c'est une contradiction et la démonstration est complète.
\end{proof}

\begin{lemma}
\label{lemma-geometrically-reduced-stein}
Supposons que $k$ soit un corps. Supposons que $X$ soit un schéma propre
géométriquement réduit sur $k$. Les éléments suivants sont équivalents
\begin{enumerate}
\item $H^0(X, \mathcal{O}_X) = k$, et
\item $X$ est géométriquement connexe.
\end{enumerate}
\end{lemma}

\begin{proof}
D'après le Lemme \ref{lemma-baby-stein}, nous avons (1)
$\Rightarrow$ (2). D'après le Lemme
\ref{lemma-proper-geometrically-reduced-global-sections}, nous avons (2) $\Rightarrow$ (1).
\end{proof}




\section{Schémas géométriquement normaux}
\label{section-geometrically-normal}

\noindent
Dans Propriétés, Définition
\ref{properties-definition-normal} nous avons défini la
notion de schéma normal. Cette notion est définie même pour
des schémas non noethériens. Ainsi, contrairement à notre
discussion des schémas géométriquement réguliers `` '' nous
considérons toutes les extensions de corps du corps de base.

\begin{definition}
\label{definition-geometrically-normal}
Supposons que $X$ soit un schéma sur le corps $k$.
\begin{enumerate}
\item Soit $x \in X$. On dit que $X$ est {\it
géométriquement normal en $x$ } si pour toute extension
de corps $k'/k$ et tout $x' \in X_{k'}$ au-dessus de $x$
l'anneau local $\mathcal{O}_{X_{k'}, x'}$ est normal.
\item On dit que $X$ est {\it géométriquement normal } sur
$k$ si $X$ est géométriquement normal en chaque $x \in X$.
\end{enumerate}
\end{definition}

\begin{lemma}
\label{lemma-geometrically-normal-at-point}
Supposons que $k$ soit un corps.
Supposons que $X$ soit un schéma
sur $k$. Soit $x \in X$. Les
énoncés suivants sont équivalents
\begin{enumerate}
\item $X$ est géométriquement normal en $x$, \item
pour chaque extension de corps finie purement
inséparable $k'$ de $k$ et $x' \in X_{k'}$ au-dessus
de $x$, l'anneau local $\mathcal{O}_{X_{k'}, x'}$
est normal, et \item l'anneau $\mathcal{O}_{X, x}$
est géométriquement normal sur $k$ (voir Algèbre,
Définition \ref{algebra-definition-geometrically-normal}).
\end{enumerate}
\end{lemma}

\begin{proof}
Il est clair que (1) implique (2). Supposons (2).
Supposons que $k'/k$ soit une extension de corps finie
purement inséparable (par exemple $k = k'$). Considérons
l'anneau $\mathcal{O}_{X, x} \otimes_k k'$. D'après
Algèbre, Lemme \ref{algebra-lemma-p-ring-map}, son spectre
est le même que le spectre de $\mathcal{O}_{X, x}$. Par
conséquent, c'est aussi un anneau local (Algèbre, Lemme
\ref{algebra-lemma-characterize-local-ring}). Il existe donc
un point unique $x' \in X_{k'}$ au-dessus de $x$ et
$\mathcal{O}_{X_{k'}, x'} \cong \mathcal{O}_{X, x} \otimes_k k'$.
Par hypothèse, c'est un anneau normal. Nous en déduisons donc (3)
d'après Algèbre, Lemme \ref{algebra-lemma-geometrically-normal}.

\medskip\noindent
Supposons (3). Supposons que $k'/k$ est une extension de
corps. Comme $\Spec(k') \to \Spec(k)$ est surjectif,
$X_{k'} \to X$ est aussi surjectif (morphismes, Lemme
\ref{morphisms-lemma-base-change-surjective}). Supposons
que $x' \in X_{k'}$ est un point quelconque situé au-dessus
de $x$. L'anneau local $\mathcal{O}_{X_{k'}, x'}$ est une
localisation de l'anneau $\mathcal{O}_{X, x} \otimes_k k'$.
Par conséquent, il est normal par hypothèse et (1) est prouvé.
\end{proof}

\begin{lemma}
\label{lemma-geometrically-normal}
Supposons que $k$ soit un corps.
Supposons que $X$ soit un schéma sur
$k$. Les éléments suivants sont équivalents
\begin{enumerate}
\item $X$ est géométriquement normal, \item $X_{k'}$
est un schéma normal pour chaque extension de corps
$k'/k$, \item $X_{k'}$ est un schéma normal pour
chaque extension de corps de type fini $k'/k$, \item
$X_{k'}$ est un schéma normal pour chaque extension
de corps finie purement inséparable $k'/k$, \item
pour chaque ouvert affine $U \subset X$ l'anneau
$\mathcal{O}_X(U)$ est géométriquement normal (voir
Algèbre, Définition
\ref{algebra-definition-geometrically-normal}), et \item $X_{k^{perf}}$ est un schéma normal.
\end{enumerate}
\end{lemma}

\begin{proof}
Supposons (1). Alors pour toute extension de corps
$k'/k$ et tout point $x' \in X_{k'}$, l'anneau
local de $X_{k'}$ en $x'$ est normal. Par
définition, cela signifie que $X_{k'}$ est normal. D'où (2).

\medskip\noindent
Il est clair que (2) implique (3) implique (4).

\medskip\noindent
Supposons (4) et soit $U \subset X$ un sous-schéma affine
ouvert. Alors $U_{k'}$ est un schéma normal pour toute
extension finie purement inséparable $k'/k$ (y compris $k =
k'$). Cela signifie que $k' \otimes_k \mathcal{O}(U)$ est un
anneau normal pour toutes les extensions finies purement
inséparables $k'/k$. Par conséquent, $\mathcal{O}(U)$ est une algèbre
géométriquement normale $k$ par définition. Ainsi, (4) implique (5).

\medskip\noindent
Supposons (5). Pour toute extension de corps $k'/k$, le
changement de base $X_{k'}$ est obtenu en collant les spectres
des anneaux $\mathcal{O}_X(U) \otimes_k k'$ où $U$ est un
ouvert affine dans $X$ (voir schémas, Section
\ref{schemes-section-fibre-products}). Donc $X_{k'}$ est normal. Ainsi (1) est vrai.

\medskip\noindent
L'équivalence de (5) et (6) découle de la
définition des algèbres géométriquement normales
et de l'équivalence (juste prouvée) de (3) et (4).
\end{proof}

\begin{lemma}
\label{lemma-geometrically-normal-upstairs}
Supposons que $k$ soit un corps. Supposons que $X$
soit un schéma sur $k$. Supposons que $k'/k$ soit
une extension de corps. Supposons que $x \in X$ soit
un point, et soit $x' \in X_{k'}$ un point situé
au-dessus de $x$. Les énoncés suivants sont équivalents
\begin{enumerate}
\item $X$ est géométriquement normal à $x$,
\item $X_{k'}$ est géométriquement normal à $x'$.
\end{enumerate}
En particulier, $X$ est géométriquement normal sur $k$ si et
seulement si $X_{k'}$ est géométriquement normal sur $k'$.
\end{lemma}

\begin{proof}
Il est clair que (1) implique (2). Supposons (2). Supposons que
$k''/k$ soit une extension de corps finie purement inséparable et
soit $x'' \in X_{k''}$ un point situé au-dessus de $x$ (en fait, il
est unique). Choisissez une extension de corps commune $k'''/k$ de
$k'/k$ et $k''/k$, voir corps, Lemme
\ref{fields-lemma-common-extension-field}. Choisissez un point $x''' \in X_{k'''}$
situé au-dessus à la fois de $x'$ et $x''$. Considérons l'application des anneaux locaux
$$
\mathcal{O}_{X_{k''}, x''} \longrightarrow \mathcal{O}_{X_{k'''}, x''''}.
$$
Ceci est un homomorphisme local plat d'anneaux et donc
fidèlement plat. D'après (2), nous voyons que l'anneau local à
droite est normal. Ainsi, d'après L'Algèbre, Lemme
\ref{algebra-lemma-descent-normal}, nous concluons que
$\mathcal{O}_{X_{k''}, x''}$ est normal. D'après le Lemme
\ref{lemma-geometrically-normal-at-point}, nous voyons que $X$ est géométriquement normal en $x$.
\end{proof}

\begin{lemma}
\label{lemma-fibre-product-normal}
Supposons que $k$ soit un corps. Supposons que $X$ soit un
schéma géométriquement normal sur $k$ et soit $Y$ un schéma
normal sur $k$. Alors $X \times_k Y$ est un schéma normal.
\end{lemma}

\begin{proof}
Ceci se réduit à l'algèbre, Lemme
\ref{algebra-lemma-geometrically-normal-tensor-normal}
par le Lemme
\ref{lemma-geometrically-normal}.
\end{proof}

\begin{lemma}
\label{lemma-base-change-normal-by-separable}
Supposons que $k$ soit un corps. Supposons que $X$ soit un schéma normal sur $k$.
Supposons que $K/k$ soit une extension de corps séparable. Alors $X_K$ est un schéma normal.
\end{lemma}

\begin{proof}
Suit du Lemme
\ref{lemma-fibre-product-normal} et du Lemme d'Algèbre
\ref{algebra-lemma-separable-field-extension-geometrically-normal}.
\end{proof}

\begin{lemma}
\label{lemma-geometrically-normal-stein}
Supposons que $k$ soit un corps. Supposons que $X$ soit un schéma propre
géométriquement normal sur $k$. Les éléments suivants sont équivalents
\begin{enumerate}
\item $H^0(X, \mathcal{O}_X) = k$, \item
$X$ est géométriquement connexe, \item
$X$ est géométriquement irréductible, et
\item $X$ est géométriquement intègre.
\end{enumerate}
\end{lemma}

\begin{proof}
D'après le lemme
\ref{lemma-geometrically-reduced-stein}, nous avons
l'équivalence de (1) et (2). Un schéma normal localement
noethérien (tel que $X_{\overline{k}}$) est une réunion
disjointe de ses composantes irréductibles (Propriétés, Lemme
\ref{properties-lemma-normal-Noetherian}). Ainsi, nous voyons que
(2) et (3) sont équivalents. Puisque $X_{\overline{k}}$ est
supposé réduit, nous voyons que (3) et (4) sont également équivalents.
\end{proof}









\section{Changement de corps et schémas localement noethériens}
\label{section-locally-Noetherian}

\noindent
Soit $X$ un schéma localement noethérien sur un corps $k$. Il n'est pas
toujours vrai que $X_{k'}$ soit aussi localement noethérien. Par
exemple, si $X = \Spec(\overline{\mathbf{Q}})$ et $k = \mathbf{Q}$,
alors $X_{\overline{\mathbf{Q}}}$ est le spectre de
$\overline{\mathbf{Q}} \otimes_{\mathbf{Q}} \overline{\mathbf{Q}}$ qui n'est
pas noethérien. (Indice : Il a trop d'idempotents). Mais si nous ne faisons
que des changements de base utilisant des extensions de corps de type fini,
alors la propriété noethérienne est préservée. (Ou si $X$ est localement de
type fini sur $k$, puisque cette propriété est préservée sous changement de base.)

\begin{lemma}
\label{lemma-locally-Noetherian-base-change}
Supposons que $k$ soit un corps. Supposons
que $X$ soit un schéma sur $k$. Supposons que
$k'/k$ soit une extension de corps de type
fini. Alors $X$ est localement noethérien si et
seulement si $X_{k'}$ est localement noethérien.
\end{lemma}

\begin{proof}
En utilisant les propriétés, le lemme
\ref{properties-lemma-locally-Noetherian} nous réduisons au cas où
$X$ est affine, disons $X = \Spec(A)$. Dans ce cas, nous devons
prouver que $A$ est noethérien si et seulement si $A_{k'}$ est
noethérien. Puisque $A \to A_{k'} = k' \otimes_k A$ est fidèlement
plat, nous voyons que si $A_{k'}$ est noethérien, alors $A$ l'est
aussi, d'après l'algèbre, le lemme
\ref{algebra-lemma-descent-Noetherian}. Réciproquement, si $A$ est noethérien, alors
$A_{k'}$ est noethérien d'après l'algèbre, le lemme \ref{algebra-lemma-Noetherian-field-extension}.
\end{proof}







\section{Schémas géométriquement réguliers}
\label{section-geometrically-regular}

\noindent
Un schéma géométriquement régulier sur un corps $k$ est un schéma
localement noethérien sur $k$ qui reste régulier après des changements
appropriés de corps de base. Un schéma de type fini sur $k$ est
géométriquement régulier si et seulement s'il est lisse sur $k$ (voir le Lemme
\ref{lemma-geometrically-regular-smooth}). La notion de régularité géométrique
est particulièrement intéressante dans les situations où la lissité ne peut pas
être utilisée, comme pour les fibres formelles (insérer référence future ici).

\medskip\noindent
Dans la définition suivante, nous nous limitons aux schémas
localement noethériens, puisque la propriété d'être un anneau
local régulier n'est définie que pour les anneaux locaux
noethériens. D'après le Lemme
\ref{lemma-locally-Noetherian-base-change} ci-dessus, si nous nous limitons
aux extensions de corps de type fini, alors cette propriété est préservée par
changement de corps de base. Ce commentaire sera utilisé sans référence
supplémentaire dans cette section. En particulier, la définition suivante a un sens.

\begin{definition}
\label{definition-geometrically-regular}
Supposons que $k$ soit un corps. Supposons que $X$ soit un schéma localement noethérien sur $k$.
\begin{enumerate}
\item Laisser $x \in X$. On dit que $X$ est {\it
géométriquement régulier en $x$ } au-dessus de $k$ si pour toute
extension de corps de type fini $k'/k$ et tout $x' \in X_{k'}$
au-dessus de $x$, l'anneau local $\mathcal{O}_{X_{k'}, x'}$ est
régulier. \item On dit que $X$ est {\it géométriquement régulier
sur $k$ } si $X$ est géométriquement régulier en tous ses points.
\end{enumerate}
\end{definition}

\noindent
Une définition similaire fonctionne pour définir
géométriquement Cohen-Macaulay, $(R_k)$, et $(S_k)$ schémas sur un
corps. Nous ajouterons une section pour ceux-ci séparément si nécessaire.

\begin{lemma}
\label{lemma-geometrically-regular-at-point}
Supposons que $k$ soit un corps.
Supposons que $X$ soit un schéma
localement noethérien sur $k$. Soit $x \in
X$. Les assertions suivantes sont équivalentes
\begin{enumerate}
\item $X$ est géométriquement régulier en $x$, \item
pour toute extension de corps finie purement
inséparable $k'$ de $k$ et $x' \in X_{k'}$ située
au-dessus de $x$ l'anneau local $\mathcal{O}_{X_{k'},
x'}$ est régulier, et \item l'anneau $\mathcal{O}_{X,
x}$ est géométriquement régulier sur $k$ (voir Algèbre,
Définition \ref{algebra-definition-geometrically-regular}).
\end{enumerate}
\end{lemma}

\begin{proof}
Il est clair que (1) implique (2). Supposons (2). Cela implique
en particulier que $\mathcal{O}_{X, x}$ est un anneau local
régulier. Supposons que $k'/k$ soit une extension de corps finie
purement inséparable. Considérons l'anneau $\mathcal{O}_{X, x}
\otimes_k k'$. D'après l'Algèbre, Lemme
\ref{algebra-lemma-p-ring-map}, son spectre est le même que le
spectre de $\mathcal{O}_{X, x}$. Par conséquent, c'est aussi un
anneau local (Algèbre, Lemme
\ref{algebra-lemma-characterize-local-ring}). Donc, il y a un point unique $x'
\in X_{k'}$ se trouvant au-dessus de $x$ et $\mathcal{O}_{X_{k'}, x'} \cong
\mathcal{O}_{X, x} \otimes_k k'$. Par hypothèse, il s'agit d'un anneau régulier.
Ainsi, nous déduisons (3) de la définition d'un anneau géométriquement régulier.

\medskip\noindent
Supposons (3). Supposons que $k'/k$ est une extension de
corps. Comme $\Spec(k') \to \Spec(k)$ est surjectif,
$X_{k'} \to X$ est aussi surjectif (morphismes, Lemme
\ref{morphisms-lemma-base-change-surjective}). Supposons que
$x' \in X_{k'}$ est un point quelconque situé au-dessus de
$x$. L'anneau local $\mathcal{O}_{X_{k'}, x'}$ est une
localisation de l'anneau $\mathcal{O}_{X, x} \otimes_k k'$.
Par conséquent, il est régulier par hypothèse et (1) est prouvé.
\end{proof}

\begin{lemma}
\label{lemma-geometrically-regular}
Supposons que $k$ soit un corps. Supposons que
$X$ soit un schéma localement noethérien sur
$k$. Les affirmations suivantes sont équivalentes
\begin{enumerate}
\item $X$ est géométriquement régulier, \item $X_{k'}$
est un schéma régulier pour toute extension de corps
de type fini $k'/k$, \item $X_{k'}$ est un schéma
régulier pour toute extension de corps finie purement
inséparable $k'/k$, \item pour tout ouvert affine $U
\subset X$ l'anneau $\mathcal{O}_X(U)$ est
géométriquement régulier (voir Algèbre, Définition
\ref{algebra-definition-geometrically-regular}), et \item il
existe un recouvrement ouvert affine $X = \bigcup U_i$ tel que
chaque $\mathcal{O}_X(U_i)$ est géométriquement régulier sur $k$.
\end{enumerate}
\end{lemma}

\begin{proof}
Supposons (1). Alors pour toute extension de corps de
type fini $k'/k$ et pour tout point $x' \in X_{k'}$,
l'anneau local de $X_{k'}$ en $x'$ est régulier. D'après
les Propriétés, Lemme
\ref{properties-lemma-characterize-regular}, cela signifie que $X_{k'}$ est régulier. D'où (2).

\medskip\noindent
Il est clair que (2) implique (3).

\medskip\noindent
Supposons (3) et soit $U \subset X$ un sous-schéma affine
ouvert. Alors $U_{k'}$ est un schéma régulier pour toute
extension finie purement inséparable $k'/k$ (y compris $k =
k'$). Cela signifie que $k' \otimes_k \mathcal{O}(U)$ est un
anneau régulier pour toutes les extensions finies purement
inséparables $k'/k$. Par conséquent, $\mathcal{O}(U)$ est une
algèbre $k$ géométriquement régulière et nous voyons que (4) est vrai.

\medskip\noindent
Il est clair que (4) implique (5). Supposons que $X = \bigcup
U_i$ soit un recouvrement affine ouvert comme dans (5). Pour toute
extension de corps $k'/k$, le changement de base $X_{k'}$ est
obtenu en recollant les spectres des anneaux $\mathcal{O}_X(U_i)
\otimes_k k'$ (voir schémas, Section
\ref{schemes-section-fibre-products}). Par conséquent, $X_{k'}$ est régulier. Donc (1) est vérifié.
\end{proof}

\begin{lemma}
\label{lemma-geometrically-regular-upstairs}
Supposons que $k$ est un corps. Supposons que $X$ est
un schéma sur $k$. Supposons que $k'/k$ est une
extension de corps de type fini. Supposons que $x \in
X$ est un point, et soit $x' \in X_{k'}$ un point situé
au-dessus de $x$. Les énoncés suivants sont équivalents
\begin{enumerate}
\item $X$ est géométriquement régulier en $x$,
\item $X_{k'}$ est géométriquement régulier en $x'$.
\end{enumerate}
En particulier, $X$ est géométriquement régulier sur $k$ si et
seulement si $X_{k'}$ est géométriquement régulier sur $k'$.
\end{lemma}

\begin{proof}
Il est clair que (1) implique (2). Supposons (2). Supposons que
$k''/k$ soit une extension de corps finie purement inséparable
et soit $x'' \in X_{k''}$ un point au-dessus de $x$ (en fait il
est unique). Nous pouvons trouver une extension de corps
commune, de type fini, $k'''/k$ de $k'/k$ et $k''/k$, voir
corps, Lemme \ref{fields-lemma-common-extension-field} et sa
démonstration, et un point $x''' \in X_{k'''}$ se trouvant au-dessus
à la fois de $x'$ et $x''$. Considérons le morphisme d’anneaux locaux
$$
\mathcal{O}_{X_{k''}, x''} \longrightarrow \mathcal{O}_{X_{k'''}, x''''}.
$$
Ceci est un homomorphisme plat local d'anneaux locaux
noethériens et donc fidèlement plat. D'après (2), nous
voyons que l'anneau local à droite est régulier. Ainsi,
selon le Lemme d'algèbre
\ref{algebra-lemma-flat-under-regular}, nous concluons que
$\mathcal{O}_{X_{k''}, x''}$ est régulier. D'après le Lemme
\ref{lemma-geometrically-regular-at-point}, nous voyons que $X$ est géométriquement régulier en $x$.
\end{proof}

\noindent
Le lemme suivant est une variante géométrique de
l'Algèbre, Lemme \ref{algebra-lemma-geometrically-regular-descent}.

\begin{lemma}
\label{lemma-flat-under-geometrically-regular}
Supposons que $k$ soit un corps. Supposons que $f : X \to Y$
soit un morphisme de schémas localement noethériens sur $k$.
Supposons que $x \in X$ soit un point et posons $y = f(x)$. Si
$X$ est géométriquement régulier en $x$ et que $f$ est plat en
$x$ alors $Y$ est géométriquement régulier en $y$. En
particulier, si $X$ est géométriquement régulier sur $k$ et que $f$
est plat et surjectif, alors $Y$ est géométriquement régulier sur $k$.
\end{lemma}

\begin{proof}
Supposons que $k'$ soit une extension finie purement
inséparable de $k$. Supposons que $f' : X_{k'} \to
Y_{k'}$ soit le changement de base de $f$. Supposons que
$x' \in X_{k'}$ soit le point unique situé au-dessus de
$x$. Si nous montrons que $Y_{k'}$ est régulier en $y' =
f'(x')$, alors $Y$ est géométriquement régulier sur $k$
en $y'$, voir le Lemme
\ref{lemma-geometrically-regular}. Par les morphismes, Lemme
\ref{morphisms-lemma-base-change-module-flat} le morphisme $X_{k'}
\to Y_{k'}$ est plat en $x'$. Par conséquent, l'homomorphisme d'anneaux
$$
\mathcal{O}_{Y_{k'}, y'}
\longrightarrow
\mathcal{O}_{X_{k'}, x'}
$$
est un homomorphisme local plat d'anneaux locaux
noethériens avec le côté droit régulier par hypothèse. Par
conséquent, le côté gauche est un anneau local régulier
selon Algèbre, Lemme \ref{algebra-lemma-flat-under-regular}.
\end{proof}

\begin{lemma}
\label{lemma-geometrically-regular-smooth}
Supposons que $k$ soit un corps. Supposons que $X$
soit un schéma localement de type fini sur $k$. Soit $x
\in X$. Alors $X$ est géométriquement régulier en $x$
si et seulement si $X \to \Spec(k)$ est lisse en $x$
(morphismes, Définition \ref{morphisms-definition-smooth}).
\end{lemma}

\begin{proof}
La question est locale autour de $x$, donc
nous pouvons supposer que $X = \Spec(A)$
pour un certain $k$-algèbre de type fini.
Soit $x$ correspondant au premier $\mathfrak p$.

\medskip\noindent
Si $A$ est lisse sur $k$ à $\mathfrak p$, alors nous pouvons
localiser $A$ et supposer que $A$ est lisse sur $k$. Dans ce cas,
$k' \otimes_k A$ est lisse sur $k'$ pour tout corps d'extension
$k'/k$, et chacun de ces anneaux noethériens est régulier selon
l'Algèbre, Lemme \ref{algebra-lemma-characterize-smooth-over-field}.

\medskip\noindent
Supposons que $X$ soit géométriquement régulier en $x$.
Considérons le corps résiduel $K := \kappa(x) =
\kappa(\mathfrak p)$ de $x$. C'est une extension de type
fini de $k$. D'après l'Algèbre, Lemme
\ref{algebra-lemma-make-separable}, il existe une
extension finie purement inséparable $k'/k$ telle que le
compositum $k'K$ soit une extension de corps séparable de
$k'$. Supposons que $\mathfrak p' \subset A' = k'
\otimes_k A$ soit un idéal premier au-dessus de $\mathfrak
p$. C'est le seul premier au-dessus de $\mathfrak p$, voir
Algèbre, Lemme \ref{algebra-lemma-p-ring-map}. Ainsi, le
corps résiduel $K' := \kappa(\mathfrak p')$ est le
compositum $k'K$. Par hypothèse, l'anneau local
$(A')_{\mathfrak p'}$ est régulier. Ainsi, d'après l'Algèbre,
Lemme \ref{algebra-lemma-separable-smooth}, nous voyons que $k'
\to A'$ est lisse en $\mathfrak p'$. Cela implique à son tour
que $k \to A$ est lisse en $\mathfrak p$ selon l'Algèbre, Lemme
\ref{algebra-lemma-smooth-field-change-local}. Le lemme est prouvé.
\end{proof}


\begin{example}
\label{example-geometrically-reduced-not-normal}
Soit $k =\mathbf{F}_p(t)$. Il est assez facile de donner un
exemple d'une variété régulière $V$ sur $k$ qui n'est pas
géométriquement réduite. Par exemple, nous pouvons prendre
$\Spec(k[x]/(x^p - t))$. En fait, il existe un exemple d'une
variété régulière $V$ qui est géométriquement réduite, mais pas même
géométriquement normale. À savoir, prenons pour $p > 2$ le schéma $V
= \Spec(k[x, y]/(y^2 - x^p + t))$. C'est une variété puisque le
polynôme $y^2 - x^p + t \in k[x, y]$ est irréductible. Le morphisme
$V \to \Spec(k)$ est lisse en tous les points sauf au point $v_0 \in
V$ correspondant à l'idéal maximal $(y, x^p - t)$ (car $2y$ est
inversible). En particulier, nous voyons que $V$ est (géométriquement)
régulière en tous les points, sauf éventuellement $v_0$. L'anneau local
$$
\mathcal{O}_{V, v_0} = \left(k[x, y]/(y^2 - x^p + t)\right)_{(y, x^p - t)}
$$
est un domaine de dimension $1$. Son idéal maximal est engendré
par l'élément $1$, à savoir $y$. Par conséquent, c'est un anneau
à valuation discrète et régulier. Soit $k' = k[t^{1/p}]$. On note
$t' = t^{1/p} \in k'$, $V' = V_{k'}$, $v'_0 \in V'$ le point
unique au-dessus de $v_0$. Au-dessus de $k'$, nous pouvons écrire
$x^p - t = (x - t')^p$, mais le polynôme $y^2 - (x - t')^p$ est
toujours irréductible et $V'$ est toujours une variété. Mais l'élément
$$
\frac{y}{x - t'} \in (\text{fraction field of }\mathcal{O}_{V', v'_0})
$$
est intégral sur $\mathcal{O}_{V', v'_0}$ (il suffit
de calculer son carré) et n'y est pas contenu, donc
$V'$ n'est pas normal sur $v'_0$. Cela conclut l'exemple.
\end{example}







\section{Changement de corps et propriété de Cohen--Macaulay}
\label{section-CM}

\noindent
Le lemme suivant dit qu'il n'a pas de sens de définir
des schémas de Cohen-Macaulay géométriques, puisque
ceux-ci seraient les mêmes que les schémas de Cohen-Macaulay.

\begin{lemma}
\label{lemma-CM-base-change}
Supposons que $X$ soit un schéma localement noethérien sur le
corps $k$. Supposons que $k'/k$ soit une extension de corps de
type fini. Supposons que $x \in X$ soit un point, et soit $x'
\in X_{k'}$ un point situé au-dessus de $x$. Alors nous avons
$$
\mathcal{O}_{X, x}\text{ is Cohen-Macaulay}
\Leftrightarrow
\mathcal{O}_{X_{k'}, x'}\text{ is Cohen-Macaulay}
$$
Si $X$ est localement de type fini sur $k$, il en
va de même pour toute extension de corps $k'/k$.
\end{lemma}

\begin{proof}
Le premier cas du lemme découle de l'Algèbre,
Lemme \ref{algebra-lemma-CM-geometrically-CM}. Le
deuxième cas du lemme est équivalent à l'Algèbre,
Lemme \ref{algebra-lemma-extend-field-CM-locus}.
\end{proof}







\section{Changement de corps et propriété de Jacobson}
\label{section-overfield}

\noindent
Un schéma localement de type fini sur un corps a beaucoup
de points fermés, à savoir qu'il est jacobien. De plus, les
corps résiduels sont des extensions finies du corps de base.

\begin{lemma}
\label{lemma-locally-finite-type-Jacobson}
Supposons que $X$ soit un schéma qui est
localement de type fini sur $k$. Alors
\begin{enumerate}
\item pour tout point fermé $x \in X$
l'extension $\kappa(x)/k$ est algébrique, et
\item $X$ est un schéma de Jacobson (Propriétés,
Définition \ref{properties-definition-jacobson} ).
\end{enumerate}
\end{lemma}

\begin{proof}
Un schéma est de Jacobson si et seulement s'il possède un
recouvrement par des ouverts affines de schémas de Jacobson,
voir Propriétés, Lemme
\ref{properties-lemma-locally-jacobson}. La propriété sur les corps
résiduels en des points fermés est également locale sur $X$. Par
conséquent, nous pouvons supposer que $X$ est affine. Dans ce cas,
le résultat est une conséquence du Nullstellensatz de Hilbert, voir
Algèbre, Théorème \ref{algebra-theorem-nullstellensatz}. Il découle
également d'une combinaison de morphismes, Lemmata
\ref{morphisms-lemma-jacobson-finite-type-points},
\ref{morphisms-lemma-Jacobson-universally-Jacobson} et \ref{morphisms-lemma-ubiquity-Jacobson-schemes}.
\end{proof}

\noindent
Il s'avère que si $X$ n'est pas localement de type fini, alors nous pouvons obtenir
le même résultat après avoir fait une extension de base de corps suffisamment grande.

\begin{lemma}
\label{lemma-make-Jacobson}
Supposons que $X$ soit un schéma sur un corps
$k$. Pour toute extension de corps $K/k$ dont la
cardinalité est suffisamment grande, nous avons
\begin{enumerate}
\item pour tout point fermé $x \in X_K$
l'extension $\kappa(x)/K$ est algébrique, et
\item $X_K$ est un schéma de Jacobson (Propriétés,
Définition \ref{properties-definition-jacobson} ).
\end{enumerate}
\end{lemma}

\begin{proof}
Choisissez un recouvrement ouvert affine $X =
\bigcup U_i$. Par l'Algèbre, Lemme
\ref{algebra-lemma-base-change-Jacobson} et les
Propriétés, Lemme
\ref{properties-lemma-affine-jacobson}, il existe des
cardinaux $\kappa_i$ tels que $U_{i, K}$ possède les
propriétés souhaitées sur $K$ si $\#(K) \geq \kappa_i$.
Posez $\kappa = \max\{\kappa_i\}$. Alors, si la cardinalité
de $K$ est supérieure à $\kappa$, nous voyons que chaque
$U_{i, K}$ satisfait les conclusions du lemme. Par
conséquent, $X_K$ est Jacobson par les Propriétés, Lemme
\ref{properties-lemma-locally-jacobson}. L'énoncé sur les corps
résiduels aux points fermés de $X_K$ découle des énoncés
correspondants pour les corps résiduels des points fermés de $U_{i, K}$.
\end{proof}





\section{Changement de corps et faisceaux inversibles amples}
\label{section-change-fields-ample}

\noindent
Le résultat suivant est typique des résultats de cette section.

\begin{lemma}
\label{lemma-ample-after-field-extension}
Supposons que $k$ est un corps. Supposons que $X$ est
un schéma sur $k$. S'il existe un faisceau inversible
ample sur $X_K$ pour une certaine extension de corps
$K/k$, alors $X$ possède un faisceau inversible ample.
\end{lemma}

\begin{proof}
Supposons que $K/k$ soit une extension de corps telle que
$X_K$ possède un faisceau inversible ample $\mathcal{L}$. Le
morphisme $X_K \to X$ est surjectif. Par conséquent, $X$ est
quasi-compact en tant qu'image d'un schéma quasi-compact
(Propriétés, Définition \ref{properties-definition-ample}).
Puisque $X_K$ est quasi-séparé (d'après Propriétés, Lemme
\ref{properties-lemma-affine-s-opens-cover-quasi-separated}), nous
voyons que $X$ est quasi-séparé : si $U, V \subset X$ sont ouverts
affines, alors $(U \cap V)_K = U_K \cap V_K$ est quasi-compact et
$(U \cap V)_K \to U \cap V$ est surjectif. Ainsi, le Lemme
\ref{schemes-lemma-characterize-quasi-separated} sur les schémas s'applique.

\medskip\noindent
Écrivez $K = \colim A_i$ comme le colimite des sous-algèbres de
$K$ qui sont de type fini sur $k$. Notons $X_i = X
\times_{\Spec(k)} \Spec(A_i)$. Puisque $X_K = \lim X_i$, nous
trouvons un $i$ et un faisceau inversible $\mathcal{L}_i$ sur
$X_i$ dont le tiré en arrière à $X_K$ est $\mathcal{L}$
(Limites, Lemme \ref{limits-lemma-descend-invertible-modules} ;
ici et ci-dessous, nous utilisons que $X$ est quasi-compact et
quasi-séparé comme montré précédemment). D'après Limites, Lemme
\ref{limits-lemma-limit-ample}, nous pouvons supposer que
$\mathcal{L}_i$ est ample après avoir éventuellement augmenté
$i$. Fixons un tel $i$ et soit $\mathfrak m \subset A_i$ un idéal
maximal. Par le Nullstellensatz de Hilbert (Algèbre, Théorème
\ref{algebra-theorem-nullstellensatz}), le corps résiduel $k' =
A_i/\mathfrak m$ est une extension finie de $k$. Ainsi $X_{k'}
\subset X_i$ est un sous-schéma fermé et possède donc un faisceau
inversible ample (Propriétés, Lemme
\ref{properties-lemma-ample-on-closed}). Puisque $X_{k'} \to X$ est localement
libre de rang fini, nous en concluons que $X$ possède un faisceau inversible
ample d’après Diviseurs, Proposition \ref{divisors-proposition-push-down-ample}.
\end{proof}

\begin{lemma}
\label{lemma-quasi-affine-after-field-extension}
Supposons que $k$ est un corps. Supposons que $X$ est un schéma sur $k$. Si $X_K$ est
quasi-affine pour une certaine extension de corps $K/k$, alors $X$ est quasi-affine.
\end{lemma}

\begin{proof}
Supposons que $K/k$ soit une extension de corps telle que $X_K$ soit
quasi-affine. Le morphisme $X_K \to X$ est surjectif. Ainsi $X$ est
quasi-compact comme l'image d'un schéma quasi-compact (Propriétés,
Définition \ref{properties-definition-quasi-affine}). Puisque $X_K$
est quasi-séparé (en tant que sous-schéma ouvert d'un schéma affine),
nous voyons que $X$ est quasi-séparé : si $U, V \subset X$ sont
ouverts affines, alors $(U \cap V)_K = U_K \cap V_K$ est quasi-compact
et $(U \cap V)_K \to U \cap V$ est surjectif. Ainsi, le lemme sur les
schémas, \ref{schemes-lemma-characterize-quasi-separated}, s'applique.

\medskip\noindent
Écrire $K = \colim A_i$ comme le colimite des
sous-algèbres de $K$ qui sont de type fini sur $k$. Noter
$X_i = X \times_{\Spec(k)} \Spec(A_i)$. Puisque $X_K = \lim
X_i$, nous trouvons un $i$ tel que $X_i$ soit quasi-affine
(Limites, Lemme \ref{limits-lemma-limit-quasi-affine} ; ici
nous utilisons que $X$ est quasi-compact et quasi-séparé
comme montré précédemment). Par le Nullstellensatz de
Hilbert (Algèbre, Théorème
\ref{algebra-theorem-nullstellensatz}), le corps résiduel $k' =
A_i/\mathfrak m$ est une extension finie de $k$. Ainsi $X_{k'}
\subset X_i$ est un sous-schéma fermé donc est quasi-affine
(Propriétés, Lemme
\ref{properties-lemma-quasi-affine-locally-closed}). Puisque $X_{k'} \to X$ est localement libre
de type fini, nous concluons par Diviseurs, Lemme \ref{divisors-lemma-push-down-quasi-affine}.
\end{proof}

\begin{lemma}
\label{lemma-quasi-projective-after-field-extension}
Supposons que $k$ soit un corps. Supposons que $X$ soit un schéma
sur $k$. Si $X_K$ est quasi-projectif sur $K$ pour une certaine
extension de corps $K/k$, alors $X$ est quasi-projectif sur $k$.
\end{lemma}

\begin{proof}
Par définition, un morphisme de schémas $g : Y \to T$ est
quasi-projectif s'il est localement de type fini, quasi-compact,
et s'il existe un $g$-faisceau inversible ample sur $Y$.
Supposons que $K/k$ soit une extension de corps telle que $X_K$
soit quasi-projectif sur $K$. Supposons que $\Spec(A) \subset X$
soit un ouvert affine. Alors $U_K$ est un sous-schéma ouvert
affine de $X_K$, donc $A_K$ est une $K$-algèbre de type fini.
Ensuite, $A$ est une $k$-algèbre de type fini d'après le Lemme
d'algèbre \ref{algebra-lemma-finite-type-descends}. Par conséquent,
$X \to \Spec(k)$ est localement de type fini. Puisque $X_K \to
\Spec(K)$ est quasi-compact, nous voyons que $X_K$ est
quasi-compact, donc $X$ est quasi-compact, donc $X \to \Spec(k)$ est
de type fini. D'après le Lemme sur les morphismes
\ref{morphisms-lemma-finite-type-over-affine-ample-very-ample}, nous
voyons que $X_K$ possède un faisceau inversible ample. Ensuite, $X$ possède
un faisceau inversible ample d'après le Lemme
\ref{lemma-ample-after-field-extension}. Par conséquent, $X \to \Spec(k)$ est
quasi-projectif par morphismes, Lemme \ref{morphisms-lemma-finite-type-over-affine-ample-very-ample}.
\end{proof}

\noindent
Le lemme suivant est un cas particulier de la Descente,
Lemme \ref{descent-lemma-descending-property-proper}.

\begin{lemma}
\label{lemma-proper-after-field-extension}
Supposons que $k$ soit un corps. Supposons que $X$ soit un
schéma sur $k$. Si $X_K$ est propre sur $K$ pour une
certaine extension de corps $K/k$, alors $X$ est propre sur $k$.
\end{lemma}

\begin{proof}
Supposons que $K/k$ soit une extension de corps telle que $X_K$ est
propre sur $K$. Rappelons que cela implique que $X_K$ est séparé et
quasi-compact (morphismes, Définition
\ref{morphisms-definition-proper}). Le morphisme $X_K \to X$ est
surjectif. Ainsi, $X$ est quasi-compact comme l'image d'un schéma
quasi-compact (Propriétés, Définition \ref{properties-definition-ample}).
Puisque $X_K$ est séparé, nous voyons que $X$ est quasi-séparé : si $U, V
\subset X$ sont ouverts affines, alors $(U \cap V)_K = U_K \cap V_K$ est
quasi-compact et $(U \cap V)_K \to U \cap V$ est surjectif. Ainsi, le
lemme des schémas \ref{schemes-lemma-characterize-quasi-separated} s'applique.

\medskip\noindent
Écrivez $K = \colim A_i$ comme le colimite des
sous-algèbres de $K$ qui sont de type fini sur $k$.
Dénotez $X_i = X \times_{\Spec(k)} \Spec(A_i)$. D'après
Limites, Lemme \ref{limits-lemma-eventually-proper}, il
existe un $i$ tel que $X_i \to \Spec(A_i)$ soit propre.
Ici, nous utilisons que $X$ est quasi-compact et
quasi-séparé comme montré précédemment. Choisissez un
idéal maximal $\mathfrak m \subset A_i$. D'après le
Nullstellensatz de Hilbert (Algèbre, Théorème
\ref{algebra-theorem-nullstellensatz}), le corps résiduel
$k' = A_i/\mathfrak m$ est une extension finie de $k$. Le
changement de base $X_{k'} \to \Spec(k')$ est propre
(morphismes, Lemme
\ref{morphisms-lemma-base-change-proper}). Puisque $k'/k$ est
fini, à la fois $X_{k'} \to X$ et la composition $X_{k'} \to
\Spec(k)$ sont également propres (morphismes, Lemmata
\ref{morphisms-lemma-finite-proper},
\ref{morphisms-lemma-base-change-proper} et
\ref{morphisms-lemma-composition-proper}). Le premier implique que $X$ est
séparé sur $k$ puisque $X_{k'}$ est séparé (morphismes, Lemme
\ref{morphisms-lemma-image-universally-closed-separated}). Le second implique que $X \to
\Spec(k)$ est propre d'après morphismes, Lemme \ref{morphisms-lemma-image-proper-is-proper}.
\end{proof}

\begin{lemma}
\label{lemma-projective-after-field-extension}
Supposons que $k$ soit un corps. Supposons que $X$ soit un
schéma sur $k$. Si $X_K$ est projectif sur $K$ pour une
certaine extension de corps $K/k$, alors $X$ est projectif sur $k$.
\end{lemma}

\begin{proof}
Un schéma sur $k$ est projectif sur $k$ si et
seulement s'il est quasi-projectif et propre sur $k$.
Voir morphismes, Lemme
\ref{morphisms-lemma-projective-is-quasi-projective-proper}.
Ainsi, le lemme découle des Lemmes
\ref{lemma-quasi-projective-after-field-extension} et \ref{lemma-proper-after-field-extension}.
\end{proof}




\section{Espaces tangents}
\label{section-tangent-spaces}

\noindent
Dans cette section, nous définissons l'espace tangent d'un morphisme de schémas
en un point de la source en utilisant des points à valeurs dans les nombres duaux.

\begin{definition}
\label{definition-dual-numbers}
Pour tout anneau $R$, les {\it nombres duaux } sur $R$
forment la $R$-algèbre notée $R[\epsilon]$. Comme
$R$-module, elle est libre de base $1$, $\epsilon$, et la
structure de $R$-algèbre est donnée par $\epsilon^2 = 0$.
\end{definition}

\noindent
Supposons que $f : X \to S$ soit un morphisme de schémas.
Supposons que $x \in X$ soit un point avec image $s =
f(x)$ dans $S$. Considérons le diagramme commutatif solide
\begin{equation}
\label{equation-tangent-space}
\vcenter{
\xymatrix{
\Spec(\kappa(x)) \ar[r] \ar[dr] \ar@/^1pc/[rr] &
\Spec(\kappa(x)[\epsilon]) \ar@{.>}[r] \ar[d]&
X \ar[d] \\
&
\Spec(\kappa(s)) \ar[r] &
S
}
}
\end{equation}
avec la flèche courbée étant le morphisme
canonique de $\Spec(\kappa(x))$ vers $X$.

\begin{lemma}
\label{lemma-tangent-space}
L'ensemble des flèches pointillées faisant commuter (
\ref{equation-tangent-space} ) possède une structure canonique d'espace vectoriel $\kappa(x)$.
\end{lemma}

\begin{proof}
Poser $\kappa = \kappa(x)$. Remarquez que nous
avons un pushout dans la catégorie des schémas
$$
\Spec(\kappa[\epsilon]) \amalg_{\Spec(\kappa)} \Spec(\kappa[\epsilon])
= \Spec(\kappa[\epsilon_1, \epsilon_2])
$$
où $\kappa[\epsilon_1, \epsilon_2]$ est l'algèbre
$\kappa$ de base $1, \epsilon_1, \epsilon_2$ et
$\epsilon_1^2 = \epsilon_1\epsilon_2 =
\epsilon_2^2 = 0$. Cela découle immédiatement du
résultat correspondant pour les anneaux et de la
description des morphismes des spectres des anneaux
locaux vers les schémas dans les schémas, Lemme
\ref{schemes-lemma-morphism-from-spec-local-ring}.
Étant donné deux flèches $\theta_1, \theta_2 :
\Spec(\kappa[\epsilon]) \to X$, nous pouvons considérer le morphisme
$$
\theta_1 + \theta_2 :
\Spec(\kappa[\epsilon]) \to
\Spec(\kappa[\epsilon_1, \epsilon_2])
\xrightarrow{\theta_1, \theta_2} X
$$
où la première flèche est donnée par $\epsilon_i \mapsto \epsilon$.
D'autre part, donné $\lambda \in \kappa$, il existe une application de
$\Spec(\kappa[\epsilon])$ vers lui-même correspondant à l'endomorphisme
d'algèbre $\kappa$ de $\kappa[\epsilon]$ qui envoie $\epsilon$ vers
$\lambda \epsilon$. Précomposer $\theta : \Spec(\kappa[\epsilon]) \to X$ par
cette application de soi-même donne $\lambda \theta$. Le lecteur peut
vérifier les axiomes d'un espace vectoriel en vérifiant l'existence de
diagrammes commutatifs appropriés de schémas. Nous omettons les détails. (Une
preuve alternative serait d'exprimer tout en termes d'anneaux locaux, puis de
vérifier les axiomes de l'espace vectoriel au niveau des homomorphismes d'anneaux.)
\end{proof}

\begin{definition}
\label{definition-tangent-space}
Supposons que $f : X \to S$ soit un morphisme de schémas. Soit $x \in X$.
L'ensemble des flèches pointillées rendant ( \ref{equation-tangent-space} )
compatible avec sa structure canonique de $\kappa(x)$-espace vectoriel est
appelé le {\it espace tangent de $X$ sur $S$ en $x$ } et nous le notons $T_{X/S,
x}$. Un élément de cet espace est appelé un {\it vecteur tangent } de $X/S$ en $x$.
\end{definition}

\noindent
Puisque les vecteurs tangents à $x \in X$ vivent dans la fibre schématique $X_s$ de
$f : X \to S$ au-dessus de $s = f(x)$, nous obtenons une identification canonique
\begin{equation}
\label{equation-tangent-space-fibre}
T_{X/S, x} = T_{X_s/s, x}
\end{equation}
Cette agréable définition faisant intervenir le
foncteur des points admet la description algébrique
suivante, qui suggère de définir {\it l'espace cotangent de
$X$ sur $S$ en $x$ } comme le $\kappa(x)$-espace vectoriel
$$
T^*_{X/S, x} = \Omega_{X/S, x} \otimes_{\mathcal{O}_{X, x}} \kappa(x)
$$
simplement parce que c'est canoniquement
$\kappa(x)$ -dual à l'espace tangent de $X$ sur $S$ en $x$.

\begin{lemma}
\label{lemma-tangent-space-cotangent-space}
Supposons que $f : X \to S$ soit un morphisme de
schémas. Soit $x \in X$. Il existe un isomorphisme canonique
$$
T_{X/S, x} = \Hom_{\mathcal{O}_{X, x}}(\Omega_{X/S, x}, \kappa(x))
$$
d'espaces vectoriels sur $\kappa(x)$.
\end{lemma}

\begin{proof}
Définir $\kappa = \kappa(x)$. Étant donné
$\theta \in T_{X/S, x}$, nous obtenons une application
$$
\theta^*\Omega_{X/S} \to
\Omega_{\Spec(\kappa[\epsilon])/\Spec(\kappa(s))} \to
\Omega_{\Spec(\kappa[\epsilon])/\Spec(\kappa)}
$$
En prenant les sections, nous obtenons une application
$\mathcal{O}_{X, x}$-linéaire $\xi_\theta : \Omega_{X/S, x}
\to \kappa \text{d}\epsilon$, c'est-à-dire un élément du
membre de droite de la formule du lemme. Pour montrer que
$\theta \mapsto \xi_\theta$ est un isomorphisme, nous
pouvons remplacer $S$ par $s$ et $X$ par la fibre schématique
$X_s$. En effet, les deux membres de la formule ne dépendent
que de la fibre schématique ; cela est clair pour $T_{X/S,
x}$ et, pour le membre de droite, voir Morphismes, lemme
\ref{morphisms-lemma-base-change-differentials}. Nous pouvons
aussi remplacer $X$ par le spectre de $\mathcal{O}_{X, x}$,
car cela ne change ni $T_{X/S, x}$ (Schémas, lemme
\ref{schemes-lemma-morphism-from-spec-local-ring}) ni $\Omega_{X/S,
x}$ (Modules, lemme \ref{modules-lemma-stalk-module-differentials}).

\medskip\noindent
Supposons que $(A, \mathfrak m, \kappa)$ soit un anneau local sur un
corps $k$. Pour terminer la preuve, nous devons montrer que toute
application $A$-linéaire $\xi : \Omega_{A/k} \to \kappa$ provient
d'une application d'algèbre $k$ unique $\varphi : A \to
\kappa[\epsilon]$ concordant avec l'application canonique $c : A \to
\kappa$ modulo $\epsilon$. Écrivons $\varphi(a) = c(a) + D(a) \epsilon$
; le lecteur voit que $a \mapsto D(a)$ est une dérivation $k$. En
utilisant la propriété universelle de $\Omega_{A/k}$, nous voyons que
chaque $D$ correspond à un $\xi$ unique et vice versa. Cela termine la preuve.
\end{proof}

\begin{lemma}
\label{lemma-tangent-space-rational-point}
Supposons que $f : X \to S$ soit un morphisme de schémas. Supposons
que $x \in X$ soit un point et soit $s = f(x) \in S$. Supposons que
$\kappa(x) = \kappa(s)$. Alors il existe des isomorphismes canoniques
$$
\mathfrak m_x/(\mathfrak m_x^2 + \mathfrak m_s\mathcal{O}_{X, x})
=
\Omega_{X/S, x} \otimes_{\mathcal{O}_{X, x}} \kappa(x)
$$
et
$$
T_{X/S, x} =
\Hom_{\kappa(x)}(
\mathfrak m_x/(\mathfrak m_x^2 + \mathfrak m_s\mathcal{O}_{X, x}),
\kappa(x))
$$
Cela fonctionne de manière plus générale si
$\kappa(x)/\kappa(s)$ est une extension algébrique séparable.
\end{lemma}

\begin{proof}
Le deuxième isomorphisme découle du premier par le Lemme
\ref{lemma-tangent-space-cotangent-space}. Pour le premier, nous
pouvons remplacer $S$ par $s$ et $X$ par $X_s$, voir morphismes,
Lemme \ref{morphisms-lemma-base-change-differentials}. Nous
pouvons également remplacer $X$ par le spectre de $\mathcal{O}_{X,
x}$, voir Modules, Lemme
\ref{modules-lemma-stalk-module-differentials}. Ainsi, nous devons montrer le fait
algébrique suivant : soit $(A, \mathfrak m, \kappa)$ un anneau local sur un corps
$k$ tel que $\kappa/k$ soit algébriquement séparable. Alors l'application canonique
$$
\mathfrak m/\mathfrak m^2
\longrightarrow
\Omega_{A/k} \otimes \kappa
$$
est un isomorphisme. Observez que $\mathfrak
m/\mathfrak m^2 = H_1(\NL_{\kappa/A})$. D'après
l'algèbre, Lemme \ref{algebra-lemma-exact-sequence-NL},
il suffit de montrer que $\Omega_{\kappa/k} = 0$ et
$H_1(\NL_{\kappa/k}) = 0$. Comme $\kappa$ est l'union de
ses extensions séparables finies dans $k$, il suffit de
prouver cela lorsque $\kappa$ est une extension séparables
finie de $k$ (Algèbre, Lemme
\ref{algebra-lemma-colimits-NL}). Dans ce cas, l'homomorphisme
d'anneaux $k \to \kappa$ est \' étale et donc $\NL_{\kappa/k} = 0$ (à peu
près par définition, voir Algèbre, Section \ref{algebra-section-etale}).
\end{proof}

\begin{lemma}
\label{lemma-map-tangent-spaces}
Supposons que $f : X \to Y$ soit un morphisme de schémas sur un schéma de
base $S$. Supposons que $x \in X$ soit un point. Posons $y = f(x)$. Si
$\kappa(y) = \kappa(x)$, alors $f$ induit une application linéaire naturelle
$$
\text{d}f : T_{X/S, x} \longrightarrow T_{Y/S, y}
$$
qui est dual à l'application linéaire $\Omega_{Y/S, y} \otimes
\kappa(y) \to \Omega_{X/S, x} \otimes \kappa(x)$ via les
identifications du Lemme \ref{lemma-tangent-space-cotangent-space}.
\end{lemma}

\begin{proof}
Omis.
\end{proof}

\begin{lemma}
\label{lemma-tangent-space-product}
Supposons que $X$, $Y$ soient des schémas sur une base $S$. Soient $x \in
X$ et $y \in Y$ avec le même point image $s \in S$ tel que $\kappa(s) =
\kappa(x)$ et $\kappa(s) = \kappa(y)$. Il existe un isomorphisme canonique
$$
T_{X \times_S Y/S, (x, y)} = T_{X/S, x} \oplus T_{Y/S, y}
$$
La application de gauche à droite est induite par les
applications sur les espaces tangents provenant des
projections $X \times_S Y \to X$ et $X \times_S Y \to Y$. La
application de droite à gauche est induite par les applications $1
\times y : X_s \to X_s \times_s Y_s$ et $x \times 1 : Y_s \to X_s
\times_s Y_s$ via l'identification
(\ref{equation-tangent-space-fibre}) des espaces tangents avec les espaces tangents des fibres.
\end{lemma}

\begin{proof}
La décomposition en somme directe découle des morphismes,
Lemme \ref{morphisms-lemma-differential-product} via le Lemme
\ref{lemma-tangent-space-rational-point}. La compatibilité avec
les applications provient du Lemme \ref{lemma-map-tangent-spaces}.
\end{proof}

\begin{lemma}
\label{lemma-injective-tangent-spaces-unramified}
Supposons que $f : X \to Y$ soit un morphisme de schémas localement de type fini sur
un schéma de base $S$. Supposons que $x \in X$ soit un point. Posons $y = f(x)$ et
supposons que $\kappa(y) = \kappa(x)$. Alors les assertions suivantes sont équivalentes
\begin{enumerate}
\item $\text{d}f : T_{X/S, x} \longrightarrow T_{Y/S,
y}$ est injectif, et \item $f$ est non ramifié en $x$.
\end{enumerate}
\end{lemma}

\begin{proof}
Le morphisme $f$ est localement de type fini par
morphismes, Lemme
\ref{morphisms-lemma-permanence-finite-type}. L'application
$\text{d}f$ est injective si et seulement si $\Omega_{Y/S, y}
\otimes \kappa(y) \to \Omega_{X/S, x} \otimes \kappa(x)$ est
surjective (Lemme \ref{lemma-map-tangent-spaces}). La suite
exacte $f^*\Omega_{Y/S} \to \Omega_{X/S} \to \Omega_{X/Y} \to
0$ (morphismes, Lemme
\ref{morphisms-lemma-triangle-differentials}) montre alors que cela se
produit si et seulement si $\Omega_{X/Y, x} \otimes \kappa(x) = 0$. D'où le
résultat découle de morphismes, Lemme \ref{morphisms-lemma-unramified-at-point}.
\end{proof}







\section{Morphismes génériquement finis}
\label{section-generically-finite}

\noindent
Dans cette section, nous revenons sur la notion de
morphisme génériquement fini de schémas telle qu'étudiée dans
morphismes, Section \ref{morphisms-section-generically-finite}.

\begin{lemma}
\label{lemma-quasi-finite-in-codim-1}
Supposons que $f : X \to Y$ soit localement de type fini. Supposons que $y
\in Y$ soit un point tel que $\mathcal{O}_{Y, y}$ soit noethérien de dimension
$\leq 1$. Supposons en outre qu'une des conditions suivantes soit satisfaite
\begin{enumerate}
\item pour chaque point générique $\eta$ d'une
composante irréductible de $X$ l'extension de corps
$\kappa(\eta)/\kappa(f(\eta))$ est finie (ou
algébrique), \item pour chaque point générique
$\eta$ d'une composante irréductible de $X$ telle
que $f(\eta) \leadsto y$ l'extension de corps
$\kappa(\eta)/\kappa(f(\eta))$ est finie (ou
algébrique), \item $f$ est quasi-finie en chaque point
générique d'une composante irréductible de $X$, \item $Y$
est localement noethérienne et $f$ est quasi-finie sur un
ensemble dense de points de $X$, \item ajouter plus ici.
\end{enumerate}
Alors $f$ est quasi-fini en chaque point de $X$ au-dessus de $y$.
\end{lemma}

\begin{proof}
La condition (4) implique que $X$ est localement noethérien
(morphismes, Lemme \ref{morphisms-lemma-finite-type-noetherian}).
L'ensemble des points en lesquels morphisme est quasi-fini est
ouvert (morphismes, Lemme
\ref{morphisms-lemma-quasi-finite-points-open}). Un ouvert dense d'un
schéma localement noethérien contient tous les points génériques des
composantes irréductibles, donc (4) implique (3). La condition (3)
implique la condition (1) selon morphismes, Lemme
\ref{morphisms-lemma-residue-field-quasi-finite}. La condition (1) implique la condition
(2). Ainsi, il suffit de prouver le lemme dans le cas où la condition (2) est vérifiée.

\medskip\noindent
Supposons que (2) soit vrai. Rappelons que
$\Spec(\mathcal{O}_{Y, y})$ est l'ensemble des points de $Y$
se spécialisant en $y$, voir schémas, Lemme
\ref{schemes-lemma-specialize-points}. Combiné avec morphismes,
Lemme \ref{morphisms-lemma-base-change-quasi-finite}, cela montre
que nous pouvons remplacer $Y$ par $\Spec(\mathcal{O}_{Y, y})$.
Ainsi nous pouvons supposer $Y = \Spec(B)$ où $B$ est un anneau
local noethérien de dimension $\leq 1$ et $y$ est le point fermé.

\medskip\noindent
Supposons que $X = \bigcup X_i$ soit les composantes irréductibles
de $X$ vues comme des sous-schémas fermés réduits. Si nous pouvons
montrer que chaque fibre $X_{i, y}$ est un espace discret, alors
$X_y = \bigcup X_{i, y}$ est également discret et nous en concluons
que $X \to Y$ est quasi-fini en tous les points de $X_y$ selon les
morphismes, Lemme
\ref{morphisms-lemma-quasi-finite-at-point-characterize}. Ainsi nous pouvons supposer que $X$ est un schéma intégral.

\medskip\noindent
Si $X \to Y$ envoie le point générique $\eta$ de $X$ à
$y$, alors $X$ est le spectre d'une extension finie de
$\kappa(y)$ et le résultat est vrai. Supposons que $X$
envoie $\eta$ à un point correspondant à un premier
minimal $\mathfrak q$ de $B$ différent de $\mathfrak m_B$.
Nous obtenons une factorisation $X \to \Spec(B/\mathfrak
q) \to \Spec(B)$. Supposons que $x \in X$ soit un point
au-dessus de $y$. Par la formule de la dimension
(morphismes, Lemme \ref{morphisms-lemma-dimension-formula}) nous avons
$$
\dim(\mathcal{O}_{X, x}) \leq \dim(B/\mathfrak q) +
\text{trdeg}_{\kappa(\mathfrak q)}(R(X)) - \text{trdeg}_{\kappa(y)} \kappa(x)
$$
Nous savons que $\dim(B/\mathfrak q) = 1$, que le point
générique de $X$ n'est pas égal à $x$ et se spécialise en $x$ et que
$R(X)$ est algébrique sur $\kappa(\mathfrak q)$. Ainsi nous obtenons
$$
1 \leq 1 - \text{trdeg}_{\kappa(y)} \kappa(x)
$$
Ainsi, chaque point $x$ de $X_y$ est fermé dans $X_y$ par
morphismes, le Lemme
\ref{morphisms-lemma-algebraic-residue-field-extension-closed-point-fibre}
et donc $X \to Y$ est quasi-fini en chaque point $x$ de $X_y$ par
morphismes, le Lemme
\ref{morphisms-lemma-quasi-finite-at-point-characterize} (ce qui implique également que $X_y$ est un espace topologique discret).
\end{proof}

\begin{lemma}
\label{lemma-finite-in-codim-1}
Supposons que $f : X \to Y$ soit un morphisme propre. Supposons que $y \in Y$
soit un point tel que $\mathcal{O}_{Y, y}$ soit noethérien de dimension
$\leq 1$. Supposons en outre qu'une des conditions suivantes soit satisfaisante.
\begin{enumerate}
\item pour tout point générique $\eta$ d'une
composante irréductible de $X$ l'extension de corps
$\kappa(\eta)/\kappa(f(\eta))$ est finie (ou
algébrique), \item pour tout point générique $\eta$
d'une composante irréductible de $X$ tel que $f(\eta)
\leadsto y$ l'extension de corps
$\kappa(\eta)/\kappa(f(\eta))$ est finie (ou algébrique),
\item $f$ est quasi-finie en tout point générique de $X$,
\item $Y$ est localement noethérien et $f$ est quasi-finie
sur un ensemble dense de points de $X$, \item ajoutez plus ici.
\end{enumerate}
Alors il existe un voisinage ouvert $V \subset
Y$ de $y$ tel que $f^{-1}(V) \to V$ soit fini.
\end{lemma}

\begin{proof}
D'après le Lemme \ref{lemma-quasi-finite-in-codim-1}, le
morphisme $f$ est quasi-fini en tout point de la fibre
$X_y$. Par conséquent, $X_y$ est un espace topologique
discret (morphismes, Lemme
\ref{morphisms-lemma-quasi-finite-at-point-characterize}). Comme $f$
est propre, la fibre $X_y$ est quasi-compacte, c'est-à-dire finie.
Ainsi, nous pouvons appliquer la Cohomologie des schémas, Lemme
\ref{coherent-lemma-proper-finite-fibre-finite-in-neighbourhood} pour conclure.
\end{proof}

\begin{lemma}
\label{lemma-modification-normal-iso-over-codimension-1}
Supposons que $X$ soit un schéma noethérien. Supposons que $f : Y \to X$ soit un
morphisme propre birationnel de schémas avec $Y$ réduit. Supposons que $U \subset
X$ soit l'ouvert maximal sur lequel $f$ est un isomorphisme. Alors $U$ contient
\begin{enumerate}
\item chaque point de codimension $0$ dans $X$, \item
chaque $x \in X$ de codimension $1$ sur $X$ tel que
$\mathcal{O}_{X, x}$ est un anneau de valuation discrète,
\item chaque $x \in X$ tel que le fibré de $Y \to X$ au-dessus
de $x$ est fini et tel que $\mathcal{O}_{X, x}$ est normal, et
\item chaque $x \in X$ tel que $f$ est quasi-fini en un $y \in
Y$ situé au-dessus de $x$ et $\mathcal{O}_{X, x}$ est normal.
\end{enumerate}
\end{lemma}

\begin{proof}
La partie (1) découle des morphismes, Lemme
\ref{morphisms-lemma-birational-isomorphism-over-dense-open}. La
partie (2) découle de la partie (3) et du Lemme
\ref{lemma-finite-in-codim-1} (et du fait que les morphismes finis ont des fibres finies).

\medskip\noindent
La partie (3) découle de la partie (4) et des morphismes,
Lemme \ref{morphisms-lemma-finite-fibre}, mais nous allons
également donner une preuve directe. Supposons que $x \in
X$ soit comme dans (3). D'après Cohomologie des schémas,
Lemme
\ref{coherent-lemma-proper-finite-fibre-finite-in-neighbourhood},
nous pouvons supposer que $f$ est fini. Nous pouvons supposer
que $X$ est affine. Cela nous ramène au cas d'un morphisme
birationnel fini de schémas affines noethériens $Y \to X$ et $x
\in X$ tels que $\mathcal{O}_{X, x}$ est un domaine normal.
Comme $\mathcal{O}_{X, x}$ est un domaine et $X$ est
noethérien, nous pouvons remplacer $X$ par un ouvert affine de
$x$ qui est intègre. Ensuite, puisque $Y \to X$ est birationnel
et $Y$ est réduit, nous voyons que $Y$ est intègre. En écrivant
$X = \Spec(A)$ et $Y = \Spec(B)$, nous voyons que $A \subset B$
est une inclusion finie de domaines ayant le même corps de
fractions. Si $\mathfrak p \subset A$ est le premier correspondant
à $x$, alors le fait que $A_\mathfrak p$ soit normal implique que
$A_\mathfrak p \subset B_\mathfrak p$ est une égalité. Puisque $B$
est un module $A$ fini, nous voyons qu'il existe un $a \in A$, $a
\not \in \mathfrak p$ tel que $A_a \to B_a$ est un isomorphisme.

\medskip\noindent
Supposons que $x \in X$ et $y \in Y$ soient comme dans (4). Après
avoir remplacé $X$ par un voisinage affine ouvert, nous pouvons
supposer $X = \Spec(A)$ et $A \subset \mathcal{O}_{X, x}$, voir
Propriétés, Lemme
\ref{properties-lemma-ring-affine-open-injective-local-ring}.
Ensuite, $A$ est un domaine et donc $X$ est intègre. Puisque $f$
est birationnel et $Y$ est réduit, il s'ensuit que $Y$ est
également intègre. Considérons l'homomorphisme d'anneaux
$\mathcal{O}_{X, x} \to \mathcal{O}_{Y, y}$. Il s'agit d'un
homomorphisme d'anneaux qui est essentiellement de type fini,
l'extension résiduelle de corps est finie, et $\dim(\mathcal{O}_{Y,
y}/\mathfrak m_x\mathcal{O}_{Y, y}) = 0$ (pour voir cela, suivez la
trace à travers les définitions des applications quasi-finies dans
morphismes, Définition \ref{morphisms-definition-quasi-finite} et
Algèbre, Définition \ref{algebra-definition-quasi-finite}). D'après
Algèbre, Lemme
\ref{algebra-lemma-essentially-finite-type-fibre-dim-zero},
$\mathcal{O}_{Y, y}$ est la localisation d'une algèbre finie
$\mathcal{O}_{X, x}$-algèbre $B$. Bien sûr, nous pouvons remplacer $B$ par
l'image de $B$ dans $\mathcal{O}_{Y, y}$ et supposer que $B$ est un
domaine ayant le même corps des fractions que $\mathcal{O}_{Y, y}$. Alors
$\mathcal{O}_{X, x} \subset B$ a le même corps des fractions que $f$ est
birationnel. Comme $\mathcal{O}_{X, x}$ est normal, nous concluons que
$\mathcal{O}_{X, x} = B$ (car fini implique intégral), en particulier, nous
voyons que $\mathcal{O}_{X, x} = \mathcal{O}_{Y, y}$. Par les morphismes,
Lemme \ref{morphisms-lemma-morphism-defined-local-ring} après avoir réduit
$X$, nous pouvons supposer qu'il existe une section $X \to Y$ de $f$
envoyant $x$ vers $y$ et induisant l'isomorphisme donné sur les anneaux
locaux. Comme $X \to Y$ est fermé (par les schémas, Lemme
\ref{schemes-lemma-section-immersion}), il en résulte nécessairement que le point
générique de $X$ est envoyé vers le point générique de $Y$ et il s'ensuit que l'image
de $X \to Y$ est $Y$. Alors $Y = X$ et nous avons prouvé ce que nous voulions montrer.
\end{proof}






\section{Variantes de la normalisation de Noether}
\label{section-noether-normalization}

\noindent
La normalisation de Noether est l'énoncé selon
lequel si $k$ est un corps et $A$ est une algèbre
$k$ de type fini de dimension $d$, alors il existe
un homomorphisme d'algèbre injectif $k$ fini $k[x_1,
\ldots, x_d] \to A$. Voir Algèbre, Lemme
\ref{algebra-lemma-Noether-normalization}.
Géométriquement, cela signifie qu'il existe un morphisme
surjectif fini $\Spec(A) \to \mathbf{A}^d_k$ au-dessus de $\Spec(k)$.

\begin{lemma}
\label{lemma-noether-normalization}
Supposons que $f : X \to S$ soit un morphisme de schémas.
Supposons que $x \in X$ avec image $s \in S$. Soit $V \subset S$
un voisinage affine ouvert de $s$. Si $f$ est localement de type
fini et $\dim_x(X_s) = d$, alors il existe un ouvert affine $U
\subset X$ avec $x \in U$ et $f(U) \subset V$ et une factorisation
$$
U \xrightarrow{\pi} \mathbf{A}^d_V \to V
$$
de $f|_U : U \to V$ de telle sorte que $\pi$ soit quasi-fini.
\end{lemma}

\begin{proof}
Cela découle de l'algèbre, lemme
\ref{algebra-lemma-quasi-finite-over-polynomial-algebra}.
\end{proof}

\begin{lemma}
\label{lemma-noether-normalization-affine}
Supposons que $f : X \to S$ soit un morphisme de type fini de schémas
affines. Soit $s \in S$. Si $\dim(X_s) = d$, alors il existe une factorisation
$$
X \xrightarrow{\pi} \mathbf{A}^d_S \to S
$$
de $f$ de telle sorte que le morphisme $\pi_s : X_s \to
\mathbf{A}^d_{\kappa(s)}$ des fibres au-dessus de $s$ soit fini.
\end{lemma}

\begin{proof}
Écrivons $S = \Spec(A)$ et $X = \Spec(B)$, et soit $A \to B$
l'homomorphisme d'anneaux correspondant à $f$. Soit
$\mathfrak p \subset A$ l'idéal premier correspondant à $s$.
Nous pouvons choisir une surjection $A[x_1, \ldots, x_r] \to
B$. Par Algèbre, lemme
\ref{algebra-lemma-Noether-normalization}, il existe des éléments
$y_1, \ldots, y_d \in A$ de la $\mathbf{Z}$-sous-algèbre de $A$
engendrée par $x_1, \ldots, x_r$ tels que l'homomorphisme de $A$-algèbres
$A[t_1, \ldots, t_d] \to B$ envoyant $t_i$ sur $y_i$ induise un
homomorphisme fini de $\kappa(\mathfrak p)$-algèbres $\kappa(\mathfrak
p)[t_1, \ldots, t_d] \to B \otimes_A \kappa(\mathfrak p)$. Ceci démontre le lemme.
\end{proof}

\begin{lemma}
\label{lemma-geometric-structure-unramified}
Supposons que $f : X \to S$ soit un morphisme de schémas.
Supposons que $x \in X$. Soit $V = \Spec(A)$ un voisinage
affine ouvert de $f(x)$ dans $S$. Si $f$ n'est pas ramifié en
$x$, alors il existe un ouvert affine $U \subset X$ avec $x \in
U$ et $f(U) \subset V$ tel que nous ayons un diagramme commutatif
$$
\xymatrix{
X \ar[d] & U \ar[l] \ar[rd] \ar[r]^-j &
\Spec(A[t]_{g'}/(g)) \ar[d] \ar[r] &
\Spec(A[t]) = \mathbf{A}^1_V \ar[ld] \\
Y & & V \ar[ll]
}
$$
où $j$ est une immersion, $g \in A[t]$ est un polynôme unitaire, et $g'$ est
la dérivée de $g$ par rapport à $t$. Si $f$ est \' étale en $x$, alors nous
pouvons choisir le diagramme de manière à ce que $j$ soit une immersion ouverte.
\end{lemma}

\begin{proof}
Le cas non ramifié est une traduction de l'Algèbre,
Proposition
\ref{algebra-proposition-unramified-locally-standard}. Dans le cas
étale \', il s'agit d'une traduction de l'Algèbre, Proposition
\ref{algebra-proposition-etale-locally-standard} ou, de manière
équivalente, cela découle des morphismes, Lemme
\ref{morphisms-lemma-etale-locally-standard-etale} bien que les énoncés diffèrent légèrement.
\end{proof}

\begin{lemma}
\label{lemma-unramfied-over-affine}
Supposons que $f : X \to S$ soit un morphisme de type fini
de schémas affines. Soit $x \in X$ avec image $s \in S$. Soit
$$
r =
\dim_{\kappa(x)} \Omega_{X/S, x} \otimes_{\mathcal{O}_{X, x}} \kappa(x) =
\dim_{\kappa(x)} \Omega_{X_s/s, x} \otimes_{\mathcal{O}_{X_s, x}} \kappa(x) =
\dim_{\kappa(x)} T_{X/S, x}
$$
Alors il existe une factorisation
$$
X \xrightarrow{\pi} \mathbf{A}^r_S \to S
$$
de $f$ de telle sorte que $\pi$ soit non ramifié en $x$.
\end{lemma}

\begin{proof}
Par les morphismes, le Lemme
\ref{morphisms-lemma-finite-type-differentials} la
première dimension est finie. La première égalité découle
du fait que la restriction de $\Omega_{X/S}$ à la fibre
est le module des différentielles des morphismes, Lemme
\ref{morphisms-lemma-base-change-differentials}. La
dernière égalité découle du Lemme
\ref{lemma-tangent-space-cotangent-space}. Ainsi, nous voyons que l'énoncé a du sens.

\medskip\noindent
Pour prouver le lemme, écrivez $S = \Spec(A)$ et $X = \Spec(B)$ et
soit $A \to B$ l'homomorphisme d'anneaux correspondant à $f$.
Supposons que $\mathfrak q \subset B$ soit l'idéal premier
correspondant à $x$. Choisissez une surjection d'algèbres
$A$-$A[x_1, \ldots, x_t] \to B$. Puisque $\Omega_{B/A}$ est généré
par $\text{d}x_1, \ldots, \text{d}x_t$, nous voyons que leurs images
dans $\Omega_{X/S, x} \otimes_{\mathcal{O}_{X, x}} \kappa(x)$
génèrent celui-ci comme un espace vectoriel $\kappa(x)$. Après
renumérotation, nous pouvons supposer que $\text{d}x_1, \ldots,
\text{d}x_r$ se mappe sur une base de $\Omega_{X/S, x}
\otimes_{\mathcal{O}_{X, x}} \kappa(x)$. Nous affirmons que $P = A[x_1,
\ldots, x_r] \to B$ est non ramifié en $\mathfrak q$. Pour le voir, il
suffit de montrer que $\Omega_{B/P, \mathfrak q} = 0$ (Algèbre, Lemme
\ref{algebra-lemma-unramified}). Notez que $\Omega_{B/P}$ est le
quotient de $\Omega_{B/A}$ par le sous-module généré par $\text{d}x_1,
\ldots, \text{d}x_r$. Par conséquent, $\Omega_{B/P, \mathfrak q}
\otimes_{B_\mathfrak q} \kappa(\mathfrak q) = 0$ par notre choix de $x_1,
\ldots, x_r$. Par le lemme de Nakayama, plus précisément Algèbre, Lemme
\ref{algebra-lemma-NAK} partie (2) qui s'applique car $\Omega_{B/P}$ est fini
(voir référence ci-dessus), nous concluons que $\Omega_{B/P, \mathfrak q} = 0$.
\end{proof}

\begin{lemma}
\label{lemma-immersion-into-affine}
Supposons que $f : X \to S$ soit un morphisme de
schémas. Supposons que $x \in X$ avec image $s
\in S$. Soit $V \subset S$ un voisinage affine
ouvert de $s$. Si $f$ est localement de type fini et
$$
r =
\dim_{\kappa(x)} \Omega_{X/S, x} \otimes_{\mathcal{O}_{X, x}} \kappa(x) =
\dim_{\kappa(x)} \Omega_{X_s/s, x} \otimes_{\mathcal{O}_{X_s, x}} \kappa(x) =
\dim_{\kappa(x)} T_{X/S, x}
$$
alors il existe
\begin{enumerate}
\item a ouvert affine $U \subset X$ avec $x
\in U$ et $f(U) \subset V$ et une factorisation
$$
U \xrightarrow{j} \mathbf{A}^{r + 1}_V \to V
$$
de $f|_U$ tel que $j$ soit une immersion, ou
\item un ouvert affine $U \subset X$ avec $x
\in U$ et $f(U) \subset V$ et une factorisation
$$
U \xrightarrow{j} D \to V
$$
de $f|_U$ de telle sorte que $j$ soit une immersion
fermée et $D \to V$ soit lisse de dimension relative $r$.
\end{enumerate}
\end{lemma}

\begin{proof}
Choisissez un ouvert affine $U \subset X$ avec $x \in U$ et
$f(U) \subset V$. Appliquez le Lemme
\ref{lemma-unramfied-over-affine} à $U \to V$ pour obtenir $U \to
\mathbf{A}^r_V \to V$ comme dans l'énoncé de ce lemme. D'après le Lemme
\ref{lemma-geometric-structure-unramified}, nous obtenons une factorisation
$$
U \xrightarrow{j} D \xrightarrow{j'} \mathbf{A}^{r + 1}_V
\xrightarrow{p} \mathbf{A}^r_V \to V
$$
où $j$ et $j'$ sont des immersions, $p$ est la
projection, et $p \circ j'$ est \' étale standard.
Ainsi, nous voyons en particulier que (1) et (2) tiennent.
\end{proof}





\section{Dimension des fibres}
\label{section-dimension-fibres}

\noindent
Nous avons déjà vu que la dimension des fibres de type fini morphismes augmente
généralement. Dans cette section, nous discutons du phénomène selon lequel en
codimension $1$, cela ne se produit pas. Plus généralement, nous discutons de l'ampleur
avec laquelle la dimension d'une fibre peut augmenter. Voici une liste de résultats connexes :
\begin{enumerate}
\item Pour un morphisme de type fini $X \to S$, l'ensemble
de $x \in X$ avec $\dim_x(X_{f(x)}) \leq d$ est ouvert,
voir Algèbre, Lemme
\ref{algebra-lemma-dimension-fibres-bounded-open-upstairs} et
morphismes, Lemme
\ref{morphisms-lemma-openness-bounded-dimension-fibres}. \item
Nous avons la formule de dimension, voir Algèbre, Lemme
\ref{algebra-lemma-dimension-formula} et morphismes, Lemme
\ref{morphisms-lemma-dimension-formula}. \item Dimension constante
des fibres pour un schéma de type fini intègre dominant un anneau de
valuation, voir Algèbre, Lemme
\ref{algebra-lemma-finite-type-domain-over-valuation-ring-dim-fibres}. \item
Si $X \to S$ est de type fini et est quasi-fini en chaque point générique de
$X$, alors $X \to S$ est quasi-fini en codimension $1$, voir Algèbre, Lemme
\ref{algebra-lemma-finite-in-codim-1} et Lemme \ref{lemma-quasi-finite-in-codim-1}.
\end{enumerate}
Le dernier résultat mentionné ci-dessus se généralise comme suit.

\begin{lemma}
\label{lemma-dimension-fibre-in-codim-1}
Supposons que $f : X \to Y$ soit localement de type fini. Supposons que $x \in
X$ soit un point dont l'image est $y \in Y$ tel que $\mathcal{O}_{Y, y}$ soit
noethérien de dimension $\leq 1$. Supposons que $d \geq 0$ soit un entier tel
que pour tout point générique $\eta$ d'une composante irréductible de $X$
contenant $x$, nous ayons $\dim_\eta(X_{f(\eta)}) = d$. Alors $\dim_x(X_y) = d$.
\end{lemma}

\begin{proof}
Rappelons que $\Spec(\mathcal{O}_{Y, y})$ est l'ensemble
des points de $Y$ se spécialisant en $y$, voir schémas,
Lemme \ref{schemes-lemma-specialize-points}. Ainsi, nous
pouvons remplacer $Y$ par $\Spec(\mathcal{O}_{Y, y})$ et
supposer $Y = \Spec(B)$ où $B$ est un anneau local
noethérien de dimension $\leq 1$ et $y$ est le point fermé. Nous
pouvons également remplacer $X$ par un voisinage affine de $x$.

\medskip\noindent
Supposons que $X = \bigcup X_i$ soit les composantes irréductibles de $X$
vues comme des sous-schémas fermés réduits. Si nous pouvons montrer que
chaque fibre $X_{i, y}$ a dimension $d$, alors $X_y = \bigcup X_{i, y}$ a
également dimension $d$. Ainsi, nous pouvons supposer que $X$ est un schéma intègre.

\medskip\noindent
Si $X \to Y$ envoie le point générique $\eta$ de $X$ sur
$y$, alors $X$ est un schéma sur $\kappa(y)$ et le
résultat est vrai par hypothèse. Supposons que $X$ envoie
$\eta$ sur un point $\xi \in Y$ correspondant à un idéal
premier minimal $\mathfrak q$ de $B$ différent de $\mathfrak
m_B$. Nous obtenons une factorisation $X \to
\Spec(B/\mathfrak q) \to \Spec(B)$. Par la formule de la dimension
(morphismes, Lemme \ref{morphisms-lemma-dimension-formula}) nous avons
$$
\dim(\mathcal{O}_{X, x}) + \text{trdeg}_{\kappa(y)} \kappa(x) \leq
\dim(B/\mathfrak q) + \text{trdeg}_{\kappa(\mathfrak q)}(R(X))
$$
Nous avons $\dim(B/\mathfrak q) = 1$. Nous avons
$\text{trdeg}_{\kappa(\mathfrak q)}(R(X)) = d$ selon
notre hypothèse que $\dim_\eta(X_\xi) = d$, voir les
morphismes, Lemme
\ref{morphisms-lemma-dimension-fibre-at-a-point}. Puisque
$\mathcal{O}_{X, x} \to \mathcal{O}_{X_s, x}$ a un noyau
(comme $\eta \mapsto \xi \not = y$) et puisque
$\mathcal{O}_{X, x}$ est un domaine noethérien, nous voyons que
$\dim(\mathcal{O}_{X, x}) > \dim(\mathcal{O}_{X_y, x})$. Nous concluons que
$$
\dim_x(X_s) =
\dim(\mathcal{O}_{X_s, x}) + \text{trdeg}_{\kappa(y)} \kappa(x) \leq d
$$
(morphismes, Lemme
\ref{morphisms-lemma-dimension-fibre-at-a-point}). D'autre part,
nous avons $\dim_x(X_s) \geq \dim_\eta(X_{f(\eta)}) = d$ par
morphismes, Lemme \ref{morphisms-lemma-openness-bounded-dimension-fibres}.
\end{proof}

\begin{lemma}
\label{lemma-dominate-valuation-ring-dimension-fibres}
Supposons que $f : X \to \Spec(R)$ soit un morphisme d'un schéma
irréductible vers le spectre d'un anneau de valuation. Si $f$ est
localement de type fini et surjectif, alors la fibre spéciale est
équidimensionnelle de dimension égale à la dimension de la fibre générique.
\end{lemma}

\begin{proof}
Nous pouvons remplacer $X$ par sa réduction car cela ne change
pas la dimension de $X$ ni celle de la fibre spéciale. Alors
$X$ est intègre et le lemme découle de l'algèbre, Lemme
\ref{algebra-lemma-finite-type-domain-over-valuation-ring-dim-fibres}.
\end{proof}

\noindent
Le lemme suivant généralise le Lemme \ref{lemma-dimension-fibre-in-codim-1}.

\begin{lemma}
\label{lemma-dimension-fibre-in-higher-codimension}
Supposons que $f : X \to Y$ soit localement de type fini. Supposons
que $x \in X$ soit un point d'image $y \in Y$ tel que $\mathcal{O}_{Y,
y}$ soit noethérien. Supposons que $d \geq 0$ soit un entier tel que
pour tout point générique $\eta$ d'une composante irréductible de $X$
qui contient $x$, nous ayons $f(\eta) \not = y$ et
$\dim_\eta(X_{f(\eta)}) = d$. Alors $\dim_x(X_y) \leq d + \dim(\mathcal{O}_{Y, y}) - 1$.
\end{lemma}

\begin{proof}
Exactement comme dans la preuve du Lemme
\ref{lemma-dimension-fibre-in-codim-1}, nous réduisons au cas $X =
\Spec(A)$ avec $A$ un domaine et $Y = \Spec(B)$ où $B$ est un anneau
local noethérien dont l'idéal maximal correspond à $y$. Après avoir
remplacé $B$ par $B/\Ker(B \to A)$, nous pouvons supposer que $B$ est un
domaine et que $B \subset A$. Ensuite, nous utilisons la formule de la
dimension (morphismes, Lemme \ref{morphisms-lemma-dimension-formula}) pour obtenir
$$
\dim(\mathcal{O}_{X, x}) + \text{trdeg}_{\kappa(y)} \kappa(x) \leq
\dim(B) + \text{trdeg}_B(A)
$$
Nous avons $\text{trdeg}_B(A) = d$ selon notre
hypothèse que $\dim_\eta(X_\xi) = d$, voir
morphismes, Lemme
\ref{morphisms-lemma-dimension-fibre-at-a-point}. Puisque
$\mathcal{O}_{X, x} \to \mathcal{O}_{X_y, x}$ a un noyau
(comme $f(\eta) \not = y$) et puisque $\mathcal{O}_{X, x}$
est un domaine noethérien, nous voyons que
$\dim(\mathcal{O}_{X, x}) > \dim(\mathcal{O}_{X_y, x})$. Nous concluons que
$$
\dim_x(X_y) =
\dim(\mathcal{O}_{X_y, x}) + \text{trdeg}_{\kappa(y)} \kappa(x)
< \dim(B) + d
$$
(égalité par morphismes, Lemme
\ref{morphisms-lemma-dimension-fibre-at-a-point}) qui prouve ce que nous voulons.
\end{proof}




\section{Schémas algébriques}
\label{section-algebraic-schemes}

\noindent
La définition suivante est tirée
de \cite[I Definition 6.4.1]{EGA}.

\begin{definition}
\label{definition-algebraic-scheme}
Supposons que $k$ soit un corps. Un {\it $k$ -schéma
algébrique } est un schéma $X$ sur $k$ tel que le morphisme de
structure $X \to \Spec(k)$ soit de type fini. Un {\it $k$ -schéma
localement algébrique } est un schéma $X$ sur $k$ tel que le
morphisme de structure $X \to \Spec(k)$ soit localement de type fini.
\end{definition}

\noindent
Notez que tout $k$ -schéma (localement) algébrique est (localement)
noethérien, voir morphismes, Lemme
\ref{morphisms-lemma-finite-type-noetherian}. La catégorie des $k$ -schémas
algébriques possède tous les produits et produits fibrés (contrairement à la catégorie
des variétés sur $k$). De même pour la catégorie des $k$ -schémas localement algébriques.

\begin{lemma}
\label{lemma-algebraic-scheme-dim-0}
Supposons que $k$ est un corps. Supposons que $X$ est un $k$ -schéma
localement algébrique de dimension $0$. Alors $X$ est une réunion
disjointe de spectres d'algèbres locales artiniennes $k$ $A$ avec
$\dim_k(A) < \infty$. Si $X$ est un $k$ -schéma algébrique de dimension
$0$, alors en plus $X$ est affine et le morphisme $X \to \Spec(k)$ est fini.
\end{lemma}

\begin{proof}
Supposons que $X$ soit un $k$ -schéma localement algébrique
de dimension $0$. Supposons que $U = \Spec(A) \subset X$
soit un ouvert affine subschéma. Comme $\dim(X) = 0$ nous
voyons que $\dim(A) = 0$. Par normalisation de Noether,
voir Algèbre, Lemme
\ref{algebra-lemma-Noether-normalization}, nous voyons qu'il
existe une injection finie $k \to A$, c'est-à-dire $\dim_k(A) <
\infty$. Par conséquent, $A$ est artinien, voir Algèbre, Lemme
\ref{algebra-lemma-finite-dimensional-algebra}. Cela implique que
$A = A_1 \times \ldots \times A_r$ est un produit d'un nombre
fini d'anneaux locaux artiniens, voir Algèbre, Lemme
\ref{algebra-lemma-artinian-finite-length}. Bien sûr, $\dim_k(A_i) <
\infty$ pour chaque $i$ puisque la somme de ces dimensions égale $\dim_k(A)$.

\medskip\noindent
Les arguments ci-dessus montrent que $X$ possède un recouvrement ouvert
dont les membres sont des espaces topologiques discrets finis. Par
conséquent, $X$ est un espace topologique discret. Il s'ensuit que $X$ est
isomorphe à l'union disjointe de ses composantes connexes, chacune étant un
singleton. Puisqu'un schéma singleton est affine, nous concluons (d'après
les résultats du paragraphe ci-dessus) que chacun de ces singletons est le
spectre d'une $k$-algèbre locale artinienne $A$ avec $\dim_k(A) < \infty$.

\medskip\noindent
Enfin, si $X$ est un $k$ -schéma algébrique de
dimension $0$, alors $X$ est quasi-compact donc
est une réunion disjointe finie $X = \Spec(A_1)
\amalg \ldots \amalg \Spec(A_r)$ donc affine (voir
schémas, Lemme
\ref{schemes-lemma-disjoint-union-affines}) et nous avons vu la
finitude de $X \to \Spec(k)$ dans le premier paragraphe de la démonstration.
\end{proof}

\noindent
Le lemme suivant rassemble quelques énoncés sur la théorie
de la dimension pour les schémas localement algébriques.

\begin{lemma}
\label{lemma-dimension-locally-algebraic}
Supposons que $k$ soit un corps. Supposons que $X$ soit un $k$-schéma localement algébrique.
\begin{enumerate}
\item
\label{item-catenary}
L'espace topologique de $X$ est
caténaire (Topologie, Définition
\ref{topology-definition-catenary} ). \item
\label{item-dimension-at-closed-point}
Pour $x \in X$ fermé, nous avons
$\dim_x(X) = \dim(\mathcal{O}_{X, x})$. \item
\label{item-dimension-irreducible}
Pour $X$ irréductible nous avons
$\dim(X) = \dim(U)$ pour tout ouvert
non vide $U \subset X$ et $\dim(X) =
\dim_x(X)$ pour tout $x \in X$. \item
\label{item-irreducible-maximal-chains}
Pour $X$ irréductible, toute chaîne de sous-ensembles
fermés irréductibles peut être étendue à une chaîne
maximale et toutes les chaînes maximales de sous-ensembles
fermés irréductibles ont une longueur égale à $\dim(X)$. \item
\label{item-dimension-irreducibles-passing-through}
Pour $x \in X$ nous avons $\dim_x(X) = \max
\dim(Z) = \min \dim(\mathcal{O}_{X, x'})$ où
le maximum est pris sur les composantes
irréductibles $Z \subset X$ contenant $x$ et
le minimum est pris sur les spécialisations $x
\leadsto x'$ avec $x'$ fermé dans $X$. \item
\label{item-dimension-irreducible-trdeg}
Si $X$ est irréductible avec le
point générique $x$, alors $\dim(X)
= \text{trdeg}_k(\kappa(x))$. \item
\label{item-immediate-specialization}
Si $x \leadsto x'$ est une
spécialisation immédiate de points de $X$,
alors nous avons
$\text{trdeg}_k(\kappa(x)) = \text{trdeg}_k(\kappa(x')) + 1$. \item
\label{item-dimension-sup-trdeg}
La dimension de $X$ est le suprémum des
nombres $\text{trdeg}_k(\kappa(x))$ où
$x$ parcourt les points génériques des
composantes irréductibles de $X$. \item
\label{item-specialization}
Si $x \leadsto x'$ est une spécialisation non triviale de points de $X$, alors
\begin{enumerate}
\item $\dim_x(X) \leq \dim_{x'}(X)$, \item
$\dim(\mathcal{O}_{X, x}) < \dim(\mathcal{O}_{X, x'})$, \item
$\text{trdeg}_k(\kappa(x)) > \text{trdeg}_k(\kappa(x'))$, et
\item toute chaîne maximale de spécialisations non triviales $x =
x_0 \leadsto x_1 \leadsto \ldots \leadsto x_n = x'$ a une
longueur $n = \text{trdeg}_k(\kappa(x)) - \text{trdeg}_k(\kappa(x'))$.
\end{enumerate}
\item
\label{item-dimension-formula}
Pour $x \in X$ nous avons
$\dim_x(X) =
\text{trdeg}_k(\kappa(x)) + \dim(\mathcal{O}_{X, x})$. \item
\label{item-immediate-specialization-local-ring}
Si $x \leadsto x'$ est une spécialisation immédiate des
points de $X$ et que $X$ est irréductible ou équidimensionnel,
alors $\dim(\mathcal{O}_{X, x'}) = \dim(\mathcal{O}_{X, x}) + 1$.
\end{enumerate}
\end{lemma}

\begin{proof}
Au lieu de s'appuyer sur les résultats plus généraux
démontrés précédemment, nous réduirons les énoncés aux
énoncés correspondants pour les $k$ -algèbres de type fini et
citerons des résultats du chapitre sur l'algèbre commutative.

\medskip\noindent
Preuve de ( \ref{item-catenary} ). Ceci est local sur $X$ par
Topologie, Lemme \ref{topology-lemma-catenary}. Ainsi, nous
pouvons supposer $X = \Spec(A)$ où $A$ est une $k$ -algèbre de
type fini. Nous devons montrer que $A$ est catenary (Algèbre,
Lemme \ref{algebra-lemma-catenary}). Nous pouvons réduire à
$k[x_1, \ldots, x_n]$ en utilisant Algèbre, Lemme
\ref{algebra-lemma-quotient-catenary} et ensuite appliquer
Algèbre, Lemme
\ref{algebra-lemma-dimension-height-polynomial-ring}. Alternativement, cela tient
parce que $k$ est Cohen-Macaulay (trivialement) et que les anneaux Cohen-Macaulay
sont universellement caténaires (Algèbre, Lemme \ref{algebra-lemma-CM-ring-catenary}).

\medskip\noindent
Preuve de ( \ref{item-dimension-at-closed-point} ). Choisissez un
voisinage affine $U = \Spec(A)$ de $x$. Alors $\dim_x(X) =
\dim_x(U)$. Par conséquent, nous nous réduisons au cas affine, qui est
Algèbre, Lemme \ref{algebra-lemma-dimension-closed-point-finite-type-field}.

\medskip\noindent
Preuve de ( \ref{item-dimension-irreducible} ). Il suffit de
montrer que deux ouverts affines non vides $U, U' \subset X$
quelconques ont la même dimension (chaîne finie de
sous-ensembles irréductibles rencontre un ouvert affine).
Choisissez un point fermé $x$ de $X$ avec $x \in U \cap U'$. Cela
est possible car $X$ est irréductible, donc $U \cap U'$ est non
vide, donc il existe un tel point fermé car $X$ est Jacobson
d'après le Lemme \ref{lemma-locally-finite-type-Jacobson}. Ensuite
$\dim(U) = \dim(\mathcal{O}_{X, x}) = \dim(U')$ selon Algèbre,
Lemme \ref{algebra-lemma-dimension-spell-it-out} (strictement
parlant, il faut remplacer $X$ par sa réduction avant d'appliquer le lemme).

\medskip\noindent
Preuve de ( \ref{item-irreducible-maximal-chains} ). Étant
donné une chaîne de sous-ensembles fermés irréductibles, nous
pouvons trouver un ouvert affine $U \subset X$ qui rencontre
le plus petit. Ainsi, l'énoncé découle de l'Algèbre, Lemme
\ref{algebra-lemma-dimension-spell-it-out} et $\dim(U) =
\dim(X)$ que nous avons vus dans ( \ref{item-dimension-irreducible} ).

\medskip\noindent
Preuve de (
\ref{item-dimension-irreducibles-passing-through} ).
Choisissez un voisinage affine $U = \Spec(A)$ de $x$. Alors
$\dim_x(X) = \dim_x(U)$. La règle $Z \mapsto Z \cap U$ est une
bijection entre les composantes irréductibles de $X$ passant par
$x$ et les composantes irréductibles de $U$ passant par $x$. De
plus, $\dim(Z \cap U) = \dim(Z)$ pour un tel $Z$ par (
\ref{item-dimension-irreducible} ). Ainsi, l'énoncé découle de Algèbre,
Lemme \ref{algebra-lemma-dimension-at-a-point-finite-type-over-field}.

\medskip\noindent
Preuve de ( \ref{item-dimension-irreducible-trdeg} ).
Par ( \ref{item-dimension-irreducible} ) cela se
réduit au cas où $X = \Spec(A)$ est affine. Dans ce
cas, cela découle de l'Algèbre, Lemme
\ref{algebra-lemma-dimension-prime-polynomial-ring} appliqué à $A_{red}$.

\medskip\noindent
Preuve de ( \ref{item-immediate-specialization} ).
Soit $Z = \overline{\{x\}} \supset Z' =
\overline{\{x'\}}$. Il en découle de (
\ref{item-irreducible-maximal-chains} ) que $Z \supset Z'$
est le début d'une chaîne maximale de sous-schémas fermés
irréductibles dans $Z$ et par conséquent $\dim(Z) = \dim(Z') +
1$. Nous concluons par ( \ref{item-dimension-irreducible-trdeg} ).

\medskip\noindent
Preuve de ( \ref{item-dimension-sup-trdeg} ). Un argument
topologique simple montre que $\dim(X) = \sup \dim(Z)$ où le
supremum est pris sur les composantes irréductibles de $X$ (indice
: utiliser Topologie, Lemme \ref{topology-lemma-irreducible} ).
Ainsi, cela découle de ( \ref{item-dimension-irreducible-trdeg} ).

\medskip\noindent
Preuve de ( \ref{item-specialization} ). La partie (a) découle du fait
que tout ouvert $U \subset X$ contenant $x'$ contient également $x$. La
partie (b) découle du fait que $\mathcal{O}_{X, x}$ est une localisation
de $\mathcal{O}_{X, x'}$, donc toute chaîne de premiers dans
$\mathcal{O}_{X, x}$ correspond à une chaîne de premiers dans $\mathcal{O}_{X,
x'}$, qui peut être prolongée en ajoutant $\mathfrak m_{x'}$ à la fin. Les
parties (c) et (d) découlent formellement de ( \ref{item-immediate-specialization} ).

\medskip\noindent
Preuve de ( \ref{item-dimension-formula} ). Choisissez un
voisinage affine $U = \Spec(A)$ de $x$. Alors $\dim_x(X) =
\dim_x(U)$. Par conséquent, nous nous réduisons au cas affine, qui est
Algèbre, Lemme \ref{algebra-lemma-dimension-at-a-point-finite-type-field}.

\medskip\noindent
Preuve de (
\ref{item-immediate-specialization-local-ring} ). Si $X$ est
équidimensionnel (Topologie, Définition
\ref{topology-definition-equidimensional}) alors $\dim(X)$ est égal à la
dimension de chaque composante irréductible de $X$, d'où $\dim_x(X) =
\dim(X) = \dim_{x'}(X)$ par (
\ref{item-dimension-irreducibles-passing-through} ). Ainsi cela découle de ( \ref{item-immediate-specialization} ).
\end{proof}

\begin{lemma}
\label{lemma-dimension-fibres-locally-algebraic}
Supposons que $k$ soit un corps. Supposons que $f : X \to
Y$ soit un morphisme de $k$-schémas localement algébriques.
\begin{enumerate}
\item Pour $y \in Y$, la fibre $X_y$ est un schéma
algébrique local sur $\kappa(y)$ donc tous les résultats
du Lemme \ref{lemma-dimension-locally-algebraic}
s'appliquent. \item Supposons que $X$ soit irréductible.
Posons $Z = \overline{f(X)}$ et $d = \dim(X) - \dim(Z)$. Alors
\begin{enumerate}
\item $\dim_x(X_{f(x)}) \geq d$ pour tous $x \in X$, \item
l'ensemble de $x \in X$ avec $\dim_x(X_{f(x)}) = d$ est ouvert
dense, \item si $\dim(\mathcal{O}_{Z, f(x)}) \geq 1$, alors
$\dim_x(X_{f(x)}) \leq d + \dim(\mathcal{O}_{Z, f(x)}) - 1$, \item
si $\dim(\mathcal{O}_{Z, f(x)}) = 1$, alors $\dim_x(X_{f(x)}) = d$,
\end{enumerate}
\item Pour $x \in X$ avec $y = f(x)$ nous
avons $\dim_x(X_y) \geq \dim_x(X) - \dim_y(Y)$.
\end{enumerate}
\end{lemma}

\begin{proof}
Le morphisme $f$ est localement de type fini par
morphismes, Lemme
\ref{morphisms-lemma-permanence-finite-type}. Par conséquent,
le changement de base $X_y \to \Spec(\kappa(y))$ est localement
de type fini. Cela prouve (1). Dans le reste de la
démonstration, nous utiliserons librement les résultats du Lemme
\ref{lemma-dimension-locally-algebraic} pour $X$, $Y$, et les fibres de $f$.

\medskip\noindent
Preuve de (2). Supposons que $\eta \in X$ soit le point générique et
posons $\xi = f(\eta)$. Alors $Z = \overline{\{\xi\}}$. Par conséquent
$$
d = \dim(X) - \dim(Z) =
\text{trdeg}_k \kappa(\eta) - \text{trdeg}_k \kappa(\xi) =
\text{trdeg}_{\kappa(\xi)} \kappa(\eta) = \dim_\eta(X_\xi)
$$
Ainsi, les parties (2)(a) et (2)(b) découlent des
morphismes, Lemme
\ref{morphisms-lemma-openness-bounded-dimension-fibres}. Les
parties (2)(c) et (2)(d) découlent des Lemmata
\ref{lemma-dimension-fibre-in-higher-codimension} et \ref{lemma-dimension-fibre-in-codim-1}.

\medskip\noindent
Preuve de (3). Supposons que $x \in X$. Soit $X' \subset X$ une
composante irréductible de $X$ passant par $x$ de dimension
$\dim_x(X)$. Alors (2) implique que $\dim_x(X_y) \geq \dim(X') -
\dim(Z')$ où $Z' \subset Y$ est l'adhérence de l'image de $X'$. Cela prouve (3).
\end{proof}

\begin{lemma}
\label{lemma-dimension-product-locally-algebraic}
\begin{slogan}
La dimension du produit est la somme des dimensions.
\end{slogan}
Supposons que $k$ soit un corps. Supposons que $X$, $Y$ soient des schémas localement algébriques $k$.
\begin{enumerate}
\item Pour $z \in X \times Y$ reposant sur $(x, y)$
nous avons $\dim_z(X \times Y) = \dim_x(X) + \dim_y(Y)$.
\item Nous avons $\dim(X \times Y) = \dim(X) + \dim(Y)$.
\end{enumerate}
\end{lemma}

\begin{proof}
Preuve de (1). Considérons la factorisation
$$
X \times Y \longrightarrow Y \longrightarrow \Spec(k)
$$
du morphisme de structure. Le premier morphisme $p : X \times Y \to
Y$ est plat en tant que changement de base du morphisme plat $X \to
\Spec(k)$ par des morphismes, Lemme
\ref{morphisms-lemma-base-change-flat}. De plus, nous avons
$\dim_z(p^{-1}(y)) = \dim_x(X)$ par des morphismes, Lemme
\ref{morphisms-lemma-dimension-fibre-after-base-change}. Par conséquent $\dim_z(X
\times Y) = \dim_x(X) + \dim_y(Y)$ par des morphismes, Lemme
\ref{morphisms-lemma-dimension-fibre-at-a-point-additive}. La partie (2) est une conséquence directe de (1).
\end{proof}






\section{Anneaux locaux complets}
\label{section-complete-local-rings}

\noindent
Quelques résultats sur les anneaux locaux complets des schémas sur corps.

\begin{lemma}
\label{lemma-complete-local-ring-structure-as-algebra}
Supposons que $k$ soit un corps. Supposons que $X$ soit un
schéma localement noethérien sur $k$. Supposons que $x \in X$ soit
un point avec le corps résiduel $\kappa$. Il existe un isomorphisme
\begin{equation}
\label{equation-complete-local-ring}
\kappa[[x_1, \ldots, x_n]]/I \longrightarrow \mathcal{O}_{X, x}^\wedge
\end{equation}
induisant l'identité sur les corps résiduels. En général, nous ne
pouvons pas choisir ( \ref{equation-complete-local-ring} ) pour
être un isomorphisme d'algèbres $k$. Cependant, si l'extension
$\kappa/k$ est séparable, alors nous pouvons choisir (
\ref{equation-complete-local-ring} ) pour être un isomorphisme d'algèbres $k$.
\end{lemma}

\begin{proof}
L'existence de l'isomorphisme est une conséquence immédiate du
théorème de structure de Cohen \footnote {. Notez que si $\kappa$ a
pour caractéristique $p$, alors le théorème dit simplement que nous
obtenons une surjection $\Lambda[[x_1, \ldots, x_n]] \to
\mathcal{O}_{X, x}^\wedge$ où $\Lambda$ est un anneau de Cohen pour
$\kappa$. Mais bien sûr, dans ce cas, l'application se factorise par
$\Lambda/p\Lambda[[x_1, \ldots, x_n]]$ et $\Lambda/p\Lambda = \kappa$.
} (Algèbre, Théorème \ref{algebra-theorem-cohen-structure-theorem}).

\medskip\noindent
Supposons que $p$ soit un nombre premier impair, notons $k =
\mathbf{F}_p(t)$, et $A = k[x, y]/(y^2 + x^p - t)$. Alors la
complétion $A^\wedge$ de $A$ dans l'idéal maximal $\mathfrak m
= (y)$ est isomorphe à $k(t^{1/p})[[z]]$ en tant qu'anneau
mais pas en tant que $k$-algèbre. La raison est que $A^\wedge$
ne contient pas d'élément dont la $p$-ième puissance est $t$
(comme le lecteur peut le voir en calculant modulo $y^2$). Cela
montre également que tout isomorphisme
(\ref{equation-complete-local-ring}) ne peut pas être un isomorphisme de $k$-algèbre.

\medskip\noindent
Si $\kappa/k$ est séparable, alors il existe un
homomorphisme d'algèbre $k$ $\kappa \to \mathcal{O}_{X,
x}^\wedge$ induisant l'identité sur les corps résiduels
selon More on Algèbre, Lemme
\ref{more-algebra-lemma-lift-residue-field}. Supposons que $f_1,
\ldots, f_n \in \mathfrak m_x$ soit générateurs. Considérons la application
$$
\kappa[[x_1, \ldots, x_n]] \longrightarrow \mathcal{O}_{X, x}^\wedge,\quad
x_i \longmapsto f_i
$$
Puisque les deux côtés sont complets $(x_1,
\ldots, x_n)$-adquement (le côté droit d'après
l'Algèbre, Lemmata
\ref{algebra-lemma-hathat-finitely-generated}), cette
application est surjective d'après Algèbre, Lemme
\ref{algebra-lemma-completion-generalities} car elle est
surjective modulo $(x_1, \ldots, x_n)$ par construction.
\end{proof}

\begin{lemma}
\label{lemma-base-change-complete-local-ring}
Supposons que $K/k$ soit une extension de corps. Supposons que
$X$ soit un $k$-schéma localement algébrique. Posons $Y = X_K$.
Supposons que $y \in Y$ soit un point avec une image $x \in X$.
Supposons que $\dim(\mathcal{O}_{X, x}) = \dim(\mathcal{O}_{Y,
y})$ et que $\kappa(x)/k$ soit séparable. Choisissez un isomorphisme
$$
\kappa(x)[[x_1, \ldots, x_n]]/(g_1, \ldots, g_m) \longrightarrow
\mathcal{O}_{X, x}^\wedge
$$
des $k$ -algèbres comme dans (
\ref{equation-complete-local-ring} ). Ensuite, nous avons un isomorphisme
$$
\kappa(y)[[x_1, \ldots, x_n]]/(g_1, \ldots, g_m) \longrightarrow
\mathcal{O}_{Y, y}^\wedge
$$
des $K$ -algèbres comme dans (
\ref{equation-complete-local-ring} ). Ici, nous utilisons $\kappa(x)
\to \kappa(y)$ pour voir $g_j$ comme une série entière sur $\kappa(y)$.
\end{lemma}

\begin{proof}
L'application d'anneau local $\mathcal{O}_{X, x}
\to \mathcal{O}_{Y, y}$ induit une application
d'anneau local $\mathcal{O}_{X, x}^\wedge \to
\mathcal{O}_{Y, y}^\wedge$. L'application induite
$$
\kappa(x) \to \kappa(x)[[x_1, \ldots, x_n]]/(g_1, \ldots, g_m)
\to \mathcal{O}_{X, x}^\wedge \to \mathcal{O}_{Y, y}^\wedge
$$
composé avec la projection vers $\kappa(y)$ est
l'homomorphisme canonique $\kappa(x) \to \kappa(y)$. Selon
le Lemme \ref{lemma-change-fields-flat}, le corps résiduel
$\kappa(y)$ est une localisation de $\kappa(x) \otimes_k K$
au noyau $\mathfrak p_0$ de $\kappa(x) \otimes_k K \to
\kappa(y)$. D'autre part, selon le Lemme
\ref{lemma-change-fields-algebraic-unramified}, l'anneau local $(\kappa(x)
\otimes_k K)_{\mathfrak p_0}$ est égal à $\kappa(y)$. Ainsi, l'application
$$
\kappa(x) \otimes_k K \to \mathcal{O}_{Y, y}^\wedge
$$
se factorise canoniquement à travers
$\kappa(y)$. Nous obtenons un diagramme commutatif
$$
\xymatrix{
\kappa(y) \ar[rr] & & \mathcal{O}_{Y, y}^\wedge \\
\kappa(x) \ar[r] \ar[u] &
\kappa(x)[[x_1, \ldots, x_n]]/(g_1, \ldots, g_m) \ar[r] &
\mathcal{O}_{X, x}^\wedge \ar[u]
}
$$
Supposons que $f_i \in \mathfrak m_x^\wedge \subset
\mathcal{O}_{X, x}^\wedge$ soit l'image de $x_i$.
Observez que $\mathfrak m_x^\wedge = (f_1, \ldots, f_n)$ en
tant que application est surjective. Considérons l'application
$$
\kappa(y)[[x_1, \ldots, x_n]] \longrightarrow \mathcal{O}_{Y, y}^\wedge,\quad
x_i \longmapsto f_i
$$
où ici $f_i$ signifie vraiment l'image de $f_i$ dans
$\mathfrak m_y^\wedge$. Puisque $\mathfrak m_x
\mathcal{O}_{Y, y} = \mathfrak m_y$ d'après le Lemme
\ref{lemma-change-fields-algebraic-unramified}, nous voyons
que le côté droit est complet par rapport à $(x_1, \ldots,
x_n)$ (utilisez Algèbre, Lemme
\ref{algebra-lemma-hathat-finitely-generated} pour voir que c'est
un anneau local complet). Puisque les deux côtés sont complets pour
la topologie $(x_1, \ldots, x_n)$-adique, notre application est
surjective par Algèbre, Lemme
\ref{algebra-lemma-completion-generalities} car elle est surjective modulo
$(x_1, \ldots, x_n)$. Bien sûr, les séries puissances $g_1, \ldots, g_m$ sont
envoyées sur zéro sous cette application, comme elles sont déjà envoyées sur
zéro dans $\mathcal{O}_{X, x}^\wedge$. Ainsi, nous avons le diagramme commutatif.
$$
\xymatrix{
\kappa(y)[[x_1, \ldots, x_n]]/(g_1, \ldots, g_m) \ar[r] &
\mathcal{O}_{Y, y}^\wedge \\
\kappa(x)[[x_1, \ldots, x_n]]/(g_1, \ldots, g_m) \ar[r] \ar[u] &
\mathcal{O}_{X, x}^\wedge \ar[u]
}
$$
Nous devons encore montrer que la flèche horizontale supérieure
est un isomorphisme. Nous savons déjà qu'elle est surjective.
Nous savons que $\mathcal{O}_{X, x} \to \mathcal{O}_{Y, y}$ est
plat (Lemme \ref{lemma-change-fields-flat}), ce qui implique
que $\mathcal{O}_{X, x}^\wedge \to \mathcal{O}_{Y, y}^\wedge$
est plat (Plus sur l'algèbre, Lemme
\ref{more-algebra-lemma-flat-completion}). Ainsi, nous pouvons
appliquer Algèbre, Lemme \ref{algebra-lemma-mod-injective} avec $R =
\kappa(x)[[x_1, \ldots, x_n]]/(g_1, \ldots, g_m)$, avec $S =
\kappa(y)[[x_1, \ldots, x_n]]/(g_1, \ldots, g_m)$, avec $M =
\mathcal{O}_{Y, y}^\wedge$ et avec $N = S$ pour conclure que l'application est injective.
\end{proof}


\section{Engendrement global}
\label{section-global-generation}

\noindent
Quelques lemmes liés à la génération globale des modules quasi-cohérents.

\begin{lemma}
\label{lemma-globally-generated-base-change}
Supposons que $X \to \Spec(A)$ soit un morphisme de schémas. Supposons que $A \subset
A'$ soit un homomorphisme d'anneaux fidèlement plat. Supposons que $\mathcal{F}$ soit
un $\mathcal{O}_X$-module quasi-cohérent. Alors $\mathcal{F}$ est globalement engendré
si et seulement si le changement de base $\mathcal{F}_{A'}$ est globalement engendré.
\end{lemma}

\begin{proof}
Plus précisément, posons $X_{A'} = X \times_{\Spec(A)}
\Spec(A')$. Soit $\mathcal{F}_{A'} = p^*\mathcal{F}$ où $p :
X_{A'} \to X$ est la projection. Par la cohomologie des
schémas, Lemme
\ref{coherent-lemma-flat-base-change-cohomology} nous avons
$H^0(X_{k'}, \mathcal{F}_{A'}) = H^0(X, \mathcal{F}) \otimes_A
A'$. Ainsi, si $s_i$, $i \in I$ sont des générateurs de $H^0(X,
\mathcal{F})$ comme module $A$, alors leurs images dans
$H^0(X_{A'}, \mathcal{F}_{A'})$ sont des générateurs de
$H^0(X_{A'}, \mathcal{F}_{A'})$ comme module $A'$. Ainsi, nous
devons montrer que l'application $\alpha : \bigoplus_{i \in I}
\mathcal{O}_X \to \mathcal{F}$, $(f_i) \mapsto \sum f_i s_i$ est
surjective si et seulement si $p^*\alpha$ est surjective. Cela,
nous pouvons le vérifier sur un ouvert affine $U = \Spec(B)$ de
$X$. Alors $\mathcal{F}|_U$ correspond à un module $B$ $M$ et
$s_i|_U$ aux éléments $x_i \in M$. Ainsi, nous devons montrer que
$\bigoplus_{i \in I} B \to M$ est surjective si et seulement si le
changement de base $\bigoplus_{i \in I} B \otimes_A A' \to M \otimes_A
A'$ est surjective. C'est vrai parce que $A \to A'$ est fidèlement plat.
\end{proof}

\begin{lemma}
\label{lemma-very-ample-vanish-at-point}
Supposons que $k$ soit un corps infini. Supposons que $X$ soit
un schéma de type fini sur $k$. Supposons que $\mathcal{L}$
soit un très faisceau inversible ample sur $X$. Supposons que
$n \geq 0$ et $x, x_1, \ldots, x_n \in X$ soient des points
avec $x$ un point $k$-rationnel, c'est-à-dire $\kappa(x) = k$,
et $x \not = x_i$ pour $i = 1, \ldots, n$. Alors il existe un $s
\in H^0(X, \mathcal{L})$ qui s'annule en $x$ mais pas en $x_i$.
\end{lemma}

\begin{proof}
Si $n = 0$ le résultat est trivial, nous supposons donc $n >
0$. Par définition d'un faisceau inversible très ample, le
lemme se réduit immédiatement au cas où $X = \mathbf{P}^r_k$
pour un certain $r > 0$ et $\mathcal{L} = \mathcal{O}_X(1)$.
Écrivons $\mathbf{P}^r_k = \text{Proj}(k[T_0, \ldots, T_r])$.
Posons $V = H^0(X, \mathcal{L}) = kT_0 \oplus \ldots \oplus
kT_r$. Puisque $x$ est un point $k$-rationnel, nous voyons que
l'ensemble $s \in V$ qui s'annule en $x$ est un sous-espace de
codimension $1$ $W \subset V$ et que $W$ génère l'idéal homogène
premier correspondant à $x$. Puisque $x_i \not = x$, le premier
homogène correspondant $\mathfrak p_i \subset k[T_0, \ldots, T_r]$
ne contient pas $W$. Comme $k$ est infini, nous voyons alors que $W
\not = \bigcup W \cap \mathfrak q_i$ et la démonstration est complète.
\end{proof}

\begin{lemma}
\label{lemma-generated-by-dim-plus-1-sections}
Supposons que $k$ soit un corps infini. Supposons que $X$ soit un $k$ -schéma
algébrique. Supposons que $\mathcal{L}$ soit un $\mathcal{O}_X$ -module
inversible. Supposons que $V \to \Gamma(X, \mathcal{L})$ soit une application
linéaire de $k$ -espaces vectoriels dont l'image génère $\mathcal{L}$. Alors il existe
un sous-espace $W \subset V$ avec $\dim_k(W) \leq \dim(X) + 1$ qui génère $\mathcal{L}$.
\end{lemma}

\begin{proof}
Tout au long de la démonstration, nous utiliserons
que pour chaque $x \in X$, l'application linéaire
$$
\psi_x : V \to \Gamma(X, \mathcal{L}) \to \mathcal{L}_x \to
\mathcal{L}_x \otimes_{\mathcal{O}_{X, x}} \kappa(x)
$$
n'est pas nul. La démonstration se fait par induction sur $\dim(X)$.

\medskip\noindent
Le cas de base est $\dim(X) = 0$. Dans ce cas, $X$ a un
nombre fini de points $X = \{x_1, \ldots, x_n\}$ (voir par
exemple le Lemme \ref{lemma-algebraic-scheme-dim-0}). Puisque
$k$ est infini, il existe un vecteur $v \in V$ tel que
$\psi_{x_i}(v) \not = 0$ pour tout $i$. Alors $W = k\cdot v$ fait l’affaire.

\medskip\noindent
Supposons $\dim(X) > 0$. Supposons que $X_i \subset X$ soit les
composantes irréductibles de dimension égale à $\dim(X)$. Comme $X$
est noethérien, il n’y en a qu’un nombre fini. Pour chaque $i$,
choisissez un point $x_i \in X_i$. Comme ci-dessus, choisissez $v
\in V$ tel que $\psi_{x_i}(v) \not = 0$ pour tout $i$. Supposons que
$Z \subset X$ soit le schéma zéro de l’image de $v$ dans $\Gamma(X,
\mathcal{L})$, voir Diviseurs, Définition
\ref{divisors-definition-zero-scheme-s}. Par construction $\dim(Z) < \dim(X)$.
Par induction, nous pouvons trouver $W \subset V$ avec $\dim(W) \leq \dim(X)$
tel que $W$ génère $\mathcal{L}|_Z$. Alors $W + k\cdot v$ génère $\mathcal{L}$.
\end{proof}






\section{Séparation des points et des vecteurs tangents}
\label{section-separating-points-tangent-vectors}

\noindent
Ceci est juste le résultat suivant.

\begin{lemma}
\label{lemma-separate-points-tangent-vectors}
Supposons que $k$ soit un corps algébriquement clos. Supposons
que $X$ soit un $k$ -schéma propre. Supposons que $\mathcal{L}$
soit un $\mathcal{O}_X$ -module inversible. Supposons que $V
\subset H^0(X, \mathcal{L})$ soit un $k$ -sous-espace vectoriel. Si
\begin{enumerate}
\item pour chaque paire de points distincts fermés $x, y \in X$ il
existe une section $s \in V$ qui s'annule en $x$ mais pas en $y$,
et \item pour chaque point fermé $x \in X$ et vecteur tangent non
nul $\theta \in T_{X/k, x}$ il existe une section $s \in V$ qui
s'annule en $x$ mais dont le tiré en arrière par $\theta$ est non nul,
\end{enumerate}
alors $\mathcal{L}$ est très ample et le
morphisme canonique $\varphi_{\mathcal{L}, V} :
X \to \mathbf{P}(V)$ est une immersion fermée.
\end{lemma}

\begin{proof}
La condition (1) implique en particulier que les éléments de $V$ génèrent
$\mathcal{L}$ sur $X$. Par conséquent, nous obtenons un morphisme canonique
$$
\varphi = \varphi_{\mathcal{L}, V} :
X
\longrightarrow
\mathbf{P}(V)
$$
par Constructions, Exemple
\ref{constructions-example-projective-space}. Le morphisme
$\varphi$ est propre par morphismes, Lemme
\ref{morphisms-lemma-image-proper-scheme-closed}. D'après (1),
l'application $\varphi$ est injective sur les points fermés (calcul
omis). En particulier, la fibre au-dessus de tout point fermé de
$\mathbf{P}(V)$ est un singleton (petit détail omis). Ainsi, nous voyons
que $\varphi$ est fini, par exemple en utilisant Cohomologie des
schémas, Lemme
\ref{coherent-lemma-proper-finite-fibre-finite-in-neighbourhood}. Pour terminer la démonstration, il suffit de montrer que l'application
$$
\varphi^\sharp :
\mathcal{O}_{\mathbf{P}(V)}
\longrightarrow
\varphi_*\mathcal{O}_X
$$
est surjective. Nous pouvons le vérifier sur les germes aux
points fermés. Soit $x \in X$ un point fermé d'image le point
fermé $p = \varphi(x) \in \mathbf{P}(V)$. Puisque
$\varphi^{-1}(\{p\}) = \{x\}$ par (1) et que $\varphi$ est
propre (donc fermé), nous voyons que $\varphi^{-1}(U)$ décrit un
système fondamental de voisinages ouverts de $x$ lorsque $U$
décrit un système fondamental de voisinages ouverts de $p$. Nous
concluons que, sur les germes en $p$, nous obtenons l'application
$$
\varphi^\sharp_x :
\mathcal{O}_{\mathbf{P}(V), p}
\longrightarrow
\mathcal{O}_{X, x}
$$
En particulier, $\mathcal{O}_{X, x}$ est un
$\mathcal{O}_{\mathbf{P}(V), p}$ -module fini. De plus, les
corps résiduels de $x$ et $p$ sont égaux à $k$ (comme $k$
est algébriquement clos, -- utilise le Nullstellensatz de
Hilbert). Enfin, la condition (2) implique que l'application
$$
T_{X/k, x} \longrightarrow T_{\mathbf{P}(V)/k, p}
$$
est injective puisque tout $\theta$ non nul dans le noyau de cette
application ne pourrait pas satisfaire la conclusion de (2). En
termes de l'application des anneaux locaux ci-dessus, cela signifie que
$$
\mathfrak m_p/\mathfrak m_p^2 \longrightarrow
\mathfrak m_x/\mathfrak m_x^2
$$
est surjective, voir le Lemme
\ref{lemma-tangent-space-rational-point}. Maintenant, la preuve est
terminée en appliquant l'Algèbre, Lemme \ref{algebra-lemma-when-surjective-local}.
\end{proof}

\begin{lemma}
\label{lemma-variant-separate-points-tangent-vectors}
Supposons que $k$ soit un corps algébriquement clos.
Supposons que $X$ soit un $k$-schéma propre. Supposons
que $\mathcal{L}$ soit un $\mathcal{O}_X$-module
inversible. Supposons que pour tout sous-schéma fermé $Z
\subset X$ de dimension $0$ et de degré $2$ sur $k$, l'application
$$
H^0(X, \mathcal{L}) \longrightarrow H^0(Z, \mathcal{L}|_Z)
$$
est surjective. Ensuite, $\mathcal{L}$ est très ample sur $X$ au-dessus de $k$.
\end{lemma}

\begin{proof}
Ceci est une reformulation du Lemme
\ref{lemma-separate-points-tangent-vectors}. À
savoir, étant donné des points distincts fermés $x, y
\in X$ en prenant $Z = x \cup y$ (considéré comme
sous-schéma fermé) nous obtenons la condition (1) du
lemme. Et étant donné un vecteur tangent non nul
$\theta \in T_{X/k, x}$ le morphisme $\theta :
\Spec(k[\epsilon]) \to X$ est une immersion fermée. En
posant $Z = \Im(\theta)$ nous obtenons la condition (2) du lemme.
\end{proof}









\section{Adhérences de produits}
\label{section-closure-of-products}

\noindent
Quelques résultats sur la relation entre l'adhérence et les produits.

\begin{lemma}
\label{lemma-closure-of-product}
Supposons que $k$ soit un corps.
Supposons que $X$, $Y$ soient des
schémas sur $k$, et soient $A \subset X$,
$B \subset Y$ des sous-ensembles. Posons
$$
AB =
\{z \in X \times_k Y \mid \text{pr}_X(z) \in A, \ \text{pr}_Y(z) \in B\}
\subset X \times_k Y
$$
Alors théoriquement nous avons
$$
\overline{A} \times_k \overline{B} = \overline{AB}
$$
\end{lemma}

\begin{proof}
L'inclusion $\overline{AB} \subset \overline{A}
\times_k \overline{B}$ est immédiate. Nous pouvons
remplacer $X$ et $Y$ par les sous-schémas fermés réduits
$\overline{A}$ et $\overline{B}$. Supposons que $W
\subset X \times_k Y$ soit un ouvert non vide. Par
morphismes, le Lemme
\ref{morphisms-lemma-scheme-over-field-universally-open}, le
sous-ensemble $U = \text{pr}_X(W)$ est ouvert non vide dans
$X$. Ainsi $A \cap U$ n'est pas vide. Prenons $a \in A \cap
U$. Désignons par $Y_a = \{a\} \times_k Y \cong
Y_{\kappa(a)}$ la fibre de $\text{pr}_X : X \times_k Y \to X$
au-dessus de $a$. Par morphismes, le Lemme
\ref{morphisms-lemma-scheme-over-field-universally-open} montre
encore que le morphisme $Y_a \to Y$ est ouvert puisque
$\Spec(\kappa(a)) \to \Spec(k)$ est universellement ouvert. Ainsi
l'ouvert non vide $W_a = W \times_{X \times_k Y} Y_a$ se projette
sur un ouvert non vide de $Y$. Nous en concluons qu'il existe un $b
\in B$ dans l'image. Ainsi $AB \cap W \not = \emptyset$ comme souhaité.
\end{proof}

\begin{lemma}
\label{lemma-closure-image-product-map}
Supposons que $k$ soit un corps. Supposons que $f : A
\to X$, $g : B \to Y$ soient des morphismes de schémas
sur $k$. Alors, d'un point de vue des ensembles, nous avons
$$
\overline{f(A)} \times_k \overline{g(B)} =
\overline{(f \times g)(A \times_k B)}
$$
\end{lemma}

\begin{proof}
Cela découle du Lemme
\ref{lemma-closure-of-product}
puisque l'image de $f \times g$ est
$f(A)g(B)$ dans la notation de ce lemme.
\end{proof}

\begin{lemma}
\label{lemma-scheme-theoretic-image-product-map}
Supposons que $k$ soit un corps. Supposons que $f : A \to X$, $g
: B \to Y$ soient des morphismes quasi-compacts de schémas sur
$k$. Supposons que $Z \subset X$ soit l'image schématique de $f$,
voir morphismes, Définition
\ref{morphisms-definition-scheme-theoretic-image}. De même, soit $Z' \subset Y$
l'image schématique de $g$. Alors $Z \times_k Z'$ est l'image schématique de $f \times g$.
\end{lemma}

\begin{proof}
Rappelez-vous que $Z$ est le plus petit sous-schéma fermé de $X$ à travers
lequel $f$ se factorise. De même pour $Z'$. Supposons que $W \subset X
\times_k Y$ soit l'image schématique de $f \times g$. Comme $f \times g$ se
factorise à travers $Z \times_k Z'$, nous voyons que $W \subset Z \times_k Z'$.

\medskip\noindent
Pour prouver l'autre inclusion, soit $U \subset X$ et $V
\subset Y$ des ouverts affines. Par les morphismes, Lemme
\ref{morphisms-lemma-quasi-compact-scheme-theoretic-image},
le schéma $Z \cap U$ est l'image schématique de
$f|_{f^{-1}(U)} : f^{-1}(U) \to U$, et de même pour $Z'
\cap V$ et $W \cap U \times_k V$. Par conséquent, nous
pouvons supposer que $X$ et $Y$ sont affines. Comme $f$ et
$g$ sont quasi-compacts, cela implique que $A = \bigcup
U_i$ est une union finie d'affines et $B = \bigcup V_j$ est
une union finie d'affines. Ensuite, nous pouvons remplacer
$A$ par $\coprod U_i$ et $B$ par $\coprod V_j$,
c'est-à-dire que nous pouvons supposer que $A$ et $B$ sont
également affines. Dans ce cas, $Z$ est défini par
$\Ker(\Gamma(X, \mathcal{O}_X) \to \Gamma(A, \mathcal{O}_A))$ et de
même pour $Z'$ et $W$. Par conséquent, le résultat découle de l'égalité
$$
\Gamma(A \times_k B, \mathcal{O}_{A \times_k B})
=
\Gamma(A, \mathcal{O}_A) \otimes_k \Gamma(B, \mathcal{O}_B)
$$
ce qui tient car $A$ et $B$ sont affines. Détails omis.
\end{proof}





\section{Schémas lisses sur des corps}
\label{section-smooth}

\noindent
Voici deux lemmes caractérisant les schémas lisses sur les corps.

\begin{lemma}
\label{lemma-char-zero-differentials-free-smooth}
Supposons que $k$ soit un corps. Supposons
que $X$ soit un schéma sur $k$. Supposons
\begin{enumerate}
\item $X$ est localement de type fini sur
$k$, \item $\Omega_{X/k}$ est localement
libre, et \item $k$ a caractéristique zéro.
\end{enumerate}
Alors la structure morphisme $X \to \Spec(k)$ est lisse.
\end{lemma}

\begin{proof}
Cela découle de l'algèbre, lemme
\ref{algebra-lemma-characteristic-zero-local-smooth}.
\end{proof}

\noindent
En caractéristique positive, il existe des schémas de type
fini non réduits dont le faisceau des différentielles est
libre, par exemple $\Spec(\mathbf{F}_p[t]/(t^p))$ sur
$\Spec(\mathbf{F}_p)$. Si le corps de base $k$ est non parfait
de caractéristique $p$, il existe des schémas réduits $X/k$
avec $\Omega_{X/k}$ libre qui ne sont pas lisses, par exemple
$\Spec(k[t]/(t^p-a)$ où $a \in k$ n'est pas une puissance $p$.

\begin{lemma}
\label{lemma-char-p-differentials-free-smooth}
Supposons que $k$ soit un corps. Supposons
que $X$ soit un schéma sur $k$. Supposons
\begin{enumerate}
\item $X$ est localement de type
fini sur $k$, \item $\Omega_{X/k}$
est localement libre, \item $X$ est
réduit, et \item $k$ est parfait.
\end{enumerate}
Alors la structure morphisme $X \to \Spec(k)$ est lisse.
\end{lemma}

\begin{proof}
Supposons que $x \in X$ soit un point. Comme $X$ est localement
noethérien (voir morphismes, Lemme
\ref{morphisms-lemma-finite-type-noetherian}), il y a un nombre
fini de composantes irréductibles $X_1, \ldots, X_n$ passant par
$x$ (voir Propriétés, Lemme
\ref{properties-lemma-Noetherian-topology} et Topologie, Lemme
\ref{topology-lemma-Noetherian}). Supposons que $\eta_i \in X_i$ soit
le point générique. Comme $X$ est réduit, nous avons $\mathcal{O}_{X,
\eta_i} = \kappa(\eta_i)$, voir Algèbre, Lemme
\ref{algebra-lemma-minimal-prime-reduced-ring}. De plus,
$\kappa(\eta_i)$ est une extension de corps de type fini du corps parfait
$k$, donc séparablement engendrée sur $k$ (voir Algèbre, Section
\ref{algebra-section-separability}). Il s'ensuit que $\Omega_{X/k, \eta_i} =
\Omega_{\kappa(\eta_i)/k}$ est libre de rang le degré de transcendance de
$\kappa(\eta_i)$ sur $k$. Par morphismes, Lemme
\ref{morphisms-lemma-dimension-fibre-at-a-point}, nous concluons que $\dim_{\eta_i}(X_i) =
\text{rank}_{\eta_i}(\Omega_{X/k})$. Puisque $x \in X_1 \cap \ldots \cap X_n$ nous voyons que
$$
\text{rank}_x(\Omega_{X/k}) = \text{rank}_{\eta_i}(\Omega_{X/k}) = \dim(X_i).
$$
Par conséquent $\dim_x(X) = \text{rank}_x(\Omega_{X/k})$, voir
Algèbre, Lemme
\ref{algebra-lemma-dimension-at-a-point-finite-type-over-field}. Il s'ensuit que $X \to \Spec(k)$
est lisse en $x$ par exemple selon Algèbre, Lemme \ref{algebra-lemma-characterize-smooth-over-field}.
\end{proof}

\begin{lemma}
\label{lemma-smooth-regular}
\begin{slogan}
Lisser un corps implique régulier
\end{slogan}
Supposons que $X \to \Spec(k)$ soit un morphisme lisse
où $k$ est un corps. Alors $X$ est un schéma régulier.
\end{lemma}

\begin{proof}
(Voir aussi le Lemme
\ref{lemma-geometrically-regular-smooth}.) D'après le
Lemme d'Algèbre
\ref{algebra-lemma-characterize-smooth-over-field}, tout
anneau local $\mathcal{O}_{X, x}$ est régulier. Et parce que
$X$ est localement de type fini sur $k$, il est localement
noethérien. Par conséquent, $X$ est régulier selon le Lemme
sur les Propriétés \ref{properties-lemma-characterize-regular}.
\end{proof}

\begin{lemma}
\label{lemma-smooth-geometrically-normal}
Supposons que $X \to \Spec(k)$ soit un morphisme lisse où
$k$ est un corps. Alors $X$ est géométriquement régulier,
géométriquement normal et géométriquement réduit sur $k$.
\end{lemma}

\begin{proof}
(Voir aussi le Lemme
\ref{lemma-geometrically-regular-smooth}.)
Supposons que $k'$ soit une extension finie
purement inséparable de $k$. Il suffit de
prouver que $X_{k'}$ est régulier, normal,
réduit, voir les Lemmata
\ref{lemma-geometrically-regular},
\ref{lemma-geometrically-normal}, et
\ref{lemma-check-only-finite-inseparable-extensions}.
Par les morphismes, Lemme
\ref{morphisms-lemma-base-change-smooth}, le morphisme $X_{k'}
\to \Spec(k')$ est lisse aussi. Par conséquent, il suffit de
montrer qu'un schéma $X$ lisse sur un corps est régulier,
normal et réduit. Nous voyons que $X$ est régulier par Lemme
\ref{lemma-smooth-regular}. Par conséquent, Propriétés, Lemme
\ref{properties-lemma-regular-normal} garantit que $X$ est normal.
\end{proof}

\begin{lemma}
\label{lemma-affine-space-over-field}
Supposons que $k$ soit un corps. Supposons que $d \geq 0$. Soit $W
\subset \mathbf{A}^d_k$ ouvert non vide. Alors il existe un point
fermé $w \in W$ tel que $k \subset \kappa(w)$ soit fini séparable.
\end{lemma}

\begin{proof}
Après un rétrécissement possible $W$, nous pouvons supposer
que $W = \mathbf{A}^d_k \setminus V(f)$ pour certains $f
\in k[x_1, \ldots, x_d]$. Si le lemme est faux, alors
$f(a_1, \ldots, a_d) = 0$ pour tous les $(a_1, \ldots, a_d) \in
(k^{sep})^d$. C'est absurde car $k^{sep}$ est un corps infini.
\end{proof}

\begin{lemma}
\label{lemma-smooth-separable-closed-points-dense}
Supposons que $k$ soit un corps. Si $X$
est lisse sur $\Spec(k)$ alors l'ensemble
$$
\{x \in X\text{ closed such that }k \subset \kappa(x)
\text{ is finite separable}\}
$$
est dense en $X$.
\end{lemma}

\begin{proof}
Il suffit de montrer que, étant donné un $X$ lisse non vide sur
$k$, il existe au moins un point fermé dont le corps résiduel
est fini séparable sur $k$. Pour le voir, choisissez un diagramme
$$
\xymatrix{
X & U \ar[l] \ar[r]^-\pi & \mathbf{A}^d_k
}
$$
avec $\pi$ \' étalé, voir morphismes, Lemme
\ref{morphisms-lemma-smooth-etale-over-affine-space}.
Le morphisme $\pi : U \to \mathbf{A}^d_k$ est ouvert,
voir morphismes, Lemme
\ref{morphisms-lemma-etale-open}. D'après le Lemme
\ref{lemma-affine-space-over-field}, nous pouvons choisir
un point fermé $w \in \pi(U)$ dont le corps résiduel est
séparable fini sur $k$. Choisissez n'importe quel $x \in
U$ avec $\pi(x) = w$. D'après morphismes, Lemme
\ref{morphisms-lemma-etale-over-field}, l'extension de
corps $\kappa(x)/\kappa(w)$ est séparable finie. Par
conséquent, $\kappa(x)/k$ est séparable finie. Le point $x$
est un point fermé de $X$ d'après morphismes, Lemme
\ref{morphisms-lemma-algebraic-residue-field-extension-closed-point-fibre}.
\end{proof}

\begin{lemma}
\label{lemma-geometrically-reduced-dense-smooth-open}
Supposons que $X$ soit un schéma réduit qui est localement de type
fini sur un corps $k$. Alors $X$ est géométriquement réduit sur $k$
si et seulement si $X$ contient un ouvert dense qui est lisse sur $k$.
\end{lemma}

\begin{proof}
Le problème est local sur $X$, par conséquent nous pouvons supposer que
$X$ est quasi-compact. Supposons que $X = X_1 \cup \ldots \cup X_n$ sont
les composantes irréductibles de $X$. Supposons d'abord que $X$ contient
un ouvert dense qui est lisse sur $k$. Alors $X$ est géométriquement
réduit en le point générique de chaque $X_i$ (par exemple selon le Lemme
\ref{lemma-smooth-geometrically-normal}). D'après le Lemme
\ref{lemma-generic-points-geometrically-reduced}, il s'ensuit que $X$ est géométriquement réduit.

\medskip\noindent
Inversement, supposons que $X$ soit géométriquement réduit.
Comme $Z = \bigcup_{i \not = j} X_i \cap X_j$ est nulle part
dense dans $X$, nous pouvons remplacer $X$ par $X \setminus Z$.
Comme $X \setminus Z$ est une réunion disjointe de schémas
irréductibles, cela nous ramène au cas où $X$ est irréductible.
Comme $X$ est irréductible et réduit, il est intègre, voir
Propriétés, Lemme \ref{properties-lemma-characterize-integral}.
Supposons que $\eta \in X$ soit son point générique. Alors le
corps des fonctions $K = k(X) = \kappa(\eta) = \mathcal{O}_{X,
\eta}$ est géométriquement réduit sur $k$, donc séparable sur
$k$, voir Algèbre, Lemme
\ref{algebra-lemma-characterize-separable-field-extensions}. Supposons
que $U = \Spec(A) \subset X$ soit un ouvert affine non vide quelconque
de sorte que $K = A_{(0)}$ soit le corps des fractions de $A$. Appliquez
Algèbre, Lemme \ref{algebra-lemma-separable-smooth} pour conclure que
$A$ est lisse en $(0)$ sur $k$. Par définition, cela signifie qu'une
certaine localisation principale de $A$ est lisse sur $k$ et nous avons gagné.
\end{proof}

\begin{lemma}
\label{lemma-dense-smooth-open-variety-over-perfect-field}
Supposons que $k$ soit un corps parfait. Supposons que $X$ soit un $k$ -schéma
réduit localement algébrique, par exemple une variété sur $k$. Alors nous avons
$$
\{x \in X \mid X \to \Spec(k)\text{ is smooth at }x\} =
\{x \in X \mid \mathcal{O}_{X, x}\text{ is regular}\}
$$
et ceci est un sous-schéma ouvert dense de $X$.
\end{lemma}

\begin{proof}
L'égalité des deux ensembles découle immédiatement de
l'Algèbre, Lemme \ref{algebra-lemma-separable-smooth} et des
définitions (voir Algèbre, Définition
\ref{algebra-definition-perfect} pour la définition d'un corps
parfait). L'ensemble est ouvert car l'ensemble des points où un
morphisme de schémas est lisse est ouvert, voir morphismes,
Définition \ref{morphisms-definition-smooth}. Enfin, nous donnons
deux arguments pour voir qu'il est dense : (1) Les points génériques
de $X$ sont dans l'ensemble car les anneaux locaux aux points
génériques sont des corps (Algèbre, Lemme
\ref{algebra-lemma-minimal-prime-reduced-ring}), donc réguliers. (2) Nous
utilisons que $X$ est géométriquement réduit d'après le Lemme
\ref{lemma-perfect-reduced} et donc le Lemme \ref{lemma-geometrically-reduced-dense-smooth-open} s'applique.
\end{proof}

\begin{lemma}
\label{lemma-flat-under-smooth}
Supposons que $k$ soit un corps. Supposons que $f : X \to Y$ soit un morphisme
de schémas localement de type fini sur $k$. Supposons que $x \in X$ soit un
point et posons $y = f(x)$. Si $X \to \Spec(k)$ est lisse en $x$ et que $f$ est
plat en $x$ alors $Y \to \Spec(k)$ est lisse en $y$. En particulier, si $X$
est lisse sur $k$ et que $f$ est plat et surjectif, alors $Y$ est lisse sur $k$.
\end{lemma}

\begin{proof}
Il suffit de montrer que $Y$ est
géométriquement régulier en $y$, voir
le Lemme
\ref{lemma-geometrically-regular-smooth}. Cela
découle du Lemme
\ref{lemma-flat-under-geometrically-regular} (et du Lemme
\ref{lemma-geometrically-regular-smooth} appliqué à $(X, x)$).
\end{proof}

\begin{lemma}
\label{lemma-variety-with-smooth-rational-point}
Soit $k$ un corps. Soit $X$ une variété sur $k$
possédant un point $k$-rationnel $x$ tel que $X$ soit
lisse en $x$. Alors $X$ est géométriquement intègre sur $k$.
\end{lemma}

\begin{proof}
Supposons que $U \subset X$ soit le lieu lisse de $X$. Par
hypothèse, $U$ n'est pas vide et est donc dense et
schématiquement dense. Alors $U_{\overline{k}} \subset
X_{\overline{k}}$ est également dense et schématiquement dense
(certains détails omis). Par conséquent, il suffit de montrer
que $U$ est géométriquement intègre. Comme $U$ possède un point
$k$-rationnel, il est géométriquement connexe d'après le Lemme
\ref{lemma-geometrically-connected-if-connected-and-point}.
D'autre part, $U_{\overline{k}}$ est réduit et normal (Lemme
\ref{lemma-smooth-geometrically-normal}). Puisqu'un schéma
Noetherien connecté et normal est intègre (Propriétés, Lemme
\ref{properties-lemma-normal-Noetherian}), la démonstration est complète.
\end{proof}

\begin{lemma}
\label{lemma-smooth-stratification}
Supposons que $X$ soit un schéma de type fini sur un
corps $k$. Il existe une extension finie purement
inséparable $k'/k$, un entier $t \geq 0$, et des sous-schémas fermés
$$
X_{k'} \supset Z_0 \supset Z_1 \supset \ldots \supset Z_t = \emptyset
$$
de sorte que $Z_0 = (X_{k'})_{red}$ et $Z_i
\setminus Z_{i + 1}$ soient lisses sur $k'$ pour tout $i$.
\end{lemma}

\begin{proof}
Nous pouvons utiliser l'induction sur $\dim(X)$. D'après le
Lemme \ref{lemma-finite-extension-geometrically-reduced}, nous
pouvons trouver une extension finie purement inséparable
$k'/k$ telle que $(X_{k'})_{red}$ soit géométriquement réduit
sur $k'$. D'après le Lemme
\ref{lemma-geometrically-reduced-dense-smooth-open}, il existe un
sous-schéma fermé de densité nulle $X' \subset (X_{k'})_{red}$ tel que
$(X_{k'})_{red} \setminus X'$ soit lisse sur $k'$. Ensuite $\dim(X') <
\dim(X)$. Par hypothèse d'induction, il existe une extension finie
purement inséparable $k''/k'$, un entier $t' \geq 0$, et des sous-schémas fermés
$$
X'_{k''} \supset Y_0 \supset Y_1 \supset \ldots \supset Y_{t'} = \emptyset
$$
tel que $Y_0 = (X'_{k''})_{red}$ et $Y_i \setminus
Y_{i + 1}$ soient lisses sur $k''$ pour tout $i$.
Ensuite, nous posons $t = t' + 1$ et nous considérons
$$
X_{k''} \supset Z_0 \supset Z_1 \supset \ldots \supset Z_t = \emptyset
$$
donné par $Z_0 = (X_{k''})_{red}$ et $Z_i = Y_{i - 1}$ pour $i > 0$ ; cela
a du sens car $X'_{k''}$ est un sous-schéma fermé de $X_{k''}$. Nous
omettons la vérification que toutes les propriétés énoncées sont satisfaites.
\end{proof}




\section{Types de variétés}
\label{section-types}

\noindent
Courte section discutant certaines propriétés globales élémentaires des variétés.

\begin{definition}
\label{definition-variety-type}
Supposons que $k$ soit un corps. Supposons que $X$ soit une variété sur $k$.
\begin{enumerate}
\item On dit que $X$ est une {\it variété affine } si $X$ est
un schéma affine. Cela équivaut à exiger que $X$ soit
isomorphe à un sous-schéma fermé de $\mathbf{A}^n_k$ pour un
certain $n$. \item On dit que $X$ est une {\it variété
projective } si le morphisme de structure $X \to \Spec(k)$ est
projectif. D'après le lemme sur les morphismes,
\ref{morphisms-lemma-characterize-locally-projective}, cela est
vrai si et seulement si $X$ est isomorphe à un sous-schéma fermé
de $\mathbf{P}^n_k$ pour un certain $n$. \item On dit que $X$ est
une {\it variété quasi-projective } si le morphisme de structure
$X \to \Spec(k)$ est quasi-projectif. D'après le lemme sur les
morphismes,
\ref{morphisms-lemma-characterize-locally-quasi-projective}, cela est vrai
si et seulement si $X$ est isomorphe à un sous-schéma localement fermé de
$\mathbf{P}^n_k$ pour un certain $n$. \item Une {\it variété propre } est une
variété telle que le morphisme $X \to \Spec(k)$ est propre. \item A {\it
variété lisse } est une variété telle que le morphisme $X \to \Spec(k)$ est lisse.
\end{enumerate}
\end{definition}

\noindent
Notez qu'une variété projective est une variété propre,
voir morphismes, Lemme
\ref{morphisms-lemma-locally-projective-proper}. De plus, une
variété affine est quasi-projective car $\mathbf{A}^n_k$ est
isomorphe à un sous-schéma ouvert de $\mathbf{P}^n_k$, voir
Constructions, Lemme \ref{constructions-lemma-standard-covering-projective-space}.

\begin{lemma}
\label{lemma-regular-functions-proper-variety}
Supposons que $X$ soit une variété propre sur $k$. Alors
\begin{enumerate}
\item $K = H^0(X, \mathcal{O}_X)$ est un corps
qui est une extension finie du corps $k$, \item si
$X$ est géométriquement réduit, alors $K/k$ est
séparable, \item si $X$ est géométriquement
irréductible, alors $K/k$ est purement inséparable,
\item si $X$ est géométriquement intègre, alors $K = k$.
\end{enumerate}
\end{lemma}

\begin{proof}
Ceci est un cas particulier du Lemme
\ref{lemma-proper-geometrically-reduced-global-sections}.
\end{proof}




\section{Normalisation}
\label{section-normalization}

\noindent
Certains problèmes associés à la normalisation.

\begin{lemma}
\label{lemma-normalization-locally-algebraic}
Supposons que $k$ soit un corps. Supposons que $X$ soit un
schéma localement algébrique sur $k$. Supposons que $\nu :
X^\nu \to X$ soit le morphisme de normalisation, voir
morphismes, Définition \ref{morphisms-definition-normalization}. Alors
\begin{enumerate}
\item $\nu$ est fini, dominant, et $X^\nu$ est une
réunion disjointe de schémas irréductibles normaux
localement algébriques sur $k$, \item $\nu$ se factorise
en $X^\nu \to X_{red} \to X$ et le premier morphisme est
le morphisme de normalisation de $X_{red}$, \item si $X$
est un schéma algébrique réduit, alors $\nu$ est
birationnel, \item si $X$ est une variété, alors $X^\nu$ est une
variété et $\nu$ est un morphisme birationnel fini de variétés.
\end{enumerate}
\end{lemma}

\begin{proof}
Puisque $X$ est localement de type fini sur un corps,
nous voyons que $X$ est localement noethérien
(morphismes, Lemme
\ref{morphisms-lemma-finite-type-noetherian}), donc chaque
ouvert quasi-compact a un nombre fini de composantes
irréductibles (Propriétés, Lemme
\ref{properties-lemma-Noetherian-irreducible-components}). Ainsi
la définition des morphismes,
\ref{morphisms-definition-normalization} s'applique. La
normalisation $X^\nu$ est toujours une réunion disjointe de schémas
intègres normaux et le morphisme de normalisation $\nu$ est toujours
dominant, voir morphismes, Lemme
\ref{morphisms-lemma-normalization-normal}. Comme $X$ est universellement Nagata
(morphismes, Lemme \ref{morphisms-lemma-ubiquity-nagata}), nous voyons que $\nu$
est fini (morphismes, Lemme \ref{morphisms-lemma-nagata-normalization}). Ainsi
$X^\nu$ est également localement algébrique. À ce stade, nous avons prouvé (1).

\medskip\noindent
La partie (2) est morphismes, Lemme \ref{morphisms-lemma-normalization-reduced}.

\medskip\noindent
La partie (3) est morphismes, Lemme \ref{morphisms-lemma-normalization-birational}.

\medskip\noindent
La partie (4) découle de (1), (2), (3), et du fait que
$X^\nu$ est séparé en tant que schéma fini sur un schéma séparé.
\end{proof}

\begin{lemma}
\label{lemma-normalize-projective}
Supposons que $k$ soit un corps. Supposons que $X$ soit un
schéma propre sur $k$. Supposons que $\nu : X^\nu \to X$ soit le
morphisme de normalisation, voir morphismes, Définition
\ref{morphisms-definition-normalization}. Alors $X^\nu$ est propre sur
$k$. Si $X$ est projectif sur $k$, alors $X^\nu$ est projectif sur $k$.
\end{lemma}

\begin{proof}
D'après le Lemme \ref{lemma-normalization-locally-algebraic}
le morphisme $\nu$ est fini. Par conséquent, $X^\nu$ est propre
sur $k$ selon les morphismes, les Lemmata
\ref{morphisms-lemma-finite-proper} et
\ref{morphisms-lemma-composition-proper}. Le morphisme $\nu$ est projectif
selon les morphismes, le Lemme \ref{morphisms-lemma-finite-projective} et
donc si $X$ est projectif sur $k$, alors $X^\nu$ est projectif sur $k$
selon les morphismes, le Lemme \ref{morphisms-lemma-composition-projective}.
\end{proof}

\begin{lemma}
\label{lemma-relative-normalization-finite}
Supposons que $k$ soit un corps. Supposons que $f : Y \to X$
soit un morphisme quasi-compact de schémas localement
algébriques sur $k$. Supposons que $X'$ soit la normalisation
de $X$ dans $Y$. Si $Y$ est réduit, alors $X' \to X$ est fini.
\end{lemma}

\begin{proof}
Puisque $Y$ est quasi-séparé (par Propriétés, Lemme
\ref{properties-lemma-locally-Noetherian-quasi-separated}
et morphismes, Lemme
\ref{morphisms-lemma-finite-type-noetherian}), le morphisme $f$
est quasi-séparé (schémas, Lemme
\ref{schemes-lemma-compose-after-separated}). Ainsi, morphismes,
Définition \ref{morphisms-definition-normalization-X-in-Y}
s'applique. Le résultat découle de morphismes, Lemme
\ref{morphisms-lemma-nagata-normalization-finite-general}. Cela utilise
que les schémas localement algébriques sont localement noethériens (et
ont donc localement un nombre fini de composantes irréductibles) et que
les schémas localement algébriques sont de Nagata (morphismes, Lemme
\ref{morphisms-lemma-ubiquity-nagata}). Quelques petits détails sont omis.
\end{proof}

\begin{lemma}
\label{lemma-finite-extension-geometrically-normal}
Supposons que $k$ soit un corps. Supposons que $X$ soit un $k$ -schéma
algébrique. Alors il existe une extension finie purement inséparable $k'/k$
telle que la normalisation $Y$ de $X_{k'}$ soit géométriquement normale sur $k'$.
\end{lemma}

\begin{proof}
Supposons que $K = k^{perf}$ soit l'adhérence parfaite.
Supposons que $Y_K$ soit la normalisation de $X_K$, voir le
Lemme \ref{lemma-normalization-locally-algebraic}. D'après
Limites, Lemme \ref{limits-lemma-descend-finite-presentation}, il
existe une sous-extension finie $K/k'/k$ et un morphisme $\nu : Y
\to X_{k'}$ de présentation finie dont le changement de base vers
$K$ est le morphisme de normalisation $\nu_K : Y_K \to X_K$. On
observe que $Y$ est géométriquement normal sur $k'$ (Lemme
\ref{lemma-geometrically-normal}). Après avoir augmenté $k'$,
nous pouvons supposer que $Y \to X_{k'}$ est fini (Limites, Lemme
\ref{limits-lemma-descend-finite-finite-presentation}). Comme
$\nu_K : Y_K \to X_K$ est le morphisme de normalisation, il induit
un morphisme birationnel $Y_K \to (X_K)_{red}$. Par conséquent, il
existe un ouvert dense $V_K \subset X_K$ tel que $\nu_K^{-1}(V_K)
\to V_K$ soit une immersion fermée (induisant un isomorphisme de
$\nu_K^{-1}(V_K)$ avec $V_{K, red}$, voir par exemple morphismes,
Lemme
\ref{morphisms-lemma-birational-isomorphism-over-dense-open}). Après avoir
augmenté $k'$, nous trouvons que $V_K$ est le changement de base d'un ouvert
dense $V \subset Y$ et que le morphisme $\nu^{-1}(V) \to V$ est une
immersion fermée (Limites, Lemmata \ref{limits-lemma-descend-opens} et
\ref{limits-lemma-descend-closed-immersion-finite-presentation}). Il s'ensuit
facilement de cela que $\nu$ est le morphisme de normalisation et la preuve est complète.
\end{proof}

\begin{lemma}
\label{lemma-normalization-and-change-of-fields}
Supposons que $k$ soit un corps. Supposons que $X$ soit un $k$ -schéma
localement algébrique. Supposons que $K/k$ soit une extension de corps.
Supposons que $\nu : X^\nu \to X$ soit la normalisation de $X$ et soit $Y^\nu
\to X_K$ la normalisation du changement de base. Alors le morphisme canonique
$$
Y^\nu \longrightarrow X^\nu \times_{\Spec(k)} \Spec(K)
$$
est un isomorphisme si $K/k$ est séparable
et un homéomorphisme universel en général.
\end{lemma}

\begin{proof}
Soit $Y = X_K$. Supposons que $X^{(0)}$, resp. \ $Y^{(0)}$ soit
l'ensemble des points génériques des composantes irréductibles
de $X$, resp. \ $Y$. Alors le morphisme de projection $\pi : Y
\to X$ satisfait $\pi(Y^{(0)}) = X^{(0)}$. Cela est vrai car
$\pi$ est surjectif, ouvert et générisant, voir morphismes, Lemme
\ref{morphisms-lemma-scheme-over-field-universally-open} et
\ref{morphisms-lemma-open-generizing}. Si nous considérons
$X^{(0)}$, resp. \ $Y^{(0)}$ comme des schémas (réduits), alors
$X^\nu$, resp. \ $Y^\nu$ est la normalisation de $X$, resp. \ $Y$
dans $X^{(0)}$, resp. \ $Y^{(0})$. Ainsi le Lemme des morphismes
\ref{morphisms-lemma-functoriality-normalization} donne un morphisme
canonique $Y^\nu \to X^\nu$ sur $Y \to X$ qui à son tour donne le
morphisme canonique du lemme par la propriété universelle du produit fibré.

\medskip\noindent
Pour prouver que ce morphisme possède les propriétés énoncées dans le
lemme, nous pouvons supposer que $X = \Spec(A)$ est affine. Supposons
que $Q(A_{red})$ soit l'anneau total des fractions de $A_{red}$. Alors
$X^\nu$ est le spectre de l'adhérence intégrale $A'$ de $A$ dans
$Q(A_{red})$, voir morphismes, lemmes
\ref{morphisms-lemma-normalization-reduced} et
\ref{morphisms-lemma-description-normalization}. De même, $Y^\nu$ est le spectre de
l'adhérence intégrale $B'$ de $A \otimes_k K$ dans $Q((A \otimes_k K)_{red})$. Il
existe une application canonique $Q(A_{red}) \to Q((A \otimes_k K)_{red})$, une
application canonique $A' \to B'$, et le morphisme du lemme correspond à l'application induite.
$$
A' \otimes_k K \longrightarrow B'
$$
des $K$ -algèbres. Le noyau est constitué d'éléments
nilpotents car le noyau de $Q(A_{red}) \otimes_k K \to Q((A
\otimes_k K)_{red})$ est l'ensemble des éléments nilpotents.

\medskip\noindent
Si $K/k$ est séparable, alors $A' \otimes_k K$ est normal
d'après le Lemme \ref{lemma-base-change-normal-by-separable}.
En particulier, il est réduit, donc $Q((A \otimes_k K)_{red})
= Q(A' \otimes_k K)$ et $B' = A' \otimes_k K$ d'après
l'Algèbre, Lemme \ref{algebra-lemma-characterize-reduced-ring-normal}.

\medskip\noindent
Supposons que $K/k$ ne soit pas séparable. Alors la
caractéristique de $k$ est $p > 0$. Nous allons montrer que pour
chaque $b \in B'$ il existe une puissance $q$ de $p$ telle que
$b^q$ se trouve dans l'image de $A' \otimes_k K$. Cela prouvera que
l'application affichée est un homéomorphisme universel d'après le
Lemme d'Algèbre \ref{algebra-lemma-p-ring-map}. Pour un $b$ donné,
il existe un sous-corps $F \subset K$ avec $F/k$ de type fini tel
que $b$ soit contenu dans $Q((A \otimes_k F)_{red})$ et soit entier
sur $A \otimes_k F$. Choisissez un polynôme unitaire $P = T^d +
\alpha_1 T^{d - 1} + \ldots + \alpha _d$ avec $P(b) = 0$ et
$\alpha_i \in A \otimes_k F$. Choisissez une base de transcendance
$t_1, \ldots, t_r$ pour $F$ sur $k$. Supposons que $F/F'/k(t_1,
\ldots, t_r)$ soit la sous-extension séparable maximale (corps, Lemme
\ref{fields-lemma-separable-first}). Comme $F/F'$ est fini purement
inséparable, il existe un $q$ tel que $\lambda^q \in F'$ pour tout
$\lambda \in F$. Alors $b^q$ est dans $Q((A \otimes_k F')_{red})$ et
satisfait le polynôme $T^d + \alpha_1^q T^{d - 1} + \ldots + \alpha
_d^q$ avec $\alpha_i^q \in A \otimes_k F'$. Par le cas séparable, nous
voyons que $b^q \in A' \otimes_k F'$ et la démonstration est complète.
\end{proof}

\begin{lemma}
\label{lemma-geometrically-normal-in-codim-1}
Supposons que $k$ soit un corps. Supposons que $X$ soit un $k$-schéma
localement algébrique. Supposons que $\nu : X^\nu \to X$ soit la
normalisation de $X$. Supposons que $x \in X$ soit un point tel que
(a) $\mathcal{O}_{X, x}$ soit réduit, (b) $\dim(\mathcal{O}_{X, x}) =
1$, et (c) pour tout $x' \in X^\nu$ avec $\nu(x') = x$ l'extension
$\kappa(x')/k$ soit séparable. Alors $X$ est géométriquement réduit en
$x$ et $X^\nu$ est géométriquement régulier en $x'$ avec $\nu(x') = x$.
\end{lemma}

\begin{proof}
Nous utiliserons les résultats du Lemme
\ref{lemma-normalization-locally-algebraic} sans autre mention.
Supposons que $x' \in X^\nu$ soit un point au-dessus de $x$. Par
la théorie de la dimension (Section
\ref{section-algebraic-schemes}), nous avons
$\dim(\mathcal{O}_{X^\nu, x'}) = 1$. Puisque $X^\nu$ est normal,
nous voyons que $\mathcal{O}_{X^\nu, x'}$ est un anneau de valuation
discrète (Propriétés, Lemme \ref{properties-lemma-criterion-normal}).
Ainsi, $\mathcal{O}_{X^\nu, x'}$ est un $k$-local régulier algèbre
dont le corps résiduel est séparable sur $k$. Par conséquent, $k \to
\mathcal{O}_{X^\nu, x'}$ est formellement lisse dans la topologie
$\mathfrak m_{x'}$-adique, voir Plus sur l'Algèbre, Lemme
\ref{more-algebra-lemma-regular-implies-fs}. Alors $\mathcal{O}_{X^\nu,
x'}$ est géométriquement régulier sur $k$ par Plus sur l'Algèbre,
Théorème \ref{more-algebra-theorem-regular-fs}. Ainsi $X^\nu$ est
géométriquement régulier en $x'$ par Lemme \ref{lemma-geometrically-regular-at-point}.

\medskip\noindent
Puisque $\mathcal{O}_{X, x}$ est réduit, la famille de applications
$\mathcal{O}_{X, x} \to \mathcal{O}_{X^\nu, x'}$ est injective.
Puisque $\mathcal{O}_{X^\nu, x'}$ est une $k$-algèbre géométriquement
réduite, il s'ensuit immédiatement que $\mathcal{O}_{X, x}$ est une
$k$-algèbre géométriquement réduite. Par conséquent, $X$ est
géométriquement réduit en $x$ selon le Lemme \ref{lemma-geometrically-reduced-at-point}.
\end{proof}









\section{Groupes de fonctions inversibles}
\label{section-units}

\noindent
Il est souvent (mais pas toujours) le cas que
$\mathcal{O}^*(X)/k^*$ est un groupe abélien de type fini
si $X$ est une variété sur $k$. Nous le montrons par une
série de lemmes. Tout repose sur le cas spécial suivant.

\begin{lemma}
\label{lemma-open-in-normal-proper}
Supposons que $k$ soit un corps algébriquement clos.
Supposons que $\overline{X}$ soit une variété propre sur
$k$. Supposons que $X \subset \overline{X}$ soit un
sous-schéma ouvert. Supposons que $X$ soit normal. Alors
$\mathcal{O}^*(X)/k^*$ est un groupe abélien de type fini.
\end{lemma}

\begin{proof}
Puisque l'énoncé ne concerne que $X$, nous pouvons remplacer
$\overline{X}$ par une variété propre différente sur $k$.
Supposons que $\nu : \overline{X}^\nu \to \overline{X}$ soit le
morphisme de normalisation. D'après le Lemme
\ref{lemma-normalization-locally-algebraic}, nous avons que $\nu$ est
fini et $\overline{X}^\nu$ est une variété. Comme $X$ est normale,
nous voyons que $\nu^{-1}(X) \to X$ est un isomorphisme (petit détail
omis). Enfin, nous voyons que $\overline{X}^\nu$ est propre sur $k$,
car un morphisme fini est propre (morphismes, Lemme
\ref{morphisms-lemma-finite-proper}) et les compositions de morphismes
propres sont propres (morphismes, Lemme
\ref{morphisms-lemma-composition-proper}). Ainsi, nous pouvons et faisons l'hypothèse que $\overline{X}$ est normal.

\medskip\noindent
Nous utiliserons sans autre mention que pour tout ouvert
affine $U$ de $\overline{X}$, l'anneau $\mathcal{O}(U)$ est
une $k$-algèbre de type fini, qui est noethérienne, intègre et
normale, voir Algèbre, Lemme
\ref{algebra-lemma-Noetherian-permanence}, Propriétés, Définition
\ref{properties-definition-integral}, Propriétés, Lemmata
\ref{properties-lemma-locally-Noetherian} et
\ref{properties-lemma-locally-normal}, morphismes, Lemme \ref{morphisms-lemma-locally-finite-type-characterize}.

\medskip\noindent
Supposons que $\xi_1, \ldots, \xi_r$ soit les points génériques du
complément de $X$ dans $\overline{X}$. Ils sont en nombre fini
puisque $\overline{X}$ possède un espace topologique sous-jacent
noethérien (voir morphismes, Lemme
\ref{morphisms-lemma-finite-type-noetherian}, Propriétés, Lemme
\ref{properties-lemma-Noetherian-topology}, et Topologie, Lemme
\ref{topology-lemma-Noetherian}). Pour chaque $i$, l’anneau local
$\mathcal{O}_i = \mathcal{O}_{X, \xi_i}$ est un domaine local normal
noethérien (comme une localisation d’un domaine normal noethérien).
Supposons que $J \subset \{1, \ldots, r\}$ soit l’ensemble des indices $i$
tels que $\dim(\mathcal{O}_i) = 1$. Pour $j \in J$, l’anneau local
$\mathcal{O}_j$ est un anneau de valuation discrète, voir Algèbre, Lemme
\ref{algebra-lemma-characterize-dvr}. Par conséquent, nous obtenons une valuation.
$$
v_j : k(\overline{X})^* \longrightarrow \mathbf{Z}
$$
avec la propriété que $v_j(f) \geq 0 \Leftrightarrow f \in \mathcal{O}_j$.

\medskip\noindent
Considérez $\mathcal{O}(X)$ comme une sous-$k$ -algèbre de $k(X)
= k(\overline{X})$. Nous affirmons que le noyau de l'application
$$
\mathcal{O}(X)^* \longrightarrow
\prod\nolimits_{j \in J} \mathbf{Z},
\quad
f \longmapsto \prod v_j(f)
$$
est $k^*$. Il est clair que cette affirmation prouve le
lemme. À savoir, supposons que $f \in \mathcal{O}(X)$ soit
un élément du noyau. Supposons que $U = \Spec(B) \subset
\overline{X}$ soit un ouvert affine quelconque. Alors $B$
est un domaine normal noethérien. Pour chaque idéal
premier de hauteur un $\mathfrak q \subset B$ avec le
point correspondant $\xi \in X$, nous voyons que soit $\xi
= \xi_j$ pour un certain $j \in J$, soit que $\xi \in X$.
La raison est que $\text{codim}(\overline{\{\xi\}},
\overline{X}) = 1$ d'après les Propriétés, le Lemme
\ref{properties-lemma-codimension-local-ring} et donc si
$\xi \in \overline{X} \setminus X$, cela doit être un point
générique de $\overline{X} \setminus X$, donc égal à un
certain $\xi_j$, $j \in J$. Nous concluons que $f \in
\mathcal{O}_{X, \xi} = B_{\mathfrak q}$ dans les deux cas
puisque $f$ est dans le noyau de l'application. Ainsi $f \in
\bigcap_{\text{ht}(\mathfrak q) = 1} B_{\mathfrak q} = B$, voir
Algèbre, Lemme
\ref{algebra-lemma-normal-domain-intersection-localizations-height-1}. En
d'autres termes, nous voyons que $f \in \Gamma(\overline{X},
\mathcal{O}_{\overline{X}})$. Mais puisque $k$ est algébriquement clos, nous
concluons que $f \in k$ d'après le Lemme \ref{lemma-regular-functions-proper-variety}.
\end{proof}

\noindent
Ensuite, nous généralisons le cas ci-dessus par quelques arguments
élémentaires, en conservant toujours le corps algébriquement clos.

\begin{lemma}
\label{lemma-units-integral-finite-type-algebraically-closed}
Supposons que $k$ est un corps algébriquement clos. Supposons
que $X$ est un schéma intègre localement de type fini sur $k$.
Alors $\mathcal{O}^*(X)/k^*$ est un groupe abélien de type fini.
\end{lemma}

\begin{proof}
Comme $X$ est intègre, l'application de restriction
$\mathcal{O}(X) \to \mathcal{O}(U)$ est injective pour
tout sous-schéma ouvert non vide $U \subset X$. Par
conséquent, nous pouvons supposer que $X$ est affine.
Choisissez une immersion fermée $X \to \mathbf{A}^n_k$
et notez $\overline{X}$ l'adhérence de $X$ dans
$\mathbf{P}^n_k$ via l'immersion usuelle $\mathbf{A}^n_k
\to \mathbf{P}^n_k$. Ainsi, nous pouvons supposer que $X$
est un ouvert affine d'une variété projective $\overline{X}$.

\medskip\noindent
Supposons que $\nu : \overline{X}^\nu \to \overline{X}$ soit
le morphisme de normalisation, voir morphismes, Définition
\ref{morphisms-definition-normalization}. Nous savons que
$\nu$ est fini, dominant, et que $\overline{X}^\nu$ est un
schéma irréductible normal, voir morphismes, Lemmata
\ref{morphisms-lemma-normalization-normal},
\ref{morphisms-lemma-Japanese-normalization}, et
\ref{morphisms-lemma-ubiquity-nagata}. Il s'ensuit que
$\overline{X}^\nu$ est une variété propre, parce que
$\overline{X} \to \Spec(k)$ est propre comme composition d'un
morphisme fini et d'un morphisme propre (voir les résultats dans
morphismes, Sections \ref{morphisms-section-proper} et
\ref{morphisms-section-integral}). Il s'ensuit également que $\nu$
est un morphisme surjectif, parce que l'image de $\nu$ est fermée
et contient le point générique de $\overline{X}$. Ainsi, en posant
$X^\nu = \nu^{-1}(X)$, nous voyons qu'il suffit de prouver le
résultat pour $X^\nu$. En d'autres termes, nous pouvons supposer que
$X$ est un ouvert non vide d'une variété propre normale
$\overline{X}$. Ce cas est traité par le Lemme \ref{lemma-open-in-normal-proper}.
\end{proof}

\noindent
Le lemme précédent implique la générisation légère suivante.

\begin{lemma}
\label{lemma-units-general-algebraically-closed}
Supposons que $k$ soit un corps algébriquement clos. Supposons
que $X$ soit un schéma réduit connexe qui est localement de type
fini sur $k$ avec un nombre fini de composantes irréductibles.
Alors $\mathcal{O}^*(X)/k^*$ est un groupe abélien de type fini.
\end{lemma}

\begin{proof}
Supposons que $X = \bigcup X_i$ soit les composantes
irréductibles. Par le Lemme
\ref{lemma-units-integral-finite-type-algebraically-closed}, nous
voyons que $\mathcal{O}(X_i)^*/k^*$ est un groupe abélien de type fini.
Supposons que $f \in \mathcal{O}(X)^*$ soit dans le noyau de l'application
$$
\mathcal{O}(X)^* \longrightarrow \prod \mathcal{O}(X_i)^*/k^*.
$$
Alors pour chaque $i$ il existe un élément $\lambda_i
\in k$ tel que $f|_{X_i} = \lambda_i$. En se
restreignant à $X_i \cap X_j$, nous concluons que
$\lambda_i = \lambda_j$ si $X_i \cap X_j \not = \emptyset$.
Puisque $X$ est connecté, nous concluons que tous les
$\lambda_i$ sont d'accord et donc que $f \in k^*$. Cela prouve que
$$
\mathcal{O}(X)^*/k^* \subset \prod \mathcal{O}(X_i)^*/k^*
$$
et le lemme en découle car à droite nous avons un
produit d’un nombre fini de groupes abéliens de type fini.
\end{proof}

\begin{lemma}
\label{lemma-integral-closure-ground-field}
Supposons que $k$ est un corps. Supposons que
$X$ est un schéma sur $k$ qui est connexe et
réduit. Alors l'adhérence intégrale de $k$
dans $\Gamma(X, \mathcal{O}_X)$ est un corps.
\end{lemma}

\begin{proof}
Supposons que $k' \subset \Gamma(X, \mathcal{O}_X)$ soit
l'adhérence intégrale de $k$. Alors $X \to \Spec(k)$ se
factorise à travers $\Spec(k')$, voir schémas, Lemme
\ref{schemes-lemma-morphism-into-affine}. Comme $X$ est réduit,
nous voyons que $k'$ n'a pas d'éléments nilpotents non nuls.
Comme $k \to k'$ est intègre, nous voyons que tout idéal
premier de $k'$ est à la fois un idéal maximal et un premier
minimal, et $\Spec(k')$ est totalement déconnecté, voir Algèbre,
Lemmata \ref{algebra-lemma-integral-no-inclusion} et
\ref{algebra-lemma-ring-with-only-minimal-primes}. Comme $X$ est
connexe, le morphisme $X \to \Spec(k')$ est constant, disons avec
pour image le point correspondant à $\mathfrak p \subset k'$.
Alors tout $f \in k'$, $f \not \in \mathfrak p$ se mappe sur un
élément inversible de $\mathcal{O}_X$. Par définition de $k'$,
cela force alors $f$ à être une unité de $k'$. Ainsi nous voyons
que $k'$ est local avec idéal maximal $\mathfrak p$, voir Algèbre,
Lemme \ref{algebra-lemma-characterize-local-ring}. Puisque nous
avons déjà vu que $k'$ est réduit, cela implique que $k'$ est un
corps, voir Algèbre, Lemme \ref{algebra-lemma-minimal-prime-reduced-ring}.
\end{proof}

\begin{proposition}
\label{proposition-units-general}
Supposons que $k$ soit un corps. Supposons que $X$ soit un schéma sur
$k$. Supposons que $X$ soit localement de type fini sur $k$, connexe,
réduit, et possède un nombre fini de composantes irréductibles. Alors
$\mathcal{O}(X)^*/k^*$ est un groupe abélien de type fini si, en plus des
conditions ci-dessus, au moins l'une des conditions suivantes est satisfaite :
\begin{enumerate}
\item l'adhérence intégrale de $k$ dans $\Gamma(X,
\mathcal{O}_X)$ est $k$, \item $X$ a un point
$k$-rationnel, ou \item $X$ est géométriquement intègre.
\end{enumerate}
\end{proposition}

\begin{proof}
Supposons que $\overline{k}$ soit une clôture algébrique de $k$.
Supposons que $Y$ soit une composante connexe de
$(X_{\overline{k}})_{red}$. Notez que le morphisme canonique $p : Y
\to X$ est ouvert (d'après morphismes, Lemme
\ref{morphisms-lemma-scheme-over-field-universally-open}) et fermé
(d'après morphismes, Lemme
\ref{morphisms-lemma-integral-universally-closed}). Ainsi, $p(Y) = X$ étant
connecté comme supposé pour $X$. En particulier, comme $X$ est réduit, cela
implique $\mathcal{O}(X) \subset \mathcal{O}(Y)$. D'après le Lemme
\ref{lemma-galois-action-irreducible-components-locally-finite-type}, nous
voyons que $Y$ possède un nombre fini de composantes irréductibles (le
corps $\overline{k}$ dans le lemme est la clôture algébrique séparable, mais
cela n'a pas d'importance car les deux changements de base
$X_{\overline{k}}$ ont les mêmes composantes irréductibles d'après le Lemme
\ref{lemma-separably-closed-field-irreducible-components}). Ainsi, le Lemme
\ref{lemma-units-general-algebraically-closed} s'applique à $Y$. Cela implique
que si $\mathcal{O}(X)^*/k^*$ n'est pas un groupe abélien de type fini, alors
il existe des éléments $f \in \mathcal{O}(X)$ et $f \not \in k$ qui se mappent
sur un élément de $\overline{k}$ via l'application $\mathcal{O}(X) \to
\mathcal{O}(Y)$. Dans ce cas, $f$ est algébrique sur $k$, donc intégral sur $k$.
Ainsi, si la condition (1) est vérifiée, cela ne peut pas se produire. Pour
terminer la démonstration, nous montrons que les conditions (2) et (3) impliquent (1).

\medskip\noindent
Supposons que $k \subset k' \subset \Gamma(X, \mathcal{O}_X)$
est l'adhérence intégrale de $k$ dans $\Gamma(X,
\mathcal{O}_X)$. D'après le lemme
\ref{lemma-integral-closure-ground-field}, nous voyons que $k'$ est
un corps. Si $e : \Spec(k) \to X$ est un point $k$-rationnel, alors
$e^\sharp : \Gamma(X, \mathcal{O}_X) \to k$ est une section de
l'application d'inclusion $k \to \Gamma(X, \mathcal{O}_X)$. En
particulier, la restriction de $e^\sharp$ à $k'$ est une application de
corps $k' \to k$ sur $k$, ce qui montre clairement que (2) implique (1).

\medskip\noindent
Si l'adhérence intégrale $k'$ de $k$ dans $\Gamma(X, \mathcal{O}_X)$
n'est pas triviale, alors nous voyons que $X$ n'est soit pas
géométriquement connexe (si $k'/k$ n'est pas purement inséparable),
soit que $X$ n'est pas géométriquement réduit (si $k'/k$ est purement
inséparable non trivial). Détails omis. Par conséquent, (3) implique (1).
\end{proof}

\begin{lemma}
\label{lemma-units-variety}
Supposons que $k$ soit un corps. Supposons que $X$
soit une variété sur $k$. Le groupe
$\mathcal{O}(X)^*/k^*$ est un groupe abélien de type fini à
condition qu'au moins une des conditions suivantes soit remplie :
\begin{enumerate}
\item $k$ est intégralement clos dans $\Gamma(X,
\mathcal{O}_X)$, \item $k$ est algébriquement clos dans
$k(X)$, \item $X$ est géométriquement intègre sur $k$, ou
\item $k$ est l'intersection `` '' des extensions de
corps $\kappa(x)/k$ où $x$ parcourt les points fermés de $x$.
\end{enumerate}
\end{lemma}

\begin{proof}
Nous voyons que (1) suffit d'après la Proposition
\ref{proposition-units-general}. Nous omettons la
vérification que chacun de (2), (3), (4) implique (1).
\end{proof}







\section{Formule de K\"unneth, I}
\label{section-kunneth}

\noindent
Dans cette section, nous prouvons la formule de K \" unneth lorsque
la base est un corps et que nous considérons la cohomologie des
modules quasi-cohérents. Pour une version plus générale, veuillez
consulter Catégories Dérivées de schémas, Section \ref{perfect-section-kunneth}.

\begin{lemma}
\label{lemma-kunneth}
Supposons que $k$ est un corps. Supposons que $X$ et $Y$ sont des
schémas sur $k$ et soient $\mathcal{F}$, respectivement \
$\mathcal{G}$, un $\mathcal{O}_X$-module quasi-cohérent,
respectivement \ $\mathcal{O}_Y$-module. Alors nous avons un isomorphisme canonique
$$
H^n(X \times_{\Spec(k)} Y, \text{pr}_1^*\mathcal{F}
\otimes_{\mathcal{O}_{X \times_{\Spec(k)} Y}} \text{pr}_2^*\mathcal{G}) =
\bigoplus\nolimits_{p + q = n}
H^p(X, \mathcal{F}) \otimes_k H^q(Y, \mathcal{G})
$$
à condition que $X$ et $Y$ soient quasi-compacts et aient un
diagonale affine \footnote { Le cas où $X$ et $Y$ sont
quasi-séparés sera discuté dans le Lemme
\ref{lemma-kunneth-general} ci-dessous. } (par exemple si $X$ et $Y$ sont séparés).
\end{lemma}

\begin{proof}
Dans cette preuve, les produits simples et les produits tensoriels sont sur $k$. En tant que applications
$$
H^p(X, \mathcal{F}) \otimes H^q(Y, \mathcal{G})
\longrightarrow
H^n(X \times Y, \text{pr}_1^*\mathcal{F}
\otimes_{\mathcal{O}_{X \times Y}} \text{pr}_2^*\mathcal{G})
$$
nous utilisons la fonctorialité de la cohomologie pour obtenir
les applications $H^p(X, \mathcal{F}) \to H^p(X \times Y,
\text{pr}_1^*\mathcal{F})$ et $H^p(Y, \mathcal{G}) \to H^p(X \times Y,
\text{pr}_2^*\mathcal{G})$, puis nous utilisons le produit en coupure
$$
\cup :
H^p(X \times Y, \text{pr}_1^*\mathcal{F})
\otimes H^q(X \times Y, \text{pr}_2^*\mathcal{G})
\longrightarrow
H^n(X \times Y, \text{pr}_1^*\mathcal{F}
\otimes_{\mathcal{O}_{X \times Y}} \text{pr}_2^*\mathcal{G})
$$
Le résultat est vrai lorsque $X$ et $Y$ sont affines en raison de la
nullité des groupes de cohomologie supérieure sur les affines
(Cohomologie des schémas, Lemme
\ref{coherent-lemma-quasi-coherent-affine-cohomology-zero}) et des définitions (des images
réciproques de modules quasi-cohérents et des produits tensoriels de modules quasi-cohérents).

\medskip\noindent
Choisissez des recouvrements affines ouverts finis
$\mathcal{U} : X = \bigcup_{i \in I} U_i$ et
$\mathcal{V} : Y = \bigcup_{j \in J} V_j$. Cela
détermine un recouvrement affine ouvert $\mathcal{W} : X
\times Y = \bigcup_{(i, j) \in I \times J} U_i \times
V_j$. Notez que $\mathcal{W}$ est un raffinement de
$\text{pr}_1^{-1}\mathcal{U}$ et de
$\text{pr}_2^{-1}\mathcal{V}$. Ainsi, par le Lemme de Cohomologie
\ref{cohomology-lemma-functoriality-cech}, nous obtenons des applications
$$
\check{\mathcal{C}}^\bullet(\mathcal{U}, \mathcal{F}) \to
\check{\mathcal{C}}^\bullet(\mathcal{W}, \text{pr}_1^*\mathcal{F})
\quad\text{and}\quad
\check{\mathcal{C}}^\bullet(\mathcal{V}, \mathcal{G}) \to
\check{\mathcal{C}}^\bullet(\mathcal{W}, \text{pr}_2^*\mathcal{G})
$$
compatible avec les applications de tirage en arrière sur la
cohomologie. En cohomologie, dans l'équation (
\ref{cohomology-equation-needs-signs} ) nous avons construit une application de complexes
$$
\text{Tot}(
\check{\mathcal{C}}^\bullet(\mathcal{W}, \text{pr}_1^*\mathcal{F})
\otimes
\check{\mathcal{C}}^\bullet(\mathcal{W}, \text{pr}_2^*\mathcal{G}))
\longrightarrow
\check{\mathcal{C}}^\bullet(\mathcal{W},
\text{pr}_1^*\mathcal{F} \otimes_{\mathcal{O}_{X \times Y}}
\text{pr}_2^*\mathcal{G})
$$
définir le produit en coupe en cohomologie. En combinant
ce qui précède, nous obtenons une application de complexes
\begin{equation}
\label{equation-kunneth-on-cech}
\text{Tot}(
\check{\mathcal{C}}^\bullet(\mathcal{U}, \mathcal{F})
\otimes
\check{\mathcal{C}}^\bullet(\mathcal{V}, \mathcal{G}))
\longrightarrow
\check{\mathcal{C}}^\bullet(\mathcal{W},
\text{pr}_1^*\mathcal{F} \otimes_{\mathcal{O}_{X \times Y}}
\text{pr}_2^*\mathcal{G})
\end{equation}
Nous avertissons le lecteur que cette application
n'est pas un isomorphisme de complexes. Rappelons que
nous pouvons calculer les cohomologies de nos
faisceaux quasi-cohérents en utilisant nos
recouvrements (Cohomologie des schémas, Lemmata
\ref{coherent-lemma-affine-diagonal} et
\ref{coherent-lemma-cech-cohomology-quasi-coherent}). Ainsi, sur la
cohomologie (\ref{equation-kunneth-on-cech}) reproduit l'application du lemme.

\medskip\noindent
Considérons une suite exacte courte $0 \to \mathcal{F}
\to \mathcal{F}' \to \mathcal{F}'' \to 0$ de modules
quasi-cohérents. Puisque la construction de (
\ref{equation-kunneth-on-cech} ) est fonctorielle en
$\mathcal{F}$ et puisque la formation des complexes { \v C }
pertinents est exacte dans la variable $\mathcal{F}$ (car nous
prenons des sections sur des ouverts affines), nous obtenons
une application entre des suites exactes courtes de complexes.
$$
\xymatrix{
\text{Tot}(
\check{\mathcal{C}}^\bullet(\mathcal{U}, \mathcal{F})
\otimes
\check{\mathcal{C}}^\bullet(\mathcal{V}, \mathcal{G})) \ar[r] \ar[d] &
\text{Tot}(
\check{\mathcal{C}}^\bullet(\mathcal{U}, \mathcal{F}')
\otimes
\check{\mathcal{C}}^\bullet(\mathcal{V}, \mathcal{G})) \ar[r] \ar[d] &
\text{Tot}(
\check{\mathcal{C}}^\bullet(\mathcal{U}, \mathcal{F}'')
\otimes
\check{\mathcal{C}}^\bullet(\mathcal{V}, \mathcal{G})) \ar[d] \\
\check{\mathcal{C}}^\bullet(\mathcal{W},
\text{pr}_1^*\mathcal{F} \otimes_{\mathcal{O}_{X \times Y}}
\text{pr}_2^*\mathcal{G}) \ar[r] &
\check{\mathcal{C}}^\bullet(\mathcal{W},
\text{pr}_1^*\mathcal{F}' \otimes_{\mathcal{O}_{X \times Y}}
\text{pr}_2^*\mathcal{G}) \ar[r] &
\check{\mathcal{C}}^\bullet(\mathcal{W},
\text{pr}_1^*\mathcal{F}'' \otimes_{\mathcal{O}_{X \times Y}}
\text{pr}_2^*\mathcal{G})
}
$$
(nous avons supprimé les zéros extérieurs). En
regardant les longues suites exactes de
cohomologie, nous constatons que si le résultat du
lemme vaut pour $2$ -hors de- $3$ de $\mathcal{F},
\mathcal{F}', \mathcal{F}''$, alors il vaut pour le troisième.

\medskip\noindent
Remarquez que $X$ a une dimension cohomologique
finie pour les modules quasi-cohérents, voir
Cohomologie des schémas, Lemme
\ref{coherent-lemma-vanishing-nr-affines}. En
utilisant l'induction sur $d(\mathcal{F}) = \max \{d
\mid H^d(X, \mathcal{F}) \not = 0\}$, nous allons
réduire au cas $d(\mathcal{F}) = 0$. Supposons
$d(\mathcal{F}) > 0$. Selon Cohomologie des schémas,
Lemme
\ref{coherent-lemma-affine-diagonal-universal-delta-functor},
nous avons vu qu'il existe un plongement $\mathcal{F} \to
\mathcal{F}'$ tel que $H^p(X, \mathcal{F}') = 0$ pour tout $p
\geq 1$. En posant $\mathcal{F}'' = \Coker(\mathcal{F} \to
\mathcal{F}')$, nous voyons que $d(\mathcal{F}'') <
d(\mathcal{F})$. Ensuite, nous pouvons appliquer le résultat du
paragraphe précédent pour voir qu'il suffit de prouver le lemme pour
$\mathcal{F}'$ et $\mathcal{F}''$, prouvant ainsi l'étape d'induction.

\medskip\noindent
En procédant de la même manière pour $\mathcal{G}$, nous
constatons que nous pouvons supposer que $\mathcal{F}$ et
$\mathcal{G}$ n'ont une cohomologie non nulle qu'au degré
$0$. Supposons que $V \subset Y$ soit un ouvert affine.
Considérons le recouvrement affine ouvert $\mathcal{U}_V : X
\times V = \bigcup_{i \in I} U_i \times V$. Il est immédiat que
$$
\check{\mathcal{C}}^\bullet(\mathcal{U}, \mathcal{F}) \otimes \mathcal{G}(V) =
\check{\mathcal{C}}^\bullet(\mathcal{U}_V,
\text{pr}_1^*\mathcal{F} \otimes_{\mathcal{O}_{X \times Y}}
\text{pr}_2^*\mathcal{G})
$$
(égalité des complexes). Nous concluons que
$$
R\text{pr}_{2, *}(\text{pr}_1^*\mathcal{F} \otimes_{\mathcal{O}_{X \times Y}}
\text{pr}_2^*\mathcal{G})
\cong
\Gamma(X, \mathcal{F}) \otimes_k \mathcal{G} \cong
\bigoplus\nolimits_{\alpha \in A} \mathcal{G}
$$
sur $Y$. Ici $A$ est une base de l’espace vectoriel $k$
sur $\Gamma(X, \mathcal{F})$. La cohomologie sur $Y$
commute avec les sommes directes (Cohomologie, Lemme
\ref{cohomology-lemma-quasi-separated-cohomology-colimit}).
En utilisant la suite spectrale de Leray pour
$\text{pr}_2$ (via Cohomologie, Lemme
\ref{cohomology-lemma-apply-Leray}), nous concluons que
$H^n(X \times Y, \text{pr}_1^*\mathcal{F}
\otimes_{\mathcal{O}_{X \times Y}} \text{pr}_2^*\mathcal{G})$
est nul pour $n > 0$ et isomorphe à $H^0(X, \mathcal{F}) \otimes
H^0(Y, \mathcal{G})$ pour $n = 0$. Cela termine la preuve (sauf
que nous devrions vérifier que l’isomorphisme est bien donné par
le produit emboîté en degré $0$ ; nous omettons la vérification).
\end{proof}

\begin{lemma}
\label{lemma-kunneth-general}
Supposons que $k$ est un corps. Supposons que $X$ et $Y$ sont des
schémas sur $k$ et soient $\mathcal{F}$, respectivement \
$\mathcal{G}$, un $\mathcal{O}_X$-module quasi-cohérent,
respectivement \ $\mathcal{O}_Y$-module. Alors nous avons un isomorphisme canonique
$$
H^n(X \times_{\Spec(k)} Y, \text{pr}_1^*\mathcal{F}
\otimes_{\mathcal{O}_{X \times_{\Spec(k)} Y}} \text{pr}_2^*\mathcal{G}) =
\bigoplus\nolimits_{p + q = n}
H^p(X, \mathcal{F}) \otimes_k H^q(Y, \mathcal{G})
$$
à condition que $X$ et $Y$ soient quasi-compacts et quasi-séparés.
\end{lemma}

\begin{proof}
Si $X$ et $Y$ sont séparés ou, plus généralement, ont une
diagonale affine, veuillez alors consulter le Lemme
\ref{lemma-kunneth} pour une meilleure preuve `` '' (la
caractéristique qu'il a par rapport à cette preuve est qu'il identifie
les applications comme des fibres produites suivies de produits cup).
Supposons que $X'$, respectivement \ $Y'$, soit l'épaississement
infinitésimal de $X$, respectivement \ $Y$, dont le faisceau de
structure est $\mathcal{O}_{X'} = \mathcal{O}_X \oplus \mathcal{F}$,
respectivement \ $\mathcal{O}_{Y'} = \mathcal{O}_Y \oplus \mathcal{G}$, où
$\mathcal{F}$, respectivement \ $\mathcal{G}$, est un idéal de carré nul. Alors
$$
\mathcal{O}_{X' \times Y'} =
\mathcal{O}_{X \times Y}
\oplus \text{pr}_1^*\mathcal{F}
\oplus \text{pr}_2^*\mathcal{G}
\oplus \text{pr}_1^*\mathcal{F}
\otimes_{\mathcal{O}_{X \times Y}} \text{pr}_2^*\mathcal{G}
$$
comme les faisceaux sur $X \times Y$. De cette
manière, nous voyons qu'il suffit de prouver que
$$
H^n(X \times Y, \mathcal{O}_{X \times Y}) =
\bigoplus\nolimits_{p + q = n}
H^p(X, \mathcal{O}_X) \otimes_k H^q(Y, \mathcal{O}_Y)
$$
pour toute paire de schémas quasi-compacts et
quasi-séparés sur $k$. Certains détails sont omis.

\medskip\noindent
Pour prouver cette affirmation, nous utilisons la cohomologie et le
changement de base sous la forme de Cohomologie des schémas, Lemme
\ref{coherent-lemma-hypercoverings}. Ce lemme nous dit qu'il existe un complexe
borné inférieurement de espaces vectoriels $k$, c'est-à-dire un complexe
$\mathcal{K}^\bullet$ de modules quasi-cohérents sur $\Spec(k)$, qui calcule
universellement la cohomologie de $Y$ sur $\Spec(k)$. En particulier, nous voyons que
$$
R\text{pr}_{1, *}(\mathcal{O}_{X \times Y}) \cong
(X \to \Spec(k))^*\mathcal{K}^\bullet
$$
dans $D(\mathcal{O}_X)$. À homotopie près, le complexe
$\mathcal{K}^\bullet$ est isomorphe à $\bigoplus_{q \geq 0}
H^q(Y, \mathcal{O}_Y)[-q]$ car cela est vrai pour tout
complexe d'espaces vectoriels sur un corps. Nous concluons que
$$
R\text{pr}_{1, *}(\mathcal{O}_{X \times Y}) \cong
\bigoplus\nolimits_{q \geq 0}
H^q(Y, \mathcal{O}_Y)[-q] \otimes_k \mathcal{O}_X
$$
dans $D(\mathcal{O}_X)$. Puis nous avons
\begin{align*}
R\Gamma(X \times Y, \mathcal{O}_{X \times Y})
& =
R\Gamma(X, R\text{pr}_{1, *}(\mathcal{O}_{X \times Y})) \\
& =
R\Gamma(X, \bigoplus\nolimits_{q \geq 0}
H^q(Y, \mathcal{O}_Y)[-q] \otimes_k \mathcal{O}_X) \\
& =
\bigoplus\nolimits_{q \geq 0}
R\Gamma(X,  H^q(Y, \mathcal{O}_Y) \otimes \mathcal{O}_X)[-q] \\
& =
\bigoplus\nolimits_{q \geq 0}
R\Gamma(X, \mathcal{O}_X) \otimes_k H^q(Y, \mathcal{O}_Y)[-q] \\
& =
\bigoplus\nolimits_{p, q \geq 0}
H^p(X, \mathcal{O}_X)[-p] \otimes_k H^q(Y, \mathcal{O}_Y)[-q]
\end{align*}
comme désiré. La première égalité par Leray pour
$\text{pr}_1$ (Cohomologie, Lemme
\ref{cohomology-lemma-before-Leray} ). La deuxième par notre
décomposition de l'image directe totale donnée ci-dessus. La
troisième parce que la cohomologie commute toujours avec les
sommes directes finies (et la cohomologie de $Y$ s'annule en
degré suffisamment grand d'après Cohomologie des schémas, Lemme
\ref{coherent-lemma-vanishing-nr-affines-quasi-separated} ). La
quatrième parce que la cohomologie sur $X$ commute avec les
sommes directes infinies d'après Cohomologie, Lemme
\ref{cohomology-lemma-quasi-separated-cohomology-colimit}. L'égalité
finale par notre remarque sur la catégorie dérivée d'un corps ci-dessus.
\end{proof}








\section{Groupes de Picard des variétés}
\label{section-picard-groups}

\noindent
Dans cette section, nous recueillons quelques résultats
élémentaires sur les groupes de Picard des variétés algébriques.

\begin{lemma}
\label{lemma-change-rings-pic-pre}
Soit $A \to B$ un homomorphisme d'anneaux fidèlement plat. Soit
$X$ un schéma quasi-compact et quasi-séparé sur $A$. Soit
$\mathcal{L}$ un $\mathcal{O}_X$-module inversible dont l'image
inverse sur $X_B$ est triviale. Alors $H^0(X, \mathcal{L})$ et $H^0(X,
\mathcal{L}^{\otimes -1})$ sont des $H^0(X, \mathcal{O}_X)$-modules
inversibles, et l'application de multiplication induit un isomorphisme
$$
H^0(X, \mathcal{L}) \otimes_{H^0(X, \mathcal{O}_X)}
H^0(X, \mathcal{L}^{\otimes -1}) \longrightarrow
H^0(X, \mathcal{O}_X)
$$
\end{lemma}

\begin{proof}
Notons $\mathcal{L}_B$ le retracement de $\mathcal{L}$ à
$X_B$. Choisissez un isomorphisme $\mathcal{L}_B \to
\mathcal{O}_{X_B}$. Posons $R = H^0(X, \mathcal{O}_X)$, $M =
H^0(X, \mathcal{L})$ et considérons $M$ comme un $R$-module.
Pour tout $\mathcal{O}_X$-module quasi-cohérent $\mathcal{F}$
avec le retracement $\mathcal{F}_B$ sur $X_B$, il existe un
isomorphisme canonique $H^0(X_B, \mathcal{F}_B) = H^0(X,
\mathcal{F}) \otimes_A B$, voir Cohomologie des schémas, Lemme
\ref{coherent-lemma-flat-base-change-cohomology}. Ainsi, nous avons
$$
M \otimes_R (R \otimes_A B) =
M \otimes_A B = H^0(X_B, \mathcal{L}_B) \cong
H^0(X_B, \mathcal{O}_{X_B}) = R \otimes_A B
$$
Puisque $R \to R \otimes_A B$ est fidèlement plat (en tant que
changement de base de l'application fidèlement plate $A \to B$), nous
concluons que $M$ est un $R$-module inversible d'après Algèbre,
Proposition \ref{algebra-proposition-ffdescent-finite-projectivity}.
De même, $N = H^0(X, \mathcal{L}^{\otimes -1})$ est un $R$-module
inversible. Pour voir que l'énoncé sur les produits tensoriels est vrai,
utilisez le fait qu'il est vrai après le tirage en arrière vers $X_B$ et
la fidélité plate de $R \to R \otimes_A B$. Quelques détails sont omis.
\end{proof}

\begin{lemma}
\label{lemma-change-rings-pic}
Supposons que $A \to B$ soit un homomorphisme d'anneaux
fidèlement plat. Supposons que $X$ soit un schéma sur $A$ tel que
\begin{enumerate}
\item $X$ est quasi-compact et quasi-séparé, et \item
$R = H^0(X, \mathcal{O}_X)$ est un semi-anneau local.
\end{enumerate}
Alors l'application de recul $\Pic(X) \to \Pic(X_B)$ est injective.
\end{lemma}

\begin{proof}
Supposons que $\mathcal{L}$ soit un $\mathcal{O}_X$-module
inversible dont le tiré en arrière $\mathcal{L}'$ vers $X_B$ est
trivial. Posons $M = H^0(X, \mathcal{L})$ et $N = H^0(X,
\mathcal{L}^{\otimes - 1})$. D'après le Lemme
\ref{lemma-change-rings-pic-pre}, les $R$-modules $M$ et $N$ sont
inversibles. Puisque $R$ est semi-local, $M \cong R$ et $N \cong R$,
voir Algèbre, Lemme \ref{algebra-lemma-locally-free-semi-local-free}.
Choisissez des générateurs $s \in M$ et $t \in N$. Alors $st \in R =
H^0(X, \mathcal{O}_X)$ est une unité selon la dernière partie du Lemme
\ref{lemma-change-rings-pic-pre}. Nous concluons que $s$ et $t$ définissent
des trivializations de $\mathcal{L}$ et $\mathcal{L}^{\otimes -1}$ sur $X$.
\end{proof}

\begin{lemma}
\label{lemma-change-fields-pic}
Supposons que $k'/k$ soit une extension de corps.
Supposons que $X$ soit un schéma sur $k$ tel que
\begin{enumerate}
\item $X$ est quasi-compact et quasi-séparé, et \item $R = H^0(X,
\mathcal{O}_X)$ est semi-local, par exemple, si $\dim_k R < \infty$.
\end{enumerate}
Alors l'application de recul $\Pic(X) \to \Pic(X_{k'})$ est injective.
\end{lemma}

\begin{proof}
Cas particulier du Lemme
\ref{lemma-change-rings-pic}. Si $\dim_k R <
\infty$, alors $R$ est artinien et donc semi-local
(Algèbre, Lemmata
\ref{algebra-lemma-finite-dimensional-algebra} et \ref{algebra-lemma-artinian-finite-nr-max}).
\end{proof}

\begin{example}
\label{example-need-condition}
Le lemme \ref{lemma-change-fields-pic} n'est pas vrai sans une certaine
condition sur le schéma $X$ sur le corps $k$. Voici un exemple. Supposons
que $k$ est un corps. Supposons que $t \in \mathbf{P}^1_k$ est un point
fermé. Posons $X = \mathbf{P}^1 \setminus \{t\}$. Alors nous avons une surjection
$$
\mathbf{Z} = \Pic(\mathbf{P}^1_k) \longrightarrow \Pic(X)
$$
La première égalité par les Diviseurs, Lemme
\ref{divisors-lemma-Pic-projective-space-UFD} et
surjective par les Diviseurs, Lemme
\ref{divisors-lemma-extend-invertible-module} (puisque
$\mathbf{P}^1_k$ est lisse de dimension $1$ sur $k$ et
donc tous ses anneaux locaux sont des anneaux de
valuation discrète). Si $\mathcal{L}$ est dans le noyau
de l'application affichée, alors $\mathcal{L} \cong
\mathcal{O}_{\mathbf{P}^1_k}(nt)$ pour un certain $n \in
\mathbf{Z}$. Nous laissons au lecteur le soin de montrer
que $\mathcal{O}_{\mathbf{P}^1_k}(t) \cong
\mathcal{O}_{\mathbf{P}^1_k}(d)$ où $d = [\kappa(t) : k]$. Par conséquent
$$
\Pic(X) = \mathbf{Z}/d\mathbf{Z}
$$
Ainsi, si $t$ n'est pas un point $k$-rationnel, alors
$d > 1$ et ce groupe de Picard sont non nuls. En
revanche, si nous étendons le corps de base $k$ à toute
extension de corps $k'$ telle qu'il existe un plongement
$k$ $\kappa(t) \to k'$, alors $\mathbf{P}^1_{k'}
\setminus X_{k'}$ possède un point $k'$-rationnel $t'$.
Par conséquent, $\mathcal{O}_{\mathbf{P}^1_{k'}}(1) =
\mathcal{O}_{\mathbf{P}^1_{k'}}(t')$ sera dans le noyau de
l'application $\mathbf{Z} \to \Pic(X_{k'})$ et il
s'ensuivra de la même manière que ci-dessus que $\Pic(X_{k'}) = 0$.
\end{example}

\noindent
Le lemme suivant nous dit que les `` modules inversibles
rationnellement équivalents '' sont isomorphes sur les variétés normales.

\begin{lemma}
\label{lemma-rational-equivalence-for-Pic}
Supposons que $k$ est un corps. Supposons que $X$ est une
variété normale sur $k$. Supposons que $U \subset
\mathbf{A}^n_k$ est un sous-schéma ouvert avec des points
$k$-rationnels $p, q \in U(k)$. Pour chaque module inversible
$\mathcal{L}$ sur $X \times_{\Spec(k)} U$, les restrictions
$\mathcal{L}|_{X \times p}$ et $\mathcal{L}|_{X \times q}$ sont isomorphes.
\end{lemma}

\begin{proof}
Les fibres de $X \times_{\Spec(k)} U \to X$ sont des
sous-schémas ouverts de l'espace affine $n$ sur un corps.
Par conséquent, ces fibres ont des groupes de Picard
triviaux d'après Diviseurs, Lemme
\ref{divisors-lemma-open-subscheme-UFD}. En appliquant Diviseurs,
Lemme \ref{divisors-lemma-in-image-pullback}, nous voyons que
$\mathcal{L}$ est le tiré en arrière d'un module inversible $\mathcal{N}$ sur $X$.
\end{proof}










\section{Unicité du corps de base}
\label{section-base-field}

\noindent
La phrase `` laisser $X$ être un schéma sur $k$ '' signifie que $X$
est un schéma qui vient équipé d'un morphisme $X \to \Spec(k)$.
Maintenant, nous pouvons demander si le corps $k$ est déterminé de
manière unique par le schéma $X$. Bien sûr, ce n'est pas le cas, car par
exemple $\mathbf{A}^1_{\mathbf{C}}$ que nous considérons habituellement
comme un schéma sur le corps $\mathbf{C}$ des nombres complexes,
pourrait également être considéré comme un schéma sur $\mathbf{Q}$. Mais
que se passe-t-il si nous demandons que le morphisme $X \to \Spec(k)$ ne
se factorise pas en tant que $X \to \Spec(k') \to \Spec(k)$ pour une
extension de corps non triviale $k'/k$ ? En d'autres termes, nous demandons
que $k$ soit d'une certaine manière maximal de sorte que $X$ vive sur $k$.

\medskip\noindent
Un exemple pour montrer que cela ne garantit
toujours pas l'unicité de $k$ est le schéma
$$
X =
\Spec\left(
\mathbf{Q}(x)[y]\left[\frac{1}{P(y)}, P \in \mathbf{Q}[y], P \not = 0\right]
\right)
$$
À première vue, cela semble être un schéma sur $\mathbf{Q}(x)$, mais
à y regarder de plus près, il est clair que c'est aussi un schéma
sur $\mathbf{Q}(y)$. De plus, les corps $\mathbf{Q}(x)$ et
$\mathbf{Q}(y)$ sont des sous-corps de $R = \Gamma(X, \mathcal{O}_X)$
maximaux parmi les sous-corps de $R$ (détails omis). En particulier,
$\mathbf{Q}(x)$ et $\mathbf{Q}(y)$ sont tous deux maximaux au sens
ci-dessus. Remarquons que les deux morphismes $X \to
\Spec(\mathbf{Q}(x))$ et $X \to \Spec(\mathbf{Q}(y))$ sont ``
essentiellement de type fini '' (c'est-à-dire que l'homomorphisme d'anneaux
correspondant est essentiellement de type fini). Ainsi $X$ est un schéma
noethérien de dimension finie, c'est-à-dire qu'il n'est pas complètement pathologique.

\medskip\noindent
Un autre problème qui peut empêcher l'unicité est que le schéma
$X$ peut ne pas être réduit. Dans ce cas, il peut y avoir de
nombreux morphismes différents de $X$ vers le spectre d'un corps
donné. À titre d'exemple explicite, considérons les nombres
duals $D = \mathbf{C}[y]/(y^2) = \mathbf{C} \oplus \epsilon
\mathbf{C}$. Étant donnée toute dérivation $\theta : \mathbf{C} \to
\mathbf{C}$ sur $\mathbf{Q}$, nous obtenons un homomorphisme d'anneaux
$$
\mathbf{C} \longrightarrow D, \quad
c \longmapsto c + \epsilon \theta(c).
$$
Le sous-domaine de $\mathbf{C}$ sur lequel toutes ces applications
sont les mêmes est la clôture algébrique de $\mathbf{Q}$. Cela
signifie que prendre l'intersection de tous les corps sur lesquels $X$
peut exister peut finir par être un très petit corps si $X$ n'est pas réduit.

\medskip\noindent
Une observation à ce sujet est la suivante : étant donné un
corps $k$ et deux sous-corps $k_1, k_2$ de $k$ tels que $k$
soit fini sur $k_1$ et sur $k_2$, alors, en général, $k$ {\it
n'est pas } fini sur $k_1 \cap k_2$. Un exemple est fourni par
le corps $k = \mathbf{Q}(t)$ et ses sous-corps $k_1 =
\mathbf{Q}(t^2)$ et $\mathbf{Q}((t + 1)^2)$. En effet, $k_1 \cap
k_2 = \mathbf{Q}$ dans ce cas. Ainsi, dans la suite, nous devons
prendre garde lorsque nous considérons des intersections de corps.

\medskip\noindent
Cela dit, nous montrons maintenant que si $X$ est
localement de type fini sur un corps, alors une
certaine unicité s'applique. Voici le résultat précis.

\begin{proposition}
\label{proposition-unique-base-field}
Supposons que $X$ soit un schéma. Supposons que $a : X \to
\Spec(k_1)$ et $b : X \to \Spec(k_2)$ soient des morphismes de
$X$ vers des spectres de corps. Supposons que $a, b$ soit
localement de type fini, et que $X$ soit réduit et connexe. Alors
nous avons $k_1' = k_2'$, où $k_i' \subset \Gamma(X, \mathcal{O}_X)$
est l'adhérence intégrale de $k_i$ dans $\Gamma(X, \mathcal{O}_X)$.
\end{proposition}

\begin{proof}
Tout d'abord, supposons que le lemme soit vrai dans le cas où
$X$ est quasi-compact (nous traiterons le cas quasi-compact
ci-dessous). Comme $X$ est localement de type fini sur un
corps, il est localement noethérien, voir morphismes, Lemme
\ref{morphisms-lemma-finite-type-noetherian}. En particulier,
cela signifie qu'il est localement connexe, les composantes
connexes des ouverts sont ouvertes, et les intersections
d'ouverts quasi-compacts sont quasi-compactes, voir
Propriétés, Lemme \ref{properties-lemma-Noetherian-topology},
Topologie, Lemme \ref{topology-lemma-locally-connected},
Topologie, Section \ref{topology-section-noetherian}, et
Topologie, Lemme
\ref{topology-lemma-constructible-Noetherian-space}. Choisissons un
recouvrement ouvert $X = \bigcup_{i \in I} U_i$ tel que chaque $U_i$ soit
quasi-compact et connexe. Pour chaque $i$, soit $K_i \subset \mathcal{O}_X(U_i)$
l'adhérence intégrale de $k_1$ et de $k_2$. Pour chaque paire $i, j \in I$ nous décomposons
$$
U_i \cap U_j = \coprod U_{i, j, l}
$$
en ses composantes connexes finies. Écrivez $K_{i, j, l}
\subset \mathcal{O}(U_{i, j, l})$ pour l'adhérence
intégrale de $k_1$ et de $k_2$. D'après le Lemme
\ref{lemma-integral-closure-ground-field}, les anneaux $K_i$
et $K_{i, j, l}$ sont des corps. Nous affirmons maintenant que
$k_1'$ et $k_2'$ sont tous deux égaux au noyau de l'application
$$
\prod K_i \longrightarrow \prod K_{i, j, l}, \quad
(x_i)_i \longmapsto x_i|_{U_{i, j, l}} - x_j|_{U_{i, j, l}}
$$
ce qui prouve ce que nous voulons. À savoir, il est clair que
$k_1'$ est contenu dans ce noyau. D'autre part, supposons que
$(x_i)_i$ soit dans le noyau. Par la condition de faisceau,
$(x_i)_i$ correspond à $f \in \mathcal{O}(X)$. Choisissons un
$i_0 \in I$ et laissons $P(T) \in k_1[T]$ être un polynôme
unitaire avec $P(x_{i_0}) = 0$. Alors nous affirmons que $P(f) =
0$, ce qui prouve que $f \in k_1$. Pour le prouver, nous devons
montrer que $P(x_i) = 0$ pour tout $i$. Choisissons $i \in I$.
Comme $X$ est connexe, il existe une suite $i_0, i_1, \ldots, i_n =
i \in I$ telle que $U_{i_t} \cap U_{i_{t + 1}} \not = \emptyset$.
Maintenant cela signifie que pour chaque $t$, il existe un $l_t$
tel que $x_{i_t}$ et $x_{i_{t + 1}}$ se projettent sur le même
élément du corps $K_{i_t, i_{t + 1}, l_t}$. Par conséquent, si
$P(x_{i_t}) = 0$, alors $P(x_{i_{t + 1}}) = 0$. Par induction, en
commençant par $P(x_{i_0}) = 0$, nous déduisons que $P(x_i) = 0$ comme désiré.

\medskip\noindent
Pour terminer la démonstration du lemme, nous prouvons
le lemme sous l’hypothèse supplémentaire que $X$ est
quasi-compact. D’après le Lemme
\ref{lemma-integral-closure-ground-field}, après avoir
remplacé $k_i$ par $k_i'$, nous pouvons supposer que $k_i$
est intégralement clos dans $\Gamma(X, \mathcal{O}_X)$.
Cela implique que $\mathcal{O}(X)^*/k_i^*$ est un groupe
abélien de type fini, voir la Proposition
\ref{proposition-units-general}. Soit $k_{12} = k_1 \cap k_2$ un
sous-anneau de $\mathcal{O}(X)$. Notez que $k_{12}$ est un corps. Puisque
$$
k_1^*/k_{12}^* \longrightarrow \mathcal{O}(X)^*/k_2^*
$$
nous voyons que $k_1^*/k_{12}^*$ est également un groupe abélien
de type fini. Par conséquent, il existe $\alpha_1, \ldots,
\alpha_n \in k_1^*$ tel que chaque élément $\lambda \in k_1$ ait la forme
$$
\lambda = c \alpha_1^{e_1} \ldots \alpha_n^{e_n}
$$
pour certains $e_i \in \mathbf{Z}$ et $c \in
k_{12}$. En particulier, l'homomorphisme d'anneaux
$$
k_{12}[x_1, \ldots, x_n, \frac{1}{x_1 \ldots x_n}] \longrightarrow k_1, \quad
x_i \longmapsto \alpha_i
$$
est surjective. D'après le Nullstellensatz de Hilbert, Algèbre,
Théorème \ref{algebra-theorem-nullstellensatz}, nous concluons
que $k_1$ est une extension finie de $k_{12}$. De la même
manière, nous concluons que $k_2$ est une extension finie de
$k_{12}$. En particulier, à la fois $k_1$ et $k_2$ sont contenus
dans l'adhérence intégrale $k_{12}'$ de $k_{12}$ dans $\Gamma(X,
\mathcal{O}_X)$. Mais puisque $k_{12}'$ est un corps selon le
Lemme \ref{lemma-integral-closure-ground-field} et que nous avons
choisi $k_i$ pour être intégralement clos dans $\Gamma(X,
\mathcal{O}_X)$, nous concluons que $k_1 = k_{12} = k_2$ comme souhaité.
\end{proof}





\section{Automorphismes}
\label{section-automorphisms}

\noindent
Une section sur les automorphismes des schémas sur les
corps. Pour certaines informations sur les automorphismes
(infinitésimaux) des courbes, voir Courbes Algébriques,
Section \ref{curves-section-vector-fields} et Modules de
Courbes, Section \ref{moduli-curves-section-finite-aut}.

\begin{lemma}
\label{lemma-automorphism}
Supposons que $X$ soit un schéma de type fini réduit sur un corps $k$.
Supposons que $f : X \to X$ soit un automorphisme sur $k$ qui induit
l’application identité sur l’espace topologique sous-jacent de $X$. Alors
\begin{enumerate}
\item $f^*\mathcal{F} \cong \mathcal{F}$ pour
chaque $\mathcal{O}_X$ -module cohérent, et
\item si $\dim(Z) > 0$ pour chaque composante
irréductible $Z \subset X$, alors $f$ est l'identité.
\end{enumerate}
\end{lemma}

\begin{proof}
La partie (1) découle de la partie (2) et du fait que les
composantes connexes de $X$ de dimension $0$ sont des spectres de corps.

\medskip\noindent
Supposons que $Z \subset X$ soit une composante irréductible considérée
comme un sous-schéma fermé intègre. De toute évidence, $f(Z) \subset Z$
et $f|_Z : Z \to Z$ sont un automorphisme sur $k$ qui induit
l'application identité sur l'espace topologique sous-jacent de $Z$.
Puisque $X$ est réduit, il suffit de montrer que les flèches $f|_Z : Z \to Z$
sont l'identité. Cela nous ramène au cas discuté dans le paragraphe suivant.

\medskip\noindent
Supposons que $X$ soit irréductible de dimension $> 0$.
Choisissez un ouvert affine non vide $U \subset X$. Puisque
$f(U) \subset U$ et puisque $U \subset X$ est schématiquement
dense, il suffit de prouver que $f|_U : U \to U$ est l'identité.

\medskip\noindent
Supposez que $X = \Spec(A)$ est affine, irréductible, de dimension $> 0$ et que
$k$ est un corps infini. Supposons que $g \in A$ n'est pas constant. L'ensemble
$$
S = \bigcup\nolimits_{\lambda \in k} V(g - \lambda)
$$
est dense dans $X$ parce que c'est l'image réciproque du
sous-ensemble dense $\mathbf{A}^1_k(k)$ par le morphisme
non constant $g : X \to \mathbf{A}^1_k$. Si $x \in S$,
alors l'image $g(x)$ de $g$ dans $\kappa(x)$ est dans
l'image de $k \to \kappa(x)$. Par conséquent, $f^\sharp :
\kappa(x) \to \kappa(x)$ fixe $g(x)$. Ainsi, l'image de
$f^\sharp(g)$ dans $\kappa(x)$ est égale à $g(x)$. Nous en concluons que
$$
S \subset V(g - f^\sharp(g))
$$
et puisque $X$ est réduit et $S$ est dense, nous concluons
$g=f^\sharp(g)$. Cela prouve $f^\sharp = \text{id}_A$ car $A$ est
engendré comme une $k$-algèbre par les éléments $g$ comme ci-dessus (détails
omis ; indice : l'ensemble des fonctions constantes est un $k$ -sous-espace
vectoriel de dimension finie de $A$). Nous concluons que $f = \text{id}_X$.

\medskip\noindent
Supposons que $X = \Spec(A)$ soit affine, irréductible, de
dimension $> 0$ et que $k$ soit un corps fini. Si pour tout
sous-schéma fermé intègre de dimension $1$, $C \subset X$, la
restriction $f|_C : C \to C$ est l'identité, alors $f$ est
l'identité. Cela nous ramène au cas où $X$ est une courbe. Une courbe
sur un corps fini possède un groupe d'automorphismes fini (détails
omis). Ainsi, $f$ a un ordre fini, disons $n$. Ensuite, nous choisissons
$g : X \to \mathbf{A}^1_k$ non constant comme ci-dessus et nous considérons
$$
S = \{x \in X\text{ closed such that }[\kappa(g(x)) : k]
\text{ is prime to }n\}
$$
En argumentant comme précédemment, nous trouvons que
$S$ est dense dans $X$. Puisque pour $x \in X$ fermé,
l'application $f^\sharp : \kappa(x) \to \kappa(x)$
est un automorphisme d'ordre divisant $n$, nous
voyons que pour $x \in S$, cet automorphisme agit
trivialement sur le sous-corps engendré par l'image de
$g$ dans $\kappa(x)$. Ainsi nous concluons que $S \subset
V(g - f^\sharp(g))$ et nous gagnons comme précédemment.
\end{proof}






\section{Caractéristiques d'Euler}
\label{section-euler}

\noindent
Dans cette section, nous prouvons quelques propriétés élémentaires des
caractères d'Euler des faisceaux cohérents sur des schémas propres sur un corps.

\begin{definition}
\label{definition-euler-characteristic}
Soit $k$ un corps. Soit $X$ un schéma propre sur $k$. Soit
$\mathcal{F}$ un $\mathcal{O}_X$-module cohérent. Dans cette
situation, la {\it caractéristique d'Euler de $\mathcal{F}$ } est l'entier
$$
\chi(X, \mathcal{F}) = \sum\nolimits_i (-1)^i \dim_k H^i(X, \mathcal{F}).
$$
Pour la justification de la formule, voir ci-dessous.
\end{definition}

\noindent
Dans le cas de la définition, seul un nombre fini des
espaces vectoriels $H^i(X, \mathcal{F})$ sont non nuls
(Cohomologie des schémas, Lemme
\ref{coherent-lemma-quasi-coherence-higher-direct-images})
et chacun de ces espaces est de dimension finie (Cohomologie
des schémas, Lemme
\ref{coherent-lemma-proper-over-affine-cohomology-finite}). Ainsi
$\chi(X, \mathcal{F}) \in \mathbf{Z}$ est bien défini. Observez que cette
définition dépend du corps $k$ et pas seulement de la paire $(X, \mathcal{F})$.

\begin{lemma}
\label{lemma-euler-characteristic-additive}
Supposons que $k$ soit un corps. Supposons que $X$ soit un schéma propre
sur $k$. Supposons que $0 \to \mathcal{F}_1 \to \mathcal{F}_2 \to
\mathcal{F}_3 \to 0$ soit une suite exacte courte de modules cohérents sur $X$. Alors
$$
\chi(X, \mathcal{F}_2) = \chi(X, \mathcal{F}_1) + \chi(X, \mathcal{F}_3)
$$
\end{lemma}

\begin{proof}
Considérons la longue suite exacte de cohomologie
$$
0 \to H^0(X, \mathcal{F}_1) \to H^0(X, \mathcal{F}_2) \to
H^0(X, \mathcal{F}_3) \to H^1(X, \mathcal{F}_1) \to \ldots
$$
associé à la suite exacte courte du lemme. Le théorème
du rang et de la nullité en algèbre linéaire montre que
$$
0 = \dim H^0(X, \mathcal{F}_1) - \dim H^0(X, \mathcal{F}_2)
+ \dim H^0(X, \mathcal{F}_3) - \dim H^1(X, \mathcal{F}_1) + \ldots
$$
Cela implique immédiatement le lemme.
\end{proof}

\begin{lemma}
\label{lemma-chi-tensor-finite}
Supposons que $k$ soit un corps. Supposons que $X$ soit un
schéma propre sur $k$. Supposons que $\mathcal{F}$ soit un
faisceau cohérent avec $\dim(\text{Supp}(\mathcal{F})) \leq 0$. Alors
\begin{enumerate}
\item $\mathcal{F}$ est généré par des sections
globales, \item $H^0(X, \mathcal{F}) =
\bigoplus_{x \in \text{Supp}(\mathcal{F})}
\mathcal{F}_x$, \item $H^i(X, \mathcal{F}) = 0$ pour
$i > 0$, \item $\chi(X, \mathcal{F}) = \dim_k
H^0(X, \mathcal{F})$, et \item $\chi(X, \mathcal{F}
\otimes \mathcal{E}) = n\chi(X, \mathcal{F})$ pour
chaque module localement libre $\mathcal{E}$ de rang $n$.
\end{enumerate}
\end{lemma}

\begin{proof}
Par la cohomologie des schémas, Lemme
\ref{coherent-lemma-coherent-support-closed}, nous voyons que
$\mathcal{F} = i_*\mathcal{G}$ où $i : Z \to X$ est l'inclusion
du support schématique de $\mathcal{F}$ et où $\mathcal{G}$ est
un $\mathcal{O}_Z$ -module cohérent. Par définition du support
schématique, l'espace topologique sous-jacent de $Z$ est
$\text{Supp}(\mathcal{F})$. Puisque la dimension de $Z$ est
$0$, nous voyons que $Z$ est affine (Propriétés, Lemme
\ref{properties-lemma-locally-Noetherian-dimension-0}). Ainsi
$\mathcal{G}$ est globalement engendré et les groupes de
cohomologie supérieurs de $\mathcal{G}$ sont nuls (Cohomologie
des schémas, Lemme
\ref{coherent-lemma-quasi-coherent-affine-cohomology-zero}). En
fait, par le Lemme \ref{lemma-algebraic-scheme-dim-0}, le schéma $Z$
est une réunion finie disjointe de spectres d'anneaux locaux
artiniens. Ainsi, correspondamment $H^0(Z, \mathcal{G}) =
\bigoplus_{z \in Z} \mathcal{G}_z$. Les cohomologies de $\mathcal{F}$
et $\mathcal{G}$ coïncident par la cohomologie des schémas, Lemme
\ref{coherent-lemma-relative-affine-cohomology}. Ainsi $H^i(X,
\mathcal{F}) = 0$ pour $i > 0$ et $H^0(X, \mathcal{F}) = H^0(Z,
\mathcal{G})$. En particulier, nous avons que (3) est vrai. Pour $z \in Z$
correspondant à $x \in \text{Supp}(\mathcal{F})$ nous avons
$\mathcal{G}_z = (i_*\mathcal{G})_x = \mathcal{F}_x$. Nous concluons que (2)
est vrai. Bien sûr, (2) implique (1). Nous avons (4) par définition de la
caractéristique d'Euler $\chi(X, \mathcal{F})$ et (3). Par la formule de
projection (Cohomologie, Lemme \ref{cohomology-lemma-projection-formula}) nous avons
$$
i_*(\mathcal{G} \otimes i^*\mathcal{E}) = \mathcal{F} \otimes \mathcal{E}.
$$
Puisque $Z$ a pour dimension $0$, le faisceau localement
libre $i^*\mathcal{E}$ est isomorphe à $\mathcal{O}_Z^{\oplus
n}$ et, en raisonnant comme ci-dessus, nous voyons que (5) tient.
\end{proof}

\begin{lemma}
\label{lemma-euler-characteristic-extend-base-field}
Supposons que $k'/k$ soit une extension de corps. Supposons que $X$ soit un
schéma propre sur $k$. Supposons que $\mathcal{F}$ soit un faisceau
cohérent sur $X$. Supposons que $\mathcal{F}'$ soit l'image réciproque de
$\mathcal{F}$ sur $X_{k'}$. Alors $\chi(X, \mathcal{F}) = \chi(X', \mathcal{F}')$.
\end{lemma}

\begin{proof}
C'est vrai parce que
$$
H^i(X_{k'}, \mathcal{F}') = H^i(X, \mathcal{F}) \otimes_k k'
$$
par changement de base plat, voir Cohomologie des
schémas, Lemme \ref{coherent-lemma-flat-base-change-cohomology}.
\end{proof}

\begin{lemma}
\label{lemma-euler-characteristic-morphism}
Supposons que $k$ soit un corps. Supposons que $f : Y \to X$ soit un morphisme de schémas
propres sur $k$. Supposons que $\mathcal{G}$ soit un $\mathcal{O}_Y$-module cohérent. Alors
$$
\chi(Y, \mathcal{G}) = \sum (-1)^i \chi(X, R^if_*\mathcal{G})
$$
\end{lemma}

\begin{proof}
La formule a du sens : les faisceaux
$R^if_*\mathcal{G}$ sont cohérents et seul un nombre
fini d'entre eux est non nul, voir Cohomologie des
schémas, Proposition
\ref{coherent-proposition-proper-pushforward-coherent} et Lemme
\ref{coherent-lemma-quasi-coherence-higher-direct-images}. D'après
Cohomologie, Lemme \ref{cohomology-lemma-Leray}, il existe une suite spectrale avec
$$
E_2^{p, q} = H^p(X, R^qf_*\mathcal{G})
$$
convergeant vers $H^{p + q}(Y, \mathcal{G})$. Par la finitude de la
cohomologie sur $X$, nous voyons qu’un nombre fini seulement de $E_2^{p, q}$ sont
non nuls et que chaque $E_2^{p, q}$ est un espace vectoriel de dimension finie.
Il en découle que la même chose est vraie pour $E_r^{p, q}$ pour $r \geq 2$ et que
$$
\sum (-1)^{p + q} \dim_k E_r^{p, q}
$$
est indépendant de $r$. Puisque pour $r$ suffisamment grand
nous avons $E_r^{p, q} = E_\infty^{p, q}$ et puisque la
convergence signifie qu'il existe une filtration sur $H^n(Y,
\mathcal{G})$ dont les morceaux gradués sont $E_\infty^{p, q}$
avec $p + q = n$ (c'est le sens de la convergence de la suite
spectrale), nous concluons. Comparez aussi avec l'Homologie plus
générale, Lemme \ref{homology-lemma-biregular-ss-relation-in-K0}.
\end{proof}








\section{Espace projectif}
\label{section-projective-space}

\noindent
Quelques résultats sur l'espace projectif sur un corps.

\begin{lemma}
\label{lemma-projective-space-smooth}
\begin{slogan}
L'espace projectif est lisse.
\end{slogan}
Supposons que $k$ soit un corps et $n \geq 0$. Alors
$\mathbf{P}^n_k$ est une variété projective lisse de dimension $n$ sur $k$.
\end{lemma}

\begin{proof}
Omis.
\end{proof}

\begin{lemma}
\label{lemma-intersection-in-affine-space}
Supposons que $k$ soit un corps et $n \geq 0$. Supposons que $X, Y
\subset \mathbf{A}^n_k$ soit des sous-ensembles fermés. Supposons que $X$
et $Y$ soient équidimensionnels, $\dim(X) = r$ et $\dim(Y) = s$. Alors
chaque composante irréductible de $X \cap Y$ a une dimension $\geq r + s - n$.
\end{lemma}

\begin{proof}
Considérons le sous-schéma fermé $X \times Y \subset
\mathbf{A}^{2n}_k$ où nous utilisons les coordonnées
$x_1, \ldots, x_n, y_1, \ldots, y_n$. Alors $X \cap Y = X
\times Y \cap V(x_1 - y_1, \ldots, x_n - y_n)$. Supposons
que $t \in X \cap Y \subset X \times Y$ soit un point
fermé. Par le Lemme
\ref{lemma-dimension-product-locally-algebraic}, nous avons $\dim_t(X
\times Y) = \dim(X) + \dim(Y)$. Ainsi $\dim(\mathcal{O}_{X \times Y,
t}) = r + s$ par le Lemme \ref{lemma-dimension-locally-algebraic}.
Par Algèbre, Lemme \ref{algebra-lemma-one-equation}, nous concluons que
$$
\dim(\mathcal{O}_{X \cap Y, t}) =
\dim(\mathcal{O}_{X \times Y, t}/(x_1 - y_1, \ldots, x_n - y_n)) \geq
r + s - n
$$
Cela implique le résultat par le Lemme \ref{lemma-dimension-locally-algebraic}.
\end{proof}

\begin{lemma}
\label{lemma-intersection-in-projective-space}
Supposons que $k$ soit un corps et $n \geq 0$. Supposons que
$X, Y \subset \mathbf{P}^n_k$ soit des sous-ensembles fermés
non vides. Si $\dim(X) = r$ et $\dim(Y) = s$ et $r + s \geq n$,
alors $X \cap Y$ est non vide et $\dim(X \cap Y) \geq r + s - n$.
\end{lemma}

\begin{proof}
Écrivez $\mathbf{A}^n = \Spec(k[x_0, \ldots, x_n])$ et
$\mathbf{P}^n = \text{Proj}(k[T_0, \ldots, T_n])$.
Considérez le morphisme $\pi : \mathbf{A}^{n + 1}
\setminus \{0\} \to \mathbf{P}^n$ qui envoie $(x_0,
\ldots, x_n)$ vers le point $[x_0 : \ldots : x_n]$. Plus
précisément, c'est le morphisme associé à la paire
$(\mathcal{O}_{\mathbf{A}^{n + 1} \setminus \{0\}}, (x_0,
\ldots, x_n))$, voir Constructions, Lemme
\ref{constructions-lemma-projective-space}. Sur l'ouvert affine standard
$D_+(T_i)$, nous obtenons le morphisme associé à l'homomorphisme d'anneaux
$$
k\left[\frac{T_0}{T_i}, \ldots, \frac{T_n}{T_i}\right]
\longrightarrow
k\left[T_0, \ldots, T_n, \frac{1}{T_i}\right] \cong
k\left[\frac{T_0}{T_i}, \ldots, \frac{T_n}{T_i}\right]
\left[T_i, \frac{1}{T_i}\right]
$$
qui est surjective et lisse de dimension relative $1$ avec des
fibres irréductibles (détails omis). Par conséquent, $\pi^{-1}(X)$
et $\pi^{-1}(Y)$ sont des sous-ensembles fermés non vides de
dimension $r + 1$ et $s + 1$. Choisissez une composante
irréductible $V \subset \pi^{-1}(X)$ de dimension $r + 1$ et une
composante irréductible $W \subset \pi^{-1}(Y)$ de dimension $s + 1$.
Observez que cela implique que $V$ et $W$ contiennent chaque fibre
de $\pi$ qu'elles rencontrent (puisque $\pi$ a des fibres
irréductibles de dimension $1$ et puisque le Lemme
\ref{lemma-dimension-fibres-locally-algebraic} dit que les fibres de $V
\to \pi(V)$ et $W \to \pi(W)$ ont pour dimension $\geq 1$). Supposons
que $\overline{V}$ et $\overline{W}$ soient l'adhérence de $V$ et $W$
dans $\mathbf{A}^{n + 1}$. Puisque $0 \in \mathbf{A}^{n + 1}$ est dans
l'adhérence de chaque fibre de $\pi$, nous voyons que $0 \in
\overline{V} \cap \overline{W}$. Par le Lemme
\ref{lemma-intersection-in-affine-space}, nous avons $\dim(\overline{V} \cap
\overline{W}) \geq r + s - n + 1$. En argumentant comme ci-dessus en utilisant à
nouveau le Lemme \ref{lemma-dimension-fibres-locally-algebraic}, nous concluons que
$\pi(V \cap W) \subset X \cap Y$ a une dimension d'au moins $r + s - n$ comme souhaité.
\end{proof}

\begin{lemma}
\label{lemma-equation-codim-1-in-projective-space}
Supposons que $k$ soit un corps. Supposons que $Z \subset
\mathbf{P}^n_k$ soit un sous-schéma fermé qui n'a pas de points intégrés
tel que chaque composante irréductible de $Z$ ait dimension $n - 1$. Alors
l'idéal $I(Z) \subset k[T_0, \ldots, T_n]$ correspondant à $Z$ est principal.
\end{lemma}

\begin{proof}
Ceci est un cas particulier des Diviseurs, Lemme
\ref{divisors-lemma-equation-codim-1-in-projective-space}.
\end{proof}




\section{Faisceaux cohérents sur l'espace projectif}
\label{section-coherent-Pn}

\noindent
Dans cette section, nous prouvons certains résultats sur la
cohomologie des faisceaux cohérents sur $\mathbf{P}^n$ sur un
corps, qui peuvent être trouvés dans \cite{Mum}. Ceux-ci seront
utiles plus tard lors de la discussion des schémas de Quot et de Hilbert.

\subsection{Préliminaires}
\label{subsection-preliminaries}

\noindent
Soit $k$ un corps, $n \geq 1$, $d \geq 1$, et soit $s \in
\Gamma(\mathbf{P}_k^n, \mathcal{O}(d))$ une section non
nulle. Dans cette section, nous noterons $\mathcal{O}(d)$ le
$d$-ième tordu du faisceau structural sur l'espace projectif
(Constructions, définitions
\ref{constructions-definition-twist} et
\ref{constructions-definition-projective-space}). Comme $\mathbf{P}^n_k$
est une variété, cette section est régulière ; ainsi $s$ est une section
régulière de $\mathcal{O}(d)$ et définit un diviseur de Cartier effectif
$H = Z(s) \subset \mathbf{P}^n_k$, voir Diviseurs, section
\ref{divisors-section-effective-Cartier-divisors}. Un tel diviseur $H$ est
appelé une {\it hypersurface } et, si $d = 1$, il est appelé un {\it hyperplan }.

\begin{lemma}
\label{lemma-hyperplane}
Supposons que $k$ soit un corps. Supposons
que $n \geq 1$. Soit $i : H \to \mathbf{P}^n_k$
un hyperplan. Alors il existe un isomorphisme
$$
\varphi : \mathbf{P}^{n - 1}_k \longrightarrow H
$$
tel que $i^*\mathcal{O}(1)$ revient à $\mathcal{O}(1)$.
\end{lemma}

\begin{proof}
Nous avons $\mathbf{P}^n_k = \text{Proj}(k[T_0,
\ldots, T_n])$. La section $s$ correspond à une
forme homogène dans $T_0, \ldots, T_n$ de degré
$1$, voir Cohomologie des schémas, Section
\ref{coherent-section-cohomology-projective-space}.
Dire $s = \sum a_i T_i$. Constructions, le Lemme
\ref{constructions-lemma-closed-in-projective-space}
montre que $H = \text{Proj}(k[T_0, \ldots, T_n]/I)$
pour l'idéal gradué $I$ défini en posant $I_d$ égal
au noyau de l'application $\Gamma(\mathbf{P}^n_k,
\mathcal{O}(d)) \to \Gamma(H, i^*\mathcal{O}(d))$.
D'après notre construction de $Z(s)$ dans Divisors,
Définition \ref{divisors-definition-zero-scheme-s},
nous voyons que sur $D_{+}(T_j)$ l'idéal de $H$ est
engendré par $\sum a_i T_i/T_j$ dans l'anneau polynomial
$k[T_0/T_j, \ldots, T_n/T_j]$. Ainsi, il est clair que
$I$ est l'idéal engendré par $\sum a_i T_i$. Remarquez que
$$
k[T_0, \ldots, T_n]/I = k[T_0, \ldots, T_n]/(\sum a_i T_i) \cong
k[S_0, \ldots, S_{n - 1}]
$$
en tant qu'anneaux gradués. Par exemple, si $a_n \not = 0$,
alors mapper $S_i$ égal à la classe de $T_i$ fonctionne.
Nous obtenons l'isomorphisme désiré par fonctorialité de
$\text{Proj}$. L'égalité des torsions de faisceaux de
structure découle par exemple des Constructions, Lemme
\ref{constructions-lemma-surjective-graded-rings-generated-degree-1-map-proj}.
\end{proof}

\begin{lemma}
\label{lemma-exact-sequence-induction}
Supposons que $k$ soit un corps infini.
Supposons que $n \geq 1$. Soit $\mathcal{F}$ un
module cohérent sur $\mathbf{P}^n_k$. Alors il
existe une section non nulle $s \in
\Gamma(\mathbf{P}^n_k, \mathcal{O}(1))$ et une suite exacte courte
$$
0 \to \mathcal{F}(-1) \to \mathcal{F} \to i_*\mathcal{G} \to 0
$$
où $i : H \to \mathbf{P}^n_k$ est l'hyperplan $H$
associé à $s$ et $\mathcal{G} = i^*\mathcal{F}$.
\end{lemma}

\begin{proof}
La application $\mathcal{F}(-1) \to \mathcal{F}$ provient
de Constructions, Équation (
\ref{constructions-equation-multiply-on-sheaf} ) avec $n =
1$, $m = -1$ et la section $s$ de $\mathcal{O}(1)$. Voyons à
quoi ressemble cette application si nous la restreignons à
$D_{+}(T_0)$. Écrivons $D_{+}(T_0) = \Spec(k[x_1, \ldots,
x_n])$ avec $x_i = T_i/T_0$. Identifions
$\mathcal{O}(1)|_{D_{+}(T_0)}$ avec $\mathcal{O}$ en utilisant la
section $T_0$. Par conséquent, si $s = \sum a_iT_i$ alors
$s|_{D_{+}(T_0)} = a_0 + \sum a_ix_i$ avec l'identification
choisie ci-dessus. De plus, supposons que
$\mathcal{F}|_{D_{+}(T_0)}$ corresponde au $k[x_1, \ldots, x_n]$-module
fini $M$. Via l'identification $\mathcal{F}(-1) = \mathcal{F} \otimes
\mathcal{O}(-1)$ et notre trivialisation choisie de $\mathcal{O}(1)$,
nous voyons que $\mathcal{F}(-1)$ correspond également à $M$. Ainsi,
restreindre $\mathcal{F}(-1) \to \mathcal{F}$ à $D_{+}(T_0)$ donne la application
$$
M \xrightarrow{a_0 + \sum a_ix_i} M
$$
Pour voir que la flèche est injective, il suffit de choisir $a_0
+ \sum a_ix_i$ en dehors de l'un quelconque des idéaux premiers
associés de $M$, voir Algèbre, Lemme
\ref{algebra-lemma-ass-zero-divisors}. D'après Algèbre, Lemme
\ref{algebra-lemma-finite-ass}, l'ensemble $\text{Ass}(M)$ des
idéaux premiers associés de $M$ est fini. Notez que pour
$\mathfrak p \in \text{Ass}(M)$, l'intersection $\mathfrak p \cap
\{a_0 + \sum a_i x_i\}$ est un sous-espace vectoriel propre de $k$.
Nous en concluons qu'il existe une famille finie de sous-espaces
vectoriels propres $V_1, \ldots, V_m \subset \Gamma(\mathbf{P}^n_k,
\mathcal{O}(1))$ tels que si nous prenons $s$ en dehors de $\bigcup
V_i$, alors la multiplication par $s$ est injective sur
$D_{+}(T_0)$. De même pour la restriction à $D_{+}(T_j)$ pour $j =
1, \ldots, n$. Comme $k$ est infini, une union finie de sous-espaces
vectoriels propres n'est jamais égale à l'espace tout entier, par
conséquent nous pouvons choisir $s$ de sorte que la application soit
injective. Le conoyau de $\mathcal{F}(-1) \to \mathcal{F}$ est annulé
par $\Im(s : \mathcal{O}(-1) \to \mathcal{O})$ qui est l'idéal
faisceau de $H$ selon Diviseurs, Définition
\ref{divisors-definition-zero-scheme-s}. Ainsi, nous obtenons $\mathcal{G}$ sur $H$ en
utilisant la cohomologie des schémas, Lemme \ref{coherent-lemma-i-star-equivalence}.
\end{proof}

\begin{remark}
\label{remark-exact-sequence-induction}
Supposons que $k$ soit un corps infini. Soit $n \geq 1$. Étant
donné un nombre fini de modules cohérents $\mathcal{F}_i$ sur
$\mathbf{P}^n_k$, nous pouvons choisir un seul $s \in
\Gamma(\mathbf{P}^n_k, \mathcal{O}(1))$ de sorte que l'énoncé du Lemme
\ref{lemma-exact-sequence-induction} fonctionne pour chacun d'eux. Pour
le prouver, il suffit d'appliquer le lemme à $\bigoplus \mathcal{F}_i$.
\end{remark}

\begin{remark}
\label{remark-exact-sequence-induction-cohomology}
Dans la situation des lemmes \ref{lemma-hyperplane} et
\ref{lemma-exact-sequence-induction}, nous avons $H \cong
\mathbf{P}^{n - 1}_k$ avec des torsions de Serre
$\mathcal{O}_H(d) = i^*\mathcal{O}_{\mathbf{P}^n_k}(d)$. Pour
chaque $d \in \mathbf{Z}$, nous avons une suite exacte courte
$$
0 \to \mathcal{F}(d - 1) \to \mathcal{F}(d) \to i_*(\mathcal{G}(d)) \to 0
$$
À savoir, le produit tensoriel par
$\mathcal{O}_{\mathbf{P}^n_k}(d)$ est un foncteur exact et,
d'après la formule de projection (Cohomologie, Lemme
\ref{cohomology-lemma-projection-formula}), nous avons
$i_*(\mathcal{G}(d)) = i_*\mathcal{G} \otimes
\mathcal{O}_{\mathbf{P}^n_k}(d)$. Nous obtenons des suites exactes longues correspondantes
$$
H^i(\mathbf{P}^n_k, \mathcal{F}(d - 1)) \to
H^i(\mathbf{P}^n_k, \mathcal{F}(d)) \to
H^i(H, \mathcal{G}(d)) \to
H^{i + 1}(\mathbf{P}^n_k, \mathcal{F}(d - 1))
$$
Cela découle de ce qui précède et du fait que nous avons
$H^i(\mathbf{P}^n_k, i_*\mathcal{G}(d)) = H^i(H,
\mathcal{G}(d))$ par la cohomologie des schémas, Lemme
\ref{coherent-lemma-relative-affine-cohomology} (les immersions fermées sont affines).
\end{remark}





\subsection{Régularité}
\label{subsection-regularity}

\noindent
Voici la définition.

\begin{definition}
\label{definition-regularity}
Supposons que $k$ est un corps. Supposons que $n \geq 0$.
Soit $\mathcal{F}$ un faisceau cohérent sur
$\mathbf{P}^n_k$. Nous disons que $\mathcal{F}$ est {\it $m$ -régulier } si
$$
H^i(\mathbf{P}^n_k, \mathcal{F}(m - i)) = 0
$$
pour $i = 1, \ldots, n$.
\end{definition}

\noindent
Notez que $\mathcal{F} = \mathcal{O}(d)$ est $m$ -régulier si
et seulement si $d \geq -m$. Cela découle du calcul des groupes
de cohomologie dans Cohomologie des schémas, Équation (
\ref{coherent-equation-identify} ). En effet, nous voyons que
$H^n(\mathbf{P}^n_k, \mathcal{O}(d)) = 0$ si et seulement si $d \geq -n$.

\begin{lemma}
\label{lemma-m-regular-extend-base-field}
Supposons que $k'/k$ soit une extension de corps. Supposons que
$n \geq 0$. Soit $\mathcal{F}$ un faisceau cohérent sur
$\mathbf{P}^n_k$. Supposons que $\mathcal{F}'$ soit le tiré en
arrière de $\mathcal{F}$ sur $\mathbf{P}^n_{k'}$. Alors $\mathcal{F}$
est $m$-régulier si et seulement si $\mathcal{F}'$ est $m$-régulier.
\end{lemma}

\begin{proof}
C'est vrai parce que
$$
H^i(\mathbf{P}^n_{k'}, \mathcal{F}') =
H^i(\mathbf{P}^n_k, \mathcal{F}) \otimes_k k'
$$
par changement de base plat, voir Cohomologie des
schémas, Lemme \ref{coherent-lemma-flat-base-change-cohomology}.
\end{proof}

\begin{lemma}
\label{lemma-m-regular}
Dans la situation du Lemme
\ref{lemma-exact-sequence-induction}, si $\mathcal{F}$ est $m$-régulier,
alors $\mathcal{G}$ est $m$-régulier sur $H \cong \mathbf{P}^{n - 1}_k$.
\end{lemma}

\begin{proof}
Rappelons que $H^i(\mathbf{P}^n_k, i_*\mathcal{G}) = H^i(H,
\mathcal{G})$ par la cohomologie des schémas, Lemme
\ref{coherent-lemma-relative-affine-cohomology}. Ainsi, nous voyons que pour $i \geq 1$ nous obtenons
$$
H^i(\mathbf{P}^n_k, \mathcal{F}(m - i)) \to
H^i(H, \mathcal{G}(m - i)) \to
H^{i + 1}(\mathbf{P}^n_k, \mathcal{F}(m - 1 - i))
$$
par la Remarque
\ref{remark-exact-sequence-induction-cohomology}. Le lemme suit.
\end{proof}

\begin{lemma}
\label{lemma-m-regular-up}
Supposons que $k$ est un corps. Supposons que $n
\geq 0$. Soit $\mathcal{F}$ un faisceau cohérent
sur $\mathbf{P}^n_k$. Si $\mathcal{F}$ est
$m$-régulier, alors $\mathcal{F}$ est $(m + 1)$-régulier.
\end{lemma}

\begin{proof}
Nous prouvons cela par induction sur $n$. Si $n = 0$ chaque
faisceau est $m$-régulier pour tout $m$, il n'y a rien à
démontrer. D'après le Lemme
\ref{lemma-m-regular-extend-base-field}, nous pouvons remplacer $k$ par
un surcorps infini et supposer que $k$ est infini. Ainsi, nous pouvons
appliquer le Lemme \ref{lemma-exact-sequence-induction}. D'après le Lemme
\ref{lemma-m-regular}, nous savons que $\mathcal{G}$ est $m$-régulier.
Par induction sur $n$, nous voyons que $\mathcal{G}$ est $(m +
1)$-régulier. En considérant la suite exacte longue de cohomologie associée à la suite
$$
0 \to \mathcal{F}(m - i) \to \mathcal{F}(m + 1 - i)
\to i_*\mathcal{G}(m + 1 - i) \to 0
$$
comme dans la Remarque
\ref{remark-exact-sequence-induction-cohomology} le lecteur déduit
facilement pour $i \geq 1$ la disparition de $H^i(\mathbf{P}^n_k,
\mathcal{F}(m + 1 - i))$ à partir de la disparition (connue) de
$H^i(\mathbf{P}^n_k, \mathcal{F}(m - i))$ et $H^i(\mathbf{P}^n_k, \mathcal{G}(m + 1 - i))$.
\end{proof}

\begin{lemma}
\label{lemma-m-regular-multiply}
Supposons que $k$ soit un corps. Supposons que $n \geq 0$. Soit
$\mathcal{F}$ un faisceau cohérent sur $\mathbf{P}^n_k$. Si
$\mathcal{F}$ est $m$-régulier, alors l'application de multiplication
$$
H^0(\mathbf{P}^n_k, \mathcal{F}(m)) \otimes_k
H^0(\mathbf{P}^n_k, \mathcal{O}(1))
\longrightarrow
H^0(\mathbf{P}^n_k, \mathcal{F}(m + 1))
$$
est surjective.
\end{lemma}

\begin{proof}
Supposons que $k'/k$ soit une extension de corps. Supposons que
$\mathcal{F}'$ soit comme dans le Lemme
\ref{lemma-m-regular-extend-base-field}. D'après la Cohomologie des schémas,
Lemme \ref{coherent-lemma-flat-base-change-cohomology}, le changement de base de
l'application linéaire du lemme vers $k'$ est la même application linéaire pour le
faisceau $\mathcal{F}'$. Puisque $k \to k'$ est fidèlement plat, il suffit de
prouver le lemme sur $k'$, c’est-à-dire que nous pouvons supposer que $k$ est infini.

\medskip\noindent
Supposons que $k$ est infini. Nous prouvons le lemme par
induction sur $n$. Le cas $n = 0$ est trivial car
$\mathcal{O}(1) \cong \mathcal{O}$ est généré par $T_0$. Pour
$n > 0$, appliquez le Lemme \ref{lemma-exact-sequence-induction}
et tensorisez la séquence par $\mathcal{O}(m + 1)$ pour obtenir
$$
0 \to \mathcal{F}(m) \xrightarrow{s} \mathcal{F}(m + 1) \to
i_*\mathcal{G}(m + 1) \to 0
$$
voir Remarque
\ref{remark-exact-sequence-induction-cohomology}. Soit $t \in
H^0(\mathbf{P}^n_k, \mathcal{F}(m + 1))$. Par induction, l'image
$\overline{t} \in H^0(H, \mathcal{G}(m + 1))$ est l'image de $\sum g_i
\otimes \overline{s}_i$ avec $\overline{s}_i \in \Gamma(H,
\mathcal{O}(1))$ et $g_i \in H^0(H, \mathcal{G}(m))$. Puisque $\mathcal{F}$
est $m$-régulier, nous avons $H^1(\mathbf{P}^n_k, \mathcal{F}(m - 1)) = 0$,
d'où la suite exacte longue de cohomologie associée à la suite exacte courte
$$
0 \to \mathcal{F}(m - 1) \xrightarrow{s} \mathcal{F}(m) \to
i_*\mathcal{G}(m) \to 0
$$
montre que nous pouvons élever $g_i$ à $f_i \in
H^0(\mathbf{P}^n_k, \mathcal{F}(m))$. Nous pouvons également
élever $\overline{s}_i$ à $s_i \in H^0(\mathbf{P}^n_k,
\mathcal{O}(1))$ (voir la preuve du Lemme \ref{lemma-hyperplane} par
exemple). Après avoir soustrait l'image de $\sum f_i \otimes s_i$ de
$t$, nous voyons que nous pouvons supposer $\overline{t} = 0$. Mais
cela signifie exactement que $t$ est l'image de $f \otimes s$ pour un
certain $f \in H^0(\mathbf{P}^n_k, \mathcal{F}(m))$ comme souhaité.
\end{proof}

\begin{lemma}
\label{lemma-m-regular-globally-generated}
Supposons que $k$ est un corps. Supposons que $n
\geq 0$. Soit $\mathcal{F}$ un faisceau cohérent
sur $\mathbf{P}^n_k$. Si $\mathcal{F}$ est
$m$-régulier, alors $\mathcal{F}(m)$ est globalement engendré.
\end{lemma}

\begin{proof}
Pour tout $d \gg 0$ le faisceau $\mathcal{F}(d)$ est globalement
engendré. Cela découle par exemple de la première partie de
Cohomologie des schémas, Lemme
\ref{coherent-lemma-coherent-projective}. Choisir $d \geq m$ de telle sorte
que $\mathcal{F}(d)$ soit globalement engendré. Choisir une base $f_1,
\ldots, f_r \in H^0(\mathbf{P}^n_k, \mathcal{F})$. D'après le Lemme
\ref{lemma-m-regular-multiply}, chaque élément $f \in H^0(\mathbf{P}^n_k,
\mathcal{F}(d))$ peut être écrit comme $f = \sum P_if_i$ pour un certain $P_i \in
k[T_0, \ldots, T_n]$ homogène de degré $d - m$. Puisque les sections $f$ génèrent
$\mathcal{F}(d)$, il en découle que les sections $f_i$ génèrent $\mathcal{F}(m)$.
\end{proof}

\subsection{Polynômes de Hilbert}
\label{subsection-hilbert}

\noindent
Le lemme suivant sera rendu obsolète par le lemme plus
général \ref{lemma-numerical-polynomial-from-euler}.

\begin{lemma}
\label{lemma-hilbert-polynomial}
Supposons que $k$ est un corps. Supposons que $n \geq 0$. Soit
$\mathcal{F}$ un faisceau cohérent sur $\mathbf{P}^n_k$. La fonction
$$
d \longmapsto \chi(\mathbf{P}^n_k, \mathcal{F}(d))
$$
est un polynôme.
\end{lemma}

\begin{proof}
Nous prouvons cela par induction sur $n$. Si $n =
0$, alors $\mathbf{P}^n_k = \Spec(k)$ et
$\mathcal{F}(d) = \mathcal{F}$. Par conséquent, dans
ce cas, la fonction est constante, c’est-à-dire un
polynôme de degré $0$. Supposons $n > 0$. Par le Lemme
\ref{lemma-euler-characteristic-extend-base-field},
nous pouvons supposer que $k$ est infini. Appliquons
le Lemme \ref{lemma-exact-sequence-induction}. En
appliquant le Lemme
\ref{lemma-euler-characteristic-additive} aux suites tordues $0 \to
\mathcal{F}(d - 1) \to \mathcal{F}(d) \to i_*\mathcal{G}(d) \to 0$, nous obtenons
$$
\chi(\mathbf{P}^n_k, \mathcal{F}(d)) -
\chi(\mathbf{P}^n_k, \mathcal{F}(d - 1)) =
\chi(H, \mathcal{G}(d))
$$
Voir la Remarque
\ref{remark-exact-sequence-induction-cohomology}. Depuis $H
\cong \mathbf{P}^{n - 1}_k$ par induction, le côté droit est un
polynôme. Le lemme se termine en notant que toute fonction $f :
\mathbf{Z} \to \mathbf{Z}$ ayant la propriété que l'application
$d \mapsto f(d) - f(d - 1)$ est un polynôme, est elle-même un
polynôme. Nous omettons la preuve de ce fait (indice : comparer
avec Algèbre, Lemme \ref{algebra-lemma-numerical-polynomial}).
\end{proof}

\begin{definition}
\label{definition-hilbert-polynomial}
Soit $k$ un corps. Soit $n \geq 0$. Soit $\mathcal{F}$
un faisceau cohérent sur $\mathbf{P}^n_k$. La fonction
$d \mapsto \chi(\mathbf{P}^n_k, \mathcal{F}(d))$ est
appelée le {\it polynôme de Hilbert } de $\mathcal{F}$.
\end{definition}

\noindent
Le polynôme de Hilbert a des coefficients dans
$\mathbf{Q}$ et pas en général dans $\mathbf{Z}$. Par exemple,
le polynôme de Hilbert de $\mathcal{O}_{\mathbf{P}^n_k}$ est
$$
d \longmapsto {d + n \choose n} = \frac{d^n}{n!} + \ldots
$$
Cela découle du lemme suivant et du fait que
$$
H^0(\mathbf{P}^n_k, \mathcal{O}_{\mathbf{P}^n_k}(d)) = k[T_0, \ldots, T_n]_d
$$
(degré $d$ partie) dont la dimension sur $k$ est ${d + n \choose n}$.

\begin{lemma}
\label{lemma-hilbert-polynomial-H0}
Supposons que $k$ soit un corps. Supposons que $n \geq 0$.
Soit $\mathcal{F}$ un faisceau cohérent sur
$\mathbf{P}^n_k$ avec le polynôme de Hilbert $P \in \mathbf{Q}[t]$. Alors
$$
P(d) = \dim_k H^0(\mathbf{P}^n_k, \mathcal{F}(d))
$$
pour tous les $d \gg 0$.
\end{lemma}

\begin{proof}
Cela découle de la disparition de la
cohomologie des torsions suffisamment élevées
de $\mathcal{F}$. Voir Cohomologie des schémas,
Lemme \ref{coherent-lemma-coherent-projective}.
\end{proof}



\subsection{Bornitude des quotients}
\label{subsection-boundedness}

\noindent
Dans cette sous-section, nous bornons la
régularité des quotients d'un faisceau cohérent donné
sur $\mathbf{P}^n$ en termes du polynôme de Hilbert.

\begin{lemma}
\label{lemma-bound-quotients-free}
Supposons que $k$ soit un corps. Soit $n \geq 0$. Soit $r
\geq 1$. Soit $P \in \mathbf{Q}[t]$. Il existe un entier $m$
dépendant de $n$, $r$ et $P$ ayant la propriété suivante : si
$$
0 \to \mathcal{K} \to \mathcal{O}^{\oplus r} \to \mathcal{F} \to 0
$$
est une suite exacte courte de faisceaux cohérents
sur $\mathbf{P}^n_k$ et $\mathcal{F}$ a pour polynôme
de Hilbert $P$, alors $\mathcal{K}$ est $m$-régulier.
\end{lemma}

\begin{proof}
Nous prouvons cela par induction sur $n$. Si $n = 0$, alors
$\mathbf{P}^n_k = \Spec(k)$ et tout module cohérent est
$0$-régulier et toute application surjective est surjective
sur les sections globales. Supposons $n > 0$. Considérons une
suite exacte comme dans le lemme. Supposons que $P' \in
\mathbf{Q}[t]$ est le polynôme $P'(t) = P(t) - P(t - 1)$.
Supposons que $m'$ est l'entier qui fonctionne pour $n - 1$,
$r$ et $P'$. D'après les lemmes
\ref{lemma-m-regular-extend-base-field} et
\ref{lemma-euler-characteristic-extend-base-field}, nous pouvons
remplacer $k$ par une extension de corps, donc nous pouvons
supposer que $k$ est infini. Appliquez le lemme
\ref{lemma-exact-sequence-induction} au faisceau cohérent
$\mathcal{F}$. Le polynôme de Hilbert de $\mathcal{F}' =
i^*\mathcal{F}$ est $P'$ (voir la preuve du lemme
\ref{lemma-hilbert-polynomial}). Puisque $i^*$ est exact à droite, nous
voyons que $\mathcal{F}'$ est un quotient de $\mathcal{O}_H^{\oplus r}
= i^*\mathcal{O}^{\oplus r}$. Ainsi, l'hypothèse d'induction s'applique
à $\mathcal{F}'$ sur $H \cong \mathbf{P}^{n - 1}_k$ (lemme
\ref{lemma-hyperplane}). Notez que la application $\mathcal{K}(-1) \to
\mathcal{K}$ est injective comme $\mathcal{K} \subset \mathcal{O}^{\oplus
r}$ et a un conoyau $i_*\mathcal{H}$ où $\mathcal{H} = i^*\mathcal{K}$.
Par le lemme du serpent (Homologie, Lemme \ref{homology-lemma-snake}) nous
obtenons un diagramme commutatif avec des colonnes et des lignes exactes.
$$
\xymatrix{
& 0 \ar[d] & 0 \ar[d] & 0 \ar[d] \\
0 \ar[r] &
\mathcal{K}(-1) \ar[r] \ar[d] &
\mathcal{O}^{\oplus r}(-1) \ar[r] \ar[d] &
\mathcal{F}(-1) \ar[d] \ar[r] & 0 \\
0 \ar[r] &
\mathcal{K} \ar[r] \ar[d] &
\mathcal{O}^{\oplus r} \ar[r] \ar[d] &
\mathcal{F} \ar[d] \ar[r] & 0\\
0 \ar[r] &
i_*\mathcal{H} \ar[r] \ar[d] &
i_*\mathcal{O}_H^{\oplus r} \ar[r] \ar[d] &
i_*\mathcal{F}' \ar[r] \ar[d] & 0 \\
& 0 & 0 & 0
}
$$
Ainsi, l'hypothèse d'induction s'applique à la suite exacte
$0 \to \mathcal{H} \to \mathcal{O}_H^{\oplus r} \to
\mathcal{F}' \to 0$ sur $H \cong \mathbf{P}^{n - 1}_k$ (Lemme
\ref{lemma-hyperplane}) et $\mathcal{H}$ est $m'$-régulier.
Rappelons que cela implique que $\mathcal{H}$ est
$d$-régulier pour tous les $d \geq m'$ (Lemme \ref{lemma-m-regular-up}).

\medskip\noindent
Soient $i \geq 2$ et $d \geq m'$. Il découle de la
suite exacte longue de cohomologie associée à la
colonne de gauche du diagramme ci-dessus et de la
nullité de $H^{i - 1}(H, \mathcal{H}(d))$ que l'application
$$
H^i(\mathbf{P}^n_k, \mathcal{K}(d - 1))
\longrightarrow
H^i(\mathbf{P}^n_k, \mathcal{K}(d))
$$
est injective. Comme ces groupes sont nuls pour
$d \gg 0$ (Cohomologie des schémas, Lemme
\ref{coherent-lemma-coherent-projective}), nous en
concluons que $H^i(\mathbf{P}^n_k, \mathcal{K}(d))$
sont nuls pour tous les $d \geq m'$ et $i \geq 2$.

\medskip\noindent
Nous devons encore contrôler $H^1$. Tout d'abord, nous observons que toutes les applications
$$
H^1(\mathbf{P}^n_k, \mathcal{K}(m' - 1)) \to
H^1(\mathbf{P}^n_k, \mathcal{K}(m')) \to
H^1(\mathbf{P}^n_k, \mathcal{K}(m' + 1)) \to \ldots
$$
sont surjectives par l'annulation de $H^1(H,
\mathcal{H}(d))$ pour $d \geq m'$. Supposons que $d > m'$ soit tel que
$$
H^1(\mathbf{P}^n_k, \mathcal{K}(d - 1))
\longrightarrow
H^1(\mathbf{P}^n_k, \mathcal{K}(d))
$$
est injective. Ensuite, $H^0(\mathbf{P}^n_k,
\mathcal{K}(d)) \to H^0(H, \mathcal{H}(d))$ est
surjective. Considérons le diagramme commutatif
$$
\xymatrix{
H^0(\mathbf{P}^n_k, \mathcal{K}(d)) \otimes_k
H^0(\mathbf{P}^n_k, \mathcal{O}(1))
\ar[r] \ar[d] &
H^0(\mathbf{P}^n_k, \mathcal{K}(d + 1)) \ar[d] \\
H^0(H, \mathcal{H}(d)) \otimes_k
H^0(H, \mathcal{O}_H(1))
\ar[r] &
H^0(H, \mathcal{H}(d + 1))
}
$$
D'après le Lemme \ref{lemma-m-regular-multiply}, nous voyons
que la flèche horizontale du bas est surjective. Par conséquent,
la flèche verticale de droite est surjective. Nous concluons que
$$
H^1(\mathbf{P}^n_k, \mathcal{K}(d))
\longrightarrow
H^1(\mathbf{P}^n_k, \mathcal{K}(d + 1))
$$
est injective. Par induction, nous voyons que
$$
H^1(\mathbf{P}^n_k, \mathcal{K}(d - 1)) \to
H^1(\mathbf{P}^n_k, \mathcal{K}(d)) \to
H^1(\mathbf{P}^n_k, \mathcal{K}(d + 1)) \to \ldots
$$
sont toutes injectives et nous concluons que
$H^1(\mathbf{P}^n_k, \mathcal{K}(d - 1)) = 0$ en
raison de la disparition éventuelle de ces groupes.
Ainsi, les dimensions des groupes $H^1(\mathbf{P}^n_k,
\mathcal{K}(d))$ pour $d \geq m'$ diminuent strictement
jusqu'à devenir nulles. Il s'ensuit que la régularité
de $\mathcal{K}$ est bornée par $m' + \dim_k
H^1(\mathbf{P}^n_k, \mathcal{K}(m'))$. D'autre part, par la
disparition des groupes de cohomologie supérieurs, nous avons
$$
\dim_k H^1(\mathbf{P}^n_k, \mathcal{K}(m')) = 
- \chi(\mathbf{P}^n_k, \mathcal{K}(m')) +
\dim_k H^0(\mathbf{P}^n_k, \mathcal{K}(m'))
$$
Notez que le $H^0$ a une dimension bornée par la dimension
de $H^0(\mathbf{P}^n_k, \mathcal{O}^{\oplus r}(m'))$ qui
est au maximum $r{n + m' \choose n}$ si $m' > 0$ et zéro
sinon. Enfin, le terme $\chi(\mathbf{P}^n_k,
\mathcal{K}(m'))$ est égal à $r{n + m' \choose n} - P(m')$. Cela
donne une borne du type souhaité, terminant la démonstration du lemme.
\end{proof}







\section{Frobenius}
\label{section-frobenius}

\noindent
Supposons que $p$ soit un nombre premier. Si $X$ est un schéma, alors nous disons
que `` $X$ a pour caractéristique $p$ '', ou que `` $X$ est de caractéristique $p$
'', ou que `` $X$ est en caractéristique $p$ '' si $p$ est nul dans $\mathcal{O}_X$.

\begin{definition}
\label{definition-absolute-frobenius}
Soit $p$ un nombre premier. Soit $X$ un schéma de caractéristique $p$.
Le {\it Frobenius absolu de $X$ } est le morphisme $F_X : X \to X$
donné par l'identité sur l'espace topologique sous-jacent et dont
$F_X^\sharp : \mathcal{O}_X \to \mathcal{O}_X$ est donné par $g \mapsto g^p$.
\end{definition}

\noindent
Cela a du sens car pour tout anneau $A$ de caractéristique $p$,
l'application $F_A : A \to A$, $a \mapsto a^p$ est un
endomorphisme d'anneau qui induit l'identité sur $\Spec(A)$. De
plus, si $A$ est local, alors $F_A$ est un homomorphisme local.
De cette manière, nous voyons que le Frobenius absolu de $X$ est
un endomorphisme de $X$ dans la catégorie des schémas. Il s'avère
que le Frobenius absolu définit une application de soi du
foncteur identité sur la catégorie des schémas en caractéristique $p$.

\begin{lemma}
\label{lemma-frobenius-endomorphism-identity}
Supposons que $p > 0$ soit un nombre premier.
Supposons que $f : X \to Y$ soit un morphisme de
schémas en caractéristique $p$. Alors le diagramme
$$
\xymatrix{
X \ar[d]_f \ar[r]_{F_X} & X \ar[d]^f \\
Y \ar[r]^{F_Y} & Y
}
$$
trajets domicile-travail.
\end{lemma}

\begin{proof}
Cela découle du fait algébrique trivial suivant : si
$\varphi : A \to B$ est un homomorphisme d'anneaux de
caractéristique $p$, alors $\varphi(a^p) = \varphi(a)^p$.
\end{proof}

\begin{lemma}
\label{lemma-frobenius}
Supposons que $p > 0$ soit un nombre premier. Supposons que $X$
soit un schéma en caractéristique $p$. Alors le frobenius absolu
$F_X : X \to X$ est un homéomorphisme universel, est intégral, et
induit des extensions de résidus de corps purement inséparables.
\end{lemma}

\begin{proof}
Cela découle des résultats correspondants pour l'endomorphisme de
Frobenius $F_A : A \to A$ d'un anneau $A$ de caractéristique $p > 0$.
Voir la discussion dans Algèbre, Section
\ref{algebra-section-universal-homeomorphism}, par exemple le Lemme \ref{algebra-lemma-p-ring-map}.
\end{proof}

\noindent
Si nous travaillons avec des schémas sur une base fixe,
alors il existe une version relative du morphisme de Frobenius.

\begin{definition}
\label{definition-relative-frobenius}
Supposons que $p > 0$ soit un nombre premier. Supposons que $S$ soit un schéma en
caractéristique $p$. Supposons que $X$ soit un schéma sur $S$. Nous définissons
$$
X^{(p)} = X^{(p/S)} = X \times_{S, F_S} S
$$
vu comme un schéma sur $S$. En appliquant le Lemme
\ref{lemma-frobenius-endomorphism-identity}, nous voyons
qu'il existe un unique morphisme $F_{X/S} : X
\longrightarrow X^{(p)}$ sur $S$ s'insérant dans le diagramme commutatif
$$
\xymatrix{
X \ar[rr]_{F_{X/S}} \ar[rrd] \ar@/^1em/[rrrr]^{F_X}
& & X^{(p)} \ar[rr] \ar[d] & & X \ar[d] \\
& & S \ar[rr]^{F_S} & & S
}
$$
où le carré de droite est cartésien. Le morphisme $F_{X/S}$
est appelé le {\it morphisme de Frobenius relatif de $X/S$ }.
\end{definition}

\noindent
Remarquez que $X \mapsto X^{(p)}$ est un foncteur ; c'est le foncteur de
changement de base pour le morphisme de Frobenius absolu $F_S : S \to S$. Nous
avons les mêmes lemmes qu'avant concernant le morphisme de Frobenius relatif.

\begin{lemma}
\label{lemma-relative-frobenius-endomorphism-identity}
Supposons que $p > 0$ soit un nombre premier. Supposons que
$S$ soit un schéma en caractéristique $p$. Supposons que $f : X
\to Y$ soit un morphisme de schémas sur $S$. Alors le diagramme
$$
\xymatrix{
X \ar[d]_f \ar[r]_{F_{X/S}} & X^{(p)} \ar[d]^{f^{(p)}} \\
Y \ar[r]^{F_{Y/S}} & Y^{(p)}
}
$$
trajets domicile-travail.
\end{lemma}

\begin{proof}
Cela découle du Lemme
\ref{lemma-frobenius-endomorphism-identity} et des définitions.
\end{proof}

\begin{lemma}
\label{lemma-relative-frobenius}
Supposons que $p > 0$ soit un nombre premier. Supposons que
$S$ soit un schéma en caractéristique $p$. Supposons que $X$
soit un schéma sur $S$. Alors le Frobenius relatif $F_{X/S} : X
\to X^{(p)}$ est un homéomorphisme universel, est intégral, et
induit des extensions de résidus de corps purement inséparables.
\end{lemma}

\begin{proof}
D'après le lemme \ref{lemma-frobenius}, les morphismes $F_X : X
\to X$ et le changement de base $h : X^{(p)} \to X$ de $F_S$ sont
des homéomorphismes universels. Puisque $h \circ F_{X/S} = F_X$,
nous concluons que $F_{X/S}$ est un homéomorphisme universel
(morphismes, Lemme
\ref{morphisms-lemma-permanence-universal-homeomorphism}). D'après les lemmes
\ref{morphisms-lemma-universal-homeomorphism} et
\ref{morphisms-lemma-universally-injective} de morphismes, nous concluons que $F_{X/S}$ possède également les autres propriétés.
\end{proof}

\begin{lemma}
\label{lemma-relative-frobenius-omega}
Supposons que $p > 0$ soit un nombre premier. Supposons que $S$ soit un schéma en
caractéristique $p$. Supposons que $X$ soit un schéma sur $S$. Alors $\Omega_{X/S} = \Omega_{X/X^{(p)}}$.
\end{lemma}

\begin{proof}
Cela se traduit par le fait algébrique suivant. Supposons
que $A \to B$ soit un homomorphisme d'anneaux de
caractéristique $p$. Posons $B' = B \otimes_{A, F_A} A$ et
considérons l'homomorphisme d'anneaux $F_{B/A} : B' \to B$,
$b \otimes a \mapsto b^pa$. Alors notre assertion est que
$\Omega_{B/A} = \Omega_{B/B'}$. C'est vrai parce que
$\text{d}(b^pa) = 0$ si $\text{d} : B \to \Omega_{B/A}$ est la
dérivation universelle et donc $\text{d}$ est une dérivation $B'$.
\end{proof}

\begin{lemma}
\label{lemma-relative-frobenius-finite}
Supposons que $p > 0$ soit un nombre premier. Supposons que $S$ soit un
schéma en caractéristique $p$. Supposons que $X$ soit un schéma sur
$S$. Si $X \to S$ est localement de type fini, alors $F_{X/S}$ est fini.
\end{lemma}

\begin{proof}
Cela se traduit par le fait algébrique suivant. Supposons que $A
\to B$ soit un homomorphisme de type fini d'anneaux de
caractéristique $p$. Posons $B' = B \otimes_{A, F_A} A$ et
considérons l'homomorphisme d'anneaux $F_{B/A} : B' \to B$, $b
\otimes a \mapsto b^pa$. Alors notre assertion est que $F_{B/A}$ est
fini. À savoir, si $x_1, \ldots, x_n \in B$ sont des générateurs sur
$A$, alors $x_i$ est intégral sur $B'$ parce que $x_i^p = F_{B/A}(x_i
\otimes 1)$. Par conséquent, $F_{B/A} : B' \to B$ est fini d'après le
Lemme d'Algèbre \ref{algebra-lemma-characterize-finite-in-terms-of-integral}.
\end{proof}

\begin{lemma}
\label{lemma-geometrically-reduced-p}
Supposons que $k$ soit un corps de caractéristique $p > 0$. Supposons que $X$ soit un
schéma sur $k$. Alors $X$ est géométriquement réduit si et seulement si $X^{(p)}$ est réduit.
\end{lemma}

\begin{proof}
Considérons le Frobenius absolu $F_k : k \to k$. Alors
$F_k(k) = k^p$, en d'autres termes, $F_k : k \to k$ est
isomorphe à l'inclusion de $k$ dans $k^{1/p}$. Ainsi, le
lemme découle du Lemme \ref{lemma-geometrically-reduced}.
\end{proof}

\begin{lemma}
\label{lemma-inseparable-deg-p-smooth}
Supposons que $k$ soit un corps de caractéristique $p > 0$. Supposons
que $X$ soit une variété sur $k$. Les énoncés suivants sont équivalents
\begin{enumerate}
\item $X^{(p)}$ est réduit, \item $X$ est
géométriquement réduit, \item il y a un
ouvert non vide $U \subset X$ lisse sur $k$.
\end{enumerate}
Dans ce cas, $X^{(p)}$ est une variété sur $k$ et $F_{X/k} : X
\to X^{(p)}$ est un morphisme dominant fini de degré $p^{\dim(X)}$.
\end{lemma}

\begin{proof}
Nous avons vu l'équivalence de (1) et (2) dans le
Lemme \ref{lemma-geometrically-reduced-p}. Nous avons
vu que (2) implique (3) dans le Lemme
\ref{lemma-geometrically-reduced-dense-smooth-open}. Si
(3) est vrai, alors $U$ est géométriquement réduit (voir
par exemple le Lemme
\ref{lemma-geometrically-regular-smooth}) et donc $X$ est
géométriquement réduit par le Lemme
\ref{lemma-generic-points-geometrically-reduced}. De cette manière, nous voyons que (1), (2) et (3) sont équivalents.

\medskip\noindent
Supposons que (1), (2) et (3) soient vérifiés. Puisque
$F_{X/k}$ est un homéomorphisme (Lemme
\ref{lemma-relative-frobenius}), nous voyons que $X^{(p)}$ est
une variété. Ensuite, $F_{X/k}$ est fini par le Lemme
\ref{lemma-relative-frobenius-finite}. Il est dominant car il
est surjectif. Pour calculer le degré (morphismes, Définition
\ref{morphisms-definition-degree}), il suffit de calculer le
degré de $F_{U/k} : U \to U^{(p)}$ (comme $F_{U/k} = F_{X/k}|_U$
par le Lemme
\ref{lemma-relative-frobenius-endomorphism-identity}). Après avoir
réduit un peu $U$, nous pouvons supposer qu'il existe un morphisme
étale \' $h : U \to \mathbf{A}^n_k$, voir morphismes, Lemme
\ref{morphisms-lemma-smooth-etale-over-affine-space}. Bien sûr $n =
\dim(U)$ parce que $\mathbf{A}^n_k \to \Spec(k)$ est lisse de
dimension relative $n$, le morphisme étale \' $h$ est lisse de
dimension relative $0$, et $U \to \Spec(k)$ est lisse de dimension
relative $\dim(U)$ et les dimensions relatives s'additionnent
correctement (morphismes, Lemme
\ref{morphisms-lemma-composition-relative-dimension-d}). Remarquons que $h$ est un
morphisme génériquement fini et dominant de variétés, et donc $\deg(h)$ est défini. D'après le
Lemme \ref{lemma-relative-frobenius-endomorphism-identity}, nous avons un diagramme commutatif
$$
\xymatrix{
U \ar[rr]_{F_{X/k}} \ar[d]_h & & U^{(p)} \ar[d]^{h^{(p)}} \\
\mathbf{A}^n_k \ar[rr]^{F_{\mathbf{A}^n_k/k}} & &
(\mathbf{A}^n_k)^{(p)}
}
$$
Puisque $h^{(p)}$ est un changement de base de $h$, il
est également \' étale et il s'ensuit que $h^{(p)}$ est
également un morphisme dominant génériquement fini de
variétés. Le degré de $h^{(p)}$ est le degré de
l'extension $k(X^{(p)})/k((\mathbf{A}^n_k)^{(p)})$, qui est
le même que le degré de l'extension
$k(X)/k(\mathbf{A}^n_k)$ parce que $h^{(p)}$ est le changement
de base de $h$ (petit détail omis). Par la multiplicativité des
degrés (morphismes, Lemme
\ref{morphisms-lemma-degree-composition}), il suffit de montrer que le
degré de $F_{\mathbf{A}^n_k/k}$ est $p^n$. Pour le voir, observez que
$(\mathbf{A}^n_k)^{(p)} = \mathbf{A}^n_k$ et que $F_{\mathbf{A}^n_k/k}$
est donné par l'application envoyant les coordonnées à leurs puissances $p$.
\end{proof}

\begin{remark}
\label{remark-n-fold-relative-frobenius}
Supposons que $p > 0$ soit un nombre premier. Supposons que $S$ soit un schéma
en caractéristique $p$. Supposons que $X$ soit un schéma sur $S$. Pour $n \geq 1$
$$
X^{(p^n)} = X^{(p^n/S)} = X \times_{S, F_S^n} S
$$
considéré comme un schéma sur $S$. Observez que $X \mapsto
X^{(p^n)}$ est un foncteur. En appliquant le Lemme
\ref{lemma-frobenius-endomorphism-identity}, nous voyons que
$F_{X/S, n} = (F_X^n, \text{id}_S) : X \longrightarrow X^{(p^n)}$
est un morphisme sur $S$ s'insérant dans le diagramme commutatif
$$
\xymatrix{
X \ar[rr]_{F_{X/S, n}} \ar[rrd] \ar@/^1em/[rrrr]^{F_X^n}
& & X^{(p^n)} \ar[rr] \ar[d] & & X \ar[d] \\
& & S \ar[rr]^{F_S^n} & & S
}
$$
où le carré de droite est cartésien. Le
morphisme $F_{X/S, n}$ est parfois appelé le {\it
morphisme de Frobenius relatif itéré $n$ fois de
$X/S$ }. Cela a du sens car nous avons la formule
$$
F_{X/S, n} =
F_{X^{(p^{n - 1})}/S} \circ \ldots \circ F_{X^{(p)}/S} \circ F_{X/S}
$$
ce qui montre que $F_{X/S, n}$ est la composition
de $n$ par rapport aux Frobenius. Puisque nous avons
$$
F_{X^{(p^m)}/S} = F_{X^{(p^{m - 1})}/S}^{(p)} = \ldots = F_{X/S}^{(p^m)}
$$
(détails omis) nous obtenons également que
$$
F_{X/S, n} =
F_{X/S}^{(p^{n - 1})} \circ \ldots \circ F_{X/S}^{(p)} \circ F_{X/S}
$$
\end{remark}




\section{Recollement d'anneaux de dimension un}
\label{section-glueing-dim-1}

\noindent
Cette section contient quelques préliminaires algébriques
pour démontrer qu'un ensemble fini de points de codimension
$1$ d'un schéma séparé est contenu dans un ouvert affine.

\begin{situation}
\label{situation-glue}
Ici, on nous donne un diagramme commutatif d'anneaux
$$
\xymatrix{
A \ar[r] & K \\
R \ar[u] \ar[r] & B \ar[u]
}
$$
où $K$ est un corps et $A$, $B$ sont des sous-anneaux de $K$ avec
le corps des fractions $K$. Enfin, $R = A \times_K B = A \cap B$.
\end{situation}

\begin{lemma}
\label{lemma-glue-valuation-ring}
Dans la situation \ref{situation-glue}, supposons que $B$ soit un anneau de valuation.
Alors pour chaque élément inversible $u$ de $A$, soit $u \in R$ soit $u^{-1} \in R$.
\end{lemma}

\begin{proof}
À savoir, si l'image $c$ de $u$ dans $K$ est
dans $B$, alors $u \in R$. Sinon, $c^{-1} \in
B$ (Algèbre, Lemme
\ref{algebra-lemma-valuation-ring-x-or-x-inverse}) et $u^{-1} \in R$.
\end{proof}

\noindent
Le lemme suivant explique la signification de la
condition `` $A \otimes B \to K$ est surjective
'' qui revient assez souvent dans ce qui suit.

\begin{lemma}
\label{lemma-glue-separated}
Dans la situation \ref{situation-glue}, supposons que $A$ est un anneau
noethérien de dimension $1$. Les assertions suivantes sont équivalentes
\begin{enumerate}
\item $A \otimes B \to K$ n'est pas surjective,
\item il existe un anneau de valuation discrète
$\mathcal{O} \subset K$ contenant à la fois $A$ et $B$.
\end{enumerate}
\end{lemma}

\begin{proof}
Il est clair que (2) implique (1). D'autre part, si $A \otimes B \to
K$ n'est pas surjective, alors l'image $C \subset K$ n'est pas un
corps, donc $C$ possède un idéal maximal non nul $\mathfrak m$.
Choisissez un anneau de valuation $\mathcal{O} \subset K$ dominant
$C_\mathfrak m$. D'après le Lemme d'Algèbre
\ref{algebra-lemma-krull-akizuki} appliqué à $A \subset \mathcal{O}$, l'anneau
$\mathcal{O}$ est noethérien. Par conséquent, $\mathcal{O}$ est un anneau de valuation
discrète selon le Lemme d'Algèbre \ref{algebra-lemma-valuation-ring-Noetherian-discrete}.
\end{proof}

\begin{lemma}
\label{lemma-semi-local}
Dans la situation \ref{situation-glue}, supposez
\begin{enumerate}
\item $A$ est un domaine semi-local noethérien de
dimension $1$, \item $B$ est un anneau de valuation discrète,
\end{enumerate}
Alors nous avons les deux possibilités suivantes
\begin{enumerate}
\item [(a)] Si $A^*$ n'est pas contenu dans $R$, alors
$\Spec(A) \to \Spec(R)$ et $\Spec(B) \to \Spec(R)$ sont
des immersions ouvertes recouvrant $\Spec(R)$ et $K = A
\otimes_R B$. \item [(b)] Si $A^*$ est contenu dans $R$,
alors $B$ domine l'un des anneaux locaux de $A$ en un
idéal maximal et $A \otimes B \to K$ n'est pas surjectif.
\end{enumerate}
\end{lemma}

\begin{proof}
L'hypothèse (a) implique qu'il existe une unité $u$ de $A$ dont
l'image dans $K$ se situe dans l'idéal maximal de $B$. Alors $u$
est un non-diviseur de zéro de $R$ et pour chaque $a \in A$ il
existe un $n$ tel que $u^n a \in R$. Il s'ensuit que $A = R_u$.

\medskip\noindent
Supposons que $\mathfrak m_A$ soit le radical de Jacobson de $A$.
Supposons que $x \in \mathfrak m_A$ soit un élément non nul. Puisque
$\dim(A) = 1$, nous voyons que $K = A_x$. Après avoir remplacé $x$ par $x^n
u^m$ pour certains $n \geq 1$ et $m \in \mathbf{Z}$, nous pouvons supposer que
$x$ se mappe sur une unité de $B$. Nous voyons que pour chaque $b \in B$ nous
avons que $x^nb$ est dans l'image de $R$ pour certains $n$. Ainsi $B = R_x$.

\medskip\noindent
Soit $z \in R$. Si $z \not \in \mathfrak m_A$ et $z$ ne se mappent pas
sur un élément de $\mathfrak m_B$, alors $z$ est inversible. Ainsi, $x
+ u$ est inversible dans $R$. Par conséquent, $\Spec(R) = D(x) \cup
D(u)$. Nous avons vu ci-dessus que $D(u) = \Spec(A)$ et $D(x) = \Spec(B)$.

\medskip\noindent
Cas (b). Si $x \in \mathfrak m_A$, alors $1 + x$ est une
unité et donc $1 + x \in R$, c'est-à-dire $x \in R$. Ainsi,
nous voyons que $\mathfrak m_A \subset R \subset A$. En
fait, dans ce cas, $A$ est intégral sur $R$. À savoir,
écrivez $A/\mathfrak m_A = \kappa_1 \times \ldots \times
\kappa_n$ comme un produit de corps. Disons que $x = (c_1,
\ldots, c_r, 0, \ldots, 0)$ est un élément avec $c_i \not = 0$. Alors
$$
x^2 - x(c_1, \ldots, c_r, 1, \ldots, 1)  = 0
$$
Puisque $R$ contient toutes les unités, nous voyons que
$A/\mathfrak m_A$ est intégral sur l'image de $R$ en lui,
et donc $A$ est intégral sur $R$. Il s'ensuit que $R \subset
A \subset B$ comme $B$ est intégralement clos. De plus, si
$x \in \mathfrak m_A$ est non nul, alors $K = A_x = \bigcup
x^{-n}A = \bigcup x^{-n}R$. Par conséquent $x^{-1} \not \in
B$, c'est-à-dire $x \in \mathfrak m_B$. Nous concluons
$\mathfrak m_A \subset \mathfrak m_B$. Ainsi, $A \cap
\mathfrak m_B$ est un idéal maximal de $A$, terminant ainsi la preuve.
\end{proof}

\begin{lemma}
\label{lemma-semi-local-dimension-one-conductor}
Supposons que $B$ soit un domaine noethérien semi-local de
dimension $1$. Supposons que $B'$ soit l'adhérence
intégrale de $B$ dans son corps des fractions. Alors $B'$
est un domaine de Dedekind semi-local. Supposons que $x$
soit un élément non nul du radical de Jacobson de $B'$. Alors
pour chaque $y \in B'$ il existe un $n$ tel que $x^n y \in B$.
\end{lemma}

\begin{proof}
Supposons que $\mathfrak m_B$ soit le radical de Jacobson
de $B$. La structure de $B'$ résulte de l'Algèbre, Lemme
\ref{algebra-lemma-integral-closure-Dedekind}. Étant
donné $x, y \in B'$ comme dans l'énoncé du lemme,
considérons le sous-anneau $B \subset A \subset B'$
engendré par $x$ et $y$. Alors $A$ est fini sur $B$
(Algèbre, Lemme
\ref{algebra-lemma-characterize-finite-in-terms-of-integral}).
Puisque le corps des fractions de $B$ et $A$ est le même, nous
voyons que le module fini $A/B$ est supporté sur l'ensemble des
points fermés de $B$. Ainsi $\mathfrak m_B^n A \subset B$ pour
un $n$ approprié. De plus, $\Spec(B') \to \Spec(A)$ est
surjectif (Algèbre, Lemme
\ref{algebra-lemma-integral-overring-surjective}), donc $A$ est également
semi-local. Il en découle également que $x$ est dans le radical de Jacobson
$\mathfrak m_A$ de $A$. Notez que $\mathfrak m_A = \sqrt{\mathfrak m_B A}$.
Ainsi $x^m y \in \mathfrak m_B A$ pour un certain $m$. Puis $x^{nm} y \in B$.
\end{proof}

\begin{lemma}
\label{lemma-semi-local-both-side}
Dans la situation \ref{situation-glue}, supposez
\begin{enumerate}
\item $A$ est un domaine semi-local noethérien de
dimension $1$, \item $B$ est un domaine semi-local
noethérien de dimension $1$, \item $A \otimes B \to K$ est surjectif.
\end{enumerate}
Alors $\Spec(A) \to \Spec(R)$ et $\Spec(B) \to \Spec(R)$ sont des
immersions ouvertes couvrant $\Spec(R)$ et $K = A \otimes_R B$.
\end{lemma}

\begin{proof}
Cas particulier : $B$ est intégralement clos dans $K$.
Cela signifie que $B$ est un anneau de Dedekind (Algèbre,
Lemme \ref{algebra-lemma-characterize-Dedekind}) d'où toutes
ses localisations en des idéaux maximaux sont des anneaux
de valuation discrète. Supposons que $\mathfrak m_1, \ldots,
\mathfrak m_r$ soit les idéaux maximaux de $B$. Nous posons
$$
R_1 = A \times_K B_{\mathfrak m_1}
$$
En observant que $A \otimes_{R_1} B_{\mathfrak m_1} \to K$ est surjectif,
nous concluons d'après le Lemme \ref{lemma-semi-local} que $A$ et
$B_{\mathfrak m_1}$ définissent des sous-schémas ouverts couvrant $\Spec(R_1)$
et que $K = A \otimes_{R_1} B_{\mathfrak m_1}$. En particulier, $R_1$ est un
anneau noethérien semi-local de dimension $1$. Par induction, nous définissons
$$
R_{i + 1} = R_i \times_K B_{\mathfrak m_{i + 1}}
$$
pour $i = 1, \ldots, r - 1$. Observez que $R = R_r$ parce
que $B = B_{\mathfrak m_1} \cap \ldots \cap B_{\mathfrak
m_r}$ (voir Algèbre, Lemme
\ref{algebra-lemma-normal-domain-intersection-localizations-height-1}).
Il découle de la procédure inductive que $R \to A$ définit
une immersion ouverte $\Spec(A) \to \Spec(R)$. D'autre part,
les idéaux maximaux $\mathfrak n_i$ de $R$ ne se trouvant pas
dans cet ouvert correspondent aux idéaux maximaux $\mathfrak
m_i$ de $B$ et en fait le homomorphisme d'anneaux $R \to B$
définit un isomorphisme $R_{\mathfrak n_i} \to B_{\mathfrak
m_i}$ (détails omis ; indice : à chaque étape nous avons ajouté
exactement un idéal maximal à $\Spec(R_i)$). Il s'ensuit que
$\Spec(B) \to \Spec(R)$ est une immersion ouverte comme souhaité.

\medskip\noindent
Cas général. Supposons que $B' \subset K$ soit l'adhérence
intégrale de $B$. Voir le Lemme
\ref{lemma-semi-local-dimension-one-conductor}. Alors le cas
spécial s'applique à $R' = A \times_K B'$. Choisissez $x \in R'$
qui n'est pas contenu dans les idéaux maximaux de $A$ et qui est
contenu dans les idéaux maximaux de $B'$ (voir Algèbre, Lemme
\ref{algebra-lemma-chinese-remainder}). D'après le Lemme
\ref{lemma-semi-local-dimension-one-conductor}, il existe un
entier $n$ tel que $x^n \in R = A \times_K B$. Remplacez $x$ par
$x^n$ afin que $x \in R$. Pour chaque $y \in R'$, il existe un
entier $n$ tel que $x^n y \in R$. D'autre part, il est clair que
$R'_x = A$. Ainsi $R_x = A$. En échangeant les rôles de $A$ et
$B$, nous trouvons également un $y \in R$ tel que $B = R_y$. Notez
que l'inversion de $x$ et $y$ ne laisse aucun premier sauf $(0)$.
Ainsi $K = R_{xy} = R_x \otimes_R R_y$. Cela termine la démonstration.
\end{proof}

\begin{lemma}
\label{lemma-glue-a-bunch-of-local-rings}
Supposons que $K$ soit un corps. Supposons que $A_1, \ldots, A_r
\subset K$ soit des semi-anneaux locaux noethériens de dimension $1$
avec corps des fractions $K$. Si $A_i \otimes A_j \to K$ est surjective
pour tout $i \not = j$, alors il existe un domaine semi-local noethérien
$A \subset K$ de dimension $1$ contenu dans $A_1, \ldots, A_r$ tel que
\begin{enumerate}
\item $A \to A_i$ induit une immersion ouverte $j_i :
\Spec(A_i) \to \Spec(A)$, \item $\Spec(A)$ est l'union
des ouverts $j_i(\Spec(A_i))$, \item chaque point fermé de
$\Spec(A)$ se trouve exactement dans l'un de ces ouverts.
\end{enumerate}
\end{lemma}

\begin{proof}
À savoir, nous pouvons prendre $A = A_1 \cap \ldots \cap
A_r$. D'abord, nous remarquons que (3), une fois que (1) et
(2) ont été prouvés, découle de l'hypothèse que $A_i \otimes
A_j \to K$ est surjectif puisque si $\mathfrak m \in
j_i(\Spec(A_i)) \cap j_j(\Spec(A_j))$, alors $A_i \otimes A_j
\to K$ se retrouve dans $A_\mathfrak m$. Pour prouver (1) et
(2), nous procédons par induction sur $r$. Si $r > 1$, par
induction, nous avons les résultats (1) et (2) pour $B = A_2 \cap
\ldots \cap A_r$. Ensuite, nous appliquons le Lemme
\ref{lemma-semi-local-both-side} pour voir qu'ils tiennent pour $A = A_1 \cap B$.
\end{proof}

\begin{lemma}
\label{lemma-create-globally-generated}
Supposons que $A$ soit un domaine avec corps des fractions $K$. Supposons que
$B_1, \ldots, B_r \subset K$ soit des domaines semi-locaux noethériens
$1$-dimensionnels dont les corps des fractions sont $K$. Si $A \otimes B_i \to
K$ sont surjectifs pour $i = 1, \ldots, r$, alors il existe un $x \in A$ tel
que $x^{-1}$ soit dans le radical de Jacobson de $B_i$ pour $i = 1, \ldots, r$.
\end{lemma}

\begin{proof}
Supposons que $B_i'$ soit l'adhérence intégrale de $B_i$ dans $K$.
Supposons que nous trouvions un $x \in A$ non nul tel que $x^{-1}$ soit
dans le radical de Jacobson de $B'_i$ pour $i = 1, \ldots, r$. Alors,
d'après le Lemme \ref{lemma-semi-local-dimension-one-conductor}, après
avoir remplacé $x$ par une puissance, nous obtenons $x^{-1} \in B_i$.
Comme $\Spec(B'_i) \to \Spec(B_i)$ est surjectif, nous voyons que
$x^{-1}$ est alors également dans le radical de Jacobson de $B_i$. Ainsi,
nous pouvons supposer que chaque $B_i$ est un domaine de Dedekind semi-local.

\medskip\noindent
Si $B_i$ n'est pas local, alors retirez $B_i$ de la liste
et ajoutez de nouveau la collection finie d'anneaux locaux
$(B_i)_\mathfrak m$. Ainsi, nous pouvons supposer que $B_i$
est un anneau de valuation discrète pour $i = 1, \ldots, r$.

\medskip\noindent
Supposons que $v_i : K \to \mathbf{Z}$, $i = 1, \ldots, r$
soient les valuations discrètes correspondantes (voir
Algèbre, Lemme \ref{algebra-lemma-characterize-Dedekind}).
Nous cherchons un $x \in A$ non nul avec $v_i(x) < 0$ pour $i
= 1, \ldots, r$. Nous allons le prouver par induction sur $r$.

\medskip\noindent
Si $r = 1$ et que le résultat est incorrect, alors $A \subset B$ et la
application $A \otimes B \to K$ n'est pas surjective, contradiction.

\medskip\noindent
Si $r > 1$, alors par induction nous pouvons trouver un $x \in A$ non nul
tel que $v_i(x) < 0$ pour $i = 1, \ldots, r - 1$. Si $v_r(x) < 0$, alors nous
avons terminé, donc nous pouvons supposer $v_r(x) \geq 0$. Par le cas de
base, nous pouvons trouver $y \in A$ non nul tel que $v_r(y) < 0$. Après avoir
remplacé $x$ par une puissance, nous pouvons supposer que $v_i(x) < v_i(y)$
pour $i = 1, \ldots, r - 1$. Alors $x + y$ est l'élément que nous recherchons.
\end{proof}

\begin{lemma}
\label{lemma-power-equal}
Supposons que $A$ soit un anneau local noethérien de
dimension $1$. Soit $L = \prod A_\mathfrak p$ où le
produit est pris sur les idéaux premiers minimaux de $A$.
Soit $a_1, a_2 \in \mathfrak m_A$ qui se projette sur le même
élément de $L$. Alors $a_1^n = a_2^n$ pour un certain $n > 0$.
\end{lemma}

\begin{proof}
Écrivez $a_1 = a_2 + x$. Ensuite, $x$ se transforme en zéro
dans $L$. Par conséquent, $x$ est un élément nilpotent de $A$
parce que $\bigcap \mathfrak p$ est le radical de $(0)$ et que
l’annihilateur $I$ de $x$ contient une puissance de l’idéal
maximal parce que $\mathfrak p \not \in V(I)$ pour tous les
premiers minimaux. Dites $x^k = 0$ et $\mathfrak m^n \subset I$. Ensuite
$$
a_1^{k + n} = a_2^{k + n} + {n + k \choose 1} a_2^{n + k - 1} x +
{n + k \choose 2} a_2^{n + k - 2} x^2 + \ldots +
{n + k \choose k - 1} a_2^{n + 1} x^{k - 1} = a_2^{n + k}
$$
à cause de $a_2 \in \mathfrak m_A$.
\end{proof}

\begin{lemma}
\label{lemma-power-works}
Supposons que $A$ soit un anneau local noethérien de dimension
$1$. Supposons que $L = \prod A_\mathfrak p$ et $I = \bigcap
\mathfrak p$ où le produit et l'intersection sont faits sur les
idéaux premiers minimaux de $A$. Soit $f \in L$ un élément de
la forme $f = i + a$ où $a \in \mathfrak m_A$ et $i \in IL$.
Alors une certaine puissance de $f$ est dans l'image de $A \to L$.
\end{lemma}

\begin{proof}
Puisque $A$ est noethérien, nous avons $I^t = 0$ pour certains
$t > 0$. Supposons que nous sachions que $f = a + i$ avec $i
\in I^kL$. Alors $f^n = a^n + na^{n - 1}i \bmod I^{k + 1}L$.
Ainsi, il suffit de montrer que $na^{n - 1}i$ est dans l'image
de $I^k \to I^kL$ pour certains $n \gg 0$. Pour voir cela,
choisissez un $g \in A$ tel que $\mathfrak m_A = \sqrt{(g)}$
(Algèbre, Lemme \ref{algebra-lemma-height-1}). Alors $L = A_g$
par exemple selon Algèbre, Proposition
\ref{algebra-proposition-dimension-zero-ring}. D'un autre côté, il existe
un $n$ tel que $a^n \in (g)$. Ainsi, nous pouvons éliminer les
dénominateurs pour les éléments de $L$ en multipliant par une puissance élevée de $a$.
\end{proof}

\begin{lemma}
\label{lemma-good-intersection}
Supposons que $A$ soit un anneau local noethérien de dimension $1$.
Supposons que $L = \prod A_\mathfrak p$ où le produit est pris sur les
premiers minimaux de $A$. Soit $K \to L$ un homomorphisme d'anneaux
intègre. Alors il existe $a \in \mathfrak m_A$ et $x \in K$ qui se
projettent sur le même élément de $L$ tel que $\mathfrak m_A = \sqrt{(a)}$.
\end{lemma}

\begin{proof}
D'après le Lemme \ref{lemma-power-works}, nous pouvons remplacer
$A$ par $A/(\bigcap \mathfrak p)$ et supposer que $A$ est un
anneau réduit (certains détails sont omis). Nous pouvons
également remplacer $K$ par l'image de $K \to L$. Alors $K$ est un
anneau réduit. L'application $\Spec(L) \to \Spec(K)$ est
surjective et fermée (détails omis). Par conséquent, $\Spec(K)$ est un
espace discret fini. Il s'ensuit que $K$ est un produit fini de corps.

\medskip\noindent
Supposons que $\mathfrak p_j$, $j = 1, \ldots, m$ soient
les idéaux premiers minimaux de $A$. Soit $L_j$ le corps
des fractions de $A_j$ de sorte que $L = \prod_{j = 1,
\ldots, m} L_j$. Supposons que $A_j$ soit la
normalisation de $A/\mathfrak p_j$. Alors $A_j$ est un
domaine de Dedekind semi-local avec au moins un idéal
maximal, voir Algèbre, Lemme
\ref{algebra-lemma-integral-closure-Dedekind}. Supposons que $n$ soit la
somme des nombres d'idéaux maximaux dans $A_1, \ldots, A_m$. Pour un
tel idéal maximal $\mathfrak m \subset A_j$, nous considérons la fonction
$$
v_{\mathfrak m} : L \to L_j \to \mathbf{Z} \cup \{\infty\}
$$
où la deuxième flèche est la valuation discrète
correspondant à l'anneau de valuation discrète
$(A_j)_{\mathfrak m}$ étendu en mappant $0$ sur $\infty$. De
cette manière, nous obtenons des fonctions $n$ $v_1, \ldots,
v_n : L \to \mathbf{Z} \cup \{\infty\}$. Nous trouverons un
élément $x \in K$ tel que $v_i(x) < 0$ pour tout $i = 1, \ldots, n$.

\medskip\noindent
Tout d'abord, nous affirmons que pour chaque $i$, il existe
un élément $x \in K$ avec $v_i(x) < 0$. À savoir, supposons
que $v_i$ corresponde à $\mathfrak m \subset A_j$. Si $v_i(x)
\geq 0$ pour tout $x \in K$, alors $K$ se mappe dans
$(A_j)_{\mathfrak m}$ à l'intérieur du corps des fractions
$L_j$ de $A_j$. L'image de $K$ dans $L_j$ est un corps sur $L_j$
et est algébrique par Algèbre, Lemme
\ref{algebra-lemma-integral-under-field}. Combiné, nous obtenons une
contradiction avec Algèbre, Lemme \ref{algebra-lemma-valuation-ring-cap-field-finite}.

\medskip\noindent
Supposons que nous ayons trouvé un élément $x \in K$ tel que
$v_1(x) < 0, \ldots, v_r(x) < 0$ pour un certain $r < n$. Si $v_{r
+ 1}(x) < 0$, alors $x$ fonctionne pour $r + 1$. Sinon, choisissez
un certain $y \in K$ avec $v_{r + 1}(y) < 0$ comme il est possible
selon le résultat du paragraphe précédent. Après avoir remplacé $x$
par $x^n$ pour un certain $n > 0$, nous pouvons supposer $v_i(x) <
v_i(y)$ pour $i = 1, \ldots, r$. Ensuite, $v_j(x + y) = v_j(x) < 0$
pour $j = 1, \ldots, r$ par les propriétés des valuations et de même
$v_{r + 1}(x + y) = v_{r + 1}(y) < 0$. En raisonnant par induction,
nous trouvons $x \in K$ avec $v_i(x) < 0$ pour $i = 1, \ldots, n$.

\medskip\noindent
En particulier, l'élément $x \in K$ a une projection non
nulle dans chaque facteur de $K$ (rappelons que $K$ est un
produit fini de corps et si un certain composant de $x$ était
nul, alors l'une des valeurs de $v_i(x)$ serait $\infty$).
Par conséquent, $x$ est inversible et $x^{-1} \in K$ est un
élément avec $\infty > v_i(x^{-1}) > 0$ pour tous les $i$. Il
découle du Lemme
\ref{lemma-semi-local-dimension-one-conductor} que pour un certain
$e < 0$, l'élément $x^e \in K$ se mappe sur un élément de
$\mathfrak m_A/\mathfrak p_j \subset A/\mathfrak p_j$ pour tous les
$j = 1, \ldots, m$. On observe que le conoyau de l'application
$\mathfrak m_A \to \prod \mathfrak m_A/\mathfrak p_j$ est annihilé
par une puissance de $\mathfrak m_A$. Ainsi, après avoir remplacé $e$
par un $e$ plus négatif, nous trouvons un élément $a \in \mathfrak
m_A$ dont l'image dans $\mathfrak m_A/\mathfrak p_j$ est égale à
l'image de $x^e$. La paire $(a, x^e)$ satisfait les conclusions du lemme.
\end{proof}

\begin{lemma}
\label{lemma-localization-semi-local}
Supposons que $A$ soit un anneau. Supposons que $\mathfrak p_1,
\ldots, \mathfrak p_r$ soit un ensemble fini de a premiers de $A$. Soit
$S = A \setminus \bigcup \mathfrak p_i$. Alors $S$ est un système
multiplicatif et $S^{-1}A$ est un semi-anneau local dont les idéaux
maximaux correspondent aux éléments maximaux de l'ensemble $\{\mathfrak p_i\}$.
\end{lemma}

\begin{proof}
Si $a, b \in A$ et $a, b \in S$, alors $a, b \not \in
\mathfrak p_i$ donc $ab \not \in \mathfrak p_i$, donc
$ab \in S$. Aussi $1 \in S$. Ainsi $S$ est un
sous-ensemble multiplicatif de $A$. D'après la
description de $\Spec(S^{-1}A)$ en Algèbre, le Lemme
\ref{algebra-lemma-spec-localization} et d'après Algèbre,
le Lemme \ref{algebra-lemma-silly}, nous voyons que les
nombres premiers de $S^{-1}A$ correspondent aux nombres
premiers de $A$ contenus dans l'un des $\mathfrak p_i$.
Ainsi les idéaux maximaux de $S^{-1}A$ correspondent
bijectivement aux éléments maximaux (p. ex. pour l'inclusion
\ ) de l'ensemble $\{\mathfrak p_1, \ldots, \mathfrak p_r\}$.
\end{proof}




\section{Schémas noethériens de dimension un}
\label{section-dimension-one}

\noindent
Le résultat principal de cette section est qu'un
schéma noethérien séparé de dimension $1$ possède un
faisceau inversible ample. Voir la Proposition
\ref{proposition-dim-1-noetherian-separated-has-ample}.

\begin{lemma}
\label{lemma-affine}
Supposons que $X$ soit un schéma dont tous les anneaux locaux sont
noethériens de dimension $\leq 1$. Supposons que $U \subset X$ soit un ouvert
rétrocompact. Notons $j : U \to X$ le morphisme d'inclusion. Alors
$R^pj_*\mathcal{F} = 0$, $p > 0$ pour tout $\mathcal{O}_U$-module quasi-cohérent $\mathcal{F}$.
\end{lemma}

\begin{proof}
Nous pouvons vérifier l'annulation de $R^pj_*\mathcal{F}$
aux germes. La formation de $R^qj_*$ commute avec le
changement de base plat, voir Cohomologie des schémas, Lemme
\ref{coherent-lemma-flat-base-change-cohomology}. Ainsi,
nous pouvons supposer que $X$ est le spectre d'un anneau
local noethérien de dimension $\leq 1$. Dans ce cas, $X$ a un
point fermé $x$ et finitement d'autres points $x_1, \ldots,
x_n$ qui se spécialisent en $x$ mais pas l'un l'autre (voir
Algèbre, Lemme
\ref{algebra-lemma-Noetherian-irreducible-components}). Si $x \in U$,
alors $U = X$ et le résultat est clair. Sinon, alors $U = \{x_1,
\ldots, x_r\}$ pour quelque $r$ après avoir éventuellement renuméroté
les points. Alors $U$ est affine (schémas, Lemme
\ref{schemes-lemma-scheme-finite-discrete-affine}). Ainsi, le résultat découle
de Cohomologie des schémas, Lemme \ref{coherent-lemma-relative-affine-vanishing}.
\end{proof}

\begin{lemma}
\label{lemma-open-in-affine-curve-affine}
Supposons que $X$ soit un schéma affine dont tous les anneaux locaux sont
noethériens de dimension $\leq 1$. Alors tout ouvert quasi-compact $U \subset X$ est affine.
\end{lemma}

\begin{proof}
Notons $j : U \to X$ le morphisme d'inclusion.
Supposons que $\mathcal{F}$ soit un
$\mathcal{O}_U$-module quasi-cohérent. D'après le
Lemme \ref{lemma-affine}, les images directes
supérieures $R^pj_*\mathcal{F}$ sont nulles. Le
$\mathcal{O}_X$-module $j_*\mathcal{F}$ est
quasi-cohérent (schémas, Lemme
\ref{schemes-lemma-push-forward-quasi-coherent}). Par
conséquent, il a des groupes de cohomologie supérieurs nuls
d'après Cohomologie des schémas, Lemme
\ref{coherent-lemma-quasi-coherent-affine-cohomology-zero}. Par
la suite spectrale de Leray Cohomologie, Lemme
\ref{cohomology-lemma-apply-Leray}, nous avons $H^p(U, \mathcal{F}) = 0$
pour tout $p > 0$. Ainsi $U$ est affine, par exemple d'après Cohomologie
des schémas, Lemme \ref{coherent-lemma-quasi-compact-h1-zero-covering}.
\end{proof}

\begin{lemma}
\label{lemma-complement-codim-1-closed-points}
Supposons que $X$ soit un schéma. Supposons que $U \subset X$ soit un ouvert. Supposons
\begin{enumerate}
\item $U$ est une ouverture rétrocompacte de
$X$, \item $X \setminus U$ est discret, et \item
pour $x \in X \setminus U$ l'anneau local
$\mathcal{O}_{X, x}$ est noethérien de dimension $\leq 1$.
\end{enumerate}
Alors (1) il existe un $\mathcal{O}_X$-module inversible
$\mathcal{L}$ et une section $s$ telles que $U = X_s$ et
(2) l'application $\Pic(X) \to \Pic(U)$ est surjective.
\end{lemma}

\begin{proof}
Supposons que $X \setminus U = \{x_i; i \in I\}$. Choisissez
des ouverts affines $U_i \subset X$ avec $x_i \in U_i$ et
$x_j \not \in U_i$ pour $j \not = i$. Cela est possible
selon la condition (2). Disons $U_i = \Spec(A_i)$. Soit
$\mathfrak m_i \subset A_i$ l'idéal maximal correspondant à
$x_i$. D'après notre hypothèse sur les anneaux locaux, il
n'existe qu'un nombre fini d'idéaux premiers $\mathfrak q
\subset \mathfrak m_i$, $\mathfrak q \not = \mathfrak m_i$
(voir Algèbre, Lemme
\ref{algebra-lemma-Noetherian-irreducible-components}). Ainsi, par
l'évitement des premiers (Algèbre, Lemme
\ref{algebra-lemma-silly}), nous pouvons trouver $f_i \in \mathfrak m_i$
non contenu dans aucun de ces premiers. Alors $V(f_i) = \{\mathfrak
m_i\} \amalg Z_i$ pour un certain fermé $Z_i \subset U_i$ car $Z_i$ est
un ouvert rétrocompact de $V(f_i)$ fermé par spécialisation, voir
Algèbre, Lemme
\ref{algebra-lemma-constructible-stable-specialization-closed}. Après avoir restreint $U_i$, nous pouvons supposer $V(f_i) = \{x_i\}$. Puis
$$
\mathcal{U} : X = U \cup \bigcup U_i
$$
est un recouvrement ouvert de $X$. Considérons le $2$-cocycle
à valeurs dans $\mathcal{O}_X^*$ donné par $f_i$ sur $U \cap
U_i$ et par $f_i/f_j$ sur $U_i \cap U_j$. Cela définit un fibré
en droites $\mathcal{L}$ tel que la section $s$ définie par
$1$ sur $U$ et $f_i$ sur $U_i$ est comme dans l'énoncé du lemme.

\medskip\noindent
Supposons que $\mathcal{N}$ soit un $\mathcal{O}_U$-module
inversible. Supposons que $N_i$ soit le module $(A_i)_{f_i}$
inversible tel que $\mathcal{N}|_{U \cap U_i}$ soit égal à
$\tilde N_i$. Observons que $(A_{\mathfrak m_i})_{f_i}$ est
un anneau artinien (comme un anneau noethérien de dimension
zéro, voir Algèbre, Lemme
\ref{algebra-lemma-Noetherian-dimension-0}). Ainsi, c'est un
produit d'anneaux locaux (Algèbre, Lemme
\ref{algebra-lemma-artinian-finite-length}) et possède donc un groupe
de Pic trivial. Ainsi, après avoir réduit $U_i$ (c'est-à-dire, après
avoir remplacé $A_i$ par $(A_i)_g$ pour certains $g \in A_i$, $g \not
\in \mathfrak m_i$), nous pouvons supposer que $N_i = (A_i)_{f_i}$,
c'est-à-dire que $\mathcal{N}|_{U \cap U_i}$ est trivial. Dans ce cas,
il est clair comment étendre $\mathcal{N}$ à un faisceau inversible sur
$X$ (en l'étendant par un module inversible trivial sur chaque $U_i$).
\end{proof}

\begin{lemma}
\label{lemma-find-globally-generated}
Supposons que $X$ soit un schéma séparé intègre. Supposons que $U \subset
X$ soit un ouvert affine non vide tel que $X \setminus U$ soit un
ensemble fini de points $x_1, \ldots, x_r$ avec $\mathcal{O}_{X, x_i}$
noethérien de dimension $1$. Alors il existe un $\mathcal{O}_X$-module
inversible globalement engendré $\mathcal{L}$ et une section $s$ tels que $U = X_s$.
\end{lemma}

\begin{proof}
Dites $U = \Spec(A)$ et laissez $K$ être le corps des fonctions de
$X$. Écrivez $B_i = \mathcal{O}_{X, x_i}$ et $\mathfrak m_i =
\mathfrak m_{x_i}$. Puisque $x_i \not \in U$ nous voyons que
l'ouvert $U \times_X \Spec(B_i)$ de $\Spec(B_i)$ n'a qu'un seul
point, c'est-à-dire $U \times_X \Spec(B_i) = \Spec(K)$. Puisque $X$
est séparé, nous trouvons que $\Spec(K)$ est un sous-schéma fermé de
$U \times \Spec(B_i)$, c'est-à-dire que l'application $A \otimes B_i
\to K$ est une surjection. Par le Lemme
\ref{lemma-create-globally-generated}, nous pouvons trouver un $f \in A$ non
nul tel que $f^{-1} \in \mathfrak m_i$ pour $i = 1, \ldots, r$. Choisissez des
ouverts $x_i \in U_i \subset X$ tels que $f^{-1} \in \mathcal{O}(U_i)$. Alors
$$
\mathcal{U} : X = U \cup \bigcup U_i
$$
est un recouvrement ouvert de $X$. Considérons le
$2$-cocycle à valeurs dans $\mathcal{O}_X^*$ donné par $f$
sur $U \cap U_i$ et par $1$ sur $U_i \cap U_j$. Cela
définit un fibré en droites $\mathcal{L}$ avec deux sections :
\begin{enumerate}
\item une section $s$ définie par $1$ sur
$U$ et $f^{-1}$ sur $U_i$ est comme dans
l'énoncé du lemme, et \item une section $t$
définie par $f$ sur $U$ et $1$ sur $U_i$.
\end{enumerate}
Notez que $X_t \supset U_1 \cup \ldots \cup U_r$. Par
conséquent, $s, t$ génère $\mathcal{L}$ et le lemme est prouvé.
\end{proof}

\begin{lemma}
\label{lemma-enough-globally-generated-ample}
Supposons que $X$ soit un schéma quasi-compact. Si pour
chaque $x \in X$ il existe une paire $(\mathcal{L}, s)$
consistant en un faisceau inversible globalement engendré
$\mathcal{L}$ et une section globale $s$ telle que $x \in X_s$ et
$X_s$ soit affine, alors $X$ possède un faisceau inversible ample.
\end{lemma}

\begin{proof}
Puisque $X$ est quasi-compact, nous pouvons trouver une
collection finie $(\mathcal{L}_i, s_i)$, $i = 1, \ldots, n$ de
paires telles que $\mathcal{L}_i$ est globalement engendré,
$X_{s_i}$ est affine et $X = \bigcup X_{s_i}$. Encore une fois,
parce que $X$ est quasi-compact, nous pouvons trouver, pour
chaque $i$, une collection finie de sections $t_{i, j}$ de
$\mathcal{L}_i$, $j = 1, \ldots, m_i$ telles que $X = \bigcup
X_{t_{i, j}}$. Posons $t_{i, 0} = s_i$. Considérons le faisceau inversible
$$
\mathcal{L} = \mathcal{L}_1 
\otimes_{\mathcal{O}_X} \ldots
\otimes_{\mathcal{O}_X} \mathcal{L}_n
$$
et les sections globales
$$
\tau_J = t_{1, j_1} \otimes \ldots \otimes t_{n, j_n}
$$
D'après les propriétés, le lemme
\ref{properties-lemma-affine-cap-s-open}, l'ouvert $X_{\tau_J}$ est affine
dès que $j_i = 0$ pour un certain $i$. Il est facile de voir que ces
ouverts recouvrent $X$. Par conséquent, $\mathcal{L}$ est ample par définition.
\end{proof}

\begin{lemma}
\label{lemma-dim-1-noetherian-integral-separated-has-ample}
Supposons que $X$ soit un schéma intègre séparé noethérien de
dimension $1$. Alors $X$ possède un faisceau inversible ample.
\end{lemma}

\begin{proof}
Choisissez un recouvrement ouvert affine $X = U_1 \cup \ldots
\cup U_n$. Comme $X$ est noethérien, chacun des ensembles $X
\setminus U_i$ est fini. Ainsi, d'après le Lemme
\ref{lemma-find-globally-generated}, nous pouvons trouver une paire
$(\mathcal{L}_i, s_i)$ composée d'un faisceau inversible globalement
engendré $\mathcal{L}_i$ et d'une section globale $s_i$ telle que
$U_i = X_{s_i}$. Nous en concluons que $X$ possède un faisceau
inversible ample d'après le Lemme \ref{lemma-enough-globally-generated-ample}.
\end{proof}

\begin{lemma}
\label{lemma-surjective-pic-birational-finite}
Supposons que $f : X \to Y$ soit un morphisme fini de
schémas. Supposons qu'il existe un ouvert $V \subset Y$ tel que
$f^{-1}(V) \to V$ soit un isomorphisme et $Y \setminus V$ soit
un espace discret. Alors tout module inversible $\mathcal{O}_X$
est le tiré en arrière d'un module inversible $\mathcal{O}_Y$.
\end{lemma}

\begin{proof}
Nous utiliserons ce $\Pic(X) = H^1(X, \mathcal{O}_X^*)$,
voir Cohomologie, Lemme
\ref{cohomology-lemma-h1-invertible}. Considérons la suite spectrale de
Leray pour le faisceau abélien $\mathcal{O}_X^*$ et $f$, voir
Cohomologie, Lemme \ref{cohomology-lemma-Leray}. Considérons l'application induite
$$
H^1(X, \mathcal{O}_X^*) \longrightarrow H^0(Y, R^1f_*\mathcal{O}_X^*)
$$
Diviseurs, le lemme
\ref{divisors-lemma-finite-trivialize-invertible-upstairs} dit exactement que
cette application est nulle. Par conséquent, Leray donne $H^1(X,
\mathcal{O}_X^*) = H^1(Y, f_*\mathcal{O}_X^*)$. Ensuite, nous considérons l'application
$$
f^\sharp : \mathcal{O}_Y^* \longrightarrow f_*\mathcal{O}_X^*
$$
Par hypothèse, le noyau et le conoyau de cette application
sont supportés sur le sous-ensemble fermé $T = Y \setminus
V$ de $Y$. Puisque $T$ est, par hypothèse, un espace
topologique discret, les groupes de cohomologie supérieurs de
tout faisceau abélien sur $Y$ supporté sur $T$ sont nuls
(cela découle de Cohomologie, Lemme
\ref{cohomology-lemma-cohomology-and-closed-immersions}, Modules,
Lemme \ref{modules-lemma-i-star-exact}, et du fait que $H^i(T,
\mathcal{F}) = 0$ pour tout $i > 0$ et tout faisceau abélien $\mathcal{F}$
sur $T$). En décomposant l'application affichée en suites exactes courtes
$$
0 \to \Ker(f^\sharp) \to \mathcal{O}_Y^* \to \Im(f^\sharp) \to 0,\quad
0 \to \Im(f^\sharp) \to f_*\mathcal{O}_X^* \to \Coker(f^\sharp) \to 0
$$
nous concluons d'abord que $H^1(Y, \mathcal{O}_Y^*) \to
H^1(Y, \Im(f^\sharp))$ est surjectif puis que $H^1(Y,
\Im(f^\sharp)) \to H^1(Y, f_*\mathcal{O}_X^*)$ est surjectif. En
combinant tout ce qui précède, nous trouvons que $H^1(Y,
\mathcal{O}_Y^*) \to H^1(X, \mathcal{O}_X^*)$ est surjectif comme souhaité.
\end{proof}

\begin{lemma}
\label{lemma-glue-invertible-sheaves}
Supposons que $X$ soit un schéma. Supposons que $Z_1, \ldots,
Z_n \subset X$ soit des sous-schémas fermés. Supposons que
$\mathcal{L}_i$ soit un faisceau inversible sur $Z_i$. Supposons que
\begin{enumerate}
\item $X$ est réduit, \item $X = \bigcup Z_i$ est
défini théoriquement, et \item $Z_i \cap Z_j$ est
un espace topologique discret pour $i \not = j$.
\end{enumerate}
Alors il existe un faisceau inversible $\mathcal{L}$ sur $X$
dont la restriction à $Z_i$ est $\mathcal{L}_i$. De plus, si
nous avons des sections $s_i \in \Gamma(Z_i, \mathcal{L}_i)$
qui ne s'annulent pas aux points de $Z_i \cap Z_j$, alors nous
pouvons choisir $\mathcal{L}$ de telle sorte qu'il existe un $s
\in \Gamma(X, \mathcal{L})$ avec $s|_{Z_i} = s_i$ pour tout $i$.
\end{lemma}

\begin{proof}
L'existence de $\mathcal{L}$ peut être déduite du Lemme
\ref{lemma-surjective-pic-birational-finite} mais nous donnerons
également une preuve directe et nous utiliserons la preuve directe
pour voir que l'énoncé concernant les sections est vrai. Posons $T =
\bigcup_{i \not = j} Z_i \cap Z_j$. Comme $X$ est réduit, nous avons
$$
X \setminus T = \bigcup (Z_i \setminus T)
$$
comme les schémas. L'hypothèse (3) implique que $T$ est un sous-ensemble
discret de $X$. Ainsi, pour chaque $t \in T$, nous pouvons trouver un ouvert
$U_t \subset X$ avec $t \in U_t$ mais $t' \not \in U_t$ pour $t' \in T$, $t'
\not = t$. En réduisant $U_t$ si nécessaire, nous pouvons supposer qu'il existe
des isomorphismes $\varphi_{t, i} : \mathcal{L}_i|_{U_t \cap Z_i} \to
\mathcal{O}_{U_t \cap Z_i}$. De plus, pour chaque $i$, choisissez un recouvrement ouvert
$$
Z_i \setminus T = \bigcup\nolimits_j U_{ij}
$$
tel qu'il existe des isomorphismes
$\varphi_{i, j} : \mathcal{L}_i|_{U_{ij}}
\cong \mathcal{O}_{U_{ij}}$. Observez que
$$
\mathcal{U} : X = \bigcup U_t \cup \bigcup U_{ij}
$$
est un recouvrement ouvert de $X$. Nous affirmons que nous pouvons
utiliser les isomorphismes $\varphi_{t, i}$ et $\varphi_{i, j}$
pour définir un $2$-cocycle à valeurs dans $\mathcal{O}_X^*$ pour ce
recouvrement qui définit $\mathcal{L}$ comme dans l’énoncé du lemme.

\medskip\noindent
À savoir, si $i \not = i'$, alors $U_{i, j} \cap U_{i', j'}
= \emptyset$ et il n'y a rien à faire. Pour $U_{i, j} \cap
U_{i, j'}$ nous avons $\mathcal{O}_X(U_{i, j} \cap U_{i, j'})
= \mathcal{O}_{Z_i}(U_{i, j} \cap U_{i, j'})$ par la
première remarque de la démonstration. Ainsi, la fonction de
transition pour $\mathcal{L}_i$ (plus précisément $\varphi_{i,
j} \circ \varphi_{i, j'}^{-1}$) définit la valeur de notre
cocycle sur cette intersection. Pour $U_t \cap U_{i, j}$, nous
pouvons faire la même chose. Enfin, pour $t \not = t'$ nous avons
$$
U_t \cap U_{t'} = \coprod (U_t \cap U_{t'}) \cap Z_i
$$
et de plus, l'intersection $U_t \cap U_{t'} \cap Z_i$ est contenue dans $Z_i
\setminus T$. Par conséquent, par le même raisonnement qu'auparavant, nous voyons que
$$
\mathcal{O}_X(U_t \cap U_{t'}) =
\prod \mathcal{O}_{Z_i}(U_t \cap U_{t'} \cap Z_i)
$$
et nous pouvons utiliser les fonctions de transition pour
$\mathcal{L}_i$ (plus précisément $\varphi_{t, i} \circ \varphi_{t',
i}^{-1}$) pour définir la valeur de notre cocycle sur $U_t \cap
U_{t'}$. Cela termine la démonstration de l'existence de $\mathcal{L}$.

\medskip\noindent
Étant donné les sections $s_i$ comme dans la dernière assertion du lemme,
dans l'argument ci-dessus, nous choisissons $U_t$ de sorte que $s_i|_{U_t
\cap Z_i}$ ne soit pas nulle et nous choisissons $\varphi_{t, i}$ de
sorte que $\varphi_{t, i}(s_i|_{U_t \cap Z_i}) = 1$. Ensuite, en utilisant
$1$ sur $U_t$ et $\varphi_{i, j}(s_i|_{U_{i, j}})$ sur $U_{i, j}$, nous
définirons une section de $\mathcal{L}$ qui se restreint à $s_i$ sur $Z_i$.
\end{proof}

\begin{remark}
\label{remark-conductor}
Supposons que $A$ soit un anneau réduit. Supposons que $I, J$ soient
des idéaux de $A$ tels que $V(I) \cup V(J) = \Spec(A)$. Posons $B =
A/J$. Alors $I \to IB$ est un isomorphisme de $A$-modules. À savoir,
nous avons $IB = I + J/J = I/(I \cap J)$ et $I \cap J$ est nul parce
que $A$ est réduit et $\Spec(A) = V(I) \cup V(J) = V(I \cap J)$. Ainsi,
pour tout $A$-module projectif $P$, nous avons également $IP = I(P/JP)$.
\end{remark}

\begin{lemma}
\label{lemma-dim-1-noetherian-reduced-separated-has-ample}
Supposons que $X$ soit un schéma séparé réduit noethérien de
dimension $1$. Alors $X$ possède un faisceau inversible ample.
\end{lemma}

\begin{proof}
Supposons que $Z_i$, $i = 1, \ldots, n$ soient les composantes
irréductibles de $X$. Nous les considérons comme des sous-schémas fermés
réduits de $X$. D'après le Lemme
\ref{lemma-dim-1-noetherian-integral-separated-has-ample}, il existe des faisceaux
inversibles amples $\mathcal{L}_i$ sur $Z_i$. Posons $T = \bigcup_{i \not = j} Z_i
\cap Z_j$. Comme $X$ est noethérien de dimension $1$, l'ensemble $T$ est fini et
consiste en des points fermés de $X$. Pour chaque $i$, nous pouvons, éventuellement
après avoir remplacé $\mathcal{L}_i$ par une puissance, choisir $s_i \in \Gamma(Z_i,
\mathcal{L}_i)$ de telle sorte que $(Z_i)_{s_i}$ soit affine et contienne $T \cap
Z_i$, voir Propriétés, Lemme \ref{properties-lemma-ample-finite-set-in-principal-affine}.

\medskip\noindent
D'après le Lemme \ref{lemma-glue-invertible-sheaves}, nous pouvons
trouver un faisceau inversible $\mathcal{L}$ sur $X$ et $s \in
\Gamma(X, \mathcal{L})$ tel que $(\mathcal{L}, s)|_{Z_i} =
(\mathcal{L}_i, s_i)$. Remarquons que $X_s$ contient $T$ et est, en
termes d'ensembles, égal aux sous-schémas affines fermés
$(Z_i)_{s_i}$. Ainsi, il est affine d'après Limites, Lemme
\ref{limits-lemma-affines-glued-in-closed-affine}. Pour finir la
démonstration, il suffit de trouver, pour tout $x \in X$, $x \not \in T$,
un entier $m > 0$ et une section $t \in \Gamma(X, \mathcal{L}^{\otimes
m})$ tels que $X_t$ soit affine et $x \in X_t$. Comme $x \not \in T$,
nous voyons que $x \in Z_i$ pour un certain $i$ unique, disons $i = 1$.
Supposons que $Z \subset X$ soit le sous-schéma fermé réduit dont l'espace
topologique sous-jacent est $Z_2 \cup \ldots \cup Z_n$. Supposons que
$\mathcal{I} \subset \mathcal{O}_X$ soit le faisceau d'idéaux de $Z$.
Notons que $\mathcal{I}_1 \subset \mathcal{O}_{Z_1}$ est l'image inverse de
ce faisceau d'idéaux sous le morphisme d'inclusion $Z_1 \to X$. Remarquons que
$$
\Gamma(X, \mathcal{I}\mathcal{L}^{\otimes m}) =
\Gamma(Z_1, \mathcal{I}_1 \mathcal{L}_1^{\otimes m})
$$
voir Remarque \ref{remark-conductor}. Ainsi, il suffit
de trouver $m > 0$ et $t \in \Gamma(Z_1, \mathcal{I}_1
\mathcal{L}_1^{\otimes m})$ avec $x \in (Z_1)_t$
affine. Puisque $\mathcal{L}_1$ est ample et que $x$
n'est pas dans $Z_1 \cap T = V(\mathcal{I}_1)$, nous
pouvons trouver une section $t_1 \in \Gamma(Z_1,
\mathcal{I}_1 \mathcal{L}_1^{\otimes m_1})$ avec $x \in
(Z_1)_{t_1}$, voir Propriétés, Proposition
\ref{properties-proposition-characterize-ample}. Puisque
$\mathcal{L}_1$ est ample, nous pouvons trouver une
section $t_2 \in \Gamma(Z_1, \mathcal{L}_1^{\otimes m_2})$
avec $x \in (Z_1)_{t_2}$ et $(Z_1)_{t_2}$ affine, voir
Propriétés, Définition \ref{properties-definition-ample}.
Posons $m = m_1 + m_2$ et $t = t_1 t_2$. Alors $t \in
\Gamma(Z_1, \mathcal{I}_1 \mathcal{L}_1^{\otimes m})$ avec $x
\in (Z_1)_t$ par construction et $(Z_1)_t$ est affine d'après
Propriétés, Lemme \ref{properties-lemma-affine-cap-s-open}.
\end{proof}

\begin{lemma}
\label{lemma-lift-line-bundle-from-reduction-dimension-1}
Supposons que $i : Z \to X$ soit une immersion fermée de
schémas. Si l'espace topologique sous-jacent de $X$ est
noethérien et $\dim(X) \leq 1$, alors $\Pic(X) \to \Pic(Z)$ est surjectif.
\end{lemma}

\begin{proof}
Considérez la suite exacte courte
$$
0 \to (1 + \mathcal{I}) \cap \mathcal{O}_X^* \to
\mathcal{O}^*_X \to i_*\mathcal{O}^*_Z \to 0
$$
des faisceaux de groupes abéliens sur $X$ où $\mathcal{I}$
est le faisceau quasi-cohérent des idéaux correspondant à
$Z$. Puisque $\dim(X) \leq 1$ nous voyons que $H^2(X,
\mathcal{F}) = 0$ pour tout faisceau abélien $\mathcal{F}$,
voir Cohomologie, Proposition
\ref{cohomology-proposition-vanishing-Noetherian}. Par
conséquent, l'application $H^1(X, \mathcal{O}^*_X) \to H^1(X,
i_*\mathcal{O}_Z^*)$ est surjective. Par Cohomologie, Lemme
\ref{cohomology-lemma-cohomology-and-closed-immersions} nous avons
$H^1(X, i_*\mathcal{O}_Z^*) = H^1(Z, \mathcal{O}_Z^*)$. Cela prouve
le lemme par Cohomologie, Lemme \ref{cohomology-lemma-h1-invertible}.
\end{proof}

\begin{proposition}
\label{proposition-dim-1-noetherian-separated-has-ample}
Supposons que $X$ soit un schéma séparé noethérien de
dimension $1$. Alors $X$ possède un faisceau inversible ample.
\end{proposition}

\begin{proof}
Supposons que $Z \subset X$ soit la réduction de $X$. D'après
le Lemme
\ref{lemma-dim-1-noetherian-reduced-separated-has-ample}, le schéma
$Z$ possède un faisceau inversible ample. Ainsi, selon le Lemme
\ref{lemma-lift-line-bundle-from-reduction-dimension-1}, il existe un
$\mathcal{O}_X$-module inversible $\mathcal{L}$ sur $X$ dont la
restriction à $Z$ est ample. Alors $\mathcal{L}$ est ample par une application
de la cohomologie des schémas, Lemme \ref{coherent-lemma-ample-on-reduction}.
\end{proof}

\begin{remark}
\label{remark-useless-generalization}
En fait, si $X$ est un schéma dont la réduction est un
schéma séparé noethérien de dimension $1$, alors $X$
possède un faisceau inversible ample. L'argument pour
le prouver est le même que la démonstration de la
Proposition
\ref{proposition-dim-1-noetherian-separated-has-ample} excepté qu'on
utilise Limites, Lemme \ref{limits-lemma-ample-on-reduction} au lieu de
Cohomologie des schémas, Lemme \ref{coherent-lemma-ample-on-reduction}.
\end{remark}

\noindent
Le lemme suivant vaut en réalité pour les morphismes séparés
quasi-finis, comme le lecteur peut le constater en utilisant le
théorème principal de Zariski (Plus sur les morphismes, Lemme
\ref{more-morphisms-lemma-quasi-finite-separated-pass-through-finite})
et le Lemme \ref{lemma-complement-codim-1-closed-points}.

\begin{lemma}
\label{lemma-surjection-on-pic-quasi-finite}
Supposons que $f : X \to Y$ soit un morphisme de schémas. Supposons
que $Y$ soit noethérien de dimension $\leq 1$, $f$ soit fini, et
qu'il existe un ouvert dense $V \subset Y$ tel que $f^{-1}(V) \to V$
soit une immersion fermée. Alors tout $\mathcal{O}_X$-module
inversible est le tiré en arrière d'un $\mathcal{O}_Y$-module inversible.
\end{lemma}

\begin{proof}
Nous factorisons $f$ en tant que $X \to Z \to Y$ où $Z$
est l'image schématique de $f$. Ensuite, $X \to Z$ est un
isomorphisme sur $V \cap Z$ et le Lemme
\ref{lemma-surjective-pic-birational-finite} s'applique.
D'autre part, le Lemme
\ref{lemma-lift-line-bundle-from-reduction-dimension-1} s'applique à $Z \to Y$. Certains détails sont omis.
\end{proof}








\section{L'invariant delta}
\label{section-delta-invariant}

\noindent
Dans cette section, nous définissons l'invariant $\delta$ d'un
point singulier sur un schéma de Nagata réduit de dimension $1$.

\begin{lemma}
\label{lemma-pre-pre-delta-invariant}
Supposons que $(A, \mathfrak m)$ soit un anneau local $1$-dimensionnel
noethérien. Soit $f \in \mathfrak m$. Les énoncés suivants sont équivalents
\begin{enumerate}
\item $\mathfrak m = \sqrt{(f)}$, \item $f$ n'est contenu dans
aucun idéal premier minimal de $A$, et \item $A_f =
\prod_{\mathfrak p\text{ minimal}} A_\mathfrak p$ en tant qu'algèbres $A$.
\end{enumerate}
Un tel $f \in \mathfrak m$ existe. Si $\text{depth}(A) = 1$ (par
exemple $A$ est réduit), alors (1) -- (3) sont également équivalents à
\begin{enumerate}
\item [(4)] $f$ est un non-diviseur de zéro, \item
[(5)] $A_f$ est l'anneau total des fractions de $A$.
\end{enumerate}
Si $A$ est réduit, alors (1) -- (5) sont également équivalents à
\begin{enumerate}
\item [(6)] $A_f$ est le produit des corps
résiduels aux idéaux premiers minimaux de $A$.
\end{enumerate}
\end{lemma}

\begin{proof}
Le spectre de $A$ a un nombre fini de nombres
premiers $\mathfrak p_1, \ldots, \mathfrak p_n$ en
plus de $\mathfrak m$ et ceux-ci sont tous minimaux,
voir Algèbre, Lemme
\ref{algebra-lemma-Noetherian-irreducible-components}.
Alors, l'équivalence de (1) et (2) découle de Algèbre,
Lemme \ref{algebra-lemma-Zariski-topology}. Clairement,
(3) implique (2). Inversement, si (2) est vrai, alors
le spectre de $A_f$ est le sous-ensemble $\{\mathfrak
p_1, \ldots, \mathfrak p_n\}$ de $\Spec(A)$ avec la
topologie induite, voir Algèbre, Lemme
\ref{algebra-lemma-spec-localization}. C'est un espace
topologique discret fini. Donc $A_f = \prod_{\mathfrak
p\text{ minimal}} A_\mathfrak p$ par Algèbre, Proposition
\ref{algebra-proposition-dimension-zero-ring}. L'existence d'un
$f$ est affirmée dans Algèbre, Lemme \ref{algebra-lemma-height-1}.

\medskip\noindent
Supposons que $A$ a une profondeur $1$. (C'est le maximum
selon l'Algèbre, Lemme \ref{algebra-lemma-bound-depth} et
cela vaut si $A$ est réduit selon l'Algèbre, Lemme
\ref{algebra-lemma-criterion-reduced}.) Alors $\mathfrak m$
n'est pas un premier associé de $A$. Chaque premier
minimal de $A$ est un premier associé (Algèbre, Proposition
\ref{algebra-proposition-minimal-primes-associated-primes}).
Ainsi, l'ensemble des éléments non nuls-diviseurs de $A$
est exactement l'ensemble des éléments non contenus dans
aucun des premiers minimaux selon l'Algèbre, Lemme
\ref{algebra-lemma-ass-zero-divisors}. Ainsi, (4) est équivalent
à (2). La partie (5) est équivalente à (3) selon l'Algèbre,
Lemme \ref{algebra-lemma-total-ring-fractions-no-embedded-points}.

\medskip\noindent
Alors $A_\mathfrak p$ est un corps pour
$\mathfrak p \subset A$ minimal, voir Algèbre,
Lemme
\ref{algebra-lemma-minimal-prime-reduced-ring}. Par conséquent, (3) est équivalent à (6).
\end{proof}

\begin{lemma}
\label{lemma-pre-delta-invariant}
Supposons que $(A, \mathfrak m)$ soit un anneau local réduit de Nagata de
dimension $1$. Supposons que $A'$ soit la clôture intégrale de $A$ dans
l'anneau total des fractions de $A$. Alors $A'$ est un anneau de Nagata
normal, $A \to A'$ est fini, et $A'/A$ a une longueur finie comme $A$-module.
\end{lemma}

\begin{proof}
Le corps total des fractions est essentiellement de type fini
sur $A$, donc $A \to A'$ est fini car $A$ est Nagata, voir
Algèbre, Lemme
\ref{algebra-lemma-nagata-in-reduced-finite-type-finite-integral-closure}.
L'anneau $A'$ est normal par exemple par Algèbre, Lemme
\ref{algebra-lemma-characterize-reduced-ring-normal} et
\ref{algebra-lemma-Noetherian-irreducible-components}. L'anneau $A'$
est Nagata par exemple par Algèbre, Lemme
\ref{algebra-lemma-quasi-finite-over-nagata}. Choisissez $f \in \mathfrak
m$ comme dans le Lemme \ref{lemma-pre-pre-delta-invariant}. Comme $A'
\subset A_f$, il est clair que $A_f = A'_f$. Par conséquent, le support du
module $A$ fini $A'/A$ est contenu dans $\{\mathfrak m\}$. Il s'ensuit qu'il
a une longueur finie par Algèbre, Lemme \ref{algebra-lemma-support-point}.
\end{proof}

\begin{definition}
\label{definition-delta-invariant-algebra}
Soit $A$ un anneau local réduit de Nagata de dimension $1$.
L'{\it invariant $\delta$ de $A$ } est $\text{length}_A(A'/A)$, où
$A'$ est défini comme dans le lemme \ref{lemma-pre-delta-invariant}.
\end{definition}

\noindent
Nous prouvons quelques lemmes sur le comportement de cet invariant.

\begin{lemma}
\label{lemma-delta-invariant-is-zero}
Supposons que $A$ soit un anneau local de Nagata réduit
de dimension $1$. L'invariant $\delta$ de $A$ est $0$ si
et seulement si $A$ est un anneau de valuation discrète.
\end{lemma}

\begin{proof}
Si $A$ est un anneau de valuation discrète, alors $A$
est normal et l'anneau $A'$ est égal à $A$.
Réciproquement, si l'invariant $\delta$ de $A$ est $0$,
alors $A$ est intégralement clos dans son anneau total de
fractions, ce qui implique que $A$ est normal (Algèbre,
Lemme
\ref{algebra-lemma-characterize-reduced-ring-normal}) et cela force $A$ à être un
anneau de valuation discrète selon Algèbre, Lemme \ref{algebra-lemma-characterize-dvr}.
\end{proof}

\begin{lemma}
\label{lemma-normalization-same-after-completion}
Supposons que $A$ soit un anneau local réduit de Nagata de
dimension $1$. Supposons que $A \to A'$ soit comme dans le
Lemme \ref{lemma-pre-delta-invariant}. Supposons que $A^h$,
$A^{sh}$, resp. \ $A^\wedge$ soit la hensélisation, la
hensélisation stricte, resp. \ complétion de $A$. Alors $A^h$,
$A^{sh}$, resp. $A^\wedge$ est un anneau local réduit de Nagata
de dimension $1$ et $A' \otimes_A A^h$, $A' \otimes_A A^{sh}$,
resp. $A' \otimes_A A^\wedge$ est l'adhérence intégrale de $A^h$,
$A^{sh}$, resp. \ $A^\wedge$ dans son anneau total des fractions.
\end{lemma}

\begin{proof}
Remarquez que $A^\wedge$ est réduit, voir Plus sur
l'Algèbre, Lemme
\ref{more-algebra-lemma-completion-reduced}. Les anneaux
$A^h$ et $A^{sh}$ sont réduits par Plus sur l'Algèbre, Lemme
\ref{more-algebra-lemma-henselization-reduced}. Les
dimensions de $A$, $A^h$, $A^{sh}$ et $A^\wedge$ sont les mêmes
selon Plus sur l'Algèbre, Lemme
\ref{more-algebra-lemma-completion-dimension} et \ref{more-algebra-lemma-henselization-dimension}.

\medskip\noindent
Rappelez-vous qu'un anneau local noethérien est Nagata si
et seulement si les fibres formelles de $A$ sont
géométriquement réduites, voir Plus sur l'algèbre, Lemme
\ref{more-algebra-lemma-Nagata-local-ring}. Cette propriété
est héritée par $A^h$ et $A^{sh}$, voir le matériel dans
Plus sur l'algèbre, Section
\ref{more-algebra-section-properties-formal-fibres} et surtout Lemme
\ref{more-algebra-lemma-henselization-P-ring}. La complétion est Nagata
selon Algèbre, Lemme \ref{algebra-lemma-Noetherian-complete-local-Nagata}.

\medskip\noindent
Nous en venons maintenant à l’énoncé sur les
clôtures intégrales. Avant de continuer, choisissons
$f \in \mathfrak m$ comme dans le Lemme
\ref{lemma-pre-pre-delta-invariant}. Alors l’image de $f$
dans $A^h$, $A^{sh}$, et $A^\wedge$ est clairement un
élément satisfaisant les propriétés (1) -- (6) dans cet anneau.

\medskip\noindent
Puisque $A \to A'$ est fini, nous voyons que $A' \otimes_A
A^h$ et $A' \otimes_A A^{sh}$ sont le produit d'anneaux
locaux henséliens finis sur $A^h$ et $A^{sh}$, voir Algèbre,
Lemme \ref{algebra-lemma-finite-over-henselian}. Chacun de
ces anneaux locaux est l'hensélisation de $A'$ en un idéal
maximal $\mathfrak m' \subset A'$ au-dessus de $\mathfrak
m$, voir Algèbre, Lemme
\ref{algebra-lemma-quasi-finite-henselization} ou
\ref{algebra-lemma-quasi-finite-strict-henselization}. Par
conséquent, ces anneaux locaux sont des domaines normaux d'après
Plus sur l'Algèbre, Lemme
\ref{more-algebra-lemma-henselization-normal}. Il s'ensuit que $A'
\otimes_A A^h$ et $A' \otimes_A A^{sh}$ sont des anneaux normaux. Puisque
$A^h \to A' \otimes_A A^h$ et $A^{sh} \to A' \otimes_A A^{sh}$ sont finis
(donc intègres) et puisque $A' \otimes_A A^h \subset (A^h)_f = Q(A^h)$ et $A'
\otimes_A A^{sh} \subset (A^{sh})_f = Q(A^{sh})$, nous concluons que $A'
\otimes_A A^h$ et $A' \otimes_A A^{sh}$ sont les clôtures intégrales souhaitées.

\medskip\noindent
Pour la complétion, nous argumentons de la même manière. D'abord, par
l'algèbre, le Lemme \ref{algebra-lemma-completion-finite-extension} nous avons
$$
A' \otimes_A A^\wedge = (A')^\wedge =
\prod\nolimits (A'_{\mathfrak m'})^\wedge
$$
Les anneaux locaux $A'_{\mathfrak m'}$ sont normaux et ont la
dimension $1$ (par l'Algèbre, Lemme
\ref{algebra-lemma-finite-in-codim-1} par exemple ou la discussion
dans Algèbre, Section
\ref{algebra-section-homomorphism-dimension}). Ainsi $A'_{\mathfrak m'}$
est un anneau de valuation discrète, voir Algèbre, Lemme
\ref{algebra-lemma-characterize-dvr}. Par conséquent, $(A'_{\mathfrak
m'})^\wedge$ est un anneau de valuation discrète selon Plus sur l'Algèbre, Lemme
\ref{more-algebra-lemma-completion-dvr}. Il s'ensuit que $A' \otimes_A A^\wedge$ est
un anneau normal et nous pouvons conclure exactement de la même manière qu'auparavant.
\end{proof}

\begin{lemma}
\label{lemma-delta-same-after-completion}
Supposons que $A$ soit un anneau local de Nagata
réduit de dimension $1$. L'invariant $\delta$ de $A$
est le même que l'invariant $\delta$ de l'hensélisation,
de l'hensélisation stricte ou de la complétion de $A$.
\end{lemma}

\begin{proof}
Faisons cela dans le cas de la complétion $B =
A^\wedge$ ; les autres cas sont prouvés de
exactement la même manière. Supposons que $A'$,
respectivement \ $B'$, soit l'adhérence intégrale de
$A$, respectivement \ $B$ dans son anneau total de
fractions. Alors $B' = A' \otimes_A B$ d'après le Lemme
\ref{lemma-normalization-same-after-completion}. Ainsi
$B'/B = (A'/A) \otimes_A B$. L'égalité découle
maintenant de l'Algèbre, Lemme
\ref{algebra-lemma-pullback-module} et du fait que $B \otimes_A \kappa_A = \kappa_B$.
\end{proof}

\begin{definition}
\label{definition-delta-invariant}
Soit $k$ un corps. Soit $X$ un $k$-schéma localement
algébrique. Soit $x \in X$ un point tel que
$\mathcal{O}_{X, x}$ soit réduit et que $\dim(\mathcal{O}_{X,
x}) = 1$. L'{\it invariant $\delta$ de $X$ en $x$ } est
l'invariant $\delta$ de $\mathcal{O}_{X, x}$ tel qu'il est
défini dans la définition \ref{definition-delta-invariant-algebra}.
\end{definition}

\noindent
Cela a du sens car l'anneau local d'un schéma
localement algébrique est Nagata selon l'Algèbre,
Proposition \ref{algebra-proposition-ubiquity-nagata}.
Bien sûr, plus généralement, nous pouvons définir ceci
chaque fois que $x \in X$ est un point d'un schéma tel que
l'anneau local $\mathcal{O}_{X, x}$ est réduit, Nagata de
dimension $1$. Il résulte du Lemme
\ref{lemma-delta-same-after-completion} que l'invariant $\delta$ de $X$ en $x$ est
$$
\delta\text{-invariant of }X\text{ at }x =
\delta\text{-invariant of }\mathcal{O}_{X, x}^h =
\delta\text{-invariant of }\mathcal{O}_{X, x}^\wedge
$$
Nous concluons que l'invariant $\delta$ est un
invariant de l'anneau local complet du point.

\begin{lemma}
\label{lemma-delta-invariant-and-change-of-fields}
Supposons que $k$ soit un corps. Supposons que $X$
soit un $k$ -schéma localement algébrique. Supposons
que $K/k$ soit une extension de corps et posons $Y =
X_K$. Soit $y \in Y$ avec image $x \in X$. Supposons
que $X$ soit géométriquement réduit en $x$ et
$\dim(\mathcal{O}_{X, x}) = \dim(\mathcal{O}_{Y, y}) = 1$. Alors
$$
\delta\text{-invariant of }X\text{ at }x \leq
\delta\text{-invariant of }Y\text{ at }y
$$
\end{lemma}

\begin{proof}
Fixez $A = \mathcal{O}_{X, x}$ et $B =
\mathcal{O}_{Y, y}$. D'après le Lemme
\ref{lemma-geometrically-reduced-at-point}, nous voyons que
$A$ est géométriquement réduit. Par conséquent, $B$ est une
localisation de $A \otimes_k K$. Supposons que $A \to A'$
soit comme dans le Lemme \ref{lemma-pre-delta-invariant}. Alors
$$
B' = B \otimes_{(A \otimes_k K)} (A' \otimes_k K)
$$
est fini sur $B$ et $B \to B'$ induit un isomorphisme sur les
anneaux totaux des fractions. À savoir, choisissez $f \in
\mathfrak m_A$ satisfaisant (1) -- (6) du Lemme
\ref{lemma-pre-pre-delta-invariant} ; puisque $\dim(B) = 1$ nous
voyons que $f \in \mathfrak m_B$ joue le même rôle pour $B$ et nous
voyons que $B_f = B'_f$ parce que $A_f = A'_f$. Supposons que $B''$ est
la clôture intégrale de $B$ dans son anneau total des fractions comme
dans le Lemme \ref{lemma-pre-delta-invariant}. Alors $B' \subset B''$.
Ainsi, l'invariant $\delta$ de $Y$ en $y$ est $\text{length}_B(B''/B)$ et
\begin{align*}
\text{length}_B(B''/B)
& \geq
\text{length}_B(B'/B) \\
& =
\text{length}_B((A'/A) \otimes_A B) \\
& =
\text{length}_B(B/\mathfrak m_A B) \text{length}_A(A'/A)
\end{align*}
par l'algèbre, le lemme
\ref{algebra-lemma-pullback-module} puisque $A \to B$ est plat
(comme une localisation de $A \to A \otimes_k K$). Puisque
$\text{length}_A(A'/A)$ est l'invariant $\delta$ de $X$ en $x$ et
puisque $\text{length}_B(B/\mathfrak m_A B) \geq 1$, le lemme est prouvé.
\end{proof}

\begin{lemma}
\label{lemma-delta-invariant-and-change-of-fields-better}
Supposons que $k$ soit un corps. Supposons que $X$ soit
un $k$-schéma localement algébrique. Supposons que
$K/k$ soit une extension de corps et posons $Y = X_K$.
Soit $y \in Y$ avec image $x \in X$. Supposons que les
hypothèses (a), (b), (c) du Lemme
\ref{lemma-geometrically-normal-in-codim-1} soient vraies pour
$x \in X$ et que $\dim(\mathcal{O}_{Y, y}) = 1$. Alors
l'invariant $\delta$ de $X$ à $x$ est l'invariant $\delta$ de $Y$ à $y$.
\end{lemma}

\begin{proof}
Poser $A = \mathcal{O}_{X, x}$ et $B = \mathcal{O}_{Y,
y}$. D'après le Lemme
\ref{lemma-geometrically-normal-in-codim-1}, nous voyons
que $A$ est géométriquement réduit. Par conséquent, $B$ est
une localisation de $A \otimes_k K$. Supposons que $A \to
A'$ soit comme dans le Lemme \ref{lemma-pre-delta-invariant}.
D'après le Lemme \ref{lemma-geometrically-normal-in-codim-1},
nous voyons que $A' \otimes_k K$ est normal. Par conséquent
$$
B' = B \otimes_{(A \otimes_k K)} (A' \otimes_k K)
$$
est normal, fini sur $B$, et $B \to B'$ induit un isomorphisme
sur les anneaux totaux des fractions. En d'autres termes,
choisissez $f \in \mathfrak m_A$ satisfaisant (1) -- (6) du
Lemme \ref{lemma-pre-pre-delta-invariant} ; puisque $\dim(B) =
1$ nous voyons que $f \in \mathfrak m_B$ joue le même rôle pour
$B$ et nous voyons que $B_f = B'_f$ parce que $A_f = A'_f$. Il
s'ensuit que $B \to B'$ est comme dans le Lemme
\ref{lemma-pre-delta-invariant} pour $B$. Ainsi, nous devons
montrer que $\text{length}_A(A'/A) = \text{length}_B(B'/B) =
\text{length}_B((A'/A) \otimes_A B)$. Puisque $A \to B$ est plat (en
tant que localisation de $A \to A \otimes_k K$) et puisque $\mathfrak
m_B = \mathfrak m_A B$ (parce que $B/\mathfrak m_A B$ est de dimension
zéro selon les remarques ci-dessus et une localisation de $K
\otimes_k \kappa(x)$ qui est réduit comme $\kappa(x)$ est séparable sur
$k$) nous concluons par Algèbre, Lemme \ref{algebra-lemma-pullback-module}.
\end{proof}





\section{Le nombre de branches}
\label{section-number-of-branches}

\noindent
Nous avons défini le nombre de branches d’un schéma à un point
dans Propriétés, Section \ref{properties-section-unibranch}.

\begin{lemma}
\label{lemma-number-of-branches}
Supposons que $X$ soit un schéma. Supposons que chaque ouvert
quasi-compact de $X$ ait un nombre fini de composantes irréductibles.
Supposons que $\nu : X^\nu \to X$ soit la normalisation de $X$. Soit $x \in X$.
\begin{enumerate}
\item Le nombre de branches de $X$ en $x$ est le
nombre d'antécédents de $x$ dans $X^\nu$. \item Le
nombre de branches géométriques de $X$ en $x$ est
$\sum_{\nu(x^\nu) = x} [\kappa(x^\nu) : \kappa(x)]_s$.
\end{enumerate}
\end{lemma}

\begin{proof}
Notez d'abord que l'hypothèse sur $X$ signifie
exactement que la normalisation est définie, voir
morphismes, Définition
\ref{morphisms-definition-normalization}. Ensuite, la
germe $A' = (\nu_*\mathcal{O}_{X^\nu})_x$ est
l'adhérence intégrale de $A = \mathcal{O}_{X, x}$ dans
l'anneau total des fractions de $A_{red}$, voir
morphismes, Lemme
\ref{morphisms-lemma-stalk-normalization}. Puisque $\nu$ est
un morphisme intégral, nous voyons que les points de $X^\nu$
au-dessus de $x$ correspondent aux idéaux premiers de $A'$
au-dessus de l'idéal maximal $\mathfrak m$ de $A$. Comme $A
\to A'$ est intégral, c'est la même chose que les idéaux
maximaux de $A'$ (Algèbre, Lemmata
\ref{algebra-lemma-integral-no-inclusion} et
\ref{algebra-lemma-integral-going-up}). Ainsi, le lemme découle maintenant de son
équivalent algébrique : Plus sur l'Algèbre, Lemme \ref{more-algebra-lemma-number-of-branches-1}.
\end{proof}

\begin{lemma}
\label{lemma-geometric-branches-and-change-of-fields}
Supposons que $k$ soit un corps. Supposons que $X$ soit un
$k$-schéma localement algébrique. Supposons que $K/k$ soit une
extension de corps. Supposons que $y \in X_K$ soit un point avec
image $x$ dans $X$. Alors, le nombre de branches géométriques de
$X$ en $x$ est le nombre de branches géométriques de $X_K$ en $y$.
\end{lemma}

\begin{proof}
Écrivez $Y = X_K$ et laissez $X^\nu$, respectivement \ $Y^\nu$ être la
normalisation de $X$, respectivement \ $Y$. Considérons le diagramme commutatif
$$
\xymatrix{
Y^\nu \ar[r] \ar[d] & X^\nu_K \ar[r] \ar[d]_{\nu_K} & X^\nu \ar[d]_\nu \\
Y \ar@{=}[r] & Y \ar[r] & X
}
$$
D'après le Lemme
\ref{lemma-normalization-and-change-of-fields}, nous voyons que
la flèche horizontale en haut à gauche est un homéomorphisme
universel. Par conséquent, elle induit des extensions
résiduelles purement inséparables de corps, voir morphismes,
Lemmata \ref{morphisms-lemma-universal-homeomorphism} et
\ref{morphisms-lemma-universally-injective}. Ainsi, le nombre de
branches géométriques de $Y$ en $y$ est $\sum_{\nu_K(y') = y}
[\kappa(y') : \kappa(y)]_s$ d'après le Lemme
\ref{lemma-number-of-branches}. De même, $\sum_{\nu(x') = x} [\kappa(x')
: \kappa(x)]_s$ est le nombre de branches géométriques de $X$ en $x$. En
utilisant schémas, Lemme \ref{schemes-lemma-points-fibre-product}, notre
énoncé découle du fait algébrique suivant : étant donnée une extension de
corps $l/\kappa$ et une extension algébrique de corps $m/\kappa$, alors
$$
\sum\nolimits_{m \otimes_\kappa l \to m'} [m' : l']_s = [m : \kappa]_s
$$
où la somme porte sur les corps quotient de $m
\otimes_\kappa l$. On peut le démontrer de manière
élémentaire, ou on peut utiliser le Lemme
\ref{lemma-separably-closed-field-connected-components} appliqué à
$$
\Spec(m \otimes_\kappa l) \times_{\Spec(l)} \Spec(\overline{l}) =
\Spec(m) \otimes_{\Spec(\kappa)} \Spec(\overline{l})
\longrightarrow
\Spec(m) \times_{\Spec(\kappa)} \Spec(\overline{\kappa})
$$
parce qu'on peut interpréter $[m : \kappa]_s$ comme
le nombre de composantes connexes du côté droit et la
somme $\sum_{m \otimes_\kappa l \to m'} [m' : l']_s$
comme le nombre de composantes connexes du côté gauche.
\end{proof}

\begin{lemma}
\label{lemma-geometrically-unibranch-and-change-of-fields}
Supposons que $k$ soit un corps. Supposons que $X$ soit un $k$ -schéma
localement algébrique. Supposons que $K/k$ soit une extension de corps. Supposons
que $y \in X_K$ soit un point avec l'image $x$ dans $X$. Alors $X$ est
géométriquement unibranche en $x$ si et seulement si $X_K$ est géométriquement unibranche en $y$.
\end{lemma}

\begin{proof}
Immédiat d'après le Lemme
\ref{lemma-geometric-branches-and-change-of-fields} et Plus sur
l'Algèbre, Lemme \ref{more-algebra-lemma-number-of-branches-1}.
\end{proof}

\begin{definition}
\label{definition-wedge}
Supposons que $A$ et $A_i$, $1 \leq i \leq n$ sont
des anneaux locaux. On dit que {\it $A$ est un bouquet
de $A_1, \ldots, A_n$ } s'il existe des isomorphismes
$$
\kappa_{A_1} \to \kappa_{A_2} \to \ldots \to \kappa_{A_n}
$$
et $A$ est isomorphe à l'anneau constitué des $n$ -uples
$(a_1, \ldots, a_n) \in A_1 \times \ldots \times A_n$
qui se projettent sur le même élément de $\kappa_{A_n}$.
\end{definition}

\noindent
Si on nous donne un anneau de base $\Lambda$ et que $A$ et
$A_i$ sont des $\Lambda$-algèbres, alors nous exigeons que
$\kappa_{A_i} \to \kappa_{A_{i + 1}}$ soit un isomorphisme de
$\Lambda$-algèbres et que $A$ soit isomorphe en tant que
$\Lambda$-algèbre à la $\Lambda$-algèbre constituée de
$n$-uplets $(a_1, \ldots, a_n) \in A_1 \times \ldots \times A_n$
qui se projettent sur le même élément de $\kappa_{A_n}$. En
particulier, si $\Lambda = k$ est un corps et si les applications
$k \to \kappa_{A_i}$ sont des isomorphismes, alors il existe un
choix unique pour les isomorphismes $\kappa_{A_i} \to \kappa_{A_{i +
1}}$ et l'on parle souvent de {\it le bouquet de $A_1, \ldots, A_n$ }.

\begin{lemma}
\label{lemma-delta-number-branches-inequality-sh}
Supposons que $(A, \mathfrak m)$ soit un anneau local
réduit de Nagata strictement hensélien de dimension $1$. Alors
$$
\delta\text{-invariant of }A \geq \text{number of geometric branches of }A - 1
$$
Si l'égalité tient, alors $A$ est un faisceau de $n \geq
1$ d'anneaux de valuation discrète strictement henséliens.
\end{lemma}

\begin{proof}
Le nombre de branches géométriques est égal au nombre
de branches de $A$ (immédiat d'après Plus sur
l'Algèbre, Définition
\ref{more-algebra-definition-number-of-branches}). Supposons que
$A \to A'$ est comme dans le Lemme
\ref{lemma-pre-delta-invariant}. On observe que le nombre de branches de
$A$ est le nombre d'idéaux maximaux de $A'$, voir Plus sur l'Algèbre,
Lemme \ref{more-algebra-lemma-number-of-branches-1}. Il existe une surjection
$$
A'/A \longrightarrow
\left(\prod\nolimits_{\mathfrak m'} \kappa(\mathfrak m')\right)/
\kappa(\mathfrak m)
$$
Puisque $\dim_{\kappa(\mathfrak m)} \prod \kappa(\mathfrak
m')$ est $\geq$ le nombre de branches, l'inégalité est évidente.

\medskip\noindent
Si l'égalité est vraie, alors $\kappa(\mathfrak m') =
\kappa(\mathfrak m)$ pour tout $\mathfrak m' \subset A'$ et la
flèche affichée ci-dessus est un isomorphisme. Puisque $A$ est
hensélien et $A \to A'$ est fini, nous voyons que $A'$ est un
produit d'anneaux locaux henséliens, voir Algèbre, Lemme
\ref{algebra-lemma-finite-over-henselian}. Les facteurs sont les
anneaux locaux $A'_{\mathfrak m'}$ et comme $A'$ est normal, ces
facteurs sont des anneaux de valuation discrète (Algèbre, Lemme
\ref{algebra-lemma-characterize-dvr}). Puisque la flèche affichée est un
isomorphisme, nous voyons que $A$ est en effet le wedge de ces anneaux locaux.
\end{proof}

\begin{lemma}
\label{lemma-delta-number-branches-inequality}
Supposons que $(A, \mathfrak m)$ soit un anneau local réduit de Nagata de dimension $1$. Alors
$$
\delta\text{-invariant of }A \geq \text{number of geometric branches of }A - 1
$$
\end{lemma}

\begin{proof}
Nous pouvons remplacer $A$ par l'hensélisation stricte de
$A$ sans modifier l'invariant $\delta$ (Lemme
\ref{lemma-delta-same-after-completion}) et sans changer le
nombre de branches géométriques de $A$ (cela découle
immédiatement de la définition, voir Plus sur l'Algèbre,
Définition \ref{more-algebra-definition-number-of-branches}). Ainsi,
nous pouvons supposer que $A$ est strictement hensélien et nous
pouvons appliquer le Lemme \ref{lemma-delta-number-branches-inequality-sh}.
\end{proof}






\section{Normalisation des schémas de dimension un}
\label{section-normalization-one-dimensional}

\noindent
Le morphisme de normalisation d’un schéma noethérien de dimension $1$ a
des propriétés exceptionnellement bonnes selon le résultat de Krull-Akizuki.

\begin{lemma}
\label{lemma-normalize-noetherian-dim-1}
Supposons que $X$ soit un schéma localement noethérien de dimension
$1$. Supposons que $\nu : X^\nu \to X$ soit la normalisation. Alors
\begin{enumerate}
\item $\nu$ est intégral, surjectif, et induit une
bijection sur les composantes irréductibles, \item
il existe une factorisation $X^\nu \to X_{red} \to
X$ et le morphisme $X^\nu \to X_{red}$ est la
normalisation de $X_{red}$, \item $X^\nu \to
X_{red}$ est birationnel, \item pour tout point
fermé $x \in X$ le germe
$(\nu_*\mathcal{O}_{X^\nu})_x$ est l'adhérence
intégrale de $\mathcal{O}_{X, x}$ dans l'anneau total
des fractions de $(\mathcal{O}_{X, x})_{red} =
\mathcal{O}_{X_{red}, x}$, \item les fibres de $\nu$ sont
finies et les extensions résiduelles de corps sont
finies, \item $X^\nu$ est une union disjointe de schémas
intègres normaux localement noethériens et chaque ouvert
affine est le spectre d'un produit fini de domaines de Dedekind.
\end{enumerate}
\end{lemma}

\begin{proof}
Beaucoup des résultats sont en fait des propriétés
générales du morphisme de normalisation, voir
morphismes, lemmes
\ref{morphisms-lemma-normalization-reduced},
\ref{morphisms-lemma-stalk-normalization},
\ref{morphisms-lemma-normalization-normal} et
\ref{morphisms-lemma-normalization-birational}. Ce qui n'est
pas clair, c'est que les fibres sont finies, que les
extensions de corps résiduelles induites sont finies, et que
$X^\nu$ ressemble localement au spectre d'un domaine de
Dedekind (et est donc noethérien). Pour voir cela, nous pouvons
supposer que $X = \Spec(A)$ est affine, noethérien, de
dimension $1$, et que $A$ est réduit. Nous pouvons alors
utiliser la description dans morphismes, lemme
\ref{morphisms-lemma-description-normalization} pour réduire au cas où
$A$ est un domaine noethérien de dimension $1$. Dans ce cas, les
propriétés désirées découlent du théorème de Krull-Akizuki sous la forme
énoncée dans Algèbre, lemme \ref{algebra-lemma-integral-closure-Dedekind}.
\end{proof}

\noindent
Bien sûr, il existe une variante du lemme
suivant dans le cas où $X$ n'est pas réduit.

\begin{lemma}
\label{lemma-prepare-delta-invariant}
Supposons que $X$ soit un schéma de Nagata réduit de dimension $1$. Supposons que
$\nu : X^\nu \to X$ soit la normalisation. Soit $x \in X$ un point fermé. Alors
\begin{enumerate}
\item $\nu : X^\nu \to X$ est fini, surjectif et
birationnel, \item $\mathcal{O}_X \subset
\nu_*\mathcal{O}_{X^\nu}$ et
$\nu_*\mathcal{O}_{X^\nu}/\mathcal{O}_X$ est une somme
directe de faisceaux skyscraper avec valeur
$\mathcal{Q}_x$ en $x$, ce qui est non nul si et
seulement si $x$ est un point singulier de $X$, \item
$A' = (\nu_*\mathcal{O}_{X^\nu})_x$ est l'adhérence
intégrale de $A = \mathcal{O}_{X, x}$ dans son anneau
total de fractions, \item $\mathcal{Q}_x = A'/A$ a une
longueur finie égale à l'invariant $\delta$ de $X$ en
$x$, \item $A'$ est un semi-anneau local qui est un
produit fini de domaines de Dedekind, \item $A^\wedge$
est un anneau local réduit noethérien complet de
dimension $1$, \item $(A')^\wedge$ est l'adhérence
intégrale de $A^\wedge$ dans son anneau total de fractions,
\item $(A')^\wedge$ est un produit fini d'anneaux de valuation
discrète complets, et \item $A'/A \cong (A')^\wedge/A^\wedge$.
\end{enumerate}
\end{lemma}

\begin{proof}
Nous pouvons et allons utiliser tous les résultats du
Lemme \ref{lemma-normalize-noetherian-dim-1}. La
finitude de $\nu$ découle des morphismes, Lemme
\ref{morphisms-lemma-nagata-normalization}. Comme $X$ est
réduit, Nagata, de dimension $1$, nous voyons que le lieu
régulier est un ouvert dense $U \subset X$ selon Plus sur
l'Algèbre, Proposition
\ref{more-algebra-proposition-ubiquity-J-2}. Comme un schéma
régulier est normal, cela montre que $\nu$ est un isomorphisme sur
$U$. Comme $\dim(X) \leq 1$, cela implique que $\nu$ n'est pas un
isomorphisme sur un ensemble discret de points fermés $x \in X$.
En particulier, nous voyons que nous avons une suite exacte courte.
$$
0 \to \mathcal{O}_X \to \nu_*\mathcal{O}_{X^\nu} \to
\bigoplus\nolimits_{x \in X \setminus U} \mathcal{Q}_x \to 0
$$
Comme nous avons la description des tiges de
$\nu_*\mathcal{O}_{X^\nu}$ par le Lemme
\ref{lemma-normalize-noetherian-dim-1}, nous concluons que $Q_x = A'/A$
a en effet une longueur égale à l'invariant $\delta$ de $X$ en $x$.
Notez que $Q_x \not = 0$ exactement lorsque $x$ est un point singulier,
par exemple selon le Lemme \ref{lemma-delta-invariant-is-zero}. La
description de $A'$ comme un produit de domaines de Dedekind semi-locaux
découle également du Lemme \ref{lemma-normalize-noetherian-dim-1}. La
relation entre $A$, $A'$ et $(A')^\wedge$ que nous avons vue dans le
Lemme \ref{lemma-normalization-same-after-completion} (et sa démonstration).
\end{proof}




\section{Recherche d'ouverts affines}
\label{section-finding-affine-opens}

\noindent
Nous continuons la discussion commencée dans Propriétés, Section
\ref{properties-section-finding-affine-opens}. Il s'avère que
nous pouvons trouver des affines contenant un ensemble fini donné
de points de codimension $1$ sur un schéma séparé. Voir la
Proposition \ref{proposition-finite-set-of-points-of-codim-1-in-affine}.

\medskip\noindent
Nous allons améliorer le lemme suivant dans Descent,
Lemme \ref{descent-lemma-characterize-open-immersion}.

\begin{lemma}
\label{lemma-characterize-open-immersion}
Supposons que $f : X \to Y$ soit un morphisme de schémas. Soit $X^0$
l'ensemble des points génériques des composantes irréductibles de $X$. Si
\begin{enumerate}
\item $f$ est séparé, \item il existe un recouvrement
ouvert $X = \bigcup U_i$ tel que $f|_{U_i} : U_i \to Y$
est une immersion ouverte, et \item si $\xi, \xi' \in
X^0$, $\xi \not = \xi'$, alors $f(\xi) \not = f(\xi')$,
\end{enumerate}
alors $f$ est une immersion ouverte.
\end{lemma}

\begin{proof}
Supposons que $y = f(x) = f(x')$. Choisissez une spécialisation $y_0 \leadsto y$
où $y_0$ est un point générique d'une composante irréductible de $Y$. Puisque
$f$ est localement sur la source un isomorphisme, nous pouvons choisir des
spécialisations $x_0 \leadsto x$ et $x'_0 \leadsto x'$ se mappant sur $y_0
\leadsto y$. Notez que $x_0, x'_0 \in X^0$. Ainsi $x_0 = x'_0$ par l'hypothèse (3).
Comme $f$ est séparé, nous concluons que $x = x'$. Donc $f$ est une immersion ouverte.
\end{proof}

\begin{lemma}
\label{lemma-local-isomorphism}
Supposons que $X \to S$ soit un morphisme de schémas.
Supposons que $x \in X$ soit un point avec image $s \in S$. Si
\begin{enumerate}
\item $\mathcal{O}_{X, x} =
\mathcal{O}_{S, s}$, \item $S$ est réduit,
\item $X \to S$ est de type fini, et \item $S$
a un nombre fini de composantes irréductibles,
\end{enumerate}
alors il existe un voisinage ouvert $U$ de
$x$ tel que $f|_U$ soit une immersion ouverte.
\end{lemma}

\begin{proof}
Nous pouvons enlever les (finitement nombreux) composantes
irréductibles de $S$ qui ne contiennent pas $s$. Nous pouvons
remplacer $S$ par un voisinage affine ouvert de $s$. Nous pouvons
remplacer $X$ par un voisinage affine ouvert de $x$. Disons $S =
\Spec(A)$ et $X = \Spec(B)$. Supposons que $\mathfrak q \subset B$,
respectivement \ $\mathfrak p \subset A$, soit l'idéal premier
correspondant à $x$, respectivement \ $s$. Comme $A$ est réduit et que
tous les idéaux premiers minimaux de $A$ sont contenus dans $\mathfrak
p$, nous voyons que $A \subset A_\mathfrak p$. Comme $X \to S$ est de
type fini, $B$ est de type fini sur $A$. Supposons que $b_1, \ldots,
b_n \in B$ soit des éléments qui engendrent $B$ sur $A$. Puisque
$A_\mathfrak p = B_\mathfrak q$, nous pouvons trouver $f \in A$, $f \not
\in \mathfrak p$ et $a_i \in A$ tels que $b_i$ et $a_i/f$ aient la même
image dans $B_\mathfrak q$. Ainsi, nous pouvons trouver $g \in B$, $g
\not \in \mathfrak q$ tels que $g(fb_i - a_i) = 0$ dans $B$. Il s'ensuit
que l'image de $A_f \to B_{fg}$ contient les images de $b_1, \ldots,
b_n$, en particulier aussi l'image de $g$. Choisissez $n \geq 0$ et $f'
\in A$ de sorte que $f'/f^n$ se projette sur l'image de $g$ dans
$B_{fg}$. Puisque $A_\mathfrak p = B_\mathfrak q$, nous voyons que $f' \not
\in \mathfrak p$. Nous concluons que $A_{ff'} \to B_{fg}$ est surjective.
Enfin, comme $A_{ff'} \subset A_\mathfrak p = B_\mathfrak q$ (voir
ci-dessus), l'application $A_{ff'} \to B_{fg}$ est injective, donc un isomorphisme.
\end{proof}

\begin{lemma}
\label{lemma-points-in-affine}
Supposons que $f : T \to X$ soit un morphisme de schémas. Supposons
que $X^0$, respectivement \ $T^0$, désignent les ensembles de
points génériques des composantes irréductibles. Soit $t_1, \ldots,
t_m \in T$ un ensemble fini de points avec des images $x_j = f(t_j)$. Si
\begin{enumerate}
\item $T$ est affine, \item $X$ est
quasi-séparé, \item $X^0$ est fini
\item $f(T^0) \subset X^0$ et $f : T^0
\to X^0$ est injectif, et \item
$\mathcal{O}_{X, x_j} = \mathcal{O}_{T, t_j}$,
\end{enumerate}
alors il existe un ouvert affine de $X$ contenant $x_1, \ldots, x_r$.
\end{lemma}

\begin{proof}
En utilisant les limites, la proposition
\ref{limits-proposition-affine} permet une réduction
immédiate au cas où $X$ et $T$ sont réduits. Détails omis.

\medskip\noindent
Supposons que $X$ et $T$ soient réduits. Nous pouvons écrire
$T = \lim_{i \in I} T_i$ comme une limite dirigée de schémas
de présentation finie sur $X$ avec des morphismes de
transition affines, voir Limites, Lemme
\ref{limits-lemma-relative-approximation}. Choisissez $i \in I$
de sorte que $T_i$ soit affine, voir Limites, Lemme
\ref{limits-lemma-limit-affine}. Disons $T_i = \Spec(R_i)$ et $T =
\Spec(R)$. Supposons que $R' \subset R$ soit l'image de $R_i \to R$.
Alors $T' = \Spec(R')$ est affine, réduit, de type fini sur $X$, et
$T \to T'$ dominant. Pour $j = 1, \ldots, r$, prenons $t'_j \in T'$
comme l'image de $t_j$. Considérons les applications de l'anneau local
$$
\mathcal{O}_{X, x_j} \to
\mathcal{O}_{T', t'_j} \to
\mathcal{O}_{T, t_j}
$$
Désignons par $(T')^0$ l'ensemble des points génériques des composantes
irréductibles de $T'$. Supposons que $\xi \leadsto t'_j$ soit une
spécialisation avec $\xi \in (T')^0$. Comme $T \to T'$ est dominant, nous
pouvons choisir $\eta \in T^0$ mappant sur $\xi$ (attention : a priori nous
ne savons pas que $\eta$ se spécialise en $t_j$). L'hypothèse (3) appliquée
à $\eta$ nous dit que l'image $\theta$ de $\xi$ dans $X$ correspond à un
idéal premier minimal de $\mathcal{O}_{X, x_j}$. En relevant $\xi$ via
l'isomorphisme de (5) nous obtenons une spécialisation $\eta' \leadsto t_j$
avec $\eta' \in T^0$ mappant sur $\theta \leadsto x_j$. L'injectivité de (4)
montre que $\eta = \eta'$. Ainsi, chaque idéal premier minimal de
$\mathcal{O}_{T', t'_j}$ est situé sous un idéal premier minimal de
$\mathcal{O}_{T, t_j}$. Nous en concluons que $\mathcal{O}_{T', t'_j} \to
\mathcal{O}_{T, t_j}$ est injectif, donc les deux applications ci-dessus sont des isomorphismes.

\medskip\noindent
D'après le Lemme \ref{lemma-local-isomorphism}, il
existe un ouvert $U \subset T'$ contenant tous les
points $t'_j$ tel que $U \to X$ soit un
isomorphisme local comme dans le Lemme
\ref{lemma-characterize-open-immersion}. D'après ce
lemme, nous voyons que $U \to X$ est une immersion
ouverte. Enfin, d'après les Propriétés, Lemme
\ref{properties-lemma-ample-finite-set-in-affine}, nous
pouvons trouver un ouvert $W \subset U \subset T'$ contenant tous
les $t'_j$. L'image de $W$ dans $X$ est l'ouvert affine désiré.
\end{proof}

\begin{lemma}
\label{lemma-finite-set-codim-1-points-in-affine}
Supposons que $X$ soit un schéma séparé intègre. Supposons que $x_1,
\ldots, x_r \in X$ soit un ensemble fini de points tel que
$\mathcal{O}_{X, x_i}$ soit noethérien de dimension $\leq 1$. Alors il
existe un sous-schéma affine ouvert de $X$ contenant tous les $x_1, \ldots, x_r$.
\end{lemma}

\begin{proof}
Supposons que $K$ soit le corps des fonctions rationnelles
de $X$. Posons $A_i = \mathcal{O}_{X, x_i}$. Alors $A_i
\subset K$ et $K$ est le corps des fractions de $A_i$.
Puisque $X$ est séparé, et $x_i \not = x_j$ il ne peut pas y
avoir un anneau de valuation $\mathcal{O} \subset K$ dominant à
la fois $A_i$ et $A_j$. À savoir, en considérant le diagramme
$$
\xymatrix{
\Spec(\mathcal{O}) \ar[r] \ar[d] & \Spec(A_1) \ar[d] \\
\Spec(A_2) \ar[r] & X
}
$$
et en appliquant le critère valuatif de séparation
(schémas, Lemme
\ref{schemes-lemma-separated-implies-valuative}) nous
obtiendrions $x_i = x_j$. Ainsi, nous voyons d'après le
Lemme \ref{lemma-glue-separated} que $A_i \otimes A_j \to K$
est surjectif pour tout $i \not = j$. D'après le Lemme
\ref{lemma-glue-a-bunch-of-local-rings}, nous voyons que $A =
A_1 \cap \ldots \cap A_r$ est un semi-anneau local
noethérien avec exactement $r$ idéaux maximaux $\mathfrak m_1,
\ldots, \mathfrak m_r$ tels que $A_i = A_{\mathfrak m_i}$. De plus,
$$
\Spec(A) = \Spec(A_1) \cup \ldots \cup \Spec(A_r)
$$
est un recouvrement ouvert et l'intersection de deux parties
quelconques de ce recouvrement est $\Spec(K)$. Ainsi, les
morphismes donnés $\Spec(A_i) \to X$ se recollent en un morphisme de schémas
$$
\Spec(A) \longrightarrow X
$$
cartographier $\mathfrak m_i$ vers $x_i$ et induire des isomorphismes
d'anneaux locaux. Ainsi, le résultat découle du Lemme \ref{lemma-points-in-affine}.
\end{proof}

\begin{lemma}
\label{lemma-extra-silly}
Supposons que $A$ soit un anneau, $I \subset A$ un idéal,
$\mathfrak p_1, \ldots, \mathfrak p_r$ des nombres premiers de
$A$, et $\overline{f} \in A/I$ un élément. Si $I \not \subset
\mathfrak p_i$ pour tout $i$, alors il existe un $f \in A$, $f
\not \in \mathfrak p_i$ qui se mappe à $\overline{f}$ dans $A/I$.
\end{lemma}

\begin{proof}
Nous pouvons supposer qu'il n'existe pas de relations d'inclusion parmi
les $\mathfrak p_i$ (en enlevant les plus petits nombres premiers).
Commencez d'abord par choisir un $f \in A$ relevant $\overline{f}$.
Supposons que $S$ soit l'ensemble $s \in \{1, \ldots, r\}$ tel que $f
\in \mathfrak p_s$. Si $S$ est vide, nous avons terminé. Sinon,
considérons l'idéal $J = I \prod_{i \not \in S} \mathfrak p_i$. Notez que
$J$ n'est pas contenu dans $\mathfrak p_s$ pour $s \in S$ car il n'y a
pas d'inclusions parmi les $\mathfrak p_i$ et que $I$ n'est contenu dans
aucun $\mathfrak p_i$. Par conséquent, nous pouvons choisir $g \in J$, $g
\not \in \mathfrak p_s$ pour $s \in S$ selon l'algèbre, Lemme
\ref{algebra-lemma-silly}. Alors $f + g$ est une solution au problème posé par le lemme.
\end{proof}

\begin{lemma}
\label{lemma-finite-set-codim-1-points-in-affine-per-component}
Supposons que $X$ est un schéma. Supposons que $T \subset X$ est un ensemble fini de points. Supposons
\begin{enumerate}
\item $X$ a un nombre fini de composantes irréductibles $Z_1,
\ldots, Z_t$, et \item $Z_i \cap T$ est contenu dans un
ouvert affine du sous-schéma réduit induit correspondant à $Z_i$.
\end{enumerate}
Alors il existe un sous-schéma affine ouvert de $X$ contenant $T$.
\end{lemma}

\begin{proof}
En utilisant les limites, la Proposition
\ref{limits-proposition-affine} conduit à une réduction
immédiate au cas où $X$ est réduit. Les détails sont omis. Dans
le reste de la démonstration, nous dotons chaque sous-ensemble
fermé de $X$ de la structure de sous-schéma fermé réduit induite.

\medskip\noindent
Nous argumentons par induction que nous pouvons trouver un
ouvert affine $U \subset Z_1 \cup \ldots \cup Z_r$
contenant $T \cap (Z_1 \cup \ldots \cup Z_r)$. Pour $r =
1$, cela tient par hypothèse. Soit $r > 1$ et soit $U
\subset Z_1 \cup \ldots \cup Z_{r - 1}$ un ouvert affine
contenant $T \cap (Z_1 \cup \ldots \cup Z_{r - 1})$.
Supposons que $V \subset X_r$ soit un ouvert affine contenant $T
\cap Z_r$ (existe par hypothèse). Alors $U \cap V$ contient $T
\cap ( Z_1 \cup \ldots \cup Z_{r - 1} ) \cap Z_r$. Par conséquent
$$
\Delta = (U \cap Z_r) \setminus (U \cap V)
$$
ne contient aucun élément de $T$. Notez que $\Delta$ est un
sous-ensemble fermé de $U$. Par l'évitation des idéaux premiers
(Algèbre, Lemme \ref{algebra-lemma-silly}), nous pouvons trouver un
ouvert standard $U'$ de $U$ contenant $T \cap U$ et évitant $\Delta$,
c'est-à-dire $U' \cap Z_r \subset U \cap V$. Après avoir remplacé
$U$ par $U'$, nous pouvons supposer que $U \cap V$ est fermé dans $U$.

\medskip\noindent
En utilisant le fait que par les mêmes arguments que ci-dessus,
l'ensemble $\Delta' = (U \cap (Z_1 \cup \ldots \cup Z_{r - 1}))
\setminus (U \cap V)$ ne contient également aucun élément de $T$, nous
trouvons un $h \in \mathcal{O}(V)$ tel que $D(h) \subset V$ contient
$T \cap V$ et tel que $U \cap D(h) \subset U \cap V$. En utilisant que
$U \cap V$ est fermé dans $U$, nous pouvons utiliser le Lemme
\ref{lemma-extra-silly} pour trouver un élément $g \in \mathcal{O}(U)$
dont la restriction à $U \cap V$ est égale à la restriction de $h$ à $U
\cap V$ et tel que $T \cap U \subset D(g)$. Ensuite, nous pouvons
remplacer $U$ par $D(g)$ et $V$ par $D(h)$ pour atteindre la situation
où $U \cap V$ est fermé à la fois dans $U$ et $V$. Dans ce cas, le
schéma $U \cup V$ est affine selon Limites, Lemme
\ref{limits-lemma-affines-glued-in-closed-affine}. Cela prouve l'étape d'induction et donc le lemme.
\end{proof}

\noindent
Voici une conclusion que nous pouvons tirer du matériel ci-dessus.

\begin{proposition}
\label{proposition-finite-set-of-points-of-codim-1-in-affine}
Supposons que $X$ soit un schéma séparé tel que chaque ouvert
quasi-compact ait un nombre fini de composantes irréductibles.
Supposons que $x_1, \ldots, x_r \in X$ soit des points tels que
$\mathcal{O}_{X, x_i}$ soit noethérien de dimension $\leq 1$. Alors il
existe un sous-schéma ouvert affine de $X$ contenant tous les $x_1, \ldots, x_r$.
\end{proposition}

\begin{proof}
Nous pouvons remplacer $X$ par un ouvert quasi-compact contenant $x_1,
\ldots, x_r$, donc nous pouvons supposer que $X$ a un nombre fini de
composantes irréductibles. Par le Lemme
\ref{lemma-finite-set-codim-1-points-in-affine-per-component}, nous revenons au cas où
$X$ est intègre. Ce cas est le Lemme \ref{lemma-finite-set-codim-1-points-in-affine}.
\end{proof}





\section{Courbes}
\label{section-curves}

\noindent
Dans le projet Stacks, nous utiliserons ce
qui suit comme notre définition d'une courbe.

\begin{definition}
\label{definition-curve}
Supposons que $k$ soit un corps. Une {\it  courbe } est une variété de dimension $1$ sur $k$.
\end{definition}

\noindent
Deux exemples standards de courbes sur $k$ sont la ligne affine
$\mathbf{A}^1_k$ et la ligne projective $\mathbf{P}^1_k$. Le schéma $X =
\Spec(k[x, y]/(f))$ est une courbe si et seulement si $f \in k[x, y]$ est irréductible.

\medskip\noindent
Notre définition d'une courbe a les mêmes problèmes que notre définition d'une
variété, voir la discussion suivant la Définition \ref{definition-variety}. De
plus, cela signifie que chaque courbe est accompagnée d'un corps de définition
spécifié. Par exemple $X = \Spec(\mathbf{C}[x])$ est une courbe sur $\mathbf{C}$
mais nous pouvons aussi la considérer comme une courbe sur $\mathbf{R}$. Le schéma
$\Spec(\mathbf{Z})$ n'est pas une courbe, même si les schémas $\Spec(\mathbf{Z})$ et
$\mathbf{A}^1_{\mathbf{F}_p}$ se comportent de manière similaire à bien des égards.

\begin{lemma}
\label{lemma-proper-minus-point}
Supposons que $X$ soit un schéma séparé et irréductible de dimension
$> 0$ sur un corps $k$. Supposons que $x \in X$ soit un point fermé.
Le sous-schéma ouvert $X \setminus \{x\}$ n'est pas propre sur $k$.
\end{lemma}

\begin{proof}
Puisque $X$ est irréductible, $U = X \setminus \{x\}$ n'est pas fermé dans
$X$. En particulier, l'immersion $U \to X$ n'est pas propre. Par les
morphismes, le Lemme \ref{morphisms-lemma-image-proper-scheme-closed} (ici
nous utilisons que $X$ est séparé), $U \to \Spec(k)$ n'est pas propre non plus.
\end{proof}

\begin{lemma}
\label{lemma-dim-1-quasi-projective}
Supposons que $X$ soit un schéma de type fini séparé sur un corps
$k$. Si $\dim(X) \leq 1$ alors $X$ est H-quasi-projectif sur $k$.
\end{lemma}

\begin{proof}
Par la Proposition
\ref{proposition-dim-1-noetherian-separated-has-ample}, le schéma $X$
possède un faisceau inversible ample $\mathcal{L}$. Par le Lemme des
morphismes \ref{morphisms-lemma-quasi-projective-finite-type-over-S}, nous
voyons que $X$ est isomorphe à un sous-schéma localement fermé de
$\mathbf{P}^n_k$ sur $\Spec(k)$. C'est la définition d'être H-quasi-projectif sur
$k$, voir Définitions des morphismes \ref{morphisms-definition-quasi-projective}.
\end{proof}

\begin{lemma}
\label{lemma-dim-1-proper-projective}
Supposons que $X$ soit un schéma propre sur un corps
$k$. Si $\dim(X) \leq 1$ alors $X$ est H-projectif sur $k$.
\end{lemma}

\begin{proof}
D'après le Lemme \ref{lemma-dim-1-quasi-projective}, nous voyons
que $X$ est un sous-schéma localement fermé de $\mathbf{P}^n_k$
pour un certain corps $k$. Puisque $X$ est propre sur $k$, il
s'ensuit que $X$ est un sous-schéma fermé de $\mathbf{P}^n_k$
(morphismes, Lemme \ref{morphisms-lemma-image-proper-scheme-closed}).
\end{proof}

\begin{lemma}
\label{lemma-dim-1-projective-completion}
Supposons que $X$ soit un schéma séparé de type fini sur
$k$. Si $\dim(X) \leq 1$, alors il existe une immersion
ouverte $j : X \to \overline{X}$ avec les propriétés suivantes
\begin{enumerate}
\item $\overline{X}$ est H-projectif sur $k$,
c'est-à-dire, $\overline{X}$ est un sous-schéma fermé
de $\mathbf{P}^d_k$ pour un certain $d$, \item $j(X)
\subset \overline{X}$ est dense et schématiquement
dense, \item $\overline{X} \setminus X = \{x_1, \ldots,
x_n\}$ pour certains points fermés $x_i \in \overline{X}$.
\end{enumerate}
\end{lemma}

\begin{proof}
D'après le Lemme \ref{lemma-dim-1-quasi-projective},
on peut supposer que $X$ est un sous-schéma localement
fermé de $\mathbf{P}^d_k$ pour un certain $d$.
Supposons que $\overline{X} \subset \mathbf{P}^d_k$
soit l'image schématique de $X \to \mathbf{P}^d_k$,
voir morphismes, Définition
\ref{morphisms-definition-scheme-theoretic-image}. La
description dans morphismes, Lemme
\ref{morphisms-lemma-quasi-compact-immersion} donne les
propriétés (1) et (2). Alors $\dim(X) = 1 \Rightarrow
\dim(\overline{X}) = 1$ par exemple en regardant les points
génériques, voir Lemme \ref{lemma-dimension-locally-algebraic}. Comme
$\overline{X}$ est noethérien, il s'ensuit alors que $\overline{X}
\setminus X = \{x_1, \ldots, x_n\}$ est un ensemble fini de points fermés.
\end{proof}

\begin{lemma}
\label{lemma-reduced-dim-1-projective-completion}
Supposons que $X$ soit un schéma séparé de type fini sur
$k$. Si $X$ est réduit et $\dim(X) \leq 1$, alors il
existe une immersion ouverte $j : X \to \overline{X}$ telle que
\begin{enumerate}
\item $\overline{X}$ est H-projectif sur $k$,
c'est-à-dire que $\overline{X}$ est un sous-schéma fermé
de $\mathbf{P}^d_k$ pour un certain $d$, \item $j(X)
\subset \overline{X}$ est dense et dense au sens
schématique, \item $\overline{X} \setminus X = \{x_1, \ldots,
x_n\}$ pour certains points fermés $x_i \in \overline{X}$,
\item les anneaux locaux $\mathcal{O}_{\overline{X}, x_i}$
sont des anneaux de valuation discrète pour $i = 1, \ldots, n$.
\end{enumerate}
\end{lemma}

\begin{proof}
Supposons que $j : X \to \overline{X}$ soit comme dans le
Lemme \ref{lemma-dim-1-projective-completion}. Considérons
la normalisation $X'$ de $\overline{X}$ dans $X$. D'après
le Lemme \ref{lemma-relative-normalization-finite}, le
morphisme $X' \to \overline{X}$ est fini. D'après le lemme
des morphismes \ref{morphisms-lemma-finite-projective}, $X'
\to \overline{X}$ est projectif. D'après le lemme des
morphismes
\ref{morphisms-lemma-projective-over-quasi-projective-is-H-projective},
nous voyons que $X' \to \overline{X}$ est H-projectif. D'après
le lemme des morphismes
\ref{morphisms-lemma-H-projective-composition}, nous voyons que $X'
\to \Spec(k)$ est H-projectif. Supposons que $\{x'_1, \ldots, x'_m\}
\subset X'$ soit l'image inverse de $\{x_1, \ldots, x_n\} =
\overline{X} \setminus X$. Alors $\dim(\mathcal{O}_{X', x'_i}) = 1$
pour tout $1 \leq i \leq m$. Par conséquent, les anneaux locaux
$\mathcal{O}_{X', x'}$ sont des anneaux à valuation discrète d'après le
lemme des morphismes
\ref{morphisms-lemma-relative-normalization-normal-codim-1}. Ensuite, $X \to X'$ et $\{x'_1, \ldots, x'_m\}$ sont comme souhaité.
\end{proof}

\begin{lemma}
\label{lemma-dim-1-projective-completion-after-insep}
Supposons que $X$ soit un schéma séparé de type fini sur $k$
avec $\dim(X) \leq 1$. Alors il existe un diagramme commutatif
$$
\xymatrix{
\overline{Y}_1 \amalg \ldots \amalg \overline{Y}_n \ar[rd] &
Y_1 \amalg \ldots \amalg Y_n \ar[r]_-\nu \ar[d] \ar[l]^j &
X_{k'} \ar[r] \ar[d] &
X \ar[d]^f \\
& \Spec(k'_1) \amalg \ldots \amalg \Spec(k'_n) \ar[r] &
\Spec(k') \ar[r] &
\Spec(k)
}
$$
de schémas avec les propriétés suivantes :
\begin{enumerate}
\item $k'/k$ est une extension finie purement inséparable
de corps, \item $\nu$ est la normalisation de $X_{k'}$,
\item $j$ est une immersion ouverte avec image dense, \item
$k'_i/k'$ est une extension finie séparable pour $i = 1,
\ldots, n$, \item $\overline{Y}_i$ est lisse, projective,
géométriquement irréductible de dimension $\leq 1$ sur $k'_i$.
\end{enumerate}
\end{lemma}

\begin{proof}
Comme nous pouvons remplacer $X$ par sa réduction, nous
pouvons et supposons que $X$ est réduit. Choisissons $X \to
\overline{X}$ comme dans le Lemme
\ref{lemma-reduced-dim-1-projective-completion}. Si nous pouvons
démontrer le lemme pour $\overline{X}$, alors le lemme s'ensuit pour $X$
(détails omis). Ainsi, nous pouvons et supposons que $X$ est projectif.

\medskip\noindent
Choisissez $k'/k$ fini purement inséparable tel que la
normalisation de $X_{k'}$ soit géométriquement normale sur
$k'$, voir le Lemme
\ref{lemma-finite-extension-geometrically-normal}. Notons $Y
= (X_{k'})^\nu$ la normalisation ; pour les propriétés de la
normalisation, voir la Section \ref{section-normalization}.
Alors $Y$ est géométriquement régulier car normal et régulier
sont les mêmes en dimension $\leq 1$, voir Propriétés, Lemme
\ref{properties-lemma-normal-dimension-1-regular}. Par
conséquent, $Y$ est lisse sur $k'$ selon le Lemme
\ref{lemma-geometrically-regular-smooth}. Supposons que $Y = Y_1
\amalg \ldots \amalg Y_n$ soit la décomposition de $Y$ en
composantes irréductibles. Posons $k'_i = \Gamma(Y_i,
\mathcal{O}_{Y_i})$. Ce sont des extensions finies séparables de $k'$
d'après le Lemme
\ref{lemma-proper-geometrically-reduced-global-sections}. La démonstration est terminée par le Lemme \ref{lemma-baby-stein}.
\end{proof}

\begin{lemma}
\label{lemma-regular-point-on-curve}
Supposons que $k$ est un corps. Supposons que $X$ est une courbe sur $k$. Supposons
que $x \in X$ est un point fermé. Nous considérons $x$ comme un sous-schéma fermé
(réduit) de $X$ avec faisceau d'idéaux $\mathcal{I}$. Les énoncés suivants sont équivalents
\begin{enumerate}
\item $\mathcal{O}_{X, x}$ est régulier, \item
$\mathcal{O}_{X, x}$ est normal, \item
$\mathcal{O}_{X, x}$ est un anneau de valuation discrète,
\item $\mathcal{I}$ est un $\mathcal{O}_X$-module
inversible, \item $x$ est un diviseur de Cartier effectif sur $X$.
\end{enumerate}
Si $k$ est parfait ou si $\kappa(x)$ est
séparable sur $k$, ceux-ci sont également équivalents à
\begin{enumerate}
\item [(6)] $X \to \Spec(k)$ est lisse à $x$.
\end{enumerate}
\end{lemma}

\begin{proof}
Puisque $X$ est une courbe, l'anneau local $\mathcal{O}_{X, x}$ est un
domaine local noethérien de dimension $1$ (Lemme
\ref{lemma-dimension-locally-algebraic}). Les parties (4) et (5) sont
équivalentes par définition et sont équivalentes au fait que $\mathcal{I}_x =
\mathfrak m_x \subset \mathcal{O}_{X, x}$ a un seul générateur (Diviseurs, Lemme
\ref{divisors-lemma-effective-Cartier-in-points}). L'équivalence de (1), (2),
(3), (4) et (5) découle donc du Lemme \ref{algebra-lemma-characterize-dvr} en
Algèbre. L'énoncé final découle du Lemme
\ref{lemma-dense-smooth-open-variety-over-perfect-field} dans le cas où $k$ est parfait. Si
$\kappa(x)/k$ est séparable, alors l'équivalence découle du Lemme \ref{algebra-lemma-separable-smooth} en Algèbre.
\end{proof}

\begin{remark}
\label{remark-smooth-projective-compactification}
Supposons que $k$ soit un corps. Supposons que $X$ soit une
courbe régulière sur $k$. D'après les lemmes
\ref{lemma-regular-point-on-curve} et
\ref{lemma-reduced-dim-1-projective-completion}, il existe une courbe
projective non singulière $\overline{X}$ qui est une compactification
de $X$, c'est-à-dire qu'il existe une immersion ouverte $j : X \to
\overline{X}$ telle que le complément consiste en un nombre fini de
points fermés. Si $k$ est parfait, alors $X$ et $\overline{X}$ sont lisses
sur $k$ et $\overline{X}$ est une compactification projective lisse de $X$.
\end{remark}

\noindent
Remarquez que si un schéma affine $X$ sur $k$ est
propre sur $k$, alors $X$ est fini sur $k$ (morphismes,
Lemme \ref{morphisms-lemma-finite-proper}) et possède
donc une dimension $0$ (Algèbre, Lemme
\ref{algebra-lemma-finite-dimensional-algebra} et
Proposition \ref{algebra-proposition-dimension-zero-ring}).
Ainsi, un schéma de dimension $> 0$ sur $k$ ne peut pas être
à la fois affine et propre sur $k$. Par conséquent, les
possibilités dans le lemme suivant sont mutuellement exclusives.

\begin{lemma}
\label{lemma-curve-affine-projective}
Supposons que $X$ soit une courbe sur $k$. Alors soit $X$
est un schéma affine, soit $X$ est H-projectif sur $k$.
\end{lemma}

\begin{proof}
Choisissez $X \to \overline{X}$ avec $\overline{X} \setminus
X = \{x_1, \ldots, x_r\}$ comme dans le Lemme
\ref{lemma-reduced-dim-1-projective-completion}. Alors
$\overline{X}$ est aussi une courbe. Si $r = 0$, alors $X =
\overline{X}$ est H-projectif sur $k$. Ainsi, nous pouvons
supposer $r \geq 1$ et notre objectif est de montrer que $X$ est
affine. D'après le Lemme \ref{lemma-open-in-affine-curve-affine}, il
suffit de montrer que $\overline{X} \setminus \{x_1\}$ est affine.
Cela nous ramène à l'affirmation énoncée dans le paragraphe suivant.

\medskip\noindent
Supposons que $X$ soit une courbe H-projective sur $k$.
Supposons que $x \in X$ soit un point fermé tel que
$\mathcal{O}_{X, x}$ soit un anneau de valuation discrète.
Affirmation : $U = X \setminus \{x\}$ est affine. D'après le
Lemme \ref{lemma-regular-point-on-curve}, le point $x$ définit un
diviseur de Cartier effectif de $X$. Pour $n \geq 1$, notons $nx =
x + \ldots + x$ la somme en $n$ exemplaires, voir Diviseurs,
Définition
\ref{divisors-definition-sum-effective-Cartier-divisors}. Notons
$\mathcal{O}_{nx}$ le faisceau structurel de $nx$ vu comme un module cohérent sur $X$.
Puisque tout module inversible sur le schéma local $nx$ est trivial, la première suite
exacte courte de Diviseurs, Remarque \ref{divisors-remark-ses-regular-section} se lit
$$
0 \to \mathcal{O}_X
\xrightarrow{1} \mathcal{O}_X(nx) \to \mathcal{O}_{nx} \to 0
$$
dans notre cas. Notez que $\dim_k H^0(X, \mathcal{O}_{nx})
\geq n$. À savoir, d'après le Lemme
\ref{lemma-chi-tensor-finite}, nous avons $H^0(X,
\mathcal{O}_{nx}) = \mathcal{O}_{X, x}/(\pi^n)$ où $\pi$ dans
$\mathcal{O}_{X, x}$ est un uniformisateur et les puissances
$\pi^i$ se mappent sur des éléments $k$ linéairement
indépendants dans $\mathcal{O}_{X, x}/(\pi^n)$ pour $i = 0, 1,
\ldots, n - 1$. Nous avons $\dim_k H^1(X, \mathcal{O}_X) <
\infty$ d'après Cohomologie des schémas, Lemme
\ref{coherent-lemma-proper-over-affine-cohomology-finite}. Si $n >
\dim_k H^1(X, \mathcal{O}_X)$, nous concluons à partir de la suite
exacte longue en cohomologie qu'il existe un $s \in \Gamma(X,
\mathcal{O}_X(nx))$ qui n'est pas une section de $\mathcal{O}_X$. Si
nous prenons $n$ minimal avec cette propriété, alors $s$ se mappent
sur un générateur de la germe $\left(\mathcal{O}_X(nx)\right)_x$,
sinon il définirait une section de $\mathcal{O}_X((n - 1)x) \subset
\mathcal{O}_X(nx)$. Pour ce $n$, nous concluons que $s_0 = 1$ et $s_1
= s$ génèrent le module inversible $\mathcal{L} = \mathcal{O}_X(nx)$.

\medskip\noindent
Considérons le morphisme correspondant $f = \varphi_{\mathcal{L},
(s_0, s_1)} : X \to \mathbf{P}^1_k$ des Constructions, section
\ref{constructions-section-projective-space}. Observons que
l'image inverse de $D_{+}(T_0)$ est $U = X \setminus \{x\}$
puisque la section $s_0$ de $\mathcal{L}$ ne s'annule qu'en $x$.
En particulier, $f$ n'est pas constante, c'est-à-dire que $\Im(f)$
a plus d'un point. Ainsi, $f$ doit envoyer le point générique
$\eta$ de $X$ au point générique de $\mathbf{P}^1_k$. Donc, si $y
\in \mathbf{P}^1_k$ est un point fermé, alors $f^{-1}(\{y\})$ est
un ensemble fermé de $X$ ne contenant pas $\eta$, donc fini. Enfin,
$f$ est propre \footnote {. À savoir, une variété H-projective est
une variété propre par morphismes, lemme
\ref{morphisms-lemma-projective-is-quasi-projective-proper}. Un
morphisme de variétés dont la source est une variété propre est un
morphisme propre par morphismes, lemme
\ref{morphisms-lemma-image-proper-scheme-closed}. }. Par la cohomologie
des schémas, Lemme
\ref{coherent-lemma-proper-finite-fibre-finite-in-neighbourhood} \footnote {
On peut éviter d'utiliser ce lemme qui repose sur le théorème des fonctions
formelles. À savoir, $X$ est projectif donc il suffit de montrer qu'un
morphisme propre $f : X \to Y$ à fibres finies entre des schémas
quasi-projectifs sur $k$ est fini. Pour ce faire, on choisit un ouvert affine
de $X$ contenant la fibre de $f$ au-dessus d'un point $y$ en utilisant le fait
que tout ensemble fini de points d'un schéma quasi-projectif sur $k$ est
contenu dans un affine. En restreignant $Y$ à un petit voisinage affine de $y$,
on se ramène au cas d'un morphisme propre entre affines. Un tel morphisme est
fini par morphismes, Lemme \ref{morphisms-lemma-integral-universally-closed}. }
on conclut que $f$ est fini. Par conséquent, $U = f^{-1}(D_{+}(T_0))$ est affine.
\end{proof}

\noindent
Le lemme suivant combiné avec le Lemme \ref{lemma-proper-minus-point}
nous indique que, étant donné un schéma séparé $X$ de dimension $1$ et
de type fini sur $k$, alors $X \setminus Z$ est affine, dès que le
sous-ensemble fermé $Z$ rencontre chaque composante irréductible de $X$.

\begin{lemma}
\label{lemma-dim-1-nonproper-affine}
Supposons que $X$ soit un schéma séparé de type fini sur
$k$. Si $\dim(X) \leq 1$ et aucune composante irréductible
de $X$ n'est propre de dimension $1$, alors $X$ est affine.
\end{lemma}

\begin{proof}
Supposons que $X = \bigcup X_i$ soit la décomposition de $X$ en
composantes irréductibles. Nous considérons $X_i$ comme un
schéma intègre (en utilisant la structure de schéma induite
réduite, voir schémas, Définition
\ref{schemes-definition-reduced-induced-scheme}). En particulier,
$X_i$ est un singleton (et donc affine) ou une courbe, donc affine
d'après le Lemme \ref{lemma-curve-affine-projective}. Ensuite,
$\coprod X_i \to X$ est fini et surjectif et $\coprod X_i$ est affine.
Ainsi, nous voyons que $X$ est affine d'après la Cohomologie des schémas,
Lemme \ref{coherent-lemma-image-affine-finite-morphism-affine-Noetherian}.
\end{proof}






\section{Degrés sur les courbes}
\label{section-divisors-curves}

\noindent
Nous commençons par définir le degré d'un faisceau
inversible et, plus généralement, d'un faisceau localement
libre sur un schéma propre de dimension $1$ sur un corps.
Dans la Section \ref{section-euler}, nous avons défini la
caractéristique d'Euler d'un faisceau cohérent
$\mathcal{F}$ sur un schéma propre $X$ sur un corps $k$ par la formule
$$
\chi(X, \mathcal{F}) = \sum (-1)^i \dim_k H^i(X, \mathcal{F}).
$$

\begin{definition}
\label{definition-degree-invertible-sheaf}
Soit $k$ un corps, soit $X$ un schéma propre de dimension
$\leq 1$ sur $k$, et soit $\mathcal{L}$ un
$\mathcal{O}_X$-module inversible. Le {\it degré } de $\mathcal{L}$ est défini par
$$
\deg(\mathcal{L}) = \chi(X, \mathcal{L}) - \chi(X, \mathcal{O}_X)
$$
Plus généralement, si $\mathcal{E}$ est un faisceau localement
libre de rang $n$ nous définissons le {\it degré } de $\mathcal{E}$ par
$$
\deg(\mathcal{E}) = \chi(X, \mathcal{E}) - n\chi(X, \mathcal{O}_X)
$$
\end{definition}

\noindent
Remarquez que cela dépend du triple $\mathcal{E}/X/k$. Si $X$ est
déconnecté et que $\mathcal{E}$ est localement libre de rang fini
(mais pas de rang constant), alors on peut modifier la définition
en sommant les degrés de la restriction de $\mathcal{E}$ aux
composantes connexes de $X$. Si $\mathcal{E}$ n'est qu'un faisceau
cohérent, il existe plusieurs façons différentes d'étendre la
définition \footnote {. Si $X$ est une courbe propre et que
$\mathcal{F}$ est un faisceau cohérent sur $X$, alors on définit souvent
le degré comme $\chi(X, \mathcal{F}) - r\chi(X, \mathcal{O}_X)$ où $r =
\dim_{\kappa(\xi)} \mathcal{F}_\xi$ est le rang de $\mathcal{F}$ au
point générique $\xi$ de $X$. }. Dans une série de lemmes, nous montrons
que cette définition possède toutes les propriétés que l'on attend du degré.

\begin{lemma}
\label{lemma-degree-base-change}
Supposons que $k'/k$ soit une extension de corps. Supposons
que $X$ soit un schéma propre de dimension $\leq 1$ sur $k$.
Supposons que $\mathcal{E}$ soit un $\mathcal{O}_X$ -module
localement libre de rang constant $n$. Alors le degré de
$\mathcal{E}/X/k$ est égal au degré de $\mathcal{E}_{k'}/X_{k'}/k'$.
\end{lemma}

\begin{proof}
Plus précisément, posons $X_{k'} = X
\times_{\Spec(k)} \Spec(k')$. Soit
$\mathcal{E}_{k'} = p^*\mathcal{E}$ où $p : X_{k'}
\to X$ est la projection. D'après la cohomologie des
schémas, Lemme
\ref{coherent-lemma-flat-base-change-cohomology} nous avons
$H^i(X_{k'}, \mathcal{E}_{k'}) = H^i(X, \mathcal{E})
\otimes_k k'$ et $H^i(X_{k'}, \mathcal{O}_{X_{k'}}) = H^i(X,
\mathcal{O}_X) \otimes_k k'$. Ainsi, nous voyons que les
caractéristiques d'Euler ne changent pas, donc le degré ne change pas.
\end{proof}

\begin{lemma}
\label{lemma-degree-additive}
Supposons que $k$ soit un corps. Supposons que $X$ soit un schéma propre
de dimension $\leq 1$ sur $k$. Supposons que $0 \to \mathcal{E}_1 \to
\mathcal{E}_2 \to \mathcal{E}_3 \to 0$ soit une suite exacte courte de
$\mathcal{O}_X$-modules localement libres, chacun de rang constant fini. Alors
$$
\deg(\mathcal{E}_2) = \deg(\mathcal{E}_1) + \deg(\mathcal{E}_3)
$$
\end{lemma}

\begin{proof}
Suit immédiatement de l'additivité des
caractéristiques d'Euler (Lemme
\ref{lemma-euler-characteristic-additive}) et de l'additivité des rangs.
\end{proof}

\begin{lemma}
\label{lemma-degree-birational-pullback}
Supposons que $k$ soit un corps. Supposons que $f : X' \to X$ soit un
morphisme birationnel de schémas propres de dimension $\leq 1$ sur $k$. Alors
$$
\deg(f^*\mathcal{E}) = \deg(\mathcal{E})
$$
pour chaque faisceau localement libre fini de rang constant.
Plus généralement, il suffit que $f$ induise une bijection
entre les composantes irréductibles de dimension $1$ et les
isomorphismes des anneaux locaux aux points génériques correspondants.
\end{lemma}

\begin{proof}
Le morphisme $f$ est propre (morphismes,
Lemme
\ref{morphisms-lemma-image-proper-scheme-closed}) et a
des fibres de dimension $\leq 0$. Par conséquent, $f$
est fini (Cohomologie des schémas, Lemme
\ref{coherent-lemma-proper-finite-fibre-finite-in-neighbourhood}). Ainsi
$$
Rf_*f^*\mathcal{E} = f_*f^*\mathcal{E} =
\mathcal{E} \otimes_{\mathcal{O}_X} f_*\mathcal{O}_{X'}
$$
Puisque $f$ induit un isomorphisme sur les anneaux locaux
aux points génériques de toutes les composantes
irréductibles de dimension $1$, nous voyons que le noyau et le conoyau
$$
0 \to \mathcal{K} \to \mathcal{O}_X \to f_*\mathcal{O}_{X'}
\to \mathcal{Q} \to 0
$$
avoir des supports de dimension $\leq 0$. Notez qu’en
tensorisant cela avec $\mathcal{E}$, il s’agit toujours d’une
suite exacte puisque $\mathcal{E}$ est localement libre. Nous obtenons
\begin{align*}
\chi(X, \mathcal{E}) - \chi(X', f^*\mathcal{E})
& =
\chi(X, \mathcal{E}) - \chi(X, f_*f^*\mathcal{E}) \\
& =
\chi(X, \mathcal{E}) - \chi(X, \mathcal{E} \otimes f_*\mathcal{O}_{X'}) \\
& =
\chi(X, \mathcal{K} \otimes \mathcal{E}) -
\chi(X, \mathcal{Q} \otimes \mathcal{E}) \\
& =
n\chi(X, \mathcal{K}) -
n\chi(X, \mathcal{Q}) \\
& =
n\chi(X, \mathcal{O}_X) - n\chi(X, f_*\mathcal{O}_{X'}) \\
& =
n\chi(X, \mathcal{O}_X) - n\chi(X', \mathcal{O}_{X'})
\end{align*}
ce qui prouve ce que nous voulons. La première égalité car $f$
est fini, voir Cohomologie des schémas, Lemme
\ref{coherent-lemma-relative-affine-cohomology}. La deuxième
égalité par la formule de projection, voir Cohomologie, Lemme
\ref{cohomology-lemma-projection-formula}. La troisième par
l'additivité des caractéristiques d'Euler, voir Lemme
\ref{lemma-euler-characteristic-additive}. La quatrième par Lemme \ref{lemma-chi-tensor-finite}.
\end{proof}

\begin{lemma}
\label{lemma-degree-on-proper-curve}
Supposons que $k$ soit un corps. Supposons que $X$ soit une courbe propre sur $k$ avec un
point générique $\xi$. Supposons que $\mathcal{E}$ soit un $\mathcal{O}_X$ -module
localement libre de rang $n$ et soit $\mathcal{F}$ un $\mathcal{O}_X$ -module cohérent. Alors
$$
\chi(X, \mathcal{E} \otimes \mathcal{F}) =
r \deg(\mathcal{E}) + n \chi(X, \mathcal{F})
$$
où $r = \dim_{\kappa(\xi)} \mathcal{F}_\xi$ est le rang de $\mathcal{F}$.
\end{lemma}

\begin{proof}
Supposons que $\mathcal{P}$ soit la propriété des faisceaux
cohérents $\mathcal{F}$ sur $X$ exprimant que la formule du lemme
est vérifiée. Nous affirmons que les hypothèses (1) et (2) de
Cohomologie des schémas, Lemme \ref{coherent-lemma-property} sont
vérifiées pour $\mathcal{P}$. À savoir, (1) est vérifiée parce que la
caractéristique d'Euler et le rang $r$ sont additifs dans les suites
exactes courtes de faisceaux cohérents. Et (2) est également vérifiée
: Si $Z = X$ alors nous pouvons prendre $\mathcal{G} = \mathcal{O}_X$
et $\mathcal{P}(\mathcal{O}_X)$ est vrai par la définition du degré.
Si $i : Z \to X$ est l'inclusion d'un point fermé, nous pouvons prendre
$\mathcal{G} = i_*\mathcal{O}_Z$ et $\mathcal{P}$ est vérifié par le
Lemme \ref{lemma-chi-tensor-finite} et le fait que $r = 0$ dans ce cas.
\end{proof}

\noindent
Soit $k$ un corps. Soit $X$ un schéma de type fini sur
$k$, de dimension $\leq 1$. Soient $C_i \subset X$, $i = 1,
\ldots, t$, ses composantes irréductibles de dimension $1$.
Nous munissons $C_i$ de la structure de schéma réduit
induite. Soit $\xi_i \in C_i$ son point générique. La {\it
multiplicité de $C_i$ dans $X$ } est définie comme la longueur
$$
m_i = \text{length}_{\mathcal{O}_{X, \xi_i}} \mathcal{O}_{X, \xi_i}
$$
Cela a du sens car $\mathcal{O}_{X, \xi_i}$ est un anneau local de
Noether de dimension zéro et a donc une longueur finie sur lui-même
(Algèbre, Proposition \ref{algebra-proposition-dimension-zero-ring}). Voir
Homologie de Chow, Section \ref{chow-section-cycle-of-closed-subscheme} pour
des informations supplémentaires. Il s'avère que le degré d'un faisceau
localement libre dépend uniquement de la restriction des composantes irréductibles.

\begin{lemma}
\label{lemma-degree-in-terms-of-components}
Supposons que $k$ soit un corps. Supposons que $X$ soit un schéma
propre de dimension $\leq 1$ sur $k$. Supposons que $\mathcal{E}$
soit un $\mathcal{O}_X$-module localement libre de rang $n$. Alors
$$
\deg(\mathcal{E}) = \sum m_i \deg(\mathcal{E}|_{C_i})
$$
où $C_i \subset X$, $i = 1, \ldots, t$ sont les composantes
irréductibles de dimension $1$ avec la structure de schéma
réduite induite et $m_i$ est la multiplicité de $C_i$ dans $X$.
\end{lemma}

\begin{proof}
Remarquez que l'énoncé a du sens car $C_i \to \Spec(k)$ est
propre de dimension $1$ ; nous utiliserons ci-dessous que tout
sous-schéma fermé de $X$ est propre sur $k$ par morphismes, les
lemmes \ref{morphisms-lemma-closed-immersion-proper} et
\ref{morphisms-lemma-composition-proper}. Considérons le
sous-schéma ouvert $U_i = X \setminus (\bigcup_{j \not = i} C_j)$
et soit $X_i \subset X$ l'adhérence schématique de $U_i$. Notez
que $X_i \cap U_i = U_i$ (schématiquement) et que $X_i \cap U_j =
\emptyset$ (au niveau des ensembles) pour $i \not = j$ ; cela
découle de la description de l'adhérence schématique dans
morphismes, lemme \ref{morphisms-lemma-quasi-compact-immersion}.
Ainsi, nous pouvons appliquer le lemme
\ref{lemma-degree-birational-pullback} au morphisme $X' = \coprod X_i \to
X$. Puisqu'il est clair que $C_i \subset X_i$ (schématiquement) et que la
multiplicité de $C_i$ dans $X_i$ est égale à la multiplicité de $C_i$ dans $X$,
nous voyons que nous nous réduisons au cas discuté dans le paragraphe suivant.

\medskip\noindent
Supposons que $X$ soit irréductible avec le point générique $\xi$.
Supposons que $C = X_{red}$ ait une multiplicité $m$. Nous devons
montrer que $\deg(\mathcal{E}) = m \deg(\mathcal{E}|_C)$. Soit
$\mathcal{I} \subset \mathcal{O}_X$ l'idéal définissant le sous-schéma
fermé $C$. Supposons que $e \geq 0$ soit minimal tel que
$\mathcal{I}^{e + 1} = 0$ (Cohomologie des schémas, Lemme
\ref{coherent-lemma-power-ideal-kills-sheaf}). Nous argumentons par
récurrence sur $e$. Si $e = 0$, alors $X = C$ et le résultat est immédiat.
Sinon, nous posons $\mathcal{F} = \mathcal{I}^e$ considéré comme un
$\mathcal{O}_C$-module cohérent (Cohomologie des schémas, Lemme
\ref{coherent-lemma-i-star-equivalence}). Supposons que $X' \subset X$ soit le
sous-schéma fermé défini par l'idéal cohérent $\mathcal{I}^e$ et soit $m'$ la
multiplicité de $C$ dans $X'$. En prenant les fibres en $\xi$ de la suite exacte courte.
$$
0 \to \mathcal{F} \to \mathcal{O}_X \to \mathcal{O}_{X'} \to 0
$$
nous trouvons (utiliser l'algèbre, les lemmes
\ref{algebra-lemma-length-additive},
\ref{algebra-lemma-dimension-is-length} et \ref{algebra-lemma-length-independent}) que
$$
m = \text{length}_{\mathcal{O}_{X, \xi}} \mathcal{O}_{X, \xi}
= \dim_{\kappa(\xi)} \mathcal{F}_\xi +
\text{length}_{\mathcal{O}_{X', \xi}} \mathcal{O}_{X', \xi}
= r + m'
$$
où $r$ est le rang de $\mathcal{F}$ en tant que faisceau cohérent sur $C$.
En tensorisant avec $\mathcal{E}$, nous obtenons une suite exacte courte
$$
0 \to \mathcal{E}|_C \otimes \mathcal{F} \to \mathcal{E} \to
\mathcal{E} \otimes \mathcal{O}_{X'} \to 0
$$
Par induction nous avons
$\deg(\mathcal{E}|_{X'}) = m'
\deg(\mathcal{E}|_C)$. Par le Lemme
\ref{lemma-degree-on-proper-curve} nous avons
$\chi(\mathcal{E}|_C \otimes \mathcal{F}) = r \deg(\mathcal{E}|_C) + n
\chi(\mathcal{F})$. En mettant tout ensemble, nous obtenons le résultat.
\end{proof}

\begin{lemma}
\label{lemma-degree-tensor-product}
Supposons que $k$ soit un corps, que $X$ soit un schéma propre de
dimension $\leq 1$ sur $k$, et que $\mathcal{E}$, $\mathcal{V}$ soient
des $\mathcal{O}_X$-modules localement libres de rang fini constant. Alors
$$
\deg(\mathcal{E} \otimes \mathcal{V}) =
\text{rank}(\mathcal{E}) \deg(\mathcal{V}) +
\text{rank}(\mathcal{V}) \deg(\mathcal{E})
$$
\end{lemma}

\begin{proof}
Par le Lemme \ref{lemma-degree-in-terms-of-components} et
l'arithmétique élémentaire, nous nous réduisons au cas d'une courbe
propre. Ce cas découle du Lemme \ref{lemma-degree-on-proper-curve}.
\end{proof}

\begin{lemma}
\label{lemma-degree-and-det}
Supposons que $k$ soit un corps, que $X$ soit un schéma
approprié de dimension $\leq 1$ sur $k$, et que $\mathcal{E}$
soit un $\mathcal{O}_X$-module localement libre de rang $n$. Alors
$$
\deg(\mathcal{E}) = \deg(\wedge^n(\mathcal{E})) = \deg(\det(\mathcal{E}))
$$
\end{lemma}

\begin{proof}
D'après le Lemme \ref{lemma-degree-in-terms-of-components} et
l'arithmétique élémentaire, nous nous ramenons au cas d'une
courbe propre. Il existe alors une modification $f : X' \to X$
telle que $f^*\mathcal{E}$ possède une filtration dont les
quotients successifs sont des modules inversibles, voir Diviseurs,
Lemme \ref{divisors-lemma-filter-after-modification}. D'après le
Lemme \ref{lemma-degree-birational-pullback}, nous pouvons travailler
sur $X'$. Ainsi, nous pouvons supposer que nous avons une filtration
$$
0 = \mathcal{E}_0 \subset \mathcal{E}_1 \subset \mathcal{E}_2 \subset
\ldots \subset \mathcal{E}_n = \mathcal{E}
$$
par des $\mathcal{O}_X$-modules localement libres avec
$\mathcal{L}_i = \mathcal{E}_i/\mathcal{E}_{i - 1}$ inversible. Par le
Lemme sur les modules \ref{modules-lemma-det-ses} et par induction,
nous trouvons $\det(\mathcal{E}) = \mathcal{L}_1 \otimes \ldots
\otimes \mathcal{L}_n$. Ainsi, l'égalité découle du Lemme
\ref{lemma-degree-tensor-product} et de l'additivité (Lemme \ref{lemma-degree-additive}).
\end{proof}

\begin{lemma}
\label{lemma-degree-effective-Cartier-divisor}
Supposons que $k$ soit un corps, supposons que $X$ soit un schéma
propre de dimension $\leq 1$ sur $k$. Supposons que $D$ soit un
diviseur de Cartier effectif sur $X$. Alors $D$ est fini sur
$\Spec(k)$ de degré $\deg(D) = \dim_k \Gamma(D, \mathcal{O}_D)$.
Pour un faisceau localement libre $\mathcal{E}$ de rang $n$, nous avons
$$
\deg(\mathcal{E}(D)) = n\deg(D) + \deg(\mathcal{E})
$$
où $\mathcal{E}(D) = \mathcal{E} \otimes_{\mathcal{O}_X} \mathcal{O}_X(D)$.
\end{lemma}

\begin{proof}
Puisque $D$ est nulle part dense dans $X$ (Diviseurs,
Lemme
\ref{divisors-lemma-complement-effective-Cartier-divisor}),
nous voyons que $\dim(D) \leq 0$. Par conséquent, $D$ est
fini sur $k$ selon le Lemme
\ref{lemma-algebraic-scheme-dim-0}. Comme $k$ est un corps,
le morphisme $D \to \Spec(k)$ est fini localement libre et a
donc un degré (morphismes, Définition
\ref{morphisms-definition-finite-locally-free}), qui est clairement
égal à $\dim_k \Gamma(D, \mathcal{O}_D)$ comme indiqué dans le
lemme. D'après Diviseurs, Définition
\ref{divisors-definition-invertible-sheaf-effective-Cartier-divisor}, il existe une suite exacte courte
$$
0 \to \mathcal{O}_X \to \mathcal{O}_X(D) \to i_*i^*\mathcal{O}_X(D) \to 0
$$
où $i : D \to X$ est l'immersion fermée. En tensorisant
avec $\mathcal{E}$, nous obtenons une suite exacte courte
$$
0 \to \mathcal{E} \to \mathcal{E}(D) \to i_*i^*\mathcal{E}(D) \to 0
$$
L'équation du lemme découle de l'additivité de la
caractéristique d'Euler (Lemme
\ref{lemma-euler-characteristic-additive}) et du Lemme \ref{lemma-chi-tensor-finite}.
\end{proof}

\begin{lemma}
\label{lemma-divisible}
Supposons que $k$ soit un corps. Supposons que $X$ soit
un schéma propre sur $k$ qui est réduit et connecté. Soit
$\kappa = H^0(X, \mathcal{O}_X)$. Alors $\kappa/k$ est
une extension finie de corps et $w = [\kappa : k]$ divise
\begin{enumerate}
\item $\deg(\mathcal{E})$ pour tous les
$\mathcal{O}_X$ -modules localement libres
$\mathcal{E}$, \item $[\kappa(x) : k]$ pour tous les
points fermés $x \in X$, et \item $\deg(D)$ pour tous
les sous-schéma fermés $D \subset X$ de dimension zéro.
\end{enumerate}
\end{lemma}

\begin{proof}
Voir le lemme
\ref{lemma-proper-geometrically-reduced-global-sections} pour
les assertions concernant $\kappa$. Pour tout
$\mathcal{O}_X$-module quasi-cohérent, les $k$-espaces vectoriels
$H^i(X, \mathcal{F})$ sont des $\kappa$-espaces vectoriels. Les
divisibilités découlent immédiatement de cet énoncé et des définitions.
\end{proof}

\begin{lemma}
\label{lemma-degree-pullback-map-proper-curves}
Supposons que $k$ soit un corps. Supposons que $f : X \to Y$ soit
un morphisme non constant de courbes propres sur $k$. Supposons que
$\mathcal{E}$ soit un $\mathcal{O}_Y$-module localement libre. Alors
$$
\deg(f^*\mathcal{E}) = \deg(X/Y) \deg(\mathcal{E})
$$
\end{lemma}

\begin{proof}
Le degré de $X$ sur $Y$ est défini dans les
morphismes, Définition
\ref{morphisms-definition-degree}. Ainsi $f_*\mathcal{O}_X$
est un $\mathcal{O}_Y$-module cohérent de rang $\deg(X/Y)$,
c’est-à-dire $\deg(X/Y) = \dim_{\kappa(\xi)}
(f_*\mathcal{O}_X)_\xi$ où $\xi$ est le point générique de $Y$. Ainsi nous obtenons
\begin{align*}
\chi(X, f^*\mathcal{E})
& =
\chi(Y, f_*f^*\mathcal{E}) \\
& =
\chi(Y, \mathcal{E} \otimes f_*\mathcal{O}_X) \\
& =
\deg(X/Y) \deg(\mathcal{E}) + n \chi(Y, f_*\mathcal{O}_X) \\
& =
\deg(X/Y) \deg(\mathcal{E}) + n \chi(X, \mathcal{O}_X)
\end{align*}
comme désiré. La première égalité car $f$ est fini, voir
Cohomologie des schémas, Lemme
\ref{coherent-lemma-relative-affine-cohomology}. La deuxième égalité par la
formule de projection, voir Cohomologie, Lemme
\ref{cohomology-lemma-projection-formula}. La troisième égalité par le Lemme \ref{lemma-degree-on-proper-curve}.
\end{proof}

\noindent
Ce qui suit est une conséquence triviale mais
importante des résultats sur les degrés ci-dessus.

\begin{lemma}
\label{lemma-check-invertible-sheaf-trivial}
Supposons que $k$ soit un corps. Supposons que $X$ soit une courbe propre sur
$k$. Supposons que $\mathcal{L}$ soit un $\mathcal{O}_X$-module inversible.
\begin{enumerate}
\item Si $\mathcal{L}$ a une section non nulle, alors
$\deg(\mathcal{L}) \geq 0$. \item Si $\mathcal{L}$ a une
section non nulle $s$ qui s'annule en un point, alors
$\deg(\mathcal{L}) > 0$. \item Si $\mathcal{L}$ et
$\mathcal{L}^{-1}$ ont des sections non nulles, alors
$\mathcal{L} \cong \mathcal{O}_X$. \item Si $\deg(\mathcal{L})
\leq 0$ et $\mathcal{L}$ ont une section non nulle, alors
$\mathcal{L} \cong \mathcal{O}_X$. \item Si $\mathcal{N} \to
\mathcal{L}$ est une application non nulle de modules
inversibles $\mathcal{O}_X$, alors $\deg(\mathcal{L}) \geq
\deg(\mathcal{N})$ et si l'égalité est réalisée, alors c'est un isomorphisme.
\end{enumerate}
\end{lemma}

\begin{proof}
Supposons que $s$ soit une section non nulle de
$\mathcal{L}$. Comme $X$ est une courbe, nous voyons que
$s$ est une section régulière. Par conséquent, il existe
un diviseur de Cartier effectif $D \subset X$ et un
isomorphisme $\mathcal{L} \to \mathcal{O}_X(D)$ envoyant
$s$ la section canonique $1$ de $\mathcal{O}_X(D)$, voir
Diviseurs, Lemme \ref{divisors-lemma-characterize-OD}.
Alors $\deg(\mathcal{L}) = \deg(D)$ d'après le Lemme
\ref{lemma-degree-effective-Cartier-divisor}. Comme
$\deg(D) \geq 0$ et $= 0$ si et seulement si $D =
\emptyset$, cela prouve (1) et (2). Dans le cas (3), nous
voyons que $\deg(\mathcal{L}) = 0$ et $D = \emptyset$. De même
pour (4). Pour voir (5), appliquez (1) et (4) au faisceau inversible
$$
\mathcal{L} \otimes_{\mathcal{O}_X} \mathcal{N}^{\otimes -1} =
\SheafHom_{\mathcal{O}_X}(\mathcal{N}, \mathcal{L})
$$
qui a le degré $\deg(\mathcal{L}) -
\deg(\mathcal{N})$ selon le Lemme \ref{lemma-degree-tensor-product}.
\end{proof}

\begin{lemma}
\label{lemma-no-sections-dual-nef}
Supposons que $k$ soit un corps. Supposons que $X$ soit un
schéma propre sur $k$ qui est réduit, connexe et
équidimensionnel de dimension $1$. Supposons que $\mathcal{L}$ soit
un $\mathcal{O}_X$-module inversible. Si $\deg(\mathcal{L}|_C) \leq
0$ pour toutes les composantes irréductibles $C$ de $X$, alors soit
$H^0(X, \mathcal{L}) = 0$ soit $\mathcal{L} \cong \mathcal{O}_X$.
\end{lemma}

\begin{proof}
Supposons que $s \in H^0(X, \mathcal{L})$ soit non
nul. Puisque $X$ est réduit, il existe une composante
irréductible $C$ de $X$ avec $s|_C \not = 0$. Mais si
$s|_C$ est non nul, alors $s$ ne s'annule nulle part
sur $C$ d'après le Lemme
\ref{lemma-check-invertible-sheaf-trivial}. Cela implique à son
tour que $s$ ne s'annule nulle part sur chaque composante
irréductible de $X$ rencontrant $C$. Puisque $X$ est connexe, nous
en concluons que $s$ ne s'annule nulle part et le lemme s'ensuit.
\end{proof}

\begin{lemma}
\label{lemma-ample-curve}
Supposons que $k$ soit un corps. Supposons que $X$ soit une courbe propre
sur $k$. Supposons que $\mathcal{L}$ soit un $\mathcal{O}_X$-module
inversible. Alors $\mathcal{L}$ est ample si et seulement si $\deg(\mathcal{L}) > 0$.
\end{lemma}

\begin{proof}
Si $\mathcal{L}$ est ample, alors il existe un $n > 0$ et
une section $s \in H^0(X, \mathcal{L}^{\otimes n})$ avec
$X_s$ affine. Comme $X$ n'est pas affine (sinon d'après les
morphismes, le Lemme \ref{morphisms-lemma-finite-proper}
$X$ serait fini), nous voyons que $s$ s'annule en un certain
point. Donc $\deg(\mathcal{L}^{\otimes n}) > 0$ d'après le
Lemme \ref{lemma-check-invertible-sheaf-trivial}. D'après le
Lemme \ref{lemma-degree-tensor-product}, nous concluons que
$\deg(\mathcal{L}) = 1/n\deg(\mathcal{L}^{\otimes n}) > 0$.

\medskip\noindent
Supposons $\deg(\mathcal{L}) > 0$. Ensuite
$$
\dim_k H^0(X, \mathcal{L}^{\otimes n}) \geq \chi(X, \mathcal{L}^n)
= n\deg(\mathcal{L}) + \chi(X, \mathcal{O}_X)
$$
croît linéairement avec $n$. Par conséquent, pour toute
collection finie de points fermés $x_1, \ldots, x_t$ de $X$, nous
pouvons trouver un $n$ tel que $\dim_k H^0(X, \mathcal{L}^{\otimes
n}) > \sum \dim_k \kappa(x_i)$. (Rappelons que par le
Nullstellensatz de Hilbert, les corps d'extension $\kappa(x_i)/k$
sont finis, voir par exemple morphismes, Lemme
\ref{morphisms-lemma-closed-point-fibre-locally-finite-type}). Par
conséquent, nous pouvons trouver un $s \in H^0(X, \mathcal{L}^{\otimes
n})$ non nul s'annulant dans $x_1, \ldots, x_t$. En particulier, si
nous choisissons $x_1, \ldots, x_t$ tel que $X \setminus \{x_1,
\ldots, x_t\}$ soit affine, alors $X_s$ est aussi affine (par exemple
par Propriétés, Lemme \ref{properties-lemma-affine-cap-s-open} bien que
si nous choisissons notre ensemble fini tel que $\mathcal{L}|_{X
\setminus \{x_1, \ldots, x_t\}}$ soit trivial, alors cela est immédiat).
La conclusion est que nous pouvons trouver un $n > 0$ et une section non
nulle $s \in H^0(X, \mathcal{L}^{\otimes n})$ tels que $X_s$ soit affine.

\medskip\noindent
Nous allons montrer que pour tout faisceau
quasi-cohérent d'idéaux $\mathcal{I}$, il existe un $m >
0$ tel que $H^1(X, \mathcal{I} \otimes
\mathcal{L}^{\otimes m})$ soit nul. Cela terminera la preuve par
la cohomologie des schémas, Lemme
\ref{coherent-lemma-vanshing-gives-ample}. Pour voir cela, nous considérons les applications
$$
\mathcal{I} \xrightarrow{s}
\mathcal{I} \otimes \mathcal{L}^{\otimes n} \xrightarrow{s}
\mathcal{I} \otimes \mathcal{L}^{\otimes 2n} \xrightarrow{s} \ldots
$$
Puisque $\mathcal{I}$ est sans torsion, ces applications
sont injectives et des isomorphismes sur $X_s$, donc les
conoyaux ont $H^1$ nul (par la cohomologie des schémas, Lemme
\ref{coherent-lemma-coherent-support-dimension-0} par exemple).
Nous en concluons que les applications des espaces vectoriels
$$
H^1(X, \mathcal{I}) \to
H^1(X, \mathcal{I} \otimes \mathcal{L}^{\otimes n}) \to
H^1(X, \mathcal{I} \otimes \mathcal{L}^{\otimes 2n}) \to \ldots
$$
sont surjectives. D'autre part, la dimension de
$H^1(X, \mathcal{I})$ est finie, et chaque élément
se projette finalement sur zéro par le Lemme de
cohomologie des schémas
\ref{coherent-lemma-section-affine-open-kills-classes}. Ainsi,
pour un certain $e > 0$, nous voyons que $H^1(X, \mathcal{I}
\otimes \mathcal{L}^{\otimes en})$ est nul. Cela termine la preuve.
\end{proof}

\begin{lemma}
\label{lemma-ampleness-in-terms-of-degrees-components}
Supposons que $k$ soit un corps. Supposons que $X$ soit un schéma propre de
dimension $\leq 1$ sur $k$. Supposons que $\mathcal{L}$ soit un $\mathcal{O}_X$
-module inversible. Supposons que $C_i \subset X$, $i = 1, \ldots, t$ soient les
composantes irréductibles de dimension $1$. Les énoncés suivants sont équivalents :
\begin{enumerate}
\item $\mathcal{L}$ est ample, et \item
$\deg(\mathcal{L}|_{C_i}) > 0$ pour $i = 1, \ldots, t$.
\end{enumerate}
\end{lemma}

\begin{proof}
Supposons que $x_1, \ldots, x_r \in X$ soit les
points isolés fermés. Pensez à $x_i =
\Spec(\kappa(x_i))$ comme un schéma. Considérez le morphisme de schémas
$$
f : C_1 \amalg \ldots \amalg C_t \amalg x_1 \amalg \ldots \amalg x_r
\longrightarrow X
$$
Ceci est un morphisme fini surjectif de schémas propres sur
$k$ (détails omis). Ainsi $\mathcal{L}$ est ample si et
seulement si $f^*\mathcal{L}$ est ample (Cohomologie des
schémas, Lemme
\ref{coherent-lemma-surjective-finite-morphism-ample}). Nous concluons ainsi par le Lemme \ref{lemma-ample-curve}.
\end{proof}

\begin{lemma}
\label{lemma-general-degree-g-line-bundle}
Supposons que $k$ soit un corps algébriquement clos.
Supposons que $X$ soit une courbe propre sur $k$. Alors il existe
\begin{enumerate}
\item un $\mathcal{O}_X$ -module inversible $\mathcal{L}$ avec
$\dim_k H^0(X, \mathcal{L}) = 1$ et $H^1(X, \mathcal{L}) = 0$,
et \item un $\mathcal{O}_X$ -module inversible $\mathcal{N}$
avec $\dim_k H^0(X, \mathcal{N}) = 0$ et $H^1(X, \mathcal{N}) = 0$.
\end{enumerate}
\end{lemma}

\begin{proof}
Choisissez une immersion fermée $i : X \to
\mathbf{P}^n_k$ (Lemme
\ref{lemma-dim-1-proper-projective}). En posant $\mathcal{L} =
i^*\mathcal{O}_{\mathbf{P}^n}(d)$ pour $d \gg 0$, nous voyons
qu'il existe un faisceau inversible $\mathcal{L}$ avec $H^0(X,
\mathcal{L}) \not = 0$ et $H^1(X, \mathcal{L}) = 0$ (voir
Cohomologie des schémas, Lemme
\ref{coherent-lemma-vanshing-gives-ample} pour l'annulation et les
références qui y sont mentionnées pour la non-annulation). Nous
terminerons la démonstration de (1) par récurrence décroissante sur $t =
\dim_k H^0(X, \mathcal{L})$. Le cas de base $t = 1$ est trivial. Supposons $t > 1$.

\medskip\noindent
Supposons que $U \subset X$ soit l'ouvert non vide des points
non singuliers étudié dans le Lemme
\ref{lemma-dense-smooth-open-variety-over-perfect-field}. Supposons
que $s \in H^0(X, \mathcal{L})$ soit non nul. Il existe un point fermé
$x \in U$ tel que $s$ ne s'annule pas dans $x$. Supposons que
$\mathcal{I}$ soit le faisceau d'idéaux de $i : x \to X$ comme dans le
Lemme \ref{lemma-regular-point-on-curve}. Considérons la suite exacte courte
$$
0 \to \mathcal{I} \otimes_{\mathcal{O}_X} \mathcal{L} \to
\mathcal{L} \to i_*i^*\mathcal{L} \to 0
$$
Remarquez que $H^0(X, i_*i^*\mathcal{L}) = H^0(x, i^*\mathcal{L})$ a
pour dimension $1$ puisque $x$ est un point $k$ rationnel ($k$ est
algébriquement clos). Puisque $s$ ne s'annule pas en $x$, nous concluons que
$$
H^0(X, \mathcal{L}) \longrightarrow H^0(X, i_*i^*\mathcal{L})
$$
est surjective. Par conséquent $\dim_k H^0(X,
\mathcal{I} \otimes_{\mathcal{O}_X} \mathcal{L}) = t - 1$.
Enfin, la longue suite exacte de cohomologie montre
également que $H^1(X, \mathcal{I} \otimes_{\mathcal{O}_X}
\mathcal{L}) = 0$, terminant ainsi la preuve de l'étape d'induction.

\medskip\noindent
Pour obtenir un faisceau inversible comme dans (2), prenez
un faisceau inversible $\mathcal{L}$ comme dans (1) et
faites l'argument du paragraphe précédent une fois de plus.
\end{proof}

\begin{lemma}
\label{lemma-vanishing-degree-2g-and-1-line-bundle}
Supposons que $k$ soit un corps algébriquement clos. Supposons que $X$
soit une courbe propre sur $k$. Posons $g = \dim_k H^1(X,
\mathcal{O}_X)$. Pour tout $\mathcal{O}_X$-module inversible $\mathcal{L}$
avec $\deg(\mathcal{L}) \geq 2g - 1$, nous avons $H^1(X, \mathcal{L}) = 0$.
\end{lemma}

\begin{proof}
Supposons que $\mathcal{N}$ soit le module inversible que
nous avons trouvé dans le Lemme
\ref{lemma-general-degree-g-line-bundle}, partie (2). Le
degré de $\mathcal{N}$ est $\chi(X, \mathcal{N}) - \chi(X,
\mathcal{O}_X) = 0 - (1 - g) = g - 1$. Donc le degré de
$\mathcal{L} \otimes \mathcal{N}^{\otimes - 1}$ est
$\deg(\mathcal{L}) - (g - 1) \geq g$. Par conséquent, $\chi(X,
\mathcal{L} \otimes \mathcal{N}^{\otimes -1}) \geq g + 1 - g =
1$. Ainsi, il existe une section globale non nulle $s$ dont le
schéma des zéros est un diviseur de Cartier effectif $D$ de degré
$\deg(\mathcal{L}) - (g - 1)$. Cela donne une suite exacte courte
$$
0 \to \mathcal{N} \xrightarrow{s} \mathcal{L} \to i_*(\mathcal{L}|_D) \to 0
$$
où $i : D \to X$ est le morphisme d'inclusion. Nous concluons
que $H^0(X, \mathcal{L})$ se projette isomorphiquement sur
$H^0(D, \mathcal{L}|_D)$ qui a une dimension
$\deg(\mathcal{L}) - (g - 1)$. Le résultat découle de la définition du degré.
\end{proof}






\section{Intersections numériques}
\label{section-num}

\noindent
Dans cette section, nous jouons avec la caractéristique
d'Euler des faisceaux cohérents sur des schémas propres pour
obtenir des nombres d'intersection numériques pour les
modules inversibles. Notre outil principal sera le lemme suivant.

\begin{lemma}
\label{lemma-numerical-polynomial-from-euler}
Supposons que $k$ est un corps. Supposons que $X$ est un schéma
propre sur $k$. Supposons que $\mathcal{F}$ est un
$\mathcal{O}_X$-module cohérent. Supposons que $\mathcal{L}_1, \ldots,
\mathcal{L}_r$ est un $\mathcal{O}_X$-module inversible. L'application
$$
(n_1, \ldots, n_r) \longmapsto
\chi(X, \mathcal{F} \otimes
\mathcal{L}_1^{\otimes n_1} \otimes \ldots \otimes
\mathcal{L}_r^{\otimes n_r})
$$
est un polynôme numérique dans $n_1, \ldots, n_r$ de degré
total au plus égal à la dimension du support de $\mathcal{F}$.
\end{lemma}

\begin{proof}
Nous le prouvons par induction sur
$\dim(\text{Supp}(\mathcal{F}))$. Si ce nombre est zéro, alors la fonction
est constante avec la valeur $\dim_k \Gamma(X, \mathcal{F})$ selon le Lemme
\ref{lemma-chi-tensor-finite}. Supposons $\dim(\text{Supp}(\mathcal{F})) > 0$.

\medskip\noindent
Si $\mathcal{F}$ a des points associés incorporés, alors
nous pouvons considérer la suite exacte courte $0 \to
\mathcal{K} \to \mathcal{F} \to \mathcal{F}' \to 0$
construite dans Diviseurs, Lemme
\ref{divisors-lemma-remove-embedded-points}. Puisque la dimension
du support de $\mathcal{K}$ est strictement moindre, le résultat
est valable pour $\mathcal{K}$ par hypothèse d'induction et avec un
degré total strictement inférieur. Par additivité de la
caractéristique d'Euler (Lemme
\ref{lemma-euler-characteristic-additive}), il suffit de prouver le résultat pour
$\mathcal{F}'$. Ainsi, nous pouvons supposer que $\mathcal{F}$ n'a pas de points associés incorporés.

\medskip\noindent
Si $i : Z \to X$ est une immersion fermée et $\mathcal{F} =
i_*\mathcal{G}$, alors nous voyons que le résultat pour $X$,
$\mathcal{F}$, $\mathcal{L}_1, \ldots, \mathcal{L}_r$ est
équivalent au résultat pour $Z$, $\mathcal{G}$, $i^*\mathcal{L}_1,
\ldots, i^*\mathcal{L}_r$ (puisque les cohomologies sont les
mêmes, voir Cohomologie des schémas, Lemme
\ref{coherent-lemma-relative-affine-cohomology}). En appliquant Diviseurs,
Lemme \ref{divisors-lemma-no-embedded-points-endos}, nous pouvons supposer
que $X$ n'a pas de composantes immergées et $X = \text{Supp}(\mathcal{F})$.

\medskip\noindent
Choisissez une section méromorphe régulière $s$ de
$\mathcal{L}_1$, voir Diviseurs, Lemme
\ref{divisors-lemma-regular-meromorphic-section-exists}. Supposons que $\mathcal{I}
\subset \mathcal{O}_X$ soit l'idéal des dénominateurs de $s$ et considérons les applications
$$
\mathcal{I}\mathcal{F} \to \mathcal{F},\quad
\mathcal{I}\mathcal{F} \to \mathcal{F} \otimes \mathcal{L}_1
$$
des diviseurs, Lemme
\ref{divisors-lemma-make-maps-regular-section}. Ceux-ci sont
injectifs et ont des conoyaux $\mathcal{Q}$, $\mathcal{Q}'$
supportés sur des sous-schémas fermés rares de $X =
\text{Supp}(\mathcal{F})$. Le produit tensoriel avec le module inversible
$\mathcal{L}_1^{\otimes n_1} \otimes \ldots \otimes \mathcal{L}_r^{\otimes
n_r}$ est exact, donc en utilisant à nouveau l'additivité, nous voyons que
\begin{align*}
&\chi(X, \mathcal{F} \otimes \mathcal{L}_1^{\otimes n_1} \otimes \ldots \otimes
\mathcal{L}_r^{\otimes n_r}) -
\chi(X, \mathcal{F} \otimes \mathcal{L}_1^{\otimes n_1 + 1}
\otimes \ldots \otimes \mathcal{L}_r^{\otimes n_r}) \\
& =
\chi(\mathcal{Q} \otimes \mathcal{L}_1^{\otimes n_1} \otimes \ldots \otimes
\mathcal{L}_r^{\otimes n_r}) -
\chi(\mathcal{Q}' \otimes \mathcal{L}_1^{\otimes n_1} \otimes \ldots \otimes
\mathcal{L}_r^{\otimes n_r})
\end{align*}
Ainsi, nous voyons que la fonction $P(n_1,
\ldots, n_r)$ du lemme a la propriété que
$$
P(n_1 + 1, n_2, \ldots, n_r) - P(n_1, \ldots, n_r)
$$
est un polynôme numérique de degré total $<$,
la dimension du support de $\mathcal{F}$. Bien
sûr, par symétrie, la même chose est vraie pour
$$
P(n_1, \ldots, n_{i - 1}, n_i + 1, n_{i + 1}, \ldots, n_r)
- P(n_1, \ldots, n_r)
$$
pour tout $i \in \{1, \ldots, r\}$. Un simple argument
arithmétique montre que $P$ est un polynôme numérique
de degré total au plus $\dim(\text{Supp}(\mathcal{F}))$.
\end{proof}

\noindent
Le lemme suivant montre approximativement que le
coefficient directeur ne dépend que de la longueur du
module cohérent aux points génériques de son support.

\begin{lemma}
\label{lemma-numerical-polynomial-leading-term}
Supposons que $k$ soit un corps. Supposons que $X$ soit un
schéma propre sur $k$. Supposons que $\mathcal{F}$ soit un
$\mathcal{O}_X$-module cohérent. Supposons que
$\mathcal{L}_1, \ldots, \mathcal{L}_r$ soit des
$\mathcal{O}_X$-modules inversibles. Supposons que $d =
\dim(\text{Supp}(\mathcal{F}))$. Soit $Z_i \subset X$ les
composantes irréductibles de $\text{Supp}(\mathcal{F})$ de dimension
$d$. Supposons que $\xi_i \in Z_i$ soit le point générique et posons
$m_i = \text{length}_{\mathcal{O}_{X, \xi_i}}(\mathcal{F}_{\xi_i})$. Alors
$$
\chi(X, \mathcal{F} \otimes \mathcal{L}_1^{\otimes n_1} \otimes \ldots \otimes
\mathcal{L}_r^{\otimes n_r}) -
\sum\nolimits_i
m_i\ \chi(Z_i, \mathcal{L}_1^{\otimes n_1} \otimes \ldots \otimes
\mathcal{L}_r^{\otimes n_r}|_{Z_i})
$$
est un polynôme numérique en $n_1, \ldots, n_r$ de degré total $< d$.
\end{lemma}

\begin{proof}
Considérons les paires $(\xi , Z)$ où $Z \subset X$ est un
sous-schéma fermé intègre de dimension $d$ et $\xi$ est son
point générique. Alors le $\mathcal{O}_{X, \xi}$-module
fini $\mathcal{F}_\xi$ a un support contenu dans $\{\xi\}$,
donc la longueur $m_Z = \text{length}_{\mathcal{O}_{X,
\xi}}(\mathcal{F}_\xi)$ est finie (Algèbre, Lemme
\ref{algebra-lemma-support-point}) et nulle sauf si $Z = Z_i$ pour
un certain $i$. Ainsi, l'expression du lemme peut s'écrire comme
$$
E(\mathcal{F}) =
\chi(X, \mathcal{F} \otimes \mathcal{L}_1^{\otimes n_1} \otimes \ldots \otimes
\mathcal{L}_r^{\otimes n_r}) -
\sum\nolimits
m_Z\ \chi(Z, \mathcal{L}_1^{\otimes n_1} \otimes \ldots \otimes
\mathcal{L}_r^{\otimes n_r}|_Z)
$$
où la somme est effectuée sur les sous-schémas fermés
intègres $Z \subset X$ de dimension $d$. L'attribution
$\mathcal{F} \mapsto E(\mathcal{F})$ est additive dans les
suites exactes courtes $0 \to \mathcal{F} \to \mathcal{F}'
\to \mathcal{F}'' \to 0$ de $\mathcal{O}_X$-modules
cohérents dont le support a une dimension $\leq d$. Cela
découle de l'additivité des caractéristiques d'Euler (Lemme
\ref{lemma-euler-characteristic-additive}) et de l'additivité
des longueurs (Algèbre, Lemme
\ref{algebra-lemma-length-additive}). Appliquons le Lemme
\ref{coherent-lemma-coherent-filter} de la Cohomologie des schémas pour trouver une filtration
$$
0 = \mathcal{F}_0 \subset \mathcal{F}_1 \subset
\ldots \subset \mathcal{F}_m = \mathcal{F}
$$
par des sous-faisceaux cohérents tels que pour
chaque $j = 1, \ldots, m$ il existe un sous-schéma
fermé intègre $V_j \subset X$ et un faisceau d'idéaux
non nul $\mathcal{I}_j \subset \mathcal{O}_{V_j}$ tel que
$$
\mathcal{F}_j/\mathcal{F}_{j - 1}
\cong (V_j \to X)_* \mathcal{I}_j
$$
Il s'ensuit que $V_j \subset \text{Supp}(\mathcal{F})$ et
donc $\dim(V_j) \leq d$. Par l'additivité que nous avons
remarquée ci-dessus, il suffit de prouver le résultat
pour chacun des sous-quotients
$\mathcal{F}_j/\mathcal{F}_{j - 1}$. Ainsi, il suffit de
prouver le résultat lorsque $\mathcal{F} = (V \to
X)_*\mathcal{I}$, où $V \subset X$ est un sous-schéma fermé
intègre de dimension $\leq d$ et $\mathcal{I} \subset
\mathcal{O}_V$ est un faisceau cohérent d'idéaux non nul. Si
$\dim(V) < d$ et plus généralement pour $\mathcal{F}$ dont le
support a une dimension $< d$, alors le premier terme dans
$E(\mathcal{F})$ a un degré total $< d$ d'après le Lemme
\ref{lemma-numerical-polynomial-from-euler} et le second terme est
nul. Si $\dim(V) = d$, alors nous pouvons utiliser la suite exacte courte
$$
0 \to (V \to X)_*\mathcal{I} \to (V \to X)_*\mathcal{O}_V
\to (V \to X)_*(\mathcal{O}_V/\mathcal{I}) \to 0
$$
Le résultat est valable pour le faisceau du
milieu parce que le seul $Z$ apparaissant dans
la somme est $Z = V$ avec $m_Z = 1$ et parce que
$$
H^i(X, ((V \to X)_*\mathcal{O}_V) \otimes 
 \mathcal{L}_1^{\otimes n_1} \otimes \ldots \otimes
\mathcal{L}_r^{\otimes n_r}) =
H^i(V,  \mathcal{L}_1^{\otimes n_1} \otimes \ldots \otimes
\mathcal{L}_r^{\otimes n_r}|_V)
$$
par la formule de projection (Cohomologie, Section
\ref{cohomology-section-projection-formula}) et
Cohomologie des schémas, Lemme
\ref{coherent-lemma-relative-affine-cohomology} ; donc
dans ce cas nous avons en réalité $E(\mathcal{F}) = 0$. Le
résultat est valable pour le faisceau de droite parce que
son support a une dimension $< d$. Ainsi, le résultat est
valable pour le faisceau de gauche et le lemme est prouvé.
\end{proof}

\begin{definition}
\label{definition-intersection-number}
Soit $k$ un corps. Soit $X$ un schéma propre sur $k$. Soit
$i : Z \to X$ un sous-schéma fermé de dimension $d$. Soient
$\mathcal{L}_1, \ldots, \mathcal{L}_d$ des
$\mathcal{O}_X$-modules inversibles. Nous définissons le {\it
nombre d'intersection } $(\mathcal{L}_1 \cdots \mathcal{L}_d \cdot
Z)$ comme le coefficient de $n_1 \ldots n_d$ dans le polynôme numérique
$$
\chi(X, i_*\mathcal{O}_Z \otimes \mathcal{L}_1^{\otimes n_1} \otimes
\ldots \otimes \mathcal{L}_d^{\otimes n_d}) =
\chi(Z, \mathcal{L}_1^{\otimes n_1} \otimes
\ldots \otimes \mathcal{L}_d^{\otimes n_d}|_Z)
$$
Dans le cas spécial de $\mathcal{L}_1 =
\ldots = \mathcal{L}_d = \mathcal{L}$,
nous écrivons $(\mathcal{L}^d \cdot Z)$.
\end{definition}

\noindent
L'égalité affichée dans la définition découle de
la formule de projection (Cohomologie, Section
\ref{cohomology-section-projection-formula}) et
de la cohomologie des schémas, Lemme
\ref{coherent-lemma-relative-affine-cohomology}. Nous
prouvons quelques lemmes pour ces nombres d'intersection.

\begin{lemma}
\label{lemma-intersection-number-integer}
Dans la situation de la Définition
\ref{definition-intersection-number}, le
nombre d'intersection $(\mathcal{L}_1
\cdots \mathcal{L}_d \cdot Z)$ est un entier.
\end{lemma}

\begin{proof}
Tout polynôme numérique de degré $e$ dans $n_1, \ldots, n_d$
peut être écrit de manière unique comme une combinaison
linéaire $\mathbf{Z}$ des fonctions ${n_1 \choose k_1}{n_2
\choose k_2} \ldots {n_d \choose k_d}$ avec $k_1 + \ldots +
k_d \leq e$. Appliquez ceci avec $e = d$. Laissé comme exercice.
\end{proof}

\begin{lemma}
\label{lemma-intersection-number-additive}
Dans la situation de la Définition
\ref{definition-intersection-number}, le nombre
d'intersection $(\mathcal{L}_1 \cdots \mathcal{L}_d
\cdot Z)$ est additif : si $\mathcal{L}_i =
\mathcal{L}_i' \otimes \mathcal{L}_i''$, alors nous avons
$$
(\mathcal{L}_1 \cdots \mathcal{L}_i \cdots \mathcal{L}_d \cdot Z) =
(\mathcal{L}_1 \cdots \mathcal{L}_i' \cdots \mathcal{L}_d \cdot Z) +
(\mathcal{L}_1 \cdots \mathcal{L}_i'' \cdots \mathcal{L}_d \cdot Z)
$$
\end{lemma}

\begin{proof}
Ceci est vrai parce que selon le Lemme
\ref{lemma-numerical-polynomial-from-euler} la fonction
$$
(n_1, \ldots, n_{i - 1}, n_i', n_i'', n_{i + 1}, \ldots, n_d)
\mapsto
\chi(Z, \mathcal{L}_1^{\otimes n_1} \otimes
\ldots \otimes (\mathcal{L}_i')^{\otimes n_i'} \otimes
(\mathcal{L}_i'')^{\otimes n_i''} \otimes \ldots \otimes
\mathcal{L}_d^{\otimes n_d}|_Z)
$$
est un polynôme numérique de degré total au plus $d$ en $d + 1$ variables.
\end{proof}

\begin{lemma}
\label{lemma-intersection-number-in-terms-of-components}
Dans la situation de la Définition
\ref{definition-intersection-number}, soit $Z_i \subset Z$ les composantes
irréductibles de dimension $d$. Soit $m_i = \text{length}_{\mathcal{O}_{X,
\xi_i}}(\mathcal{O}_{Z, \xi_i})$ où $\xi_i \in Z_i$ est le point générique. Alors
$$
(\mathcal{L}_1 \cdots \mathcal{L}_d \cdot Z) =
\sum m_i(\mathcal{L}_1 \cdots \mathcal{L}_d \cdot Z_i)
$$
\end{lemma}

\begin{proof}
Immédiat d'après le Lemme
\ref{lemma-numerical-polynomial-leading-term} et les définitions.
\end{proof}

\begin{lemma}
\label{lemma-intersection-number-and-pullback}
Supposons que $k$ soit un corps. Supposons que $f : Y \to X$ soit un
morphisme de schémas propres sur $k$. Supposons que $Z \subset Y$ soit
un sous-schéma fermé intègre de dimension $d$ et soit $\mathcal{L}_1,
\ldots, \mathcal{L}_d$ des $\mathcal{O}_X$-modules inversibles. Alors
$$
(f^*\mathcal{L}_1 \cdots f^*\mathcal{L}_d \cdot Z) =
\deg(f|_Z : Z \to f(Z)) (\mathcal{L}_1 \cdots \mathcal{L}_d \cdot f(Z))
$$
où $\deg(Z \to f(Z))$ est comme dans
morphismes, Définition
\ref{morphisms-definition-degree} ou $0$ si $\dim(f(Z)) < d$.
\end{lemma}

\begin{proof}
Le côté gauche est calculé en utilisant le
coefficient de $n_1 \ldots n_d$ dans la fonction
$$
\chi(Y, \mathcal{O}_Z \otimes f^*\mathcal{L}_1^{\otimes n_1} \otimes
\ldots \otimes f^*\mathcal{L}_d^{\otimes n_d}) =
\sum (-1)^i
\chi(X, R^if_*\mathcal{O}_Z \otimes
\mathcal{L}_1^{\otimes n_1} \otimes \ldots \otimes
\mathcal{L}_d^{\otimes n_d})
$$
L'égalité découle du Lemme
\ref{lemma-euler-characteristic-morphism} et de la formule
de projection (Cohomologie, Lemme
\ref{cohomology-lemma-projection-formula}). Si $f(Z)$ a
dimension $< d$, alors le côté droit est un polynôme de
degré total $< d$ selon le Lemme
\ref{lemma-numerical-polynomial-from-euler} et le résultat
est vrai. Supposons $\dim(f(Z)) = d$. Supposons que $\xi \in
f(Z)$ soit le point générique. Par la théorie de la dimension
(voir les Lemmata \ref{lemma-dimension-locally-algebraic} et
\ref{lemma-dimension-fibres-locally-algebraic}), le point
générique de $Z$ est l'unique point de $Z$ se projetant sur
$\xi$. Alors $f : Z \to f(Z)$ est fini au-dessus d'un ouvert
non vide de $f(Z)$, voir morphismes, Lemme
\ref{morphisms-lemma-generically-finite}. Ainsi $\deg(f : Z \to
f(Z))$ est défini et en fait il est égal à la longueur du germe
de $f_*\mathcal{O}_Z$ en $\xi$ au-dessus de $\mathcal{O}_{X,
\xi}$. De plus, le germe de $R^if_*\mathcal{O}_X$ en $\xi$ est
nul pour $i > 0$ car nous venons de voir que $f|_Z$ est fini
dans un voisinage de $\xi$ (de sorte que Cohomologie des
schémas, Lemme \ref{coherent-lemma-finite-pushforward-coherent}
donne l'annulation). Ainsi les termes $\chi(X, R^if_*\mathcal{O}_Z
\otimes \mathcal{L}_1^{\otimes n_1} \otimes \ldots \otimes
\mathcal{L}_d^{\otimes n_d})$ avec $i > 0$ ont un degré total de $< d$ et
$$
\chi(X, f_*\mathcal{O}_Z \otimes
\mathcal{L}_1^{\otimes n_1} \otimes \ldots \otimes
\mathcal{L}_d^{\otimes n_d})
=
\deg(f : Z \to f(Z)) \chi(f(Z),
\mathcal{L}_1^{\otimes n_1} \otimes \ldots \otimes
\mathcal{L}_d^{\otimes n_d}|_{f(Z)})
$$
modulo un polynôme de degré total $< d$
selon le Lemme
\ref{lemma-numerical-polynomial-leading-term}. Le résultat souhaité s'ensuit.
\end{proof}

\begin{lemma}
\label{lemma-numerical-intersection-effective-Cartier-divisor}
Supposons que $k$ soit un corps. Supposons que $X$ soit propre sur $k$.
Supposons que $Z \subset X$ soit un sous-schéma fermé de dimension $d$.
Supposons que $\mathcal{L}_1, \ldots, \mathcal{L}_d$ soit des
$\mathcal{O}_X$-modules inversibles. Supposons qu'il existe un diviseur de
Cartier effectif $D \subset Z$ tel que $\mathcal{L}_1|_Z \cong \mathcal{O}_Z(D)$. Alors
$$
(\mathcal{L}_1 \cdots \mathcal{L}_d \cdot Z) =
(\mathcal{L}_2 \cdots \mathcal{L}_d \cdot D)
$$
\end{lemma}

\begin{proof}
Nous pouvons remplacer $X$ par $Z$ et $\mathcal{L}_i$ par
$\mathcal{L}_i|_Z$. Ainsi, nous pouvons supposer $X = Z$ et
$\mathcal{L}_1 = \mathcal{O}_X(D)$. Ensuite, $\mathcal{L}_1^{-1}$ est le
faisceau idéal de $D$ et nous pouvons considérer la suite exacte courte
$$
0 \to \mathcal{L}_1^{\otimes -1} \to \mathcal{O}_X \to \mathcal{O}_D \to 0
$$
Définir $P(n_1, \ldots, n_d) =
\chi(X, \mathcal{L}_1^{\otimes
n_1} \otimes \ldots \otimes
\mathcal{L}_d^{\otimes n_d})$ et
$Q(n_1, \ldots, n_d) = \chi(D,
\mathcal{L}_1^{\otimes n_1}
\otimes \ldots \otimes
\mathcal{L}_d^{\otimes n_d}|_D)$. Nous
concluons à partir de l'additivité que
$$
P(n_1, \ldots, n_d) - P(n_1 - 1, n_2, \ldots, n_d) =
Q(n_1, \ldots, n_d)
$$
Parce que le degré total de $P$ est au plus $d$, nous
voyons que le coefficient de $n_1 \ldots n_d$ dans $P$
est égal au coefficient de $n_2 \ldots n_d$ dans $Q$.
\end{proof}

\begin{lemma}
\label{lemma-ample-positive}
Supposons que $k$ soit un corps. Supposons que $X$ soit propre sur $k$. Supposons que
$Z \subset X$ soit un sous-schéma fermé de dimension $d$. Si $\mathcal{L}_1, \ldots,
\mathcal{L}_d$ est ample, alors $(\mathcal{L}_1 \cdots \mathcal{L}_d \cdot Z)$ est positif.
\end{lemma}

\begin{proof}
Nous allons le prouver par induction sur $d$. Le cas $d = 0$
découle du Lemme \ref{lemma-chi-tensor-finite}. Supposons $d
> 0$. Par le Lemme
\ref{lemma-intersection-number-in-terms-of-components}, nous
pouvons supposer que $Z$ est un sous-schéma fermé intègre. En
fait, nous pouvons remplacer $X$ par $Z$ et $\mathcal{L}_i$
par $\mathcal{L}_i|_Z$ pour réduire au cas où $Z = X$ est une
variété propre de dimension $d$. D'après le Lemme
\ref{lemma-intersection-number-additive}, nous pouvons
remplacer $\mathcal{L}_1$ par une puissance tensorielle
positive. Ainsi, nous pouvons supposer qu'il existe une section
non nulle $s \in \Gamma(X, \mathcal{L}_1)$ telle que $X_s$ soit
affine (ici nous utilisons la définition de faisceau inversible
ample, voir Propriétés, Définition
\ref{properties-definition-ample}). Observez que $X$ n'est pas
affine car être propre et affine implique fini (morphismes, Lemme
\ref{morphisms-lemma-finite-proper}), ce qui contredit $d > 0$. Il
s'ensuit que $s$ a un schéma des zéros $Z(s) \subset X$ non vide.
Puisque $X$ est une variété, $s$ est une section régulière de
$\mathcal{L}_1$, donc $Z(s)$ est un diviseur de Cartier effectif,
ainsi $Z(s)$ a une codimension $1$ dans $X$, et donc $Z(s)$ a une
dimension $d - 1$ (ici nous utilisons le matériel des Diviseurs,
Sections \ref{divisors-section-effective-Cartier-divisors},
\ref{divisors-section-effective-Cartier-invertible} et
\ref{divisors-section-Noetherian-effective-Cartier} et de la théorie de la
dimension comme dans le Lemme \ref{lemma-dimension-locally-algebraic}). D'après
le Lemme \ref{lemma-numerical-intersection-effective-Cartier-divisor}, nous avons
$$
(\mathcal{L}_1 \cdots \mathcal{L}_d \cdot X) =
(\mathcal{L}_2 \cdots \mathcal{L}_d \cdot Z(s))
$$
Par induction, le côté droit est positif et la preuve est complète.
\end{proof}

\begin{definition}
\label{definition-degree}
Soit $k$ un corps. Soit $X$ un schéma propre sur
$k$. Soit $\mathcal{L}$ un $\mathcal{O}_X$-module
inversible ample. Pour tout sous-schéma fermé, le
{\it degré de $Z$ par rapport à $\mathcal{L}$ },
noté $\deg_\mathcal{L}(Z)$, est le nombre
d'intersection $(\mathcal{L}^d \cdot Z)$, où $d = \dim(Z)$.
\end{definition}

\noindent
D'après le Lemme \ref{lemma-ample-positive}, le degré d'un sous-schéma est
toujours un entier positif. Nous notons que $\deg_\mathcal{L}(Z) = d$ si et seulement si
$$
\chi(Z, \mathcal{L}^{\otimes n}|_Z) = \frac{d}{\dim(Z)!} n^{\dim(Z)} + l.o.t
$$
comme on peut le voir en utilisant cela
$$
(n_1 + \ldots + n_{\dim(Z)})^{\dim(Z)} =
\dim(Z)!\ n_1 \ldots n_{\dim(Z)} + \text{other terms}
$$

\begin{lemma}
\label{lemma-degree-finite-morphism-in-terms-degrees}
Supposons que $k$ soit un corps. Supposons que $f :
Y \to X$ soit un morphisme dominant fini de
variétés propres sur $k$. Supposons que $\mathcal{L}$
soit un $\mathcal{O}_X$-module inversible ample. Alors
$$
\deg_{f^*\mathcal{L}}(Y) = \deg(f) \deg_\mathcal{L}(X)
$$
où $\deg(f)$ est comme dans morphismes,
Définition \ref{morphisms-definition-degree}.
\end{lemma}

\begin{proof}
L'énoncé a du sens parce que $f^*\mathcal{L}$ est ample par
morphismes, Lemme
\ref{morphisms-lemma-pullback-ample-tensor-relatively-ample}. Cela dit, le
résultat est un cas particulier du Lemme \ref{lemma-intersection-number-and-pullback}.
\end{proof}

\noindent
Enfin, nous relions le nombre d'intersections avec une
courbe à la notion de degrés des modules inversibles sur les
courbes introduite à la Section \ref{section-divisors-curves}.

\begin{lemma}
\label{lemma-intersection-numbers-and-degrees-on-curves}
Supposons que $k$ soit un corps. Supposons que $X$ soit un
schéma propre sur $k$. Supposons que $Z \subset X$ soit un
sous-schéma fermé de dimension $\leq 1$. Supposons que
$\mathcal{L}$ soit un $\mathcal{O}_X$-module inversible. Alors
$$
(\mathcal{L} \cdot Z) = \deg(\mathcal{L}|_Z)
$$
où $\deg(\mathcal{L}|_Z)$ est comme dans la
Définition
\ref{definition-degree-invertible-sheaf}. Si $\mathcal{L}$
est ample, alors $\deg_\mathcal{L}(Z) = \deg(\mathcal{L}|_Z)$.
\end{lemma}

\begin{proof}
Cela découle du fait que la fonction $n \mapsto \chi(Z,
\mathcal{L}|_Z^{\otimes n})$ a un degré $1$ et que, par conséquent,
le coefficient directeur est la différence des valeurs consécutives.
\end{proof}

\begin{proposition}[Riemann-Roch asymptotique]
\label{proposition-asymptotic-riemann-roch}
Supposons que $k$ soit un corps. Supposons que $X$ soit un
schéma propre sur $k$ de dimension $d$. Supposons que
$\mathcal{L}$ soit un $\mathcal{O}_X$-module inversible ample. Alors
$$
\dim_k \Gamma(X, \mathcal{L}^{\otimes n}) \sim c n^d + l.o.t.
$$
où $c = \deg_\mathcal{L}(X)/d!$ est une constante positive.
\end{proposition}

\begin{proof}
Cela découle des définitions, du Lemme
\ref{lemma-ample-positive}, et de l'annulation de
la cohomologie supérieure dans Cohomologie des
schémas, Lemme \ref{coherent-lemma-vanshing-gives-ample}.
\end{proof}






\section{Dimension de plongement}
\label{section-embedding-dimension}

\noindent
Il existe plusieurs manières de définir la dimension d'intrusion, mais
pour des points fermés sur des schémas algébriques sur des corps
algébriquement clos, toutes les définitions sont équivalentes à ce qui suit.

\begin{definition}
\label{definition-embed-dim}
Soit $k$ un corps algébriquement clos. Soit $X$ un
$k$-schéma localement algébrique et soit $x \in X$ un
point fermé. La {\it dimension de plongement de $X$
en $x$ } est $\dim_k \mathfrak m_x/\mathfrak m_x^2$.
\end{definition}

\noindent
Faits sur la dimension d'encastrement. Supposons que $k, X,
x$ soit comme dans la Définition \ref{definition-embed-dim}.
\begin{enumerate}
\item La dimension d'immersion de $X$ en $x$
est la dimension de l'espace tangent $T_{X/k,
x}$ (Définition \ref{definition-tangent-space})
en tant qu'espace vectoriel $k$. \item La
dimension d'immersion de $X$ en $x$ est le plus
petit entier $d \geq 0$ tel qu'il existe une surjection
$$
k[[x_1, \ldots, x_d]] \longrightarrow \mathcal{O}_{X, x}^\wedge
$$
des algèbres $k$. \item La dimension d'immersion de $X$
en $x$ est le plus petit entier $d \geq 0$ tel qu'il
existe un voisinage ouvert $U \subset X$ de $x$ et une
immersion fermée $U \to Y$ où $Y$ est une variété lisse
de dimension $d$ sur $k$. \item La dimension d'immersion
de $X$ en $x$ est le plus petit entier $d \geq 0$ tel
qu'il existe un voisinage ouvert $U \subset X$ de $x$ et
un morphisme non ramifié $U \to \mathbf{A}^d_k$. \item Si
l'on nous donne une immersion fermée $X \to Y$ avec $Y$
lisse sur $k$, alors la dimension d'immersion de $X$ en
$x$ est le plus petit entier $d \geq 0$ tel qu'il existe
un sous-schéma fermé $Z \subset Y$ avec $X \subset Z$,
avec $Z \to \Spec(k)$ lisse en $x$, et avec $\dim_x(Z) = d$.
\end{enumerate}
Si nous avons jamais besoin de ceux-ci, nous
formulerons un résultat précis et fournirons une preuve.

\medskip\noindent
Corps de base non algébriquement clos ou points non fermés. Supposons
que $k$ soit un corps et que $X$ soit un $k$-schéma localement
algébrique. Si $x \in X$ est un point, alors nous avons plusieurs options
pour la dimension d’immersion de $X$ en $x$. À savoir, nous pourrions utiliser
\begin{enumerate}
\item $\dim_{\kappa(x)}(\mathfrak m_x/\mathfrak m_x^2)$,
\item $\dim_{\kappa(x)}(T_{X/k, x}) =
\dim_{\kappa(x)}(\Omega_{X/k, x} \otimes_{\mathcal{O}_{X, x}}
\kappa(x))$ (Lemme \ref{lemma-tangent-space-cotangent-space}),
\item le plus petit entier $d \geq 0$ tel qu'il existe un
voisinage ouvert $U \subset X$ de $x$ et une immersion fermée
$U \to Y$ où $Y$ est une variété lisse de dimension $d$ sur $k$.
\end{enumerate}
En caractéristique zéro (1) $=$ (2) si $x$ est un
point fermé ; de manière plus générale, cela vaut si
$\kappa(x)$ est algébrique séparable sur $k$, voir le
Lemme \ref{lemma-tangent-space-rational-point}. Il
semble que la définition géométrique (3) corresponde
le plus étroitement à l'intuition géométrique que
l'expression `` dimension d'immersion '' évoque.
Puisque l'on peut montrer que (3) et (2) définissent
le même nombre (cela découle du Lemme
\ref{lemma-immersion-into-affine}), c'est ce que nous
utiliserons. Dans notre terminologie, nous préciserons que nous
prenons la dimension d'immersion relativement au corps de base.

\begin{definition}
\label{definition-embedding-dimension}
Supposons que $k$ soit un corps. Supposons que $X$ soit un $k$ -schéma
localement algébrique. Supposons que $x \in X$ soit un point. La {\it
dimension de plongement de $X/k$ en $x$ } est $\dim_{\kappa(x)}(T_{X/k, x})$.
\end{definition}

\noindent
Si $(A, \mathfrak m, \kappa)$ est un anneau local noethérien, la
{\it dimension de plongement de $A$ } est parfois définie comme la
dimension de $\mathfrak m/\mathfrak m^2$ sur $\kappa$. Ci-dessus,
nous avons vu que si $A$ est donné comme une algèbre sur un corps
$k$, il peut être préférable d'utiliser $\dim_\kappa(\Omega_{A/k}
\otimes_A \kappa)$. Appelons cette quantité la {\it dimension de
plongement de $A/k$ }. Avec cette terminologie en place, nous avons
$$
\text{embed dim of }X/k\text{ at }x =
\text{embed dim of }\mathcal{O}_{X, x}/k =
\text{embed dim of }\mathcal{O}_{X, x}^\wedge/k
$$
si $k, X, x$ sont comme dans la Définition \ref{definition-embedding-dimension}.






\section{Théorèmes de Bertini}
\label{section-bertini}

\noindent
Dans cette section, nous prouvons des résultats du type :
étant donnée une variété projective lisse $X$ sur un corps
$k$, il existe un diviseur ample $H \subset X$ qui est lisse.

\begin{lemma}
\label{lemma-generate-over-complement}
Supposons que $k$ soit un corps. Supposons que $X$ soit un
schéma propre sur $k$. Supposons que $\mathcal{L}$ soit un
$\mathcal{O}_X$-module inversible ample. Supposons que $Z \subset
X$ soit un sous-schéma fermé. Alors il existe un entier $n_0$ tel
que pour tout $n \geq n_0$ le noyau $V_n$ de $\Gamma(X,
\mathcal{L}^{\otimes n}) \to \Gamma(Z, \mathcal{L}^{\otimes n}|_Z)$
engendre $\mathcal{L}^{\otimes n}|_{X \setminus Z}$ et le morphisme canonique
$$
X \setminus Z \longrightarrow \mathbf{P}(V_n)
$$
est une immersion de schémas sur $k$.
\end{lemma}

\begin{proof}
Supposons que $\mathcal{I} \subset \mathcal{O}_X$ soit le
faisceau d'idéaux quasi-cohérent de $Z$. Remarquez que via
l'inclusion $\mathcal{I} \otimes_{\mathcal{O}_X}
\mathcal{L}^{\otimes n} \subset \mathcal{L}^{\otimes n}$ nous
avons $V_n = \Gamma(X, \mathcal{I} \otimes_{\mathcal{O}_X}
\mathcal{L}^{\otimes n})$. Choisissez $n_1$ de sorte que pour $n
\geq n_1$ le faisceau $\mathcal{I} \otimes \mathcal{L}^{\otimes
n}$ soit globalement engendré, voir Propriétés, Proposition
\ref{properties-proposition-characterize-ample}. Il s'ensuit que $V_n$
engendre $\mathcal{L}^{\otimes n}|_{X \setminus Z}$ pour $n \geq n_1$.

\medskip\noindent
Pour $n \geq n_1$, notons $\psi_n : V_n \to
\Gamma(X \setminus Z, \mathcal{L}^{\otimes
n}|_{X \setminus Z})$ l'application de
restriction. Nous obtenons un morphisme canonique
$$
\varphi = \varphi_{\mathcal{L}^{\otimes n}|_{X \setminus Z}, \psi_n} :
X \setminus Z
\longrightarrow
\mathbf{P}(V_n)
$$
par Constructions, Exemple
\ref{constructions-example-projective-space}. Choisissez $n_2$ de sorte
que pour tout $n \geq n_2$ le faisceau inversible $\mathcal{L}^{\otimes n}$
soit très ample sur $X$. Nous affirmons que $n_0 = n_1 + n_2$ fonctionne.

\medskip\noindent
Preuve de l'affirmation. Disons $n \geq n_0$ et écrivons $n =
n_1 + n'$. Pour $x \in X \setminus Z$, nous pouvons choisir $s_1
\in V_1$ ne s'annulant pas en $x$. Posons $V' = \Gamma(X,
\mathcal{L}^{\otimes n'})$. Par notre choix de $n$ et $n'$, nous
voyons que le morphisme correspondant $\varphi' : X \to
\mathbf{P}(V')$ est une immersion fermée. Ainsi, si nous
choisissons $s' \in \Gamma(X, \mathcal{L}^{\otimes n'})$ ne
s'annulant pas en $x$, alors $X_{s'} = (\varphi')^{-1}(D_+(s'))$
(voir Constructions, Lemme
\ref{constructions-lemma-invertible-map-into-proj}) est affine et $X_{s'}
\to D_+(s')$ est une immersion fermée. Ensuite, $s = s_1 \otimes s' \in
V_n$ ne s'annule pas en $x$. Si $D_+(s) \subset \mathbf{P}(V_n)$ désigne
l'espace affine ouvert correspondant de notre espace projectif, alors
$\varphi^{-1}(D_+(s)) = X_s \subset X \setminus Z$ (voir la référence
ci-dessus). L'ouvert $X_s = X_{s'} \cap X_{s_1}$ est affine, voir Propriétés,
Lemme \ref{properties-lemma-affine-cap-s-open}. Considérons le homomorphisme d'anneaux
$$
\text{Sym}(V)_{(s)} \longrightarrow \mathcal{O}_X(X_s)
$$
définissant le morphisme $X_s \to D_+(s)$. Parce que $X_{s'}
\to D_+(s')$ est une immersion fermée, les images des éléments
$$
\frac{s_1 \otimes t'}{s_1 \otimes s'}
$$
où $t' \in V'$ génère l'image de
$\mathcal{O}_X(X_{s'}) \to \mathcal{O}_X(X_s)$. Puisque
$X_s \to X_{s'}$ est une immersion ouverte, cela implique
que $X_s \to D_+(s)$ est une immersion de schémas affines
(voir ci-dessous). Ainsi $\varphi_n$ est une immersion par
morphismes, Lemme \ref{morphisms-lemma-check-immersion}.

\medskip\noindent
Supposons que $a : A' \to A$ et $c : B \to A$ soient des
homomorphismes d'anneaux tels que $\Spec(a)$ soit une
immersion et $\Im(a) \subset \Im(c)$. Posons $B' = A' \times_A
B$ avec les projections $b : B' \to B$ et $c' : B' \to A'$. Par
hypothèse, $c'$ est surjectif et donc $\Spec(c')$ est une
immersion fermée. D'où $\Spec(c') \circ \Spec(a)$ est une
immersion (schémas, Lemme
\ref{schemes-lemma-composition-immersion}). Alors $\Spec(c)$ doit être une
immersion car il factorise l'immersion $\Spec(c') \circ \Spec(a) = \Spec(b)
\circ \Spec(c)$, voir morphismes, Lemme \ref{morphisms-lemma-immersion-permanence}.
\end{proof}

\begin{situation}
\label{situation-family-divisors}
Supposons que $k$ soit un corps, que $X$ soit un
schéma sur $k$, que $\mathcal{L}$ soit un
$\mathcal{O}_X$-module inversible, que $V$ soit un espace
vectoriel $k$ de dimension finie, et que $\psi : V \to
\Gamma(X, \mathcal{L})$ soit une application linéaire $k$.
Disons $\dim(V) = r$ et nous avons une base $v_1, \ldots,
v_r$ de $V$. Alors nous obtenons un diviseur universel `` ''
$$
H_{univ} = Z(s_{univ}) \subset  \mathbf{A}^r \times_k X
$$
comme le schéma zéro (Diviseurs, Définition
\ref{divisors-definition-zero-scheme-s}) de la section
$$
s_{univ} = \sum\nolimits_{i = 1, \ldots, r} x_i \psi(v_i) \in
\Gamma(\mathbf{A}^r \times_k X, \text{pr}_2^*\mathcal{L})
$$
Pour une extension de corps $k'/k$, les $k'$-points $v \in
\mathbf{A}^r_k(k')$ correspondent à des vecteurs $(a_1, \ldots,
a_r)$ d'éléments de $k'$. Ainsi, nous pouvons, d'une part,
considérer $v$ comme l'élément $v = \sum_{i = 1, \ldots, r} a_i
v_i \in V \otimes_k k'$ et, d'autre part, associer à $v$ la section
$$
\psi(v) = \sum\nolimits_{i = 1, \ldots, r} a_i \psi(v_i) \in
\Gamma(X_{k'}, \mathcal{L}|_{X_{k'}})
$$
Avec cette notation, il est clair que la
fibre de $H_{univ}$ sur $v \in V \otimes k'$
est le schéma nul de $\psi(v)$. En formule :
$$
H_v = H_{univ, v} = Z(\psi(v))
$$
Nous noterons cette valeur commune par $H_v$ comme indiqué. Enfin, dans
cette situation, laissons $P$ être une propriété des vecteurs $v \in V
\otimes_k k'$ pour $k'/k$ une extension de corps arbitraire \footnote {.
Par exemple, nous pourrions considérer la condition que $H_v$ est lisse
sur $k'$, ou géométriquement irréductible sur $k'$. }. Nous disons que
$P$ {\it est vérifiée en général } $v \in V \otimes_k k'$ s'il existe un
ouvert de Zariski non vide $U \subset \mathbf{A}^r_k$ tel que si $v$
correspond à un point $k'$ de $U$ pour tout $k'/k$ alors $P(v)$ est vérifié.
\end{situation}

\begin{lemma}
\label{lemma-bertini}
Dans la situation \ref{situation-family-divisors}, supposez
\begin{enumerate}
\item $X$ est lisse sur $k$, \item
l'image de $\psi : V \to \Gamma(X,
\mathcal{L})$ génère $\mathcal{L}$,
\item le morphisme correspondant
$\varphi_{\mathcal{L}, \psi} : X
\to \mathbf{P}(V)$ est une immersion.
\end{enumerate}
Alors pour le $v \in V \otimes_k k'$
général, le schéma $H_v$ est lisse sur $k'$.
\end{lemma}

\begin{proof}
(Nous observons que $X$ est séparé et de type fini en
tant que sous-schéma localement fermé d'un espace
projectif.) Utilisons la notation introduite au-dessus
de l'énoncé du lemme. Nous considérons les projections
$$
\xymatrix{
\mathbf{A}^r_k \times_k X \ar[d] &
H_{univ} \ar[l] \ar[ld]^p \ar[r] \ar[rd]_q &
\mathbf{A}^r_k \times_k X \ar[d] \\
X & &
\mathbf{A}^r_k
}
$$
Supposons que $\Sigma \subset H_{univ}$ soit le lieu singulier du
morphisme $q : H_{univ} \to \mathbf{A}^r_k$, c’est-à-dire
l’ensemble des points où $q$ n’est pas lisse. Alors $\Sigma$ est
fermé car le lieu lisse d’un morphisme est ouvert par définition.
Puisque la fibre d’un morphisme lisse est lisse, il suffit de
prouver $q(\Sigma)$ est contenue dans un sous-ensemble fermé
propre de $\mathbf{A}^r_k$. Puisque $\Sigma$ (avec structure de
schéma induite réduite) est un schéma de type fini sur $k$ il
suffit de prouver $\dim(\Sigma) < r$ Cela découle du Lemme
\ref{lemma-dimension-fibres-locally-algebraic}. Puisque les
dimensions ne changent pas en remplaçant $k$ par un corps plus grand
(morphismes, Lemme
\ref{morphisms-lemma-dimension-fibre-after-base-change} ), on peut
supposer et supposer que $k$ est algébriquement clos. Par la théorie de la
dimension (Lemme \ref{lemma-dimension-fibres-locally-algebraic} ), il
suffit de démontrer que pour $x \in X \setminus Z$ fermé, nous avons
$p^{-1}(\{x\}) \cap \Sigma$ a dimension $< r - \dim(X')$ où $X'$ est le
composante unique irréductible de $X$ contenant $x$. Comme $X$ est lisse sur
$k$ et $x$ est un point fermé, nous avons $\dim(X') = \dim \mathfrak
m_x/\mathfrak m_x^2$ (morphismes, Lemme
\ref{morphisms-lemma-smooth-omega-finite-locally-free} et Algèbre, Lemme \ref{algebra-lemma-rank-omega}). Ainsi, nous gagnons si
$$
\dim p^{-1}(x) \cap \Sigma < r - \dim \mathfrak m_x/\mathfrak m_x^2
$$
pour tous $x \in X$ fermés.

\medskip\noindent
Puisque $V$ génère globalement $\mathcal{L}$, pour chaque composante
irréductible $X'$ de $X$, il existe un ouvert de Zariski non vide
de $\mathbf{A}^r$ tel que les fibres de $q$ au-dessus de cet ouvert
ne contiennent pas $X'$. (Par exemple, si $x' \in X'$ est un point
fermé, alors nous pouvons prendre l'ouvert correspondant à ces
vecteurs $v \in V$ tels que $\psi(v)$ ne s'annule pas en $x'$. Cet
ouvert sera le complément d'un hyperplan dans $\mathbf{A}^r_k$.)
Supposons que $U \subset \mathbf{A}^r$ soit l'intersection (non vide)
de ces ouverts. Alors les fibres de $q^{-1}(U) \to U$ sont des
diviseurs de Cartier effectifs sur les fibres de $U \times_k X \to U$
(car une section non nulle d'un module inversible sur un schéma
intègre est une section régulière). Par conséquent, le morphisme
$q^{-1}(U) \to U$ est plat selon le Lemme sur les Diviseurs,
\ref{divisors-lemma-fibre-Cartier}. Ainsi pour $x \in X$ fermé et $v \in V
= \mathbf{A}^r_k(k)$, si $(x, v) \in H_{univ}$, c'est-à-dire, si $x \in
H_v$ alors $q$ est lisse en $(x, v)$ si et seulement si la fibre $H_v$ est
lisse en $x$, voir morphismes, Lemme \ref{morphisms-lemma-smooth-at-point}.

\medskip\noindent
Considérez l'image $\psi(v)_x$ dans la tige
$\mathcal{L}_x$ de la section correspondant à $v \in V$. Nous avons
$$
x \in H_v \Leftrightarrow \psi(v)_x \in \mathfrak m_x\mathcal{L}_x
$$
Si c'est vrai, alors nous avons
$$
H_v\text{ singular at }x \Leftrightarrow
\psi(v)_x \in \mathfrak m_x^2\mathcal{L}_x
$$
À savoir, $\psi(v)_x$ n'est pas contenu dans
$\mathfrak m_x^2\mathcal{L}_x$
$\Leftrightarrow$ l'équation locale pour $H_v
\subset X$ en $x$ n'est pas contenue dans
$\mathfrak m_x^2$ $\Leftrightarrow$
$\mathcal{O}_{H_v, x}$ est régulière (Algèbre,
Lemme \ref{algebra-lemma-regular-ring-CM})
$\Leftrightarrow$ $H_v$ est lisse en $x$ sur $k$
(Algèbre, Lemme
\ref{algebra-lemma-separable-smooth}). Nous concluons
que les points fermés de $p^{-1}(x) \cap \Sigma$
correspondent à ceux de $v \in V$ tels que $\psi(v)_x
\in \mathfrak m_x^2\mathcal{L}_x$. Cependant, comme
$\varphi_{\mathcal{L}, \psi}$ est une immersion, l'application
$$
V \longrightarrow \mathcal{L}_x/\mathfrak m_x^2\mathcal{L}_x
$$
est surjective (petit détail omis). D'après ce qui précède, les
points fermés du lieu $p^{-1}(x) \cap \Sigma$ vus comme
sous-espace de $V$ sont le noyau de cette application et ont donc la
dimension $r - \dim \mathfrak m_x/\mathfrak m_x^2 - 1$ comme souhaité.
\end{proof}






\section{Enriques--Severi--Zariski}
\label{section-vanishing-negative}

\noindent
Dans cette section, nous prouvons certains résultats
de la forme : tordre par un module inversible `` très
négatif '' annule la cohomologie de bas degré. Nous en
déduisons également la connexité d'une section
hyper-surface d'un schéma propre normal de dimension $\geq 2$.

\begin{lemma}
\label{lemma-vanishin-h0-negative}
Supposons que $k$ soit un corps. Supposons que $X$ soit un
schéma propre sur $k$. Supposons que $\mathcal{L}$ soit un
$\mathcal{O}_X$-inversible ample. Supposons que $\mathcal{F}$
soit un $\mathcal{O}_X$ cohérent. Si $\text{Ass}(\mathcal{F})$
ne contient aucun point fermé, alors $\Gamma(X, \mathcal{F}
\otimes_{\mathcal{O}_X} \mathcal{L}^{\otimes n}) = 0$ pour $n \ll 0$.
\end{lemma}

\begin{proof}
Pour un $\mathcal{O}_X$-module cohérent
$\mathcal{F}$, soit $\mathcal{P}(\mathcal{F})$ la
propriété : il existe un $n_0 \in \mathbf{Z}$ tel
que pour tout $n \leq n_0$, chaque section $s$ de
$\mathcal{F} \otimes_{\mathcal{O}_X}
\mathcal{L}^{\otimes n}$ a un support ne consistant
qu’en des points fermés. Depuis
$\text{Ass}(\mathcal{F}) = \text{Ass}(\mathcal{F}
\otimes_{\mathcal{O}_X} \mathcal{L}^{\otimes n})$, nous
voyons qu’il suffit de prouver que $\mathcal{P}$ vaut
pour tous les modules cohérents sur $X$. Pour ce faire,
nous allons prouver que les conditions (1), (2) et (3)
du Lemme
\ref{coherent-lemma-property-higher-rank-cohomological} de la cohomologie des schémas sont satisfaites.

\medskip\noindent
Pour voir la condition (1), supposons que
$$
0 \to \mathcal{F}_1 \to \mathcal{F} \to \mathcal{F}_2 \to 0
$$
est une suite exacte courte de $\mathcal{O}_X$-modules
cohérents telle que nous ayons $\mathcal{P}$ pour
$\mathcal{F}_i$, $i = 1, 2$. Supposons que $n_1, n_2$ soit les
coupures que nous trouvons. Supposons que $\mathcal{F}'_2
\subset \mathcal{F}_2$ soit le sous-module cohérent maximal
dont le support est un ensemble fini de points fermés.
Supposons que $\mathcal{I} \subset \mathcal{O}_X$ soit
l'annulateur de $\mathcal{F}'_2$. Puisque $\mathcal{L}$ est
ample, nous pouvons trouver un $e > 0$ tel que $\mathcal{I}
\otimes_{\mathcal{O}_X} \mathcal{L}^{\otimes e}$ soit globalement
engendré. Posons $n_0 = \min(n_2, n_1 - e)$. Supposons que $n
\leq n_0$ et que $t$ soit une section globale de $\mathcal{F}
\otimes \mathcal{L}^{\otimes n}$. L'image de $t$ dans
$\mathcal{F}_2 \otimes \mathcal{L}^{\otimes n}$ tombe dans
$\mathcal{F}'_2 \otimes \mathcal{L}^{\otimes n}$ parce que $n \leq
n_2$. Ainsi, pour tout $s \in \Gamma(X, \mathcal{I}
\otimes_{\mathcal{O}_X} \mathcal{L}^{\otimes e})$, le produit $t
\otimes s$ se trouve dans $\mathcal{F}_1 \otimes \mathcal{L}^{\otimes
n + e}$. Par conséquent, $t \otimes s$ a un support contenu dans
l'ensemble fini de points fermés dans $\text{Ass}(\mathcal{F}_1)$
parce que $n + e \leq n_1$. Puisque par notre choix de $e$ nous
pouvons choisir $s$ inversible en tout point ne se trouvant pas dans le
support de $\mathcal{F}'_2$, nous en concluons que le support de $t$
est contenu dans l'union de l'ensemble fini de points fermés dans
$\text{Ass}(\mathcal{F}_1)$ et de l'ensemble fini de points fermés dans
$\text{Ass}(\mathcal{F}_2)$. Cela termine la démonstration de la condition (1).

\medskip\noindent
La condition (2) est immédiate.

\medskip\noindent
Pour la condition (3), nous choisissons $\mathcal{G} =
\mathcal{O}_Z$. Dans ce cas, si $Z$ est un point fermé de
$X$, alors il n'y a rien à montrer. Si $\dim(Z) > 0$,
alors nous montrerons que $\Gamma(Z, \mathcal{L}^{\otimes
n}|_Z) = 0$ pour $n < 0$. C'est-à-dire, soit $s$ une
section non nulle d'une puissance négative de
$\mathcal{L}|_Z$. Choisissez une section non nulle $t$ d'une
puissance positive de $\mathcal{L}|_Z$ (cela est possible
car $\mathcal{L}$ est ample, voir Propriétés, Proposition
\ref{properties-proposition-characterize-ample}). Alors
$s^{\deg(t)} \otimes t^{\deg(s)}$ est une section globale non
nulle de $\mathcal{O}_Z$ (parce que $Z$ est intègre) et donc
une unité (Lemme
\ref{lemma-proper-geometrically-reduced-global-sections}). Cela
implique que $t$ est une section trivialise d'une puissance positive
de $\mathcal{L}$. Ainsi, la fonction $n \mapsto \dim_k \Gamma(X,
\mathcal{L}^{\otimes n})$ est bornée sur un ensemble infini d'entiers
positifs ce qui contredit le théorème de Riemann-Roch asymptotique
(Proposition \ref{proposition-asymptotic-riemann-roch}) puisque $\dim(Z) > 0$.
\end{proof}

\begin{lemma}[Enriques--Severi--Zariski]
\label{lemma-vanishin-h1-negative}
Supposons que $k$ soit un corps. Supposons que $X$ soit un
schéma propre sur $k$. Supposons que $\mathcal{L}$ soit un
$\mathcal{O}_X$-module inversible ample. Supposons que
$\mathcal{F}$ soit un $\mathcal{O}_X$-module cohérent.
Supposons que pour $x \in X$ fermé nous ayons
$\text{depth}(\mathcal{F}_x) \geq 2$. Alors $H^1(X, \mathcal{F}
\otimes_{\mathcal{O}_X} \mathcal{L}^{\otimes m}) = 0$ pour $m \ll 0$.
\end{lemma}

\begin{proof}
Choisissez une immersion fermée $i : X \to
\mathbf{P}^n_k$ telle que $i^*\mathcal{O}(1) \cong
\mathcal{L}^{\otimes e}$ pour un certain $e > 0$ (voir
morphismes, Lemme
\ref{morphisms-lemma-finite-type-over-affine-ample-very-ample}). Il suffit alors de prouver le lemme pour
$$
\mathcal{G} =
i_*(\mathcal{F} \oplus \mathcal{F} \otimes \mathcal{L} \oplus \ldots
\oplus \mathcal{F} \otimes \mathcal{L}^{\otimes e - 1})
\quad\text{and}\quad \mathcal{O}(1)
$$
sur $\mathbf{P}^n_k$. À savoir, nous avons
$$
H^1(\mathbf{P}^n_k, \mathcal{G}(m)) =
\bigoplus\nolimits_{j = 0, \ldots, e - 1}
H^1(X, \mathcal{F} \otimes \mathcal{L}^{\otimes j + me})
$$
par la cohomologie des schémas, Lemme
\ref{coherent-lemma-relative-affine-cohomology}. De
plus, si $y \in \mathbf{P}^n_k$ est un point fermé
alors $\text{depth}(\mathcal{G}_y) = \infty$ si $y
\not \in i(X)$ et $\text{depth}(\mathcal{G}_y) =
\text{depth}(\mathcal{F}_x)$ si $y = i(x)$ car dans ce
cas $\mathcal{G}_y \cong \mathcal{F}_x^{\oplus e}$
comme module sur $\mathcal{O}_{\mathbf{P}^n_k, x}$ et
nous pouvons utiliser par exemple Algèbre, Lemme
\ref{algebra-lemma-depth-goes-down-finite} pour obtenir l'égalité.

\medskip\noindent
Supposons que $X = \mathbf{P}^n_k$, que $\mathcal{L} =
\mathcal{O}(1)$ et que $k$ soit infini. Choisissons $s \in
H^0(\mathbf{P}^1_k, \mathcal{O}(1))$ qui détermine une suite exacte
$$
0 \to \mathcal{F}(-1) \xrightarrow{s} \mathcal{F} \to \mathcal{G} \to 0
$$
comme dans le Lemme \ref{lemma-exact-sequence-induction}.
Puisque l'application $\mathcal{F}(-1) \to \mathcal{F}$ est
affine localement donnée par la multiplication par un
non-diviseur de zéro sur $\mathcal{F}$, nous voyons que pour
$x \in \mathbf{P}^n_k$ fermé nous avons
$\text{depth}(\mathcal{G}_x) \geq 1$, voir l'Algèbre, Lemme
\ref{algebra-lemma-depth-drops-by-one}. Par conséquent, d'après le
Lemme \ref{lemma-vanishin-h0-negative}, nous avons
$H^0(\mathcal{G}(m)) = 0$ pour $m \ll 0$. En examinant la suite exacte
longue de la cohomologie après tordage (voir Remarque
\ref{remark-exact-sequence-induction-cohomology}), nous trouvons que la suite de nombres
$$
\dim H^1(\mathbf{P}^n_k, \mathcal{F}(m))
$$
se stabilise pour $m \leq m_0$ pour un certain entier
$m_0$. Supposons que $N$ soit la dimension commune de ces
espaces pour $m \leq m_0$. Nous devons montrer $N = 0$.

\medskip\noindent
Pour $d > 0$ et $m \leq m_0$, considérez l'application bilinéaire
$$
H^0(\mathbf{P}^n_k, \mathcal{O}(d)) \times
H^1(\mathbf{P}^n_k, \mathcal{F}(m - d))
\longrightarrow
H^1(\mathbf{P}^n_k, \mathcal{F}(m))
$$
Selon l'algèbre linéaire, il existe un sous-espace de
codimension $\leq N^2$ $V_m \subset
H^0(\mathbf{P}^n_k, \mathcal{O}(d))$ tel que la
multiplication par $s' \in V_m$ annule $H^1(\mathbf{P}^n_k,
\mathcal{F}(m - d))$. Observez que pour $m' < m \leq m_0$ le diagramme
$$
\xymatrix{
H^0(\mathbf{P}^n_k, \mathcal{O}(d)) \times
H^1(\mathbf{P}^n_k, \mathcal{F}(m' - d)) \ar[r] \ar[d]^{1 \times s^{m' - m}} &
H^1(\mathbf{P}^n_k, \mathcal{F}(m')) \ar[d]^{s^{m' - m}}\\
H^0(\mathbf{P}^n_k, \mathcal{O}(d)) \times
H^1(\mathbf{P}^n_k, \mathcal{F}(m - d)) \ar[r] &
H^1(\mathbf{P}^n_k, \mathcal{F}(m))
}
$$
commute avec les isomorphismes en allant verticalement. Ainsi $V_m = V$ est
indépendant de $m \leq m_0$. Pour $x \in \text{Ass}(\mathcal{F})$, on pose
$Z = \overline{\{x\}}$. Pour $d$ suffisamment grand, l'application linéaire
$$
H^0(\mathbf{P}^n_k, \mathcal{O}(d)) \to H^0(Z, \mathcal{O}(d)|_Z)
$$
a le rang $> N^2$ parce que $\dim(Z) \geq 1$ (par exemple, cela
découle du Riemann-Roch asymptotique et de l'ampleness $\mathcal{O}(1)$
; détails omis). Par conséquent, nous pouvons trouver $s' \in V$ de
sorte que $s'$ ne s'annule en aucun point associé de $\mathcal{F}$
(utiliser que l'ensemble des points associés est fini). Ensuite, nous obtenons
$$
0 \to \mathcal{F}(-d) \xrightarrow{s'} \mathcal{F} \to \mathcal{G}' \to 0
$$
et comme précédemment, nous concluons comme avant que
la multiplication par $s'$ sur $H^1(\mathbf{P}^n_k,
\mathcal{F}(m - d))$ est injective pour $m \ll 0$. Cela
contredit le choix de $s'$ sauf $N = 0$ comme souhaité.

\medskip\noindent
Nous devons encore traiter le cas où $k$ est fini.
Dans ce cas, soit $K/k$ toute extension algébrique
infinie de corps. Notons $\mathcal{F}_K$ et
$\mathcal{L}_K$ les rétractations de $\mathcal{F}$ et
$\mathcal{L}$ vers $X_K = \Spec(K) \times_{\Spec(k)} X$. Nous avons
$$
H^1(X_K, \mathcal{F}_K \otimes \mathcal{L}_K^{\otimes m}) =
H^1(X, \mathcal{F} \otimes \mathcal{L}^{\otimes m}) \otimes_k K
$$
par la cohomologie des schémas, Lemme
\ref{coherent-lemma-flat-base-change-cohomology}. D'autre part,
un point fermé $x_K$ de $X_K$ se mappe à un point fermé $x$ de
$X$ parce que $K/k$ est une extension algébrique.
L'homomorphisme d'anneaux $\mathcal{O}_{X, x} \to \mathcal{O}_{X_K,
x_K}$ est plat (Lemme \ref{lemma-change-fields-flat}). Ainsi nous avons
$$
\text{depth}(\mathcal{F}_{x_K}) =
\text{depth}(\mathcal{F}_x \otimes_{\mathcal{O}_{X, x}}
\mathcal{O}_{X_K, x_K}) \geq
\text{depth}(\mathcal{F}_x)
$$
par l'Algèbre, le Lemme
\ref{algebra-lemma-apply-grothendieck-module} (en fait l'égalité tient
ici mais nous n'en avons pas besoin). Par conséquent, le résultat sur
$k$ découle du résultat sur le corps infini $K$ et la preuve est complète.
\end{proof}

\begin{lemma}
\label{lemma-connectedness-ample-divisor}
Supposons que $k$ soit un corps. Supposons que $X$ soit un schéma
propre sur $k$. Supposons que $\mathcal{L}$ soit un $\mathcal{O}_X$
-module inversible ample. Soit $s \in \Gamma(X, \mathcal{L})$. Supposons
\begin{enumerate}
\item $s$ est une section régulière
(Diviseurs, Définition
\ref{divisors-definition-regular-section}), \item pour
chaque point fermé $x \in X$ nous avons
$\text{depth}(\mathcal{O}_{X, x}) \geq 2$, et \item $X$ est connexe.
\end{enumerate}
Ensuite, le schéma zéro $Z(s)$ de $s$ est connecté.
\end{lemma}

\begin{proof}
Puisque $s$ est une section régulière, il en est de même de $s^n \in
\Gamma(X, \mathcal{L}^{\otimes n})$ pour tout $n > 1$. De plus, le
morphisme d'inclusion $Z(s) \to Z(s^n)$ est une bijection sur les espaces
topologiques sous-jacents. Ainsi, si $Z(s)$ est déconnecté, il en est de
même de $Z(s^n)$. Considérons maintenant la suite exacte courte canonique
$$
0 \to \mathcal{L}^{\otimes -n} \xrightarrow{s^n}
\mathcal{O}_X \to \mathcal{O}_{Z(s^n)} \to 0
$$
Considérons la $k$-algèbre $R_n = \Gamma(X,
\mathcal{O}_{Z(s^n)})$. Si $Z(s)$ est déconnecté, c'est-à-dire si
$Z(s^n)$ est déconnecté, alors soit $R_n$ est nul dans le cas
$Z(s^n) = \emptyset$, soit $R_n$ contient un idempotent non trivial
dans le cas $Z(s^n) = U \amalg V$, avec $U, V \subset Z(s^n)$
ouverts et non vides (le lecteur peut consulter le lemme
\ref{lemma-proper-geometrically-reduced-global-sections}). Ainsi,
l'application $\Gamma(X, \mathcal{O}_X) \to R_n$ ne peut pas être un
isomorphisme. Il s'ensuit que $H^0(X, \mathcal{L}^{\otimes -n})$ ou
$H^1(X, \mathcal{L}^{\otimes -n})$ est non nul pour une infinité
d'entiers positifs $n$. Cela contredit le lemme
\ref{lemma-vanishin-h0-negative} ou le lemme \ref{lemma-vanishin-h1-negative}, et la preuve est complète.
\end{proof}








\begin{multicols}{2}[\section{Autres chapitres}]
\noindent
Préliminaires
\begin{enumerate}
\item \hyperref[introduction-section-phantom]{Introduction}
\item \hyperref[conventions-section-phantom]{Conventions}
\item \hyperref[sets-section-phantom]{Théorie des ensembles}
\item \hyperref[categories-section-phantom]{Catégories}
\item \hyperref[topology-section-phantom]{Topologie}
\item \hyperref[sheaves-section-phantom]{Faisceaux sur les espaces}
\item \hyperref[sites-section-phantom]{Sites et faisceaux}
\item \hyperref[stacks-section-phantom]{Champs}
\item \hyperref[fields-section-phantom]{Corps}
\item \hyperref[algebra-section-phantom]{Algèbre commutative}
\item \hyperref[brauer-section-phantom]{Groupes de Brauer}
\item \hyperref[homology-section-phantom]{Algèbre homologique}
\item \hyperref[derived-section-phantom]{Catégories dérivées}
\item \hyperref[simplicial-section-phantom]{Méthodes simpliciales}
\item \hyperref[more-algebra-section-phantom]{Compléments d'algèbre}
\item \hyperref[smoothing-section-phantom]{Lissage des morphismes d'anneaux}
\item \hyperref[modules-section-phantom]{Faisceaux de modules}
\item \hyperref[sites-modules-section-phantom]{Modules sur les sites}
\item \hyperref[injectives-section-phantom]{Objets injectifs}
\item \hyperref[cohomology-section-phantom]{Cohomologie des faisceaux}
\item \hyperref[sites-cohomology-section-phantom]{Cohomologie sur les sites}
\item \hyperref[dga-section-phantom]{Algèbre différentielle graduée}
\item \hyperref[dpa-section-phantom]{Algèbre à puissances divisées}
\item \hyperref[sdga-section-phantom]{Faisceaux différentiels gradués}
\item \hyperref[hypercovering-section-phantom]{Hyperrecouvrements}
\end{enumerate}
Schémas
\begin{enumerate}
\setcounter{enumi}{25}
\item \hyperref[schemes-section-phantom]{Schémas}
\item \hyperref[constructions-section-phantom]{Constructions de schémas}
\item \hyperref[properties-section-phantom]{Propriétés des schémas}
\item \hyperref[morphisms-section-phantom]{Morphismes de schémas}
\item \hyperref[coherent-section-phantom]{Cohomologie des schémas}
\item \hyperref[divisors-section-phantom]{Diviseurs}
\item \hyperref[limits-section-phantom]{Limites de schémas}
\item \hyperref[varieties-section-phantom]{Variétés}
\item \hyperref[topologies-section-phantom]{Topologies sur les schémas}
\item \hyperref[descent-section-phantom]{Descente}
\item \hyperref[perfect-section-phantom]{Catégories dérivées de schémas}
\item \hyperref[more-morphisms-section-phantom]{Compléments sur les morphismes}
\item \hyperref[flat-section-phantom]{Compléments sur la platitude}
\item \hyperref[groupoids-section-phantom]{Schémas en groupoïdes}
\item \hyperref[more-groupoids-section-phantom]{Compléments sur les schémas en groupoïdes}
\item \hyperref[etale-section-phantom]{Morphismes étales de schémas}
\end{enumerate}
Thèmes de théorie des schémas
\begin{enumerate}
\setcounter{enumi}{41}
\item \hyperref[chow-section-phantom]{Homologie de Chow}
\item \hyperref[intersection-section-phantom]{Théorie de l'intersection}
\item \hyperref[pic-section-phantom]{Schémas de Picard des courbes}
\item \hyperref[weil-section-phantom]{Théories cohomologiques de Weil}
\item \hyperref[adequate-section-phantom]{Modules adéquats}
\item \hyperref[dualizing-section-phantom]{Complexes dualisants}
\item \hyperref[duality-section-phantom]{Dualité pour les schémas}
\item \hyperref[discriminant-section-phantom]{Discriminants et différentes}
\item \hyperref[derham-section-phantom]{Cohomologie de de Rham}
\item \hyperref[local-cohomology-section-phantom]{Cohomologie locale}
\item \hyperref[algebraization-section-phantom]{Géométrie algébrique et formelle}
\item \hyperref[curves-section-phantom]{Courbes algébriques}
\item \hyperref[resolve-section-phantom]{Résolution des surfaces}
\item \hyperref[models-section-phantom]{Réduction semi-stable}
\item \hyperref[functors-section-phantom]{Foncteurs et morphismes}
\item \hyperref[equiv-section-phantom]{Catégories dérivées de variétés}
\item \hyperref[pione-section-phantom]{Groupes fondamentaux de schémas}
\item \hyperref[etale-cohomology-section-phantom]{Cohomologie étale}
\item \hyperref[crystalline-section-phantom]{Cohomologie cristalline}
\item \hyperref[proetale-section-phantom]{Cohomologie pro-étale}
\item \hyperref[relative-cycles-section-phantom]{Cycles relatifs}
\item \hyperref[more-etale-section-phantom]{Compléments de cohomologie étale}
\item \hyperref[trace-section-phantom]{La formule des traces}
\end{enumerate}
Espaces algébriques
\begin{enumerate}
\setcounter{enumi}{64}
\item \hyperref[spaces-section-phantom]{Espaces algébriques}
\item \hyperref[spaces-properties-section-phantom]{Propriétés des espaces algébriques}
\item \hyperref[spaces-morphisms-section-phantom]{Morphismes d'espaces algébriques}
\item \hyperref[decent-spaces-section-phantom]{Espaces algébriques décents}
\item \hyperref[spaces-cohomology-section-phantom]{Cohomologie des espaces algébriques}
\item \hyperref[spaces-limits-section-phantom]{Limites d'espaces algébriques}
\item \hyperref[spaces-divisors-section-phantom]{Diviseurs sur les espaces algébriques}
\item \hyperref[spaces-over-fields-section-phantom]{Espaces algébriques sur des corps}
\item \hyperref[spaces-topologies-section-phantom]{Topologies sur les espaces algébriques}
\item \hyperref[spaces-descent-section-phantom]{Descente et espaces algébriques}
\item \hyperref[spaces-perfect-section-phantom]{Catégories dérivées d'espaces}
\item \hyperref[spaces-more-morphisms-section-phantom]{Compléments sur les morphismes d'espaces}
\item \hyperref[spaces-flat-section-phantom]{Platitude sur les espaces algébriques}
\item \hyperref[spaces-groupoids-section-phantom]{Groupoïdes dans les espaces algébriques}
\item \hyperref[spaces-more-groupoids-section-phantom]{Compléments sur les groupoïdes dans les espaces}
\item \hyperref[bootstrap-section-phantom]{Amorçage}
\item \hyperref[spaces-pushouts-section-phantom]{Sommes amalgamées d'espaces algébriques}
\end{enumerate}
Thèmes de géométrie
\begin{enumerate}
\setcounter{enumi}{81}
\item \hyperref[spaces-chow-section-phantom]{Groupes de Chow des espaces}
\item \hyperref[groupoids-quotients-section-phantom]{Quotients de groupoïdes}
\item \hyperref[spaces-more-cohomology-section-phantom]{Compléments sur la cohomologie des espaces}
\item \hyperref[spaces-simplicial-section-phantom]{Espaces simpliciaux}
\item \hyperref[spaces-duality-section-phantom]{Dualité pour les espaces}
\item \hyperref[formal-spaces-section-phantom]{Espaces algébriques formels}
\item \hyperref[restricted-section-phantom]{Algébrisation des espaces formels}
\item \hyperref[spaces-resolve-section-phantom]{Résolution des surfaces revisitée}
\end{enumerate}
Théorie des déformations
\begin{enumerate}
\setcounter{enumi}{89}
\item \hyperref[formal-defos-section-phantom]{Théorie formelle des déformations}
\item \hyperref[defos-section-phantom]{Théorie des déformations}
\item \hyperref[cotangent-section-phantom]{Le complexe cotangent}
\item \hyperref[examples-defos-section-phantom]{Problèmes de déformation}
\end{enumerate}
Champs algébriques
\begin{enumerate}
\setcounter{enumi}{93}
\item \hyperref[algebraic-section-phantom]{Champs algébriques}
\item \hyperref[examples-stacks-section-phantom]{Exemples de champs}
\item \hyperref[stacks-sheaves-section-phantom]{Faisceaux sur les champs algébriques}
\item \hyperref[criteria-section-phantom]{Critères de représentabilité}
\item \hyperref[artin-section-phantom]{Axiomes d'Artin}
\item \hyperref[quot-section-phantom]{Espaces de Quot et de Hilbert}
\item \hyperref[stacks-properties-section-phantom]{Propriétés des champs algébriques}
\item \hyperref[stacks-morphisms-section-phantom]{Morphismes de champs algébriques}
\item \hyperref[stacks-limits-section-phantom]{Limites de champs algébriques}
\item \hyperref[stacks-cohomology-section-phantom]{Cohomologie des champs algébriques}
\item \hyperref[stacks-perfect-section-phantom]{Catégories dérivées de champs}
\item \hyperref[stacks-introduction-section-phantom]{Introduction aux champs algébriques}
\item \hyperref[stacks-more-morphisms-section-phantom]{Compléments sur les morphismes de champs}
\item \hyperref[stacks-geometry-section-phantom]{La géométrie des champs}
\end{enumerate}
Thèmes de théorie des espaces de modules
\begin{enumerate}
\setcounter{enumi}{107}
\item \hyperref[moduli-section-phantom]{Champs de modules}
\item \hyperref[moduli-curves-section-phantom]{Espaces de modules de courbes}
\end{enumerate}
Divers
\begin{enumerate}
\setcounter{enumi}{109}
\item \hyperref[examples-section-phantom]{Exemples}
\item \hyperref[exercises-section-phantom]{Exercices}
\item \hyperref[guide-section-phantom]{Guide bibliographique}
\item \hyperref[desirables-section-phantom]{Desiderata}
\item \hyperref[coding-section-phantom]{Style de programmation}
\item \hyperref[obsolete-section-phantom]{Obsolète}
\item \hyperref[fdl-section-phantom]{Licence de documentation libre GNU}
\item \hyperref[index-section-phantom]{Index généré automatiquement}
\end{enumerate}
\end{multicols}


\bibliography{my}
\bibliographystyle{amsalpha}

\end{document}
